\documentclass{article}
\usepackage{geometry}
\geometry{margin=1.5cm, vmargin={0pt,1cm}}
\setlength{\topmargin}{-1cm}
\setlength{\headheight}{12.5pt}
\setlength{\paperheight}{29.7cm}
\setlength{\textheight}{25.3cm}

% useful packages.
% \usepackage[UTF8]{ctex}
\usepackage{fancyhdr}
\usepackage{layout}
\usepackage{luacode}
\usepackage{tikz}
\usetikzlibrary{calc}

% Matplotlib PGF exports may emit \mathdefault in tick labels.
\providecommand{\mathdefault}[1]{#1}

% import common macros
% % % % % % % % % % % % % % % % % % % % % % % % % % % % % % % %
% Common Macros for LaTeX
% % % % % % % % % % % % % % % % % % % % % % % % % % % % % % % %

% Useful packages
\usepackage{amsfonts}
\usepackage{amsmath}
\usepackage{amssymb}
\usepackage{amsthm}
\usepackage{enumerate}
\usepackage{graphicx}
\usepackage{multicol}
\usepackage[ruled,linesnumbered]{algorithm2e}

% Math operators and common symbols
\newcommand{\dif}{\mathrm{d}}
\newcommand{\innerProd}[1]{\left\langle #1 \right\rangle}
\newcommand{\difFrac}[2]{\frac{\dif #1}{\dif #2}}
\newcommand{\pdfFrac}[2]{\frac{\partial #1}{\partial #2}}
\newcommand{\T}{\mathrm{T}}
\newcommand{\norm}[1]{\left\| #1 \right\|}
\newcommand{\abs}[1]{\left| #1 \right|}

% Vector and Matrix shortcuts
\newcommand{\vecTwo}[2]{
  \begin{bmatrix}
    #1 \\
    #2
\end{bmatrix}}
\newcommand{\vecThree}[3]{
  \begin{bmatrix}
    #1 \\
    #2 \\
    #3
\end{bmatrix}}
\newcommand{\vecTwoH}[2]{
  \begin{bmatrix}
    #1 & #2
  \end{bmatrix}
}
\newcommand{\vecThreeH}[3]{
  \begin{bmatrix}
    #1 & #2 & #3
  \end{bmatrix}
}

% Common fields
\newcommand{\R}{\mathbb{R}}
\newcommand{\C}{\mathbb{C}}
\newcommand{\Z}{\mathbb{Z}}
\newcommand{\N}{\mathbb{N}}

% Solutions and proofs
\newenvironment{solution}{
\begin{proof}[解]}{
\end{proof}}

% Utility to get environment variables (requires lualatex)
\newcommand{\getEnv}[2]{\directlua{tex.print(os.getenv("#1") or "#2")}}


% Name and uID
\newcommand{\name}{\getEnv{NAME}{NAME}}
\newcommand{\uid}{\getEnv{UID}{UID}}
\newcommand{\course}{Complex Analysis}
\newcommand{\hwNum}{1}

\begin{document}

\pagestyle{fancy}
\fancyhead{}
\lhead{\zhstr{\name} (\uid)}
\chead{\course\ \#\hwNum}
\rhead{\today}

\section{Textbook Exercise 2 on Page 24}

Let $\langle \cdot, \cdot \rangle$ denote the usual inner product in
$\mathbb{R}^2$. In other words, if $Z = (x_1, y_1)$ and $W = (x_2, y_2)$, then
\[
  \langle Z, W \rangle = x_1 x_2 + y_1 y_2.
\]
Similarly, we may define a Hermitian inner product $(\cdot, \cdot)$
in $\mathbb{C}$ by
\[
  (z, w) = z \bar{w}.
\]
The term Hermitian is used to describe the fact that $(\cdot, \cdot)$
is not symmetric, but rather satisfies the relation
\[
  (z, w) = \overline{(w, z)} \quad \text{for all } z, w \in \mathbb{C}.
\]
Show that
\[
  \langle z, w \rangle = \frac{1}{2} [(z, w) + (w, z)] = \text{Re}(z, w),
\]
where we use the usual identification $z = x + iy \in \mathbb{C}$
with $(x, y) \in \mathbb{R}^2$.

\begin{proof}
  Let $z = x_1 + i y_1$ and $w = x_2 + i y_2$, where $x_1, y_1, x_2,
  y_2 \in \R$.
  Then $\bar{w} = x_2 - i y_2$, and we compute
  \[
    (z, w) = z\bar{w} = (x_1 + i y_1)(x_2 - i y_2)
    = (x_1 x_2 + y_1 y_2) + i (y_1 x_2 - x_1 y_2).
  \]
  Taking the real part yields
  \[
    \Re(z, w) = x_1 x_2 + y_1 y_2 = \langle z, w \rangle,
  \]
  where we identify $z = (x_1, y_1)$ and $w = (x_2, y_2)$ with points in $\R^2$.

  Moreover, note that
  \[
    \overline{(z, w)} = \overline{z\bar{w}} = \bar{z}w = w\bar{z} = (w, z).
  \]
  Since $\Re(z, w) = \frac{1}{2}\big[(z, w) + \overline{(z, w)}\big]$
  for any complex number,
  we obtain
  \[
    \langle z, w \rangle = \Re(z, w)
    = \frac{1}{2}\big[(z, w) + (w, z)\big].
  \]
\end{proof}

\section{Textbook Exercise 3 on Page 25}

With $\omega = s e^{i\varphi}$, where $s \geq 0$ and $\varphi \in
\mathbb{R}$, solve the equation $z^n = \omega$ in $\mathbb{C}$ where
$n$ is a natural number. How many solutions are there?

\begin{solution}
  Since $s \ge 0$, there exists a unique nonnegative real number $r \ge 0$
  such that $r^n = s$. Write $z$ in polar form as $z = r e^{i\theta}$.
  Then
  \[
    z^n = (r e^{i\theta})^n = r^n e^{in\theta} = s e^{i\varphi} = \omega.
  \]
  Equating the moduli and arguments, we obtain $r^n = s$ (which holds by
  construction) and
  \[
    n\theta = \varphi + 2k\pi, \quad k \in \mathbb{Z}.
  \]
  Thus
  \[
    \theta = \frac{\varphi + 2k\pi}{n}, \quad k = 0, 1, \dots, n-1.
  \]
  For $k = 0, 1, \dots, n-1$, we obtain $n$ distinct solutions:
  \[
    z_k = r e^{i(\varphi + 2k\pi)/n}, \quad k = 0, 1, \dots, n-1.
  \]
  There are exactly $n$ solutions in total.
\end{solution}

\section{Exercise 1.40*}

\begin{definition}[Definition 1.32: Real Numbers]
  Real numbers, denoted by $\mathbb{R}$, are equivalence classes that
  partition the set of all Rational Cauchy Sequences (RCSs), i.e.,
  \[
    \mathbb{R} := \{ \{s_n\} : \{s_n\} \text{ is a RCS} \}.
  \]
\end{definition}

\begin{theorem}[Theorem 1.39: Cauchy 1821]
  A sequence of real numbers is convergent (with a real number as its
  limit) if and only if it is a Cauchy sequence.
\end{theorem}

Prove Theorem 1.39 via the viewpoint of Definition 1.32.

\begin{proof}
  ($\Rightarrow$) Suppose $\{r_n\}$ converges to $r \in \R$. For any
  $\varepsilon > 0$, there exists $N$ such that $|r_n - r| < \varepsilon/2$
  for all $n > N$. Then for all $m, n > N$,
  \[
    |r_n - r_m| \le |r_n - r| + |r - r_m| < \varepsilon,
  \]
  so $\{r_n\}$ is a Cauchy sequence.

  ($\Leftarrow$) Let $\{r_n\}$ be a Cauchy sequence in $\R$. Write each
  real number $r_n$ as an equivalence class of rational Cauchy sequences
  $r_n = [\{r_{n,k}\}]$, where $r_{n,k} \in \Q$ and $\{r_{n,k}\}_{k\ge1}$
  is a rational Cauchy sequence. Define the rational sequence $y_n =
  r_{n,n}$.

  \textbf{Step 1: $\{y_n\}$ is a rational Cauchy sequence.} Since
  $\{r_n\}$ is Cauchy in $\R$, for any $\varepsilon > 0$ there exists
  $N_1 \in \N$ such that $|r_n - r_m| < \varepsilon$ for all $m, n > N_1$.
  Fix $n, m > N_1$. For any $k \in \N$, the triangle inequality gives
  \[
    |r_{n,k} - r_{m,k}| \le |r_{n,k} - r_n| + |r_n - r_m| + |r_m - r_{m,k}|.
  \]
  Since each $\{r_{n,k}\}$ converges to $r_n$, there exists $K$ such that
  $|r_{n,k} - r_n| < \varepsilon$ and $|r_{m,k} - r_m| < \varepsilon$ for
  all $k > K$. Taking $k > \max\{K, N_1\}$, we obtain
  \[
    |r_{n,k} - r_{m,k}| < 3\varepsilon.
  \]
  Letting $k \to \infty$ and using the uniqueness of limits, we get
  $|r_n - r_m| \le 3\varepsilon$. Hence $\{y_n\}$ is a rational Cauchy
  sequence.

  \textbf{Step 2: Let $y = [\{y_n\}] \in \R$. Then $r_n \to y$.} For any
  $\varepsilon > 0$, since $\{r_n\}$ is Cauchy, there exists $N$ such that
  $|r_n - r_m| < \varepsilon$ for all $n, m > N$. By the argument in Step
  1, for sufficiently large $k$ we have
  \[
    |r_{n,k} - r_n| < \varepsilon \quad \text{for all } n > N.
  \]
  Thus for $n > N$,
  \[
    |r_n - y| = \lim_{k\to\infty} |r_{n,k} - y_k| \le \varepsilon.
  \]
  Therefore $r_n \to y$, so $\{r_n\}$ converges.
\end{proof}

\section{Exercise 1.56}

\begin{corollary}[Corollary 1.55]
  $\mathbb{R}$ is an ordered field while $\mathbb{C}$ is not.
\end{corollary}

Prove Corollary 1.55.

\begin{solution}
  \textbf{$\R$ is ordered.} Define $P = \{r \in \R : r > 0\}$, where
  $r > 0$ means that $r$ is the equivalence class of a rational
  Cauchy sequence $\{r_k\}$ such that $r_k > 0$ for sufficiently
  large $k$. Then $P$ satisfies the required properties of a positive
  cone.
  \textbf{$\C$ is not ordered.} Suppose $\C$ is an ordered field. Then
  either $i > 0$ or $-i > 0$. If $i > 0$, then $-1 = i \cdot i < 0$.
  If $-i > 0$, then $-1 = (-i) \cdot (-i) < 0$. In either case, we
  have a contradiction with ODF-1. Hence $\C$ is not ordered.
\end{solution}

\section{Exercise 1.60}

\begin{remark}[Formula 1.59: Solving Generic Quartics]
  The generic quartic equation $ax^4 + bx^3 + cx^2 + dx + e = 0$
  (WLOG $a=1$) is solved by:
  \begin{enumerate}
    \item Moving terms to get $x^4 + bx^3 = -cx^2 - dx - e$.
    \item Adding $(\frac{b}{2}x)^2$ to both sides to complete the
      square on the LHS.
    \item Introducing a parameter $y$ and adding $(x^2 +
      \frac{b}{2}x)y + \frac{1}{4}y^2$ to both sides.
    \item Choosing $y$ such that the RHS becomes a perfect square (by
      solving its cubic discriminant).
  \end{enumerate}
\end{remark}

Use Formula 1.59 to solve the equation
\[
  x^4 + 2x^3 + 2x + 1 = 0.
\]

\begin{solution}
  Move the quadratic, linear, and constant terms to the RHS to get
  \[
    x^4 + 2x^3 = -2x - 1.
  \]
  Adding $(\frac{2}{2}x)^2 = x^2$ to both sides yields
  \[
    (x^2 + x)^2 = x^2 - 2x - 1.
  \]
  Introducing a parameter $y$ and adding $(x^2 + x)y +
  \frac{1}{4}y^2$ to both sides gives
  \[
    (x^2 + x + \frac{y}{2})^2 = (x^2 - 2x - 1) + (x^2 + x)y + \frac{1}{4}y^2
    = (y+1) x^2 + (y - 2)x + \frac{1}{4}y^2 - 1.
  \]
  We want to choose $y$ such that the RHS becomes a perfect square,
  which means that:
  \[
    \begin{aligned}
      &(y-2) = 2 \sqrt{(y+1)(\frac{1}{4}y^2-1)} \\
      \Rightarrow& y^2 - 4y + 4 = 4(y+1)(\frac{1}{4}y^2-1) = y^3+y^2-4y-4. \\
      \Rightarrow& y^3 - 8 = 0. \\
      \Rightarrow& y = 2 \ \text{or} \pm 2\omega, \quad \text{where }
      \omega = e^{2\pi i/3}.
    \end{aligned}
  \]

  We take $y = 2$ to get
  \[
    (x^2 + x + 1)^2 = 3x^2.
  \]
  Thus, \[
    (x^2 + 1)(x^2 + 2x + 1) = 0
  \]
  and the solutions are $x = \pm i$ and $x = -1$.

  To check that these are indeed solutions, we can directly substitute them
  into the original equation:
  \begin{itemize}
    \item For $x = i$, we have $i^4 + 2i^3 + 2i + 1 = 1 - 2i + 2i + 1 = 0$.
    \item For $x = -i$, we have $(-i)^4 + 2(-i)^3 + 2(-i) + 1 = 1 +
      2i - 2i + 1 = 0$.
    \item For $x = -1$, we have $(-1)^4 + 2(-1)^3 + 2(-1) + 1 = 1 - 2
      - 2 + 1 = 0$.
  \end{itemize}

\end{solution}

\section{Exercise 1.67}
Construct four squares for a quadrilateral, as shown in [a]. Prove
that the line-segments joining the centers of opposite squares are
perpendicular and of equal length; see [b].

% The TikZ code is generated by *gemini-3-flash*.
\begin{figure}[h]
  \centering
  \begin{tikzpicture}[scale=1.2]
    % [a] Version
    \begin{scope}
      \node at (-1.5, 3.5) {[a]};
      % Define the quadrilateral vertices
      \coordinate (O) at (0, 0);
      \coordinate (B) at (1.5, 0.5);
      \coordinate (C) at (1.75, 2.25);
      \coordinate (D) at (-0.2, 1.8);

      % Draw the main quadrilateral (filled)
      \filldraw[fill=white!80, draw=black, thick] (O) -- (B) -- (C)
      -- (D) -- cycle;

      % Define helper for drawing squares and finding center
      % Side OA (here O to D)
      \coordinate (S1_1) at ($(O)!1!90:(D)$);
      \coordinate (S1_2) at ($(D)!1!-90:(O)$);
      \filldraw[fill=black!20!white!60, draw=black] (O) -- (D) --
      (S1_2) -- (S1_1) -- cycle;
      \coordinate (C1) at ($(O)!0.5!(S1_2)$);

      % Side AB (here D to C)
      \coordinate (S2_1) at ($(D)!1!90:(C)$);
      \coordinate (S2_2) at ($(C)!1!-90:(D)$);
      \filldraw[fill=black!20!white!60, draw=black] (D) -- (C) --
      (S2_2) -- (S2_1) -- cycle;
      \coordinate (C2) at ($(D)!0.5!(S2_2)$);

      % Side BC (here C to B)
      \coordinate (S3_1) at ($(C)!1!90:(B)$);
      \coordinate (S3_2) at ($(B)!1!-90:(C)$);
      \filldraw[fill=black!20!white!60, draw=black] (C) -- (B) --
      (S3_2) -- (S3_1) -- cycle;
      \coordinate (C3) at ($(C)!0.5!(S3_2)$);

      % Side CO (here B to O)
      \coordinate (S4_1) at ($(B)!1!90:(O)$);
      \coordinate (S4_2) at ($(O)!1!-90:(B)$);
      \filldraw[fill=black!20!white!60, draw=black] (B) -- (O) --
      (S4_2) -- (S4_1) -- cycle;
      \coordinate (C4) at ($(B)!0.5!(S4_2)$);

      % Dots for vertices and labels
      \filldraw[fill=black, draw=white] (O) circle (1pt) node[left,
      black] {$0$};
      \filldraw[fill=black, draw=white] (C1) circle (1.5pt)
      node[above, black] {$p$};
      \filldraw[fill=black, draw=white] (C2) circle (1.5pt)
      node[right, black] {$q$};
      \filldraw[fill=black, draw=white] (C3) circle (1.5pt)
      node[below, black] {$r$};
      \filldraw[fill=black, draw=white] (C4) circle (1.5pt)
      node[left, black] {$s$};

      % Lines and labels from image
      \draw[black, thick] (C1) -- (C3);
      \draw[black, thick] (C2) -- (C4);

      % Right angle marker
      \begin{scope}[shift={($(intersection of C1--C3 and C2--C4)$)}]
        \draw[black] (0, 0.2) -- (0.2, 0.2) -- (0.2, 0);
      \end{scope}

      % Vectors/labels
      \draw[->, >=stealth, black, thick] (O) -- ($(O)!0.5!(D)$)
      node[midway, right] {$2a$};
      \draw[->, >=stealth, black, thick] (D) -- ($(D)!0.5!(C)$)
      node[midway, below] {$2b$};
      \draw[->, >=stealth, black, thick] (C) -- ($(C)!0.5!(B)$)
      node[midway, left] {$2c$};
      \draw[->, >=stealth, black, thick] (B) -- ($(B)!0.5!(O)$)
      node[midway, above] {$2d$};
    \end{scope}

    % [b] Version
    \begin{scope}[xshift=6.5cm]
      \node at (-1.5, 3.5) {[b]};
      % Define the same vertices as [a]
      \coordinate (O) at (0, 0);
      \coordinate (B) at (1.5, 0.5);
      \coordinate (C) at (1.75, 2.25);
      \coordinate (D) at (-0.2, 1.8);

      % Draw the two squares involved: square on OD and square on OB
      % (or equivalent from [a])
      \filldraw[fill=black!20!white!60, draw=black] (O) -- (D) --
      ($(D)!1!-90:(O)$)-- ($(O)!1!90:(D)$) -- cycle;
      \filldraw[fill=black!20!white!60, draw=black] (B) -- (O) --
      ($(O)!1!-90:(B)$) -- ($(B)!1!90:(O)$) -- cycle;

      % Get centers p and s
      \coordinate (p) at ($(O)!0.5!($(D)!1!-90:(O)$)$);
      \coordinate (s) at ($(B)!0.5!($(O)!1!-90:(B)$)$);

      % Draw the relevant part of the quadrilateral
      \draw[black, thick] (D) -- (O) -- (B);
      \draw[black, thin] (D) -- (B);

      % Midpoint m of DB (This is the pivot in the visual proof)
      \coordinate (m) at ($(D)!0.5!(B)$);

      % Draw dashed lines connecting centers to m
      \draw[black, dashed, thick] (p) -- (m);
      \draw[black, dashed, thick] (s) -- (m);

      % Dots and labels
      \filldraw[fill=black, draw=black] (p) circle (1.5pt) node[above
      left, black] {$p$};
      \filldraw[fill=black, draw=black] (s) circle (1.5pt) node[below
      left, black] {$s$};
      \filldraw[fill=black, draw=black] (m) circle (1.5pt)
      node[right, black] {$m$};

      % Right angle marker at m
      \begin{scope}[shift={(m)}]
        \path let \p1=($(p)-(m)$), \n1={atan2(\y1,\x1)} in
        \pgfextra{\xdef\mAngle{\n1}};
        \begin{scope}[rotate=\mAngle]
          \draw[black] (0.2, 0) -- (0.2, 0.2) -- (0, 0.2);
        \end{scope}
      \end{scope}
    \end{scope}
  \end{tikzpicture}
\end{figure}

\begin{solution}
  Let the four edges of the quadrilateral be $2a, 2b, 2c, 2d$ as
  shown in the figure, and the four vertices be $O, D, C, B$ in
  counterclockwise order. For convenience, we set $O = 0$.
  Then we have:
  \[
    \begin{aligned}
      p &= a + ia \\
      q &= 2a + b + ib \\
      r &= 2a + 2b + c + ic \\
      s &= 2a + 2b + 2c + d + id \\
      0 &= a+b+c+d
    \end{aligned}
  \]
  Let $X = p - r$ and $Y = s - q$. Then we want to show that $X$ and
  $Y$ are perpendicular and of equal length, i.e., $Y = iX$.
  \[
    \begin{aligned}
      X &= p-r = -a - 2b - c + i(a - c) = -b+d+i(a-c) \\
      Y &= s-q = b + 2c + d + i (d - b) = c-a+i(d-b)
    \end{aligned}
  \]
  Thus \[
    iX = i(-b+d+i(a-c)) = c-a + i(d-b) = Y.
  \]
  So $X$ and $Y$ are perpendicular and of equal length.

\end{solution}

\section{Exercise 1.91}

\begin{lemma}[Lemma 1.90: Properties of the Exponential Function]
  The exponential function satisfies:
  \begin{itemize}
    \item \textbf{(EXP-1):} Its restriction to $\mathbb{R}$ is
      monotonically increasing and positive.
    \item \textbf{(EXP-2):} There exists a positive number $\pi$ such
      that $e^{\frac{\pi}{2} i} = i$ and $e^z = 1$ iff $z \in 2\pi i
      \mathbb{Z}$.
    \item \textbf{(EXP-3):} $\exp$ is periodic with period $2\pi i$.
    \item \textbf{(EXP-4):} $\theta \mapsto e^{i\theta}$ maps
      $\mathbb{R}$ onto the unit circle.
    \item \textbf{(EXP-5):} $w \in \mathbb{C} \setminus \{0\}$
      implies $w = e^z$ for some $z \in \mathbb{C}$.
  \end{itemize}
\end{lemma}

Prove Lemma 1.90 using the power series definition
$e^z = \sum_{n=0}^{\infty} \frac{z^n}{n!}$. \\
\textit{Hint: define $\cos(\theta) := \text{Re}(e^{i\theta})$,
  $\sin(\theta) := \text{Im}(e^{i\theta})$, and look for the zeros of
these functions.}

\begin{proof}
  \textbf{(EXP-1):} For real $x \in \R$, all terms in the series are
  positive for $x > 0$, so $e^x > 0$. For $x_1 < x_2$, we have
  $e^{x_1} < e^{x_2}$ since each term in $e^{x_2} - e^{x_1}$ is positive.

  \textbf{(EXP-2):} Define $\cos(\theta) = \Re(e^{i\theta})$ and
  $\sin(\theta) = \Im(e^{i\theta})$. Substituting $z = i\theta$ into the
  power series:
  \[
    e^{i\theta} = \sum_{n=0}^{\infty} \frac{(i\theta)^n}{n!}
    = \sum_{k=0}^{\infty} \frac{(-1)^k \theta^{2k}}{(2k)!}
    + i \sum_{k=0}^{\infty} \frac{(-1)^k \theta^{2k+1}}{(2k+1)!}
    = \cos\theta + i\sin\theta.
    \tag{1}
  \]
  From (1), we have $\cos^2\theta + \sin^2\theta = |e^{i\theta}|^2
  = e^{i\theta}e^{-i\theta} = 1$. Using $e^{i(\theta+\phi)} =
  e^{i\theta}e^{i\phi}$ and comparing real and imaginary parts yields
  the addition formulas for cosine and sine.

  To define $\pi$, we show that $\cos\theta$ has a unique positive zero.
  From the power series, $\cos\theta = 1 - \frac{\theta^2}{2!} +
  \frac{\theta^4}{4!} - \cdots$.
  For $0 < \theta \le 1$, we have $\cos\theta \ge 1 -
  \frac{\theta^2}{2} \ge \frac{1}{2}$.
  For $\theta = 2$, the series alternates with decreasing terms starting
  from $1 - 2^2/2! = -1$, so by the alternating series test, $\cos 2 < 0$.
  Since $\cos$ is continuous on $[1, 2]$, there exists
  $\pi_0 \in (2, 4)$ such that $\cos(\pi_0/2) = 0$. We define $\pi = \pi_0$.
  Then $\sin(\pi/2) = \sqrt{1 - \cos^2(\pi/2)} = 1$, so $e^{i\pi/2} = i$.

  Now we show $e^z = 1$ iff $z \in 2\pi i\mathbb{Z}$. First, from
  $e^{i\pi/2} = i$
  and $e^{z+w} = e^z e^w$, we have:
  \[
    e^{i\pi} = (e^{i\pi/2})^2 = i^2 = -1, \quad
    e^{i2\pi} = (e^{i\pi})^2 = (-1)^2 = 1.
  \]
  Hence $e^{2k\pi i} = 1$ for all $k \in \mathbb{Z}$.

  Conversely, suppose $e^z = 1$. Write $z = x + iy$ with $x, y \in \R$.
  Then $e^z = e^x e^{iy} = 1$ implies $e^x = 1$, so $x = 0$. Thus $e^{iy} = 1$.
  Let $y = 2k\pi + r$ with $k \in \mathbb{Z}$ and $r \in [0, 2\pi)$. Then
  $e^{iy} = e^{i2k\pi} e^{ir} = e^{ir} = 1$. So it suffices to show that
  the only solution to $e^{ir} = 1$ in $[0, 2\pi)$ is $r = 0$.

  From the power series, $e^{i\theta}$ traces a continuous path on the
  unit circle. We have shown $e^{i0} = 1$, $e^{i\pi/2} = i$, $e^{i\pi} = -1$,
  $e^{i3\pi/2} = -i$, and $e^{i2\pi} = 1$. By continuity, the equation
  $e^{ir} = 1$ can only hold at $r = 0$ within $[0, 2\pi)$. Thus $y = 2k\pi$.

  \textbf{(EXP-3):} Since $e^{2\pi i} = 1$ from (EXP-2), we have
  $e^{z + 2\pi i} = e^z e^{2\pi i} = e^z$. Conversely, if $e^{z_1} = e^{z_2}$,
  then $e^{z_1 - z_2} = 1$, so $z_1 - z_2 = 2\pi i k$ for some $k \in \Z$.

  \textbf{(EXP-4):} For any $\theta \in \R$, $|e^{i\theta}|^2 = 1$, so
  $e^{i\theta}$ lies on the unit circle. Conversely, any $w$ on the unit
  circle can be written as $w = e^{i\theta}$ for some $\theta \in \R$.

  \textbf{(EXP-5):} For $w \in \C \setminus \{0\}$, write $w = re^{i\theta}$
  with $r = |w| > 0$. Then $w = e^{\ln r + i\theta}$, where $\ln r$ is
  the real logarithm of $r$.
\end{proof}

\section{Exercise 1.99}
Prove $\tan(3\theta) = \frac{3T - T^3}{1 - 3T^2}$ where $T = \tan \theta$.

\begin{proof}
  First, we prove the addition formula for tangent. For any
  $\alpha, \beta \in \R$ with $\cos\alpha\cos\beta \neq 0$, we have
  \[
    \tan(\alpha + \beta)
    = \frac{\sin(\alpha + \beta)}{\cos(\alpha + \beta)}
    = \frac{\sin\alpha\cos\beta + \cos\alpha\sin\beta}
    {\cos\alpha\cos\beta - \sin\alpha\sin\beta}
    = \frac{\tan\alpha + \tan\beta}{1 - \tan\alpha\tan\beta}.
    \tag{1}
  \]

  Now compute $\tan(2\theta)$ using (1) with $\alpha = \beta = \theta$:
  \[
    \tan(2\theta) = \frac{2\tan\theta}{1 - \tan^2\theta}.
    \tag{2}
  \]

  Next, compute $\tan(3\theta) = \tan(2\theta + \theta)$ using (1) again:
  \[
    \tan(3\theta)
    = \frac{\tan(2\theta) + \tan\theta}{1 - \tan(2\theta)\tan\theta}.
  \]

  Substituting (2) into the above and letting $T = \tan\theta$:
  \[
    \begin{aligned}
      \tan(3\theta)
      &= \frac{\dfrac{2T}{1 - T^2} + T}{1 - \dfrac{2T}{1 - T^2} \cdot T} \\
      &= \frac{\dfrac{2T + T(1 - T^2)}{1 - T^2}}
      {\dfrac{1 - T^2 - 2T^2}{1 - T^2}} \\
      &= \frac{2T + T - T^3}{1 - T^2 - 2T^2} \\
      &= \frac{3T - T^3}{1 - 3T^2}.
    \end{aligned}
  \]
  This completes the proof.
\end{proof}

\section{Exercise 1.105}

\begin{corollary}[Corollary 1.101: Relations between Trigonometric
  and Hyperbolic Functions]
  Hyperbolic functions have period $2\pi i$ and satisfy:
  \begin{align*}
    \cosh(z+w)       & = \cosh z \cosh w + \sinh z \sinh w \\
    \sinh(z+w)       & = \sinh z \cosh w + \cosh z \sinh w \\
    \cosh z          & = \cos(iz), \quad \sinh z = -i \sin(iz) \\
    \cos z           & = \cosh(iz), \quad \sin z = -i \sinh(iz) \\
    \cosh^2 z - \sinh^2 z & = 1
  \end{align*}
\end{corollary}

What are the shapes of images of straight lines $x = \text{const}$
and $y = \text{const}$ under $z = x + iy \mapsto \cos(z)$? Draw a
dual-plane plot similar to that in Example 1.85 to answer this
question. (Hint: use Corollary 1.101)

\begin{solution}
  Using Corollary 1.101, we have
  \[
    \cos(z) = \cos(x + iy) = \cos x \cos(iy) - \sin x \sin(iy) = \cos x \cosh y - i \sin x \sinh y.
  \]
  Let $w = u + iv = \cos(z)$. Then:
  \[
    u = \cos x \cosh y, \quad v = -\sin x \sinh y.
  \]
  
  \begin{enumerate}
    \item \textbf{Vertical lines ($x = c$)}:
    If $c \neq n\pi/2$, we have $\frac{u^2}{\cos^2 c} - \frac{v^2}{\sin^2 c} = \cosh^2 y - \sinh^2 y = 1$, which represents a family of \textbf{hyperbolas}.
    \item \textbf{Horizontal lines ($y = d$)}:
    If $d > 0$, we have $\frac{u^2}{\cosh^2 d} + \frac{v^2}{\sinh^2 d} = \cos^2 x + \sin^2 x = 1$, which represents a family of \textbf{ellipses}.
  \end{enumerate}

  The dual-plane plot below (Figure \ref{fig:mapping_cos}) visualizes these conformal mappings.

  \begin{figure}[htbp]
    \centering
    \resizebox{0.9\textwidth}{!}{%% Creator: Matplotlib, PGF backend
%%
%% To include the figure in your LaTeX document, write
%%   \input{<filename>.pgf}
%%
%% Make sure the required packages are loaded in your preamble
%%   \usepackage{pgf}
%%
%% Also ensure that all the required font packages are loaded; for instance,
%% the lmodern package is sometimes necessary when using math font.
%%   \usepackage{lmodern}
%%
%% Figures using additional raster images can only be included by \input if
%% they are in the same directory as the main LaTeX file. For loading figures
%% from other directories you can use the `import` package
%%   \usepackage{import}
%%
%% and then include the figures with
%%   \import{<path to file>}{<filename>.pgf}
%%
%% Matplotlib used the following preamble
%%   \def\mathdefault#1{#1}
%%   \everymath=\expandafter{\the\everymath\displaystyle}
%%   \IfFileExists{scrextend.sty}{
%%     \usepackage[fontsize=10.000000pt]{scrextend}
%%   }{
%%     \renewcommand{\normalsize}{\fontsize{10.000000}{12.000000}\selectfont}
%%     \normalsize
%%   }
%%   \usepackage{amsmath}
%%   \usepackage{amssymb}
%%   \def\mathdefault#1{#1}
%%   \makeatletter\@ifpackageloaded{underscore}{}{\usepackage[strings]{underscore}}\makeatother
%%
\begingroup%
\makeatletter%
\begin{pgfpicture}%
\pgfpathrectangle{\pgfpointorigin}{\pgfqpoint{10.000000in}{4.500000in}}%
\pgfusepath{use as bounding box, clip}%
\begin{pgfscope}%
\pgfsetbuttcap%
\pgfsetmiterjoin%
\definecolor{currentfill}{rgb}{1.000000,1.000000,1.000000}%
\pgfsetfillcolor{currentfill}%
\pgfsetlinewidth{0.000000pt}%
\definecolor{currentstroke}{rgb}{1.000000,1.000000,1.000000}%
\pgfsetstrokecolor{currentstroke}%
\pgfsetdash{}{0pt}%
\pgfpathmoveto{\pgfqpoint{0.000000in}{0.000000in}}%
\pgfpathlineto{\pgfqpoint{10.000000in}{0.000000in}}%
\pgfpathlineto{\pgfqpoint{10.000000in}{4.500000in}}%
\pgfpathlineto{\pgfqpoint{0.000000in}{4.500000in}}%
\pgfpathlineto{\pgfqpoint{0.000000in}{0.000000in}}%
\pgfpathclose%
\pgfusepath{fill}%
\end{pgfscope}%
\begin{pgfscope}%
\pgfsetbuttcap%
\pgfsetmiterjoin%
\definecolor{currentfill}{rgb}{1.000000,1.000000,1.000000}%
\pgfsetfillcolor{currentfill}%
\pgfsetlinewidth{0.000000pt}%
\definecolor{currentstroke}{rgb}{0.000000,0.000000,0.000000}%
\pgfsetstrokecolor{currentstroke}%
\pgfsetstrokeopacity{0.000000}%
\pgfsetdash{}{0pt}%
\pgfpathmoveto{\pgfqpoint{0.495251in}{1.287254in}}%
\pgfpathlineto{\pgfqpoint{4.932639in}{1.287254in}}%
\pgfpathlineto{\pgfqpoint{4.932639in}{3.212746in}}%
\pgfpathlineto{\pgfqpoint{0.495251in}{3.212746in}}%
\pgfpathlineto{\pgfqpoint{0.495251in}{1.287254in}}%
\pgfpathclose%
\pgfusepath{fill}%
\end{pgfscope}%
\begin{pgfscope}%
\pgfpathrectangle{\pgfqpoint{0.495251in}{1.287254in}}{\pgfqpoint{4.437388in}{1.925493in}}%
\pgfusepath{clip}%
\pgfsetrectcap%
\pgfsetroundjoin%
\pgfsetlinewidth{0.803000pt}%
\definecolor{currentstroke}{rgb}{0.690196,0.690196,0.690196}%
\pgfsetstrokecolor{currentstroke}%
\pgfsetstrokeopacity{0.300000}%
\pgfsetdash{}{0pt}%
\pgfpathmoveto{\pgfqpoint{0.628043in}{1.287254in}}%
\pgfpathlineto{\pgfqpoint{0.628043in}{3.212746in}}%
\pgfusepath{stroke}%
\end{pgfscope}%
\begin{pgfscope}%
\pgfsetbuttcap%
\pgfsetroundjoin%
\definecolor{currentfill}{rgb}{0.000000,0.000000,0.000000}%
\pgfsetfillcolor{currentfill}%
\pgfsetlinewidth{0.803000pt}%
\definecolor{currentstroke}{rgb}{0.000000,0.000000,0.000000}%
\pgfsetstrokecolor{currentstroke}%
\pgfsetdash{}{0pt}%
\pgfsys@defobject{currentmarker}{\pgfqpoint{0.000000in}{-0.048611in}}{\pgfqpoint{0.000000in}{0.000000in}}{%
\pgfpathmoveto{\pgfqpoint{0.000000in}{0.000000in}}%
\pgfpathlineto{\pgfqpoint{0.000000in}{-0.048611in}}%
\pgfusepath{stroke,fill}%
}%
\begin{pgfscope}%
\pgfsys@transformshift{0.628043in}{1.287254in}%
\pgfsys@useobject{currentmarker}{}%
\end{pgfscope}%
\end{pgfscope}%
\begin{pgfscope}%
\definecolor{textcolor}{rgb}{0.000000,0.000000,0.000000}%
\pgfsetstrokecolor{textcolor}%
\pgfsetfillcolor{textcolor}%
\pgftext[x=0.628043in,y=1.190031in,,top]{\color{textcolor}{\rmfamily\fontsize{10.000000}{12.000000}\selectfont\catcode`\^=\active\def^{\ifmmode\sp\else\^{}\fi}\catcode`\%=\active\def%{\%}$-\pi$}}%
\end{pgfscope}%
\begin{pgfscope}%
\pgfpathrectangle{\pgfqpoint{0.495251in}{1.287254in}}{\pgfqpoint{4.437388in}{1.925493in}}%
\pgfusepath{clip}%
\pgfsetrectcap%
\pgfsetroundjoin%
\pgfsetlinewidth{0.803000pt}%
\definecolor{currentstroke}{rgb}{0.690196,0.690196,0.690196}%
\pgfsetstrokecolor{currentstroke}%
\pgfsetstrokeopacity{0.300000}%
\pgfsetdash{}{0pt}%
\pgfpathmoveto{\pgfqpoint{1.670994in}{1.287254in}}%
\pgfpathlineto{\pgfqpoint{1.670994in}{3.212746in}}%
\pgfusepath{stroke}%
\end{pgfscope}%
\begin{pgfscope}%
\pgfsetbuttcap%
\pgfsetroundjoin%
\definecolor{currentfill}{rgb}{0.000000,0.000000,0.000000}%
\pgfsetfillcolor{currentfill}%
\pgfsetlinewidth{0.803000pt}%
\definecolor{currentstroke}{rgb}{0.000000,0.000000,0.000000}%
\pgfsetstrokecolor{currentstroke}%
\pgfsetdash{}{0pt}%
\pgfsys@defobject{currentmarker}{\pgfqpoint{0.000000in}{-0.048611in}}{\pgfqpoint{0.000000in}{0.000000in}}{%
\pgfpathmoveto{\pgfqpoint{0.000000in}{0.000000in}}%
\pgfpathlineto{\pgfqpoint{0.000000in}{-0.048611in}}%
\pgfusepath{stroke,fill}%
}%
\begin{pgfscope}%
\pgfsys@transformshift{1.670994in}{1.287254in}%
\pgfsys@useobject{currentmarker}{}%
\end{pgfscope}%
\end{pgfscope}%
\begin{pgfscope}%
\definecolor{textcolor}{rgb}{0.000000,0.000000,0.000000}%
\pgfsetstrokecolor{textcolor}%
\pgfsetfillcolor{textcolor}%
\pgftext[x=1.670994in,y=1.190031in,,top]{\color{textcolor}{\rmfamily\fontsize{10.000000}{12.000000}\selectfont\catcode`\^=\active\def^{\ifmmode\sp\else\^{}\fi}\catcode`\%=\active\def%{\%}$-\pi/2$}}%
\end{pgfscope}%
\begin{pgfscope}%
\pgfpathrectangle{\pgfqpoint{0.495251in}{1.287254in}}{\pgfqpoint{4.437388in}{1.925493in}}%
\pgfusepath{clip}%
\pgfsetrectcap%
\pgfsetroundjoin%
\pgfsetlinewidth{0.803000pt}%
\definecolor{currentstroke}{rgb}{0.690196,0.690196,0.690196}%
\pgfsetstrokecolor{currentstroke}%
\pgfsetstrokeopacity{0.300000}%
\pgfsetdash{}{0pt}%
\pgfpathmoveto{\pgfqpoint{2.713945in}{1.287254in}}%
\pgfpathlineto{\pgfqpoint{2.713945in}{3.212746in}}%
\pgfusepath{stroke}%
\end{pgfscope}%
\begin{pgfscope}%
\pgfsetbuttcap%
\pgfsetroundjoin%
\definecolor{currentfill}{rgb}{0.000000,0.000000,0.000000}%
\pgfsetfillcolor{currentfill}%
\pgfsetlinewidth{0.803000pt}%
\definecolor{currentstroke}{rgb}{0.000000,0.000000,0.000000}%
\pgfsetstrokecolor{currentstroke}%
\pgfsetdash{}{0pt}%
\pgfsys@defobject{currentmarker}{\pgfqpoint{0.000000in}{-0.048611in}}{\pgfqpoint{0.000000in}{0.000000in}}{%
\pgfpathmoveto{\pgfqpoint{0.000000in}{0.000000in}}%
\pgfpathlineto{\pgfqpoint{0.000000in}{-0.048611in}}%
\pgfusepath{stroke,fill}%
}%
\begin{pgfscope}%
\pgfsys@transformshift{2.713945in}{1.287254in}%
\pgfsys@useobject{currentmarker}{}%
\end{pgfscope}%
\end{pgfscope}%
\begin{pgfscope}%
\definecolor{textcolor}{rgb}{0.000000,0.000000,0.000000}%
\pgfsetstrokecolor{textcolor}%
\pgfsetfillcolor{textcolor}%
\pgftext[x=2.713945in,y=1.190031in,,top]{\color{textcolor}{\rmfamily\fontsize{10.000000}{12.000000}\selectfont\catcode`\^=\active\def^{\ifmmode\sp\else\^{}\fi}\catcode`\%=\active\def%{\%}$0$}}%
\end{pgfscope}%
\begin{pgfscope}%
\pgfpathrectangle{\pgfqpoint{0.495251in}{1.287254in}}{\pgfqpoint{4.437388in}{1.925493in}}%
\pgfusepath{clip}%
\pgfsetrectcap%
\pgfsetroundjoin%
\pgfsetlinewidth{0.803000pt}%
\definecolor{currentstroke}{rgb}{0.690196,0.690196,0.690196}%
\pgfsetstrokecolor{currentstroke}%
\pgfsetstrokeopacity{0.300000}%
\pgfsetdash{}{0pt}%
\pgfpathmoveto{\pgfqpoint{3.756896in}{1.287254in}}%
\pgfpathlineto{\pgfqpoint{3.756896in}{3.212746in}}%
\pgfusepath{stroke}%
\end{pgfscope}%
\begin{pgfscope}%
\pgfsetbuttcap%
\pgfsetroundjoin%
\definecolor{currentfill}{rgb}{0.000000,0.000000,0.000000}%
\pgfsetfillcolor{currentfill}%
\pgfsetlinewidth{0.803000pt}%
\definecolor{currentstroke}{rgb}{0.000000,0.000000,0.000000}%
\pgfsetstrokecolor{currentstroke}%
\pgfsetdash{}{0pt}%
\pgfsys@defobject{currentmarker}{\pgfqpoint{0.000000in}{-0.048611in}}{\pgfqpoint{0.000000in}{0.000000in}}{%
\pgfpathmoveto{\pgfqpoint{0.000000in}{0.000000in}}%
\pgfpathlineto{\pgfqpoint{0.000000in}{-0.048611in}}%
\pgfusepath{stroke,fill}%
}%
\begin{pgfscope}%
\pgfsys@transformshift{3.756896in}{1.287254in}%
\pgfsys@useobject{currentmarker}{}%
\end{pgfscope}%
\end{pgfscope}%
\begin{pgfscope}%
\definecolor{textcolor}{rgb}{0.000000,0.000000,0.000000}%
\pgfsetstrokecolor{textcolor}%
\pgfsetfillcolor{textcolor}%
\pgftext[x=3.756896in,y=1.190031in,,top]{\color{textcolor}{\rmfamily\fontsize{10.000000}{12.000000}\selectfont\catcode`\^=\active\def^{\ifmmode\sp\else\^{}\fi}\catcode`\%=\active\def%{\%}$\pi/2$}}%
\end{pgfscope}%
\begin{pgfscope}%
\pgfpathrectangle{\pgfqpoint{0.495251in}{1.287254in}}{\pgfqpoint{4.437388in}{1.925493in}}%
\pgfusepath{clip}%
\pgfsetrectcap%
\pgfsetroundjoin%
\pgfsetlinewidth{0.803000pt}%
\definecolor{currentstroke}{rgb}{0.690196,0.690196,0.690196}%
\pgfsetstrokecolor{currentstroke}%
\pgfsetstrokeopacity{0.300000}%
\pgfsetdash{}{0pt}%
\pgfpathmoveto{\pgfqpoint{4.799846in}{1.287254in}}%
\pgfpathlineto{\pgfqpoint{4.799846in}{3.212746in}}%
\pgfusepath{stroke}%
\end{pgfscope}%
\begin{pgfscope}%
\pgfsetbuttcap%
\pgfsetroundjoin%
\definecolor{currentfill}{rgb}{0.000000,0.000000,0.000000}%
\pgfsetfillcolor{currentfill}%
\pgfsetlinewidth{0.803000pt}%
\definecolor{currentstroke}{rgb}{0.000000,0.000000,0.000000}%
\pgfsetstrokecolor{currentstroke}%
\pgfsetdash{}{0pt}%
\pgfsys@defobject{currentmarker}{\pgfqpoint{0.000000in}{-0.048611in}}{\pgfqpoint{0.000000in}{0.000000in}}{%
\pgfpathmoveto{\pgfqpoint{0.000000in}{0.000000in}}%
\pgfpathlineto{\pgfqpoint{0.000000in}{-0.048611in}}%
\pgfusepath{stroke,fill}%
}%
\begin{pgfscope}%
\pgfsys@transformshift{4.799846in}{1.287254in}%
\pgfsys@useobject{currentmarker}{}%
\end{pgfscope}%
\end{pgfscope}%
\begin{pgfscope}%
\definecolor{textcolor}{rgb}{0.000000,0.000000,0.000000}%
\pgfsetstrokecolor{textcolor}%
\pgfsetfillcolor{textcolor}%
\pgftext[x=4.799846in,y=1.190031in,,top]{\color{textcolor}{\rmfamily\fontsize{10.000000}{12.000000}\selectfont\catcode`\^=\active\def^{\ifmmode\sp\else\^{}\fi}\catcode`\%=\active\def%{\%}$\pi$}}%
\end{pgfscope}%
\begin{pgfscope}%
\definecolor{textcolor}{rgb}{0.000000,0.000000,0.000000}%
\pgfsetstrokecolor{textcolor}%
\pgfsetfillcolor{textcolor}%
\pgftext[x=2.713945in,y=0.995587in,,top]{\color{textcolor}{\rmfamily\fontsize{10.000000}{12.000000}\selectfont\catcode`\^=\active\def^{\ifmmode\sp\else\^{}\fi}\catcode`\%=\active\def%{\%}$x$}}%
\end{pgfscope}%
\begin{pgfscope}%
\pgfpathrectangle{\pgfqpoint{0.495251in}{1.287254in}}{\pgfqpoint{4.437388in}{1.925493in}}%
\pgfusepath{clip}%
\pgfsetrectcap%
\pgfsetroundjoin%
\pgfsetlinewidth{0.803000pt}%
\definecolor{currentstroke}{rgb}{0.690196,0.690196,0.690196}%
\pgfsetstrokecolor{currentstroke}%
\pgfsetstrokeopacity{0.300000}%
\pgfsetdash{}{0pt}%
\pgfpathmoveto{\pgfqpoint{0.495251in}{1.420046in}}%
\pgfpathlineto{\pgfqpoint{4.932639in}{1.420046in}}%
\pgfusepath{stroke}%
\end{pgfscope}%
\begin{pgfscope}%
\pgfsetbuttcap%
\pgfsetroundjoin%
\definecolor{currentfill}{rgb}{0.000000,0.000000,0.000000}%
\pgfsetfillcolor{currentfill}%
\pgfsetlinewidth{0.803000pt}%
\definecolor{currentstroke}{rgb}{0.000000,0.000000,0.000000}%
\pgfsetstrokecolor{currentstroke}%
\pgfsetdash{}{0pt}%
\pgfsys@defobject{currentmarker}{\pgfqpoint{-0.048611in}{0.000000in}}{\pgfqpoint{-0.000000in}{0.000000in}}{%
\pgfpathmoveto{\pgfqpoint{-0.000000in}{0.000000in}}%
\pgfpathlineto{\pgfqpoint{-0.048611in}{0.000000in}}%
\pgfusepath{stroke,fill}%
}%
\begin{pgfscope}%
\pgfsys@transformshift{0.495251in}{1.420046in}%
\pgfsys@useobject{currentmarker}{}%
\end{pgfscope}%
\end{pgfscope}%
\begin{pgfscope}%
\definecolor{textcolor}{rgb}{0.000000,0.000000,0.000000}%
\pgfsetstrokecolor{textcolor}%
\pgfsetfillcolor{textcolor}%
\pgftext[x=0.220559in, y=1.371852in, left, base]{\color{textcolor}{\rmfamily\fontsize{10.000000}{12.000000}\selectfont\catcode`\^=\active\def^{\ifmmode\sp\else\^{}\fi}\catcode`\%=\active\def%{\%}$\mathdefault{0.0}$}}%
\end{pgfscope}%
\begin{pgfscope}%
\pgfpathrectangle{\pgfqpoint{0.495251in}{1.287254in}}{\pgfqpoint{4.437388in}{1.925493in}}%
\pgfusepath{clip}%
\pgfsetrectcap%
\pgfsetroundjoin%
\pgfsetlinewidth{0.803000pt}%
\definecolor{currentstroke}{rgb}{0.690196,0.690196,0.690196}%
\pgfsetstrokecolor{currentstroke}%
\pgfsetstrokeopacity{0.300000}%
\pgfsetdash{}{0pt}%
\pgfpathmoveto{\pgfqpoint{0.495251in}{1.752028in}}%
\pgfpathlineto{\pgfqpoint{4.932639in}{1.752028in}}%
\pgfusepath{stroke}%
\end{pgfscope}%
\begin{pgfscope}%
\pgfsetbuttcap%
\pgfsetroundjoin%
\definecolor{currentfill}{rgb}{0.000000,0.000000,0.000000}%
\pgfsetfillcolor{currentfill}%
\pgfsetlinewidth{0.803000pt}%
\definecolor{currentstroke}{rgb}{0.000000,0.000000,0.000000}%
\pgfsetstrokecolor{currentstroke}%
\pgfsetdash{}{0pt}%
\pgfsys@defobject{currentmarker}{\pgfqpoint{-0.048611in}{0.000000in}}{\pgfqpoint{-0.000000in}{0.000000in}}{%
\pgfpathmoveto{\pgfqpoint{-0.000000in}{0.000000in}}%
\pgfpathlineto{\pgfqpoint{-0.048611in}{0.000000in}}%
\pgfusepath{stroke,fill}%
}%
\begin{pgfscope}%
\pgfsys@transformshift{0.495251in}{1.752028in}%
\pgfsys@useobject{currentmarker}{}%
\end{pgfscope}%
\end{pgfscope}%
\begin{pgfscope}%
\definecolor{textcolor}{rgb}{0.000000,0.000000,0.000000}%
\pgfsetstrokecolor{textcolor}%
\pgfsetfillcolor{textcolor}%
\pgftext[x=0.220559in, y=1.703833in, left, base]{\color{textcolor}{\rmfamily\fontsize{10.000000}{12.000000}\selectfont\catcode`\^=\active\def^{\ifmmode\sp\else\^{}\fi}\catcode`\%=\active\def%{\%}$\mathdefault{0.5}$}}%
\end{pgfscope}%
\begin{pgfscope}%
\pgfpathrectangle{\pgfqpoint{0.495251in}{1.287254in}}{\pgfqpoint{4.437388in}{1.925493in}}%
\pgfusepath{clip}%
\pgfsetrectcap%
\pgfsetroundjoin%
\pgfsetlinewidth{0.803000pt}%
\definecolor{currentstroke}{rgb}{0.690196,0.690196,0.690196}%
\pgfsetstrokecolor{currentstroke}%
\pgfsetstrokeopacity{0.300000}%
\pgfsetdash{}{0pt}%
\pgfpathmoveto{\pgfqpoint{0.495251in}{2.084009in}}%
\pgfpathlineto{\pgfqpoint{4.932639in}{2.084009in}}%
\pgfusepath{stroke}%
\end{pgfscope}%
\begin{pgfscope}%
\pgfsetbuttcap%
\pgfsetroundjoin%
\definecolor{currentfill}{rgb}{0.000000,0.000000,0.000000}%
\pgfsetfillcolor{currentfill}%
\pgfsetlinewidth{0.803000pt}%
\definecolor{currentstroke}{rgb}{0.000000,0.000000,0.000000}%
\pgfsetstrokecolor{currentstroke}%
\pgfsetdash{}{0pt}%
\pgfsys@defobject{currentmarker}{\pgfqpoint{-0.048611in}{0.000000in}}{\pgfqpoint{-0.000000in}{0.000000in}}{%
\pgfpathmoveto{\pgfqpoint{-0.000000in}{0.000000in}}%
\pgfpathlineto{\pgfqpoint{-0.048611in}{0.000000in}}%
\pgfusepath{stroke,fill}%
}%
\begin{pgfscope}%
\pgfsys@transformshift{0.495251in}{2.084009in}%
\pgfsys@useobject{currentmarker}{}%
\end{pgfscope}%
\end{pgfscope}%
\begin{pgfscope}%
\definecolor{textcolor}{rgb}{0.000000,0.000000,0.000000}%
\pgfsetstrokecolor{textcolor}%
\pgfsetfillcolor{textcolor}%
\pgftext[x=0.220559in, y=2.035815in, left, base]{\color{textcolor}{\rmfamily\fontsize{10.000000}{12.000000}\selectfont\catcode`\^=\active\def^{\ifmmode\sp\else\^{}\fi}\catcode`\%=\active\def%{\%}$\mathdefault{1.0}$}}%
\end{pgfscope}%
\begin{pgfscope}%
\pgfpathrectangle{\pgfqpoint{0.495251in}{1.287254in}}{\pgfqpoint{4.437388in}{1.925493in}}%
\pgfusepath{clip}%
\pgfsetrectcap%
\pgfsetroundjoin%
\pgfsetlinewidth{0.803000pt}%
\definecolor{currentstroke}{rgb}{0.690196,0.690196,0.690196}%
\pgfsetstrokecolor{currentstroke}%
\pgfsetstrokeopacity{0.300000}%
\pgfsetdash{}{0pt}%
\pgfpathmoveto{\pgfqpoint{0.495251in}{2.415991in}}%
\pgfpathlineto{\pgfqpoint{4.932639in}{2.415991in}}%
\pgfusepath{stroke}%
\end{pgfscope}%
\begin{pgfscope}%
\pgfsetbuttcap%
\pgfsetroundjoin%
\definecolor{currentfill}{rgb}{0.000000,0.000000,0.000000}%
\pgfsetfillcolor{currentfill}%
\pgfsetlinewidth{0.803000pt}%
\definecolor{currentstroke}{rgb}{0.000000,0.000000,0.000000}%
\pgfsetstrokecolor{currentstroke}%
\pgfsetdash{}{0pt}%
\pgfsys@defobject{currentmarker}{\pgfqpoint{-0.048611in}{0.000000in}}{\pgfqpoint{-0.000000in}{0.000000in}}{%
\pgfpathmoveto{\pgfqpoint{-0.000000in}{0.000000in}}%
\pgfpathlineto{\pgfqpoint{-0.048611in}{0.000000in}}%
\pgfusepath{stroke,fill}%
}%
\begin{pgfscope}%
\pgfsys@transformshift{0.495251in}{2.415991in}%
\pgfsys@useobject{currentmarker}{}%
\end{pgfscope}%
\end{pgfscope}%
\begin{pgfscope}%
\definecolor{textcolor}{rgb}{0.000000,0.000000,0.000000}%
\pgfsetstrokecolor{textcolor}%
\pgfsetfillcolor{textcolor}%
\pgftext[x=0.220559in, y=2.367796in, left, base]{\color{textcolor}{\rmfamily\fontsize{10.000000}{12.000000}\selectfont\catcode`\^=\active\def^{\ifmmode\sp\else\^{}\fi}\catcode`\%=\active\def%{\%}$\mathdefault{1.5}$}}%
\end{pgfscope}%
\begin{pgfscope}%
\pgfpathrectangle{\pgfqpoint{0.495251in}{1.287254in}}{\pgfqpoint{4.437388in}{1.925493in}}%
\pgfusepath{clip}%
\pgfsetrectcap%
\pgfsetroundjoin%
\pgfsetlinewidth{0.803000pt}%
\definecolor{currentstroke}{rgb}{0.690196,0.690196,0.690196}%
\pgfsetstrokecolor{currentstroke}%
\pgfsetstrokeopacity{0.300000}%
\pgfsetdash{}{0pt}%
\pgfpathmoveto{\pgfqpoint{0.495251in}{2.747972in}}%
\pgfpathlineto{\pgfqpoint{4.932639in}{2.747972in}}%
\pgfusepath{stroke}%
\end{pgfscope}%
\begin{pgfscope}%
\pgfsetbuttcap%
\pgfsetroundjoin%
\definecolor{currentfill}{rgb}{0.000000,0.000000,0.000000}%
\pgfsetfillcolor{currentfill}%
\pgfsetlinewidth{0.803000pt}%
\definecolor{currentstroke}{rgb}{0.000000,0.000000,0.000000}%
\pgfsetstrokecolor{currentstroke}%
\pgfsetdash{}{0pt}%
\pgfsys@defobject{currentmarker}{\pgfqpoint{-0.048611in}{0.000000in}}{\pgfqpoint{-0.000000in}{0.000000in}}{%
\pgfpathmoveto{\pgfqpoint{-0.000000in}{0.000000in}}%
\pgfpathlineto{\pgfqpoint{-0.048611in}{0.000000in}}%
\pgfusepath{stroke,fill}%
}%
\begin{pgfscope}%
\pgfsys@transformshift{0.495251in}{2.747972in}%
\pgfsys@useobject{currentmarker}{}%
\end{pgfscope}%
\end{pgfscope}%
\begin{pgfscope}%
\definecolor{textcolor}{rgb}{0.000000,0.000000,0.000000}%
\pgfsetstrokecolor{textcolor}%
\pgfsetfillcolor{textcolor}%
\pgftext[x=0.220559in, y=2.699778in, left, base]{\color{textcolor}{\rmfamily\fontsize{10.000000}{12.000000}\selectfont\catcode`\^=\active\def^{\ifmmode\sp\else\^{}\fi}\catcode`\%=\active\def%{\%}$\mathdefault{2.0}$}}%
\end{pgfscope}%
\begin{pgfscope}%
\pgfpathrectangle{\pgfqpoint{0.495251in}{1.287254in}}{\pgfqpoint{4.437388in}{1.925493in}}%
\pgfusepath{clip}%
\pgfsetrectcap%
\pgfsetroundjoin%
\pgfsetlinewidth{0.803000pt}%
\definecolor{currentstroke}{rgb}{0.690196,0.690196,0.690196}%
\pgfsetstrokecolor{currentstroke}%
\pgfsetstrokeopacity{0.300000}%
\pgfsetdash{}{0pt}%
\pgfpathmoveto{\pgfqpoint{0.495251in}{3.079954in}}%
\pgfpathlineto{\pgfqpoint{4.932639in}{3.079954in}}%
\pgfusepath{stroke}%
\end{pgfscope}%
\begin{pgfscope}%
\pgfsetbuttcap%
\pgfsetroundjoin%
\definecolor{currentfill}{rgb}{0.000000,0.000000,0.000000}%
\pgfsetfillcolor{currentfill}%
\pgfsetlinewidth{0.803000pt}%
\definecolor{currentstroke}{rgb}{0.000000,0.000000,0.000000}%
\pgfsetstrokecolor{currentstroke}%
\pgfsetdash{}{0pt}%
\pgfsys@defobject{currentmarker}{\pgfqpoint{-0.048611in}{0.000000in}}{\pgfqpoint{-0.000000in}{0.000000in}}{%
\pgfpathmoveto{\pgfqpoint{-0.000000in}{0.000000in}}%
\pgfpathlineto{\pgfqpoint{-0.048611in}{0.000000in}}%
\pgfusepath{stroke,fill}%
}%
\begin{pgfscope}%
\pgfsys@transformshift{0.495251in}{3.079954in}%
\pgfsys@useobject{currentmarker}{}%
\end{pgfscope}%
\end{pgfscope}%
\begin{pgfscope}%
\definecolor{textcolor}{rgb}{0.000000,0.000000,0.000000}%
\pgfsetstrokecolor{textcolor}%
\pgfsetfillcolor{textcolor}%
\pgftext[x=0.220559in, y=3.031759in, left, base]{\color{textcolor}{\rmfamily\fontsize{10.000000}{12.000000}\selectfont\catcode`\^=\active\def^{\ifmmode\sp\else\^{}\fi}\catcode`\%=\active\def%{\%}$\mathdefault{2.5}$}}%
\end{pgfscope}%
\begin{pgfscope}%
\definecolor{textcolor}{rgb}{0.000000,0.000000,0.000000}%
\pgfsetstrokecolor{textcolor}%
\pgfsetfillcolor{textcolor}%
\pgftext[x=0.165003in,y=2.250000in,,bottom,rotate=90.000000]{\color{textcolor}{\rmfamily\fontsize{10.000000}{12.000000}\selectfont\catcode`\^=\active\def^{\ifmmode\sp\else\^{}\fi}\catcode`\%=\active\def%{\%}$y$}}%
\end{pgfscope}%
\begin{pgfscope}%
\pgfpathrectangle{\pgfqpoint{0.495251in}{1.287254in}}{\pgfqpoint{4.437388in}{1.925493in}}%
\pgfusepath{clip}%
\pgfsetrectcap%
\pgfsetroundjoin%
\pgfsetlinewidth{1.505625pt}%
\definecolor{currentstroke}{rgb}{0.267004,0.004874,0.329415}%
\pgfsetstrokecolor{currentstroke}%
\pgfsetstrokeopacity{0.800000}%
\pgfsetdash{}{0pt}%
\pgfpathmoveto{\pgfqpoint{0.628043in}{1.420046in}}%
\pgfpathlineto{\pgfqpoint{0.628043in}{3.079954in}}%
\pgfpathlineto{\pgfqpoint{0.628043in}{3.079954in}}%
\pgfusepath{stroke}%
\end{pgfscope}%
\begin{pgfscope}%
\pgfpathrectangle{\pgfqpoint{0.495251in}{1.287254in}}{\pgfqpoint{4.437388in}{1.925493in}}%
\pgfusepath{clip}%
\pgfsetrectcap%
\pgfsetroundjoin%
\pgfsetlinewidth{1.505625pt}%
\definecolor{currentstroke}{rgb}{0.283229,0.120777,0.440584}%
\pgfsetstrokecolor{currentstroke}%
\pgfsetstrokeopacity{0.800000}%
\pgfsetdash{}{0pt}%
\pgfpathmoveto{\pgfqpoint{0.975694in}{1.420046in}}%
\pgfpathlineto{\pgfqpoint{0.975694in}{3.079954in}}%
\pgfpathlineto{\pgfqpoint{0.975694in}{3.079954in}}%
\pgfusepath{stroke}%
\end{pgfscope}%
\begin{pgfscope}%
\pgfpathrectangle{\pgfqpoint{0.495251in}{1.287254in}}{\pgfqpoint{4.437388in}{1.925493in}}%
\pgfusepath{clip}%
\pgfsetrectcap%
\pgfsetroundjoin%
\pgfsetlinewidth{1.505625pt}%
\definecolor{currentstroke}{rgb}{0.267968,0.223549,0.512008}%
\pgfsetstrokecolor{currentstroke}%
\pgfsetstrokeopacity{0.800000}%
\pgfsetdash{}{0pt}%
\pgfpathmoveto{\pgfqpoint{1.323344in}{1.420046in}}%
\pgfpathlineto{\pgfqpoint{1.323344in}{3.079954in}}%
\pgfpathlineto{\pgfqpoint{1.323344in}{3.079954in}}%
\pgfusepath{stroke}%
\end{pgfscope}%
\begin{pgfscope}%
\pgfpathrectangle{\pgfqpoint{0.495251in}{1.287254in}}{\pgfqpoint{4.437388in}{1.925493in}}%
\pgfusepath{clip}%
\pgfsetrectcap%
\pgfsetroundjoin%
\pgfsetlinewidth{1.505625pt}%
\definecolor{currentstroke}{rgb}{0.229739,0.322361,0.545706}%
\pgfsetstrokecolor{currentstroke}%
\pgfsetstrokeopacity{0.800000}%
\pgfsetdash{}{0pt}%
\pgfpathmoveto{\pgfqpoint{1.670994in}{1.420046in}}%
\pgfpathlineto{\pgfqpoint{1.670994in}{3.079954in}}%
\pgfpathlineto{\pgfqpoint{1.670994in}{3.079954in}}%
\pgfusepath{stroke}%
\end{pgfscope}%
\begin{pgfscope}%
\pgfpathrectangle{\pgfqpoint{0.495251in}{1.287254in}}{\pgfqpoint{4.437388in}{1.925493in}}%
\pgfusepath{clip}%
\pgfsetrectcap%
\pgfsetroundjoin%
\pgfsetlinewidth{1.505625pt}%
\definecolor{currentstroke}{rgb}{0.190631,0.407061,0.556089}%
\pgfsetstrokecolor{currentstroke}%
\pgfsetstrokeopacity{0.800000}%
\pgfsetdash{}{0pt}%
\pgfpathmoveto{\pgfqpoint{2.018644in}{1.420046in}}%
\pgfpathlineto{\pgfqpoint{2.018644in}{3.079954in}}%
\pgfpathlineto{\pgfqpoint{2.018644in}{3.079954in}}%
\pgfusepath{stroke}%
\end{pgfscope}%
\begin{pgfscope}%
\pgfpathrectangle{\pgfqpoint{0.495251in}{1.287254in}}{\pgfqpoint{4.437388in}{1.925493in}}%
\pgfusepath{clip}%
\pgfsetrectcap%
\pgfsetroundjoin%
\pgfsetlinewidth{1.505625pt}%
\definecolor{currentstroke}{rgb}{0.157729,0.485932,0.558013}%
\pgfsetstrokecolor{currentstroke}%
\pgfsetstrokeopacity{0.800000}%
\pgfsetdash{}{0pt}%
\pgfpathmoveto{\pgfqpoint{2.366295in}{1.420046in}}%
\pgfpathlineto{\pgfqpoint{2.366295in}{3.079954in}}%
\pgfpathlineto{\pgfqpoint{2.366295in}{3.079954in}}%
\pgfusepath{stroke}%
\end{pgfscope}%
\begin{pgfscope}%
\pgfpathrectangle{\pgfqpoint{0.495251in}{1.287254in}}{\pgfqpoint{4.437388in}{1.925493in}}%
\pgfusepath{clip}%
\pgfsetrectcap%
\pgfsetroundjoin%
\pgfsetlinewidth{1.505625pt}%
\definecolor{currentstroke}{rgb}{0.127568,0.566949,0.550556}%
\pgfsetstrokecolor{currentstroke}%
\pgfsetstrokeopacity{0.800000}%
\pgfsetdash{}{0pt}%
\pgfpathmoveto{\pgfqpoint{2.713945in}{1.420046in}}%
\pgfpathlineto{\pgfqpoint{2.713945in}{3.079954in}}%
\pgfpathlineto{\pgfqpoint{2.713945in}{3.079954in}}%
\pgfusepath{stroke}%
\end{pgfscope}%
\begin{pgfscope}%
\pgfpathrectangle{\pgfqpoint{0.495251in}{1.287254in}}{\pgfqpoint{4.437388in}{1.925493in}}%
\pgfusepath{clip}%
\pgfsetrectcap%
\pgfsetroundjoin%
\pgfsetlinewidth{1.505625pt}%
\definecolor{currentstroke}{rgb}{0.126326,0.644107,0.525311}%
\pgfsetstrokecolor{currentstroke}%
\pgfsetstrokeopacity{0.800000}%
\pgfsetdash{}{0pt}%
\pgfpathmoveto{\pgfqpoint{3.061595in}{1.420046in}}%
\pgfpathlineto{\pgfqpoint{3.061595in}{3.079954in}}%
\pgfpathlineto{\pgfqpoint{3.061595in}{3.079954in}}%
\pgfusepath{stroke}%
\end{pgfscope}%
\begin{pgfscope}%
\pgfpathrectangle{\pgfqpoint{0.495251in}{1.287254in}}{\pgfqpoint{4.437388in}{1.925493in}}%
\pgfusepath{clip}%
\pgfsetrectcap%
\pgfsetroundjoin%
\pgfsetlinewidth{1.505625pt}%
\definecolor{currentstroke}{rgb}{0.208030,0.718701,0.472873}%
\pgfsetstrokecolor{currentstroke}%
\pgfsetstrokeopacity{0.800000}%
\pgfsetdash{}{0pt}%
\pgfpathmoveto{\pgfqpoint{3.409245in}{1.420046in}}%
\pgfpathlineto{\pgfqpoint{3.409245in}{3.079954in}}%
\pgfpathlineto{\pgfqpoint{3.409245in}{3.079954in}}%
\pgfusepath{stroke}%
\end{pgfscope}%
\begin{pgfscope}%
\pgfpathrectangle{\pgfqpoint{0.495251in}{1.287254in}}{\pgfqpoint{4.437388in}{1.925493in}}%
\pgfusepath{clip}%
\pgfsetrectcap%
\pgfsetroundjoin%
\pgfsetlinewidth{1.505625pt}%
\definecolor{currentstroke}{rgb}{0.369214,0.788888,0.382914}%
\pgfsetstrokecolor{currentstroke}%
\pgfsetstrokeopacity{0.800000}%
\pgfsetdash{}{0pt}%
\pgfpathmoveto{\pgfqpoint{3.756896in}{1.420046in}}%
\pgfpathlineto{\pgfqpoint{3.756896in}{3.079954in}}%
\pgfpathlineto{\pgfqpoint{3.756896in}{3.079954in}}%
\pgfusepath{stroke}%
\end{pgfscope}%
\begin{pgfscope}%
\pgfpathrectangle{\pgfqpoint{0.495251in}{1.287254in}}{\pgfqpoint{4.437388in}{1.925493in}}%
\pgfusepath{clip}%
\pgfsetrectcap%
\pgfsetroundjoin%
\pgfsetlinewidth{1.505625pt}%
\definecolor{currentstroke}{rgb}{0.565498,0.842430,0.262877}%
\pgfsetstrokecolor{currentstroke}%
\pgfsetstrokeopacity{0.800000}%
\pgfsetdash{}{0pt}%
\pgfpathmoveto{\pgfqpoint{4.104546in}{1.420046in}}%
\pgfpathlineto{\pgfqpoint{4.104546in}{3.079954in}}%
\pgfpathlineto{\pgfqpoint{4.104546in}{3.079954in}}%
\pgfusepath{stroke}%
\end{pgfscope}%
\begin{pgfscope}%
\pgfpathrectangle{\pgfqpoint{0.495251in}{1.287254in}}{\pgfqpoint{4.437388in}{1.925493in}}%
\pgfusepath{clip}%
\pgfsetrectcap%
\pgfsetroundjoin%
\pgfsetlinewidth{1.505625pt}%
\definecolor{currentstroke}{rgb}{0.783315,0.879285,0.125405}%
\pgfsetstrokecolor{currentstroke}%
\pgfsetstrokeopacity{0.800000}%
\pgfsetdash{}{0pt}%
\pgfpathmoveto{\pgfqpoint{4.452196in}{1.420046in}}%
\pgfpathlineto{\pgfqpoint{4.452196in}{3.079954in}}%
\pgfpathlineto{\pgfqpoint{4.452196in}{3.079954in}}%
\pgfusepath{stroke}%
\end{pgfscope}%
\begin{pgfscope}%
\pgfpathrectangle{\pgfqpoint{0.495251in}{1.287254in}}{\pgfqpoint{4.437388in}{1.925493in}}%
\pgfusepath{clip}%
\pgfsetrectcap%
\pgfsetroundjoin%
\pgfsetlinewidth{1.505625pt}%
\definecolor{currentstroke}{rgb}{0.993248,0.906157,0.143936}%
\pgfsetstrokecolor{currentstroke}%
\pgfsetstrokeopacity{0.800000}%
\pgfsetdash{}{0pt}%
\pgfpathmoveto{\pgfqpoint{4.799846in}{1.420046in}}%
\pgfpathlineto{\pgfqpoint{4.799846in}{3.079954in}}%
\pgfpathlineto{\pgfqpoint{4.799846in}{3.079954in}}%
\pgfusepath{stroke}%
\end{pgfscope}%
\begin{pgfscope}%
\pgfpathrectangle{\pgfqpoint{0.495251in}{1.287254in}}{\pgfqpoint{4.437388in}{1.925493in}}%
\pgfusepath{clip}%
\pgfsetbuttcap%
\pgfsetroundjoin%
\pgfsetlinewidth{1.505625pt}%
\definecolor{currentstroke}{rgb}{0.050383,0.029803,0.527975}%
\pgfsetstrokecolor{currentstroke}%
\pgfsetstrokeopacity{0.800000}%
\pgfsetdash{{5.550000pt}{2.400000pt}}{0.000000pt}%
\pgfpathmoveto{\pgfqpoint{0.628043in}{1.420046in}}%
\pgfpathlineto{\pgfqpoint{4.799846in}{1.420046in}}%
\pgfpathlineto{\pgfqpoint{4.799846in}{1.420046in}}%
\pgfusepath{stroke}%
\end{pgfscope}%
\begin{pgfscope}%
\pgfpathrectangle{\pgfqpoint{0.495251in}{1.287254in}}{\pgfqpoint{4.437388in}{1.925493in}}%
\pgfusepath{clip}%
\pgfsetbuttcap%
\pgfsetroundjoin%
\pgfsetlinewidth{1.505625pt}%
\definecolor{currentstroke}{rgb}{0.299855,0.009561,0.631624}%
\pgfsetstrokecolor{currentstroke}%
\pgfsetstrokeopacity{0.800000}%
\pgfsetdash{{5.550000pt}{2.400000pt}}{0.000000pt}%
\pgfpathmoveto{\pgfqpoint{0.628043in}{1.627535in}}%
\pgfpathlineto{\pgfqpoint{4.799846in}{1.627535in}}%
\pgfpathlineto{\pgfqpoint{4.799846in}{1.627535in}}%
\pgfusepath{stroke}%
\end{pgfscope}%
\begin{pgfscope}%
\pgfpathrectangle{\pgfqpoint{0.495251in}{1.287254in}}{\pgfqpoint{4.437388in}{1.925493in}}%
\pgfusepath{clip}%
\pgfsetbuttcap%
\pgfsetroundjoin%
\pgfsetlinewidth{1.505625pt}%
\definecolor{currentstroke}{rgb}{0.494877,0.011990,0.657865}%
\pgfsetstrokecolor{currentstroke}%
\pgfsetstrokeopacity{0.800000}%
\pgfsetdash{{5.550000pt}{2.400000pt}}{0.000000pt}%
\pgfpathmoveto{\pgfqpoint{0.628043in}{1.835023in}}%
\pgfpathlineto{\pgfqpoint{4.799846in}{1.835023in}}%
\pgfpathlineto{\pgfqpoint{4.799846in}{1.835023in}}%
\pgfusepath{stroke}%
\end{pgfscope}%
\begin{pgfscope}%
\pgfpathrectangle{\pgfqpoint{0.495251in}{1.287254in}}{\pgfqpoint{4.437388in}{1.925493in}}%
\pgfusepath{clip}%
\pgfsetbuttcap%
\pgfsetroundjoin%
\pgfsetlinewidth{1.505625pt}%
\definecolor{currentstroke}{rgb}{0.665129,0.138566,0.585582}%
\pgfsetstrokecolor{currentstroke}%
\pgfsetstrokeopacity{0.800000}%
\pgfsetdash{{5.550000pt}{2.400000pt}}{0.000000pt}%
\pgfpathmoveto{\pgfqpoint{0.628043in}{2.042512in}}%
\pgfpathlineto{\pgfqpoint{4.799846in}{2.042512in}}%
\pgfpathlineto{\pgfqpoint{4.799846in}{2.042512in}}%
\pgfusepath{stroke}%
\end{pgfscope}%
\begin{pgfscope}%
\pgfpathrectangle{\pgfqpoint{0.495251in}{1.287254in}}{\pgfqpoint{4.437388in}{1.925493in}}%
\pgfusepath{clip}%
\pgfsetbuttcap%
\pgfsetroundjoin%
\pgfsetlinewidth{1.505625pt}%
\definecolor{currentstroke}{rgb}{0.798216,0.280197,0.469538}%
\pgfsetstrokecolor{currentstroke}%
\pgfsetstrokeopacity{0.800000}%
\pgfsetdash{{5.550000pt}{2.400000pt}}{0.000000pt}%
\pgfpathmoveto{\pgfqpoint{0.628043in}{2.250000in}}%
\pgfpathlineto{\pgfqpoint{4.799846in}{2.250000in}}%
\pgfpathlineto{\pgfqpoint{4.799846in}{2.250000in}}%
\pgfusepath{stroke}%
\end{pgfscope}%
\begin{pgfscope}%
\pgfpathrectangle{\pgfqpoint{0.495251in}{1.287254in}}{\pgfqpoint{4.437388in}{1.925493in}}%
\pgfusepath{clip}%
\pgfsetbuttcap%
\pgfsetroundjoin%
\pgfsetlinewidth{1.505625pt}%
\definecolor{currentstroke}{rgb}{0.901807,0.425087,0.359688}%
\pgfsetstrokecolor{currentstroke}%
\pgfsetstrokeopacity{0.800000}%
\pgfsetdash{{5.550000pt}{2.400000pt}}{0.000000pt}%
\pgfpathmoveto{\pgfqpoint{0.628043in}{2.457488in}}%
\pgfpathlineto{\pgfqpoint{4.799846in}{2.457488in}}%
\pgfpathlineto{\pgfqpoint{4.799846in}{2.457488in}}%
\pgfusepath{stroke}%
\end{pgfscope}%
\begin{pgfscope}%
\pgfpathrectangle{\pgfqpoint{0.495251in}{1.287254in}}{\pgfqpoint{4.437388in}{1.925493in}}%
\pgfusepath{clip}%
\pgfsetbuttcap%
\pgfsetroundjoin%
\pgfsetlinewidth{1.505625pt}%
\definecolor{currentstroke}{rgb}{0.973416,0.585761,0.251540}%
\pgfsetstrokecolor{currentstroke}%
\pgfsetstrokeopacity{0.800000}%
\pgfsetdash{{5.550000pt}{2.400000pt}}{0.000000pt}%
\pgfpathmoveto{\pgfqpoint{0.628043in}{2.664977in}}%
\pgfpathlineto{\pgfqpoint{4.799846in}{2.664977in}}%
\pgfpathlineto{\pgfqpoint{4.799846in}{2.664977in}}%
\pgfusepath{stroke}%
\end{pgfscope}%
\begin{pgfscope}%
\pgfpathrectangle{\pgfqpoint{0.495251in}{1.287254in}}{\pgfqpoint{4.437388in}{1.925493in}}%
\pgfusepath{clip}%
\pgfsetbuttcap%
\pgfsetroundjoin%
\pgfsetlinewidth{1.505625pt}%
\definecolor{currentstroke}{rgb}{0.993033,0.771720,0.154808}%
\pgfsetstrokecolor{currentstroke}%
\pgfsetstrokeopacity{0.800000}%
\pgfsetdash{{5.550000pt}{2.400000pt}}{0.000000pt}%
\pgfpathmoveto{\pgfqpoint{0.628043in}{2.872465in}}%
\pgfpathlineto{\pgfqpoint{4.799846in}{2.872465in}}%
\pgfpathlineto{\pgfqpoint{4.799846in}{2.872465in}}%
\pgfusepath{stroke}%
\end{pgfscope}%
\begin{pgfscope}%
\pgfpathrectangle{\pgfqpoint{0.495251in}{1.287254in}}{\pgfqpoint{4.437388in}{1.925493in}}%
\pgfusepath{clip}%
\pgfsetbuttcap%
\pgfsetroundjoin%
\pgfsetlinewidth{1.505625pt}%
\definecolor{currentstroke}{rgb}{0.940015,0.975158,0.131326}%
\pgfsetstrokecolor{currentstroke}%
\pgfsetstrokeopacity{0.800000}%
\pgfsetdash{{5.550000pt}{2.400000pt}}{0.000000pt}%
\pgfpathmoveto{\pgfqpoint{0.628043in}{3.079954in}}%
\pgfpathlineto{\pgfqpoint{4.799846in}{3.079954in}}%
\pgfpathlineto{\pgfqpoint{4.799846in}{3.079954in}}%
\pgfusepath{stroke}%
\end{pgfscope}%
\begin{pgfscope}%
\pgfpathrectangle{\pgfqpoint{0.495251in}{1.287254in}}{\pgfqpoint{4.437388in}{1.925493in}}%
\pgfusepath{clip}%
\pgfsetbuttcap%
\pgfsetroundjoin%
\definecolor{currentfill}{rgb}{0.000000,0.000000,0.000000}%
\pgfsetfillcolor{currentfill}%
\pgfsetfillopacity{0.900000}%
\pgfsetlinewidth{1.003750pt}%
\definecolor{currentstroke}{rgb}{0.000000,0.000000,0.000000}%
\pgfsetstrokecolor{currentstroke}%
\pgfsetstrokeopacity{0.900000}%
\pgfsetdash{}{0pt}%
\pgfsys@defobject{currentmarker}{\pgfqpoint{-0.001736in}{-0.001736in}}{\pgfqpoint{0.001736in}{0.001736in}}{%
\pgfpathmoveto{\pgfqpoint{0.000000in}{-0.001736in}}%
\pgfpathcurveto{\pgfqpoint{0.000460in}{-0.001736in}}{\pgfqpoint{0.000902in}{-0.001553in}}{\pgfqpoint{0.001228in}{-0.001228in}}%
\pgfpathcurveto{\pgfqpoint{0.001553in}{-0.000902in}}{\pgfqpoint{0.001736in}{-0.000460in}}{\pgfqpoint{0.001736in}{0.000000in}}%
\pgfpathcurveto{\pgfqpoint{0.001736in}{0.000460in}}{\pgfqpoint{0.001553in}{0.000902in}}{\pgfqpoint{0.001228in}{0.001228in}}%
\pgfpathcurveto{\pgfqpoint{0.000902in}{0.001553in}}{\pgfqpoint{0.000460in}{0.001736in}}{\pgfqpoint{0.000000in}{0.001736in}}%
\pgfpathcurveto{\pgfqpoint{-0.000460in}{0.001736in}}{\pgfqpoint{-0.000902in}{0.001553in}}{\pgfqpoint{-0.001228in}{0.001228in}}%
\pgfpathcurveto{\pgfqpoint{-0.001553in}{0.000902in}}{\pgfqpoint{-0.001736in}{0.000460in}}{\pgfqpoint{-0.001736in}{0.000000in}}%
\pgfpathcurveto{\pgfqpoint{-0.001736in}{-0.000460in}}{\pgfqpoint{-0.001553in}{-0.000902in}}{\pgfqpoint{-0.001228in}{-0.001228in}}%
\pgfpathcurveto{\pgfqpoint{-0.000902in}{-0.001553in}}{\pgfqpoint{-0.000460in}{-0.001736in}}{\pgfqpoint{0.000000in}{-0.001736in}}%
\pgfpathlineto{\pgfqpoint{0.000000in}{-0.001736in}}%
\pgfpathclose%
\pgfusepath{stroke,fill}%
}%
\begin{pgfscope}%
\pgfsys@transformshift{2.850210in}{1.753696in}%
\pgfsys@useobject{currentmarker}{}%
\end{pgfscope}%
\begin{pgfscope}%
\pgfsys@transformshift{2.871174in}{1.753696in}%
\pgfsys@useobject{currentmarker}{}%
\end{pgfscope}%
\begin{pgfscope}%
\pgfsys@transformshift{2.892137in}{1.753696in}%
\pgfsys@useobject{currentmarker}{}%
\end{pgfscope}%
\begin{pgfscope}%
\pgfsys@transformshift{2.913101in}{1.753696in}%
\pgfsys@useobject{currentmarker}{}%
\end{pgfscope}%
\begin{pgfscope}%
\pgfsys@transformshift{2.934065in}{1.753696in}%
\pgfsys@useobject{currentmarker}{}%
\end{pgfscope}%
\begin{pgfscope}%
\pgfsys@transformshift{2.955029in}{1.753696in}%
\pgfsys@useobject{currentmarker}{}%
\end{pgfscope}%
\begin{pgfscope}%
\pgfsys@transformshift{2.975993in}{1.753696in}%
\pgfsys@useobject{currentmarker}{}%
\end{pgfscope}%
\begin{pgfscope}%
\pgfsys@transformshift{2.996957in}{1.753696in}%
\pgfsys@useobject{currentmarker}{}%
\end{pgfscope}%
\begin{pgfscope}%
\pgfsys@transformshift{3.017920in}{1.753696in}%
\pgfsys@useobject{currentmarker}{}%
\end{pgfscope}%
\begin{pgfscope}%
\pgfsys@transformshift{3.038884in}{1.753696in}%
\pgfsys@useobject{currentmarker}{}%
\end{pgfscope}%
\begin{pgfscope}%
\pgfsys@transformshift{3.059848in}{1.753696in}%
\pgfsys@useobject{currentmarker}{}%
\end{pgfscope}%
\begin{pgfscope}%
\pgfsys@transformshift{3.080812in}{1.753696in}%
\pgfsys@useobject{currentmarker}{}%
\end{pgfscope}%
\begin{pgfscope}%
\pgfsys@transformshift{3.101776in}{1.753696in}%
\pgfsys@useobject{currentmarker}{}%
\end{pgfscope}%
\begin{pgfscope}%
\pgfsys@transformshift{3.122740in}{1.753696in}%
\pgfsys@useobject{currentmarker}{}%
\end{pgfscope}%
\begin{pgfscope}%
\pgfsys@transformshift{3.143703in}{1.753696in}%
\pgfsys@useobject{currentmarker}{}%
\end{pgfscope}%
\begin{pgfscope}%
\pgfsys@transformshift{3.164667in}{1.753696in}%
\pgfsys@useobject{currentmarker}{}%
\end{pgfscope}%
\begin{pgfscope}%
\pgfsys@transformshift{3.185631in}{1.753696in}%
\pgfsys@useobject{currentmarker}{}%
\end{pgfscope}%
\begin{pgfscope}%
\pgfsys@transformshift{3.206595in}{1.753696in}%
\pgfsys@useobject{currentmarker}{}%
\end{pgfscope}%
\begin{pgfscope}%
\pgfsys@transformshift{3.227559in}{1.753696in}%
\pgfsys@useobject{currentmarker}{}%
\end{pgfscope}%
\begin{pgfscope}%
\pgfsys@transformshift{2.850210in}{1.762037in}%
\pgfsys@useobject{currentmarker}{}%
\end{pgfscope}%
\begin{pgfscope}%
\pgfsys@transformshift{2.871174in}{1.762037in}%
\pgfsys@useobject{currentmarker}{}%
\end{pgfscope}%
\begin{pgfscope}%
\pgfsys@transformshift{2.892137in}{1.762037in}%
\pgfsys@useobject{currentmarker}{}%
\end{pgfscope}%
\begin{pgfscope}%
\pgfsys@transformshift{2.913101in}{1.762037in}%
\pgfsys@useobject{currentmarker}{}%
\end{pgfscope}%
\begin{pgfscope}%
\pgfsys@transformshift{2.934065in}{1.762037in}%
\pgfsys@useobject{currentmarker}{}%
\end{pgfscope}%
\begin{pgfscope}%
\pgfsys@transformshift{2.955029in}{1.762037in}%
\pgfsys@useobject{currentmarker}{}%
\end{pgfscope}%
\begin{pgfscope}%
\pgfsys@transformshift{2.975993in}{1.762037in}%
\pgfsys@useobject{currentmarker}{}%
\end{pgfscope}%
\begin{pgfscope}%
\pgfsys@transformshift{2.996957in}{1.762037in}%
\pgfsys@useobject{currentmarker}{}%
\end{pgfscope}%
\begin{pgfscope}%
\pgfsys@transformshift{3.017920in}{1.762037in}%
\pgfsys@useobject{currentmarker}{}%
\end{pgfscope}%
\begin{pgfscope}%
\pgfsys@transformshift{3.038884in}{1.762037in}%
\pgfsys@useobject{currentmarker}{}%
\end{pgfscope}%
\begin{pgfscope}%
\pgfsys@transformshift{3.059848in}{1.762037in}%
\pgfsys@useobject{currentmarker}{}%
\end{pgfscope}%
\begin{pgfscope}%
\pgfsys@transformshift{3.080812in}{1.762037in}%
\pgfsys@useobject{currentmarker}{}%
\end{pgfscope}%
\begin{pgfscope}%
\pgfsys@transformshift{3.101776in}{1.762037in}%
\pgfsys@useobject{currentmarker}{}%
\end{pgfscope}%
\begin{pgfscope}%
\pgfsys@transformshift{3.122740in}{1.762037in}%
\pgfsys@useobject{currentmarker}{}%
\end{pgfscope}%
\begin{pgfscope}%
\pgfsys@transformshift{3.143703in}{1.762037in}%
\pgfsys@useobject{currentmarker}{}%
\end{pgfscope}%
\begin{pgfscope}%
\pgfsys@transformshift{3.164667in}{1.762037in}%
\pgfsys@useobject{currentmarker}{}%
\end{pgfscope}%
\begin{pgfscope}%
\pgfsys@transformshift{3.185631in}{1.762037in}%
\pgfsys@useobject{currentmarker}{}%
\end{pgfscope}%
\begin{pgfscope}%
\pgfsys@transformshift{3.206595in}{1.762037in}%
\pgfsys@useobject{currentmarker}{}%
\end{pgfscope}%
\begin{pgfscope}%
\pgfsys@transformshift{3.227559in}{1.762037in}%
\pgfsys@useobject{currentmarker}{}%
\end{pgfscope}%
\begin{pgfscope}%
\pgfsys@transformshift{2.850210in}{1.770378in}%
\pgfsys@useobject{currentmarker}{}%
\end{pgfscope}%
\begin{pgfscope}%
\pgfsys@transformshift{2.871174in}{1.770378in}%
\pgfsys@useobject{currentmarker}{}%
\end{pgfscope}%
\begin{pgfscope}%
\pgfsys@transformshift{2.892137in}{1.770378in}%
\pgfsys@useobject{currentmarker}{}%
\end{pgfscope}%
\begin{pgfscope}%
\pgfsys@transformshift{2.913101in}{1.770378in}%
\pgfsys@useobject{currentmarker}{}%
\end{pgfscope}%
\begin{pgfscope}%
\pgfsys@transformshift{2.934065in}{1.770378in}%
\pgfsys@useobject{currentmarker}{}%
\end{pgfscope}%
\begin{pgfscope}%
\pgfsys@transformshift{2.955029in}{1.770378in}%
\pgfsys@useobject{currentmarker}{}%
\end{pgfscope}%
\begin{pgfscope}%
\pgfsys@transformshift{2.975993in}{1.770378in}%
\pgfsys@useobject{currentmarker}{}%
\end{pgfscope}%
\begin{pgfscope}%
\pgfsys@transformshift{2.996957in}{1.770378in}%
\pgfsys@useobject{currentmarker}{}%
\end{pgfscope}%
\begin{pgfscope}%
\pgfsys@transformshift{3.017920in}{1.770378in}%
\pgfsys@useobject{currentmarker}{}%
\end{pgfscope}%
\begin{pgfscope}%
\pgfsys@transformshift{3.038884in}{1.770378in}%
\pgfsys@useobject{currentmarker}{}%
\end{pgfscope}%
\begin{pgfscope}%
\pgfsys@transformshift{3.059848in}{1.770378in}%
\pgfsys@useobject{currentmarker}{}%
\end{pgfscope}%
\begin{pgfscope}%
\pgfsys@transformshift{3.080812in}{1.770378in}%
\pgfsys@useobject{currentmarker}{}%
\end{pgfscope}%
\begin{pgfscope}%
\pgfsys@transformshift{3.101776in}{1.770378in}%
\pgfsys@useobject{currentmarker}{}%
\end{pgfscope}%
\begin{pgfscope}%
\pgfsys@transformshift{3.122740in}{1.770378in}%
\pgfsys@useobject{currentmarker}{}%
\end{pgfscope}%
\begin{pgfscope}%
\pgfsys@transformshift{3.143703in}{1.770378in}%
\pgfsys@useobject{currentmarker}{}%
\end{pgfscope}%
\begin{pgfscope}%
\pgfsys@transformshift{3.164667in}{1.770378in}%
\pgfsys@useobject{currentmarker}{}%
\end{pgfscope}%
\begin{pgfscope}%
\pgfsys@transformshift{3.185631in}{1.770378in}%
\pgfsys@useobject{currentmarker}{}%
\end{pgfscope}%
\begin{pgfscope}%
\pgfsys@transformshift{3.206595in}{1.770378in}%
\pgfsys@useobject{currentmarker}{}%
\end{pgfscope}%
\begin{pgfscope}%
\pgfsys@transformshift{3.227559in}{1.770378in}%
\pgfsys@useobject{currentmarker}{}%
\end{pgfscope}%
\begin{pgfscope}%
\pgfsys@transformshift{2.850210in}{1.778720in}%
\pgfsys@useobject{currentmarker}{}%
\end{pgfscope}%
\begin{pgfscope}%
\pgfsys@transformshift{2.871174in}{1.778720in}%
\pgfsys@useobject{currentmarker}{}%
\end{pgfscope}%
\begin{pgfscope}%
\pgfsys@transformshift{2.892137in}{1.778720in}%
\pgfsys@useobject{currentmarker}{}%
\end{pgfscope}%
\begin{pgfscope}%
\pgfsys@transformshift{2.913101in}{1.778720in}%
\pgfsys@useobject{currentmarker}{}%
\end{pgfscope}%
\begin{pgfscope}%
\pgfsys@transformshift{2.934065in}{1.778720in}%
\pgfsys@useobject{currentmarker}{}%
\end{pgfscope}%
\begin{pgfscope}%
\pgfsys@transformshift{2.955029in}{1.778720in}%
\pgfsys@useobject{currentmarker}{}%
\end{pgfscope}%
\begin{pgfscope}%
\pgfsys@transformshift{2.975993in}{1.778720in}%
\pgfsys@useobject{currentmarker}{}%
\end{pgfscope}%
\begin{pgfscope}%
\pgfsys@transformshift{2.996957in}{1.778720in}%
\pgfsys@useobject{currentmarker}{}%
\end{pgfscope}%
\begin{pgfscope}%
\pgfsys@transformshift{3.017920in}{1.778720in}%
\pgfsys@useobject{currentmarker}{}%
\end{pgfscope}%
\begin{pgfscope}%
\pgfsys@transformshift{3.038884in}{1.778720in}%
\pgfsys@useobject{currentmarker}{}%
\end{pgfscope}%
\begin{pgfscope}%
\pgfsys@transformshift{3.059848in}{1.778720in}%
\pgfsys@useobject{currentmarker}{}%
\end{pgfscope}%
\begin{pgfscope}%
\pgfsys@transformshift{3.080812in}{1.778720in}%
\pgfsys@useobject{currentmarker}{}%
\end{pgfscope}%
\begin{pgfscope}%
\pgfsys@transformshift{3.101776in}{1.778720in}%
\pgfsys@useobject{currentmarker}{}%
\end{pgfscope}%
\begin{pgfscope}%
\pgfsys@transformshift{3.122740in}{1.778720in}%
\pgfsys@useobject{currentmarker}{}%
\end{pgfscope}%
\begin{pgfscope}%
\pgfsys@transformshift{3.143703in}{1.778720in}%
\pgfsys@useobject{currentmarker}{}%
\end{pgfscope}%
\begin{pgfscope}%
\pgfsys@transformshift{3.164667in}{1.778720in}%
\pgfsys@useobject{currentmarker}{}%
\end{pgfscope}%
\begin{pgfscope}%
\pgfsys@transformshift{3.185631in}{1.778720in}%
\pgfsys@useobject{currentmarker}{}%
\end{pgfscope}%
\begin{pgfscope}%
\pgfsys@transformshift{3.206595in}{1.778720in}%
\pgfsys@useobject{currentmarker}{}%
\end{pgfscope}%
\begin{pgfscope}%
\pgfsys@transformshift{3.227559in}{1.778720in}%
\pgfsys@useobject{currentmarker}{}%
\end{pgfscope}%
\begin{pgfscope}%
\pgfsys@transformshift{2.850210in}{1.787061in}%
\pgfsys@useobject{currentmarker}{}%
\end{pgfscope}%
\begin{pgfscope}%
\pgfsys@transformshift{2.871174in}{1.787061in}%
\pgfsys@useobject{currentmarker}{}%
\end{pgfscope}%
\begin{pgfscope}%
\pgfsys@transformshift{2.892137in}{1.787061in}%
\pgfsys@useobject{currentmarker}{}%
\end{pgfscope}%
\begin{pgfscope}%
\pgfsys@transformshift{2.913101in}{1.787061in}%
\pgfsys@useobject{currentmarker}{}%
\end{pgfscope}%
\begin{pgfscope}%
\pgfsys@transformshift{2.934065in}{1.787061in}%
\pgfsys@useobject{currentmarker}{}%
\end{pgfscope}%
\begin{pgfscope}%
\pgfsys@transformshift{2.955029in}{1.787061in}%
\pgfsys@useobject{currentmarker}{}%
\end{pgfscope}%
\begin{pgfscope}%
\pgfsys@transformshift{2.975993in}{1.787061in}%
\pgfsys@useobject{currentmarker}{}%
\end{pgfscope}%
\begin{pgfscope}%
\pgfsys@transformshift{2.996957in}{1.787061in}%
\pgfsys@useobject{currentmarker}{}%
\end{pgfscope}%
\begin{pgfscope}%
\pgfsys@transformshift{3.017920in}{1.787061in}%
\pgfsys@useobject{currentmarker}{}%
\end{pgfscope}%
\begin{pgfscope}%
\pgfsys@transformshift{3.038884in}{1.787061in}%
\pgfsys@useobject{currentmarker}{}%
\end{pgfscope}%
\begin{pgfscope}%
\pgfsys@transformshift{3.059848in}{1.787061in}%
\pgfsys@useobject{currentmarker}{}%
\end{pgfscope}%
\begin{pgfscope}%
\pgfsys@transformshift{3.080812in}{1.787061in}%
\pgfsys@useobject{currentmarker}{}%
\end{pgfscope}%
\begin{pgfscope}%
\pgfsys@transformshift{3.101776in}{1.787061in}%
\pgfsys@useobject{currentmarker}{}%
\end{pgfscope}%
\begin{pgfscope}%
\pgfsys@transformshift{3.122740in}{1.787061in}%
\pgfsys@useobject{currentmarker}{}%
\end{pgfscope}%
\begin{pgfscope}%
\pgfsys@transformshift{3.143703in}{1.787061in}%
\pgfsys@useobject{currentmarker}{}%
\end{pgfscope}%
\begin{pgfscope}%
\pgfsys@transformshift{3.164667in}{1.787061in}%
\pgfsys@useobject{currentmarker}{}%
\end{pgfscope}%
\begin{pgfscope}%
\pgfsys@transformshift{3.185631in}{1.787061in}%
\pgfsys@useobject{currentmarker}{}%
\end{pgfscope}%
\begin{pgfscope}%
\pgfsys@transformshift{3.206595in}{1.787061in}%
\pgfsys@useobject{currentmarker}{}%
\end{pgfscope}%
\begin{pgfscope}%
\pgfsys@transformshift{3.227559in}{1.787061in}%
\pgfsys@useobject{currentmarker}{}%
\end{pgfscope}%
\begin{pgfscope}%
\pgfsys@transformshift{2.850210in}{1.795402in}%
\pgfsys@useobject{currentmarker}{}%
\end{pgfscope}%
\begin{pgfscope}%
\pgfsys@transformshift{2.871174in}{1.795402in}%
\pgfsys@useobject{currentmarker}{}%
\end{pgfscope}%
\begin{pgfscope}%
\pgfsys@transformshift{2.892137in}{1.795402in}%
\pgfsys@useobject{currentmarker}{}%
\end{pgfscope}%
\begin{pgfscope}%
\pgfsys@transformshift{2.913101in}{1.795402in}%
\pgfsys@useobject{currentmarker}{}%
\end{pgfscope}%
\begin{pgfscope}%
\pgfsys@transformshift{2.934065in}{1.795402in}%
\pgfsys@useobject{currentmarker}{}%
\end{pgfscope}%
\begin{pgfscope}%
\pgfsys@transformshift{2.955029in}{1.795402in}%
\pgfsys@useobject{currentmarker}{}%
\end{pgfscope}%
\begin{pgfscope}%
\pgfsys@transformshift{2.975993in}{1.795402in}%
\pgfsys@useobject{currentmarker}{}%
\end{pgfscope}%
\begin{pgfscope}%
\pgfsys@transformshift{2.996957in}{1.795402in}%
\pgfsys@useobject{currentmarker}{}%
\end{pgfscope}%
\begin{pgfscope}%
\pgfsys@transformshift{3.017920in}{1.795402in}%
\pgfsys@useobject{currentmarker}{}%
\end{pgfscope}%
\begin{pgfscope}%
\pgfsys@transformshift{3.038884in}{1.795402in}%
\pgfsys@useobject{currentmarker}{}%
\end{pgfscope}%
\begin{pgfscope}%
\pgfsys@transformshift{3.059848in}{1.795402in}%
\pgfsys@useobject{currentmarker}{}%
\end{pgfscope}%
\begin{pgfscope}%
\pgfsys@transformshift{3.080812in}{1.795402in}%
\pgfsys@useobject{currentmarker}{}%
\end{pgfscope}%
\begin{pgfscope}%
\pgfsys@transformshift{3.101776in}{1.795402in}%
\pgfsys@useobject{currentmarker}{}%
\end{pgfscope}%
\begin{pgfscope}%
\pgfsys@transformshift{3.122740in}{1.795402in}%
\pgfsys@useobject{currentmarker}{}%
\end{pgfscope}%
\begin{pgfscope}%
\pgfsys@transformshift{3.143703in}{1.795402in}%
\pgfsys@useobject{currentmarker}{}%
\end{pgfscope}%
\begin{pgfscope}%
\pgfsys@transformshift{3.164667in}{1.795402in}%
\pgfsys@useobject{currentmarker}{}%
\end{pgfscope}%
\begin{pgfscope}%
\pgfsys@transformshift{3.185631in}{1.795402in}%
\pgfsys@useobject{currentmarker}{}%
\end{pgfscope}%
\begin{pgfscope}%
\pgfsys@transformshift{3.206595in}{1.795402in}%
\pgfsys@useobject{currentmarker}{}%
\end{pgfscope}%
\begin{pgfscope}%
\pgfsys@transformshift{3.227559in}{1.795402in}%
\pgfsys@useobject{currentmarker}{}%
\end{pgfscope}%
\begin{pgfscope}%
\pgfsys@transformshift{2.850210in}{1.803743in}%
\pgfsys@useobject{currentmarker}{}%
\end{pgfscope}%
\begin{pgfscope}%
\pgfsys@transformshift{2.871174in}{1.803743in}%
\pgfsys@useobject{currentmarker}{}%
\end{pgfscope}%
\begin{pgfscope}%
\pgfsys@transformshift{2.892137in}{1.803743in}%
\pgfsys@useobject{currentmarker}{}%
\end{pgfscope}%
\begin{pgfscope}%
\pgfsys@transformshift{2.913101in}{1.803743in}%
\pgfsys@useobject{currentmarker}{}%
\end{pgfscope}%
\begin{pgfscope}%
\pgfsys@transformshift{2.934065in}{1.803743in}%
\pgfsys@useobject{currentmarker}{}%
\end{pgfscope}%
\begin{pgfscope}%
\pgfsys@transformshift{2.955029in}{1.803743in}%
\pgfsys@useobject{currentmarker}{}%
\end{pgfscope}%
\begin{pgfscope}%
\pgfsys@transformshift{2.975993in}{1.803743in}%
\pgfsys@useobject{currentmarker}{}%
\end{pgfscope}%
\begin{pgfscope}%
\pgfsys@transformshift{2.996957in}{1.803743in}%
\pgfsys@useobject{currentmarker}{}%
\end{pgfscope}%
\begin{pgfscope}%
\pgfsys@transformshift{3.017920in}{1.803743in}%
\pgfsys@useobject{currentmarker}{}%
\end{pgfscope}%
\begin{pgfscope}%
\pgfsys@transformshift{3.038884in}{1.803743in}%
\pgfsys@useobject{currentmarker}{}%
\end{pgfscope}%
\begin{pgfscope}%
\pgfsys@transformshift{3.059848in}{1.803743in}%
\pgfsys@useobject{currentmarker}{}%
\end{pgfscope}%
\begin{pgfscope}%
\pgfsys@transformshift{3.080812in}{1.803743in}%
\pgfsys@useobject{currentmarker}{}%
\end{pgfscope}%
\begin{pgfscope}%
\pgfsys@transformshift{3.101776in}{1.803743in}%
\pgfsys@useobject{currentmarker}{}%
\end{pgfscope}%
\begin{pgfscope}%
\pgfsys@transformshift{3.122740in}{1.803743in}%
\pgfsys@useobject{currentmarker}{}%
\end{pgfscope}%
\begin{pgfscope}%
\pgfsys@transformshift{3.143703in}{1.803743in}%
\pgfsys@useobject{currentmarker}{}%
\end{pgfscope}%
\begin{pgfscope}%
\pgfsys@transformshift{3.164667in}{1.803743in}%
\pgfsys@useobject{currentmarker}{}%
\end{pgfscope}%
\begin{pgfscope}%
\pgfsys@transformshift{3.185631in}{1.803743in}%
\pgfsys@useobject{currentmarker}{}%
\end{pgfscope}%
\begin{pgfscope}%
\pgfsys@transformshift{3.206595in}{1.803743in}%
\pgfsys@useobject{currentmarker}{}%
\end{pgfscope}%
\begin{pgfscope}%
\pgfsys@transformshift{3.227559in}{1.803743in}%
\pgfsys@useobject{currentmarker}{}%
\end{pgfscope}%
\begin{pgfscope}%
\pgfsys@transformshift{2.850210in}{1.812085in}%
\pgfsys@useobject{currentmarker}{}%
\end{pgfscope}%
\begin{pgfscope}%
\pgfsys@transformshift{2.871174in}{1.812085in}%
\pgfsys@useobject{currentmarker}{}%
\end{pgfscope}%
\begin{pgfscope}%
\pgfsys@transformshift{2.892137in}{1.812085in}%
\pgfsys@useobject{currentmarker}{}%
\end{pgfscope}%
\begin{pgfscope}%
\pgfsys@transformshift{2.913101in}{1.812085in}%
\pgfsys@useobject{currentmarker}{}%
\end{pgfscope}%
\begin{pgfscope}%
\pgfsys@transformshift{2.934065in}{1.812085in}%
\pgfsys@useobject{currentmarker}{}%
\end{pgfscope}%
\begin{pgfscope}%
\pgfsys@transformshift{2.955029in}{1.812085in}%
\pgfsys@useobject{currentmarker}{}%
\end{pgfscope}%
\begin{pgfscope}%
\pgfsys@transformshift{2.975993in}{1.812085in}%
\pgfsys@useobject{currentmarker}{}%
\end{pgfscope}%
\begin{pgfscope}%
\pgfsys@transformshift{2.996957in}{1.812085in}%
\pgfsys@useobject{currentmarker}{}%
\end{pgfscope}%
\begin{pgfscope}%
\pgfsys@transformshift{3.017920in}{1.812085in}%
\pgfsys@useobject{currentmarker}{}%
\end{pgfscope}%
\begin{pgfscope}%
\pgfsys@transformshift{3.038884in}{1.812085in}%
\pgfsys@useobject{currentmarker}{}%
\end{pgfscope}%
\begin{pgfscope}%
\pgfsys@transformshift{3.059848in}{1.812085in}%
\pgfsys@useobject{currentmarker}{}%
\end{pgfscope}%
\begin{pgfscope}%
\pgfsys@transformshift{3.080812in}{1.812085in}%
\pgfsys@useobject{currentmarker}{}%
\end{pgfscope}%
\begin{pgfscope}%
\pgfsys@transformshift{3.101776in}{1.812085in}%
\pgfsys@useobject{currentmarker}{}%
\end{pgfscope}%
\begin{pgfscope}%
\pgfsys@transformshift{3.122740in}{1.812085in}%
\pgfsys@useobject{currentmarker}{}%
\end{pgfscope}%
\begin{pgfscope}%
\pgfsys@transformshift{3.143703in}{1.812085in}%
\pgfsys@useobject{currentmarker}{}%
\end{pgfscope}%
\begin{pgfscope}%
\pgfsys@transformshift{3.164667in}{1.812085in}%
\pgfsys@useobject{currentmarker}{}%
\end{pgfscope}%
\begin{pgfscope}%
\pgfsys@transformshift{3.185631in}{1.812085in}%
\pgfsys@useobject{currentmarker}{}%
\end{pgfscope}%
\begin{pgfscope}%
\pgfsys@transformshift{3.206595in}{1.812085in}%
\pgfsys@useobject{currentmarker}{}%
\end{pgfscope}%
\begin{pgfscope}%
\pgfsys@transformshift{3.227559in}{1.812085in}%
\pgfsys@useobject{currentmarker}{}%
\end{pgfscope}%
\begin{pgfscope}%
\pgfsys@transformshift{2.850210in}{1.820426in}%
\pgfsys@useobject{currentmarker}{}%
\end{pgfscope}%
\begin{pgfscope}%
\pgfsys@transformshift{2.871174in}{1.820426in}%
\pgfsys@useobject{currentmarker}{}%
\end{pgfscope}%
\begin{pgfscope}%
\pgfsys@transformshift{2.892137in}{1.820426in}%
\pgfsys@useobject{currentmarker}{}%
\end{pgfscope}%
\begin{pgfscope}%
\pgfsys@transformshift{2.913101in}{1.820426in}%
\pgfsys@useobject{currentmarker}{}%
\end{pgfscope}%
\begin{pgfscope}%
\pgfsys@transformshift{2.934065in}{1.820426in}%
\pgfsys@useobject{currentmarker}{}%
\end{pgfscope}%
\begin{pgfscope}%
\pgfsys@transformshift{2.955029in}{1.820426in}%
\pgfsys@useobject{currentmarker}{}%
\end{pgfscope}%
\begin{pgfscope}%
\pgfsys@transformshift{2.975993in}{1.820426in}%
\pgfsys@useobject{currentmarker}{}%
\end{pgfscope}%
\begin{pgfscope}%
\pgfsys@transformshift{2.996957in}{1.820426in}%
\pgfsys@useobject{currentmarker}{}%
\end{pgfscope}%
\begin{pgfscope}%
\pgfsys@transformshift{3.017920in}{1.820426in}%
\pgfsys@useobject{currentmarker}{}%
\end{pgfscope}%
\begin{pgfscope}%
\pgfsys@transformshift{3.038884in}{1.820426in}%
\pgfsys@useobject{currentmarker}{}%
\end{pgfscope}%
\begin{pgfscope}%
\pgfsys@transformshift{3.059848in}{1.820426in}%
\pgfsys@useobject{currentmarker}{}%
\end{pgfscope}%
\begin{pgfscope}%
\pgfsys@transformshift{3.080812in}{1.820426in}%
\pgfsys@useobject{currentmarker}{}%
\end{pgfscope}%
\begin{pgfscope}%
\pgfsys@transformshift{3.101776in}{1.820426in}%
\pgfsys@useobject{currentmarker}{}%
\end{pgfscope}%
\begin{pgfscope}%
\pgfsys@transformshift{3.122740in}{1.820426in}%
\pgfsys@useobject{currentmarker}{}%
\end{pgfscope}%
\begin{pgfscope}%
\pgfsys@transformshift{3.143703in}{1.820426in}%
\pgfsys@useobject{currentmarker}{}%
\end{pgfscope}%
\begin{pgfscope}%
\pgfsys@transformshift{3.164667in}{1.820426in}%
\pgfsys@useobject{currentmarker}{}%
\end{pgfscope}%
\begin{pgfscope}%
\pgfsys@transformshift{3.185631in}{1.820426in}%
\pgfsys@useobject{currentmarker}{}%
\end{pgfscope}%
\begin{pgfscope}%
\pgfsys@transformshift{3.206595in}{1.820426in}%
\pgfsys@useobject{currentmarker}{}%
\end{pgfscope}%
\begin{pgfscope}%
\pgfsys@transformshift{3.227559in}{1.820426in}%
\pgfsys@useobject{currentmarker}{}%
\end{pgfscope}%
\begin{pgfscope}%
\pgfsys@transformshift{2.850210in}{1.828767in}%
\pgfsys@useobject{currentmarker}{}%
\end{pgfscope}%
\begin{pgfscope}%
\pgfsys@transformshift{2.871174in}{1.828767in}%
\pgfsys@useobject{currentmarker}{}%
\end{pgfscope}%
\begin{pgfscope}%
\pgfsys@transformshift{2.892137in}{1.828767in}%
\pgfsys@useobject{currentmarker}{}%
\end{pgfscope}%
\begin{pgfscope}%
\pgfsys@transformshift{2.913101in}{1.828767in}%
\pgfsys@useobject{currentmarker}{}%
\end{pgfscope}%
\begin{pgfscope}%
\pgfsys@transformshift{2.934065in}{1.828767in}%
\pgfsys@useobject{currentmarker}{}%
\end{pgfscope}%
\begin{pgfscope}%
\pgfsys@transformshift{2.955029in}{1.828767in}%
\pgfsys@useobject{currentmarker}{}%
\end{pgfscope}%
\begin{pgfscope}%
\pgfsys@transformshift{2.975993in}{1.828767in}%
\pgfsys@useobject{currentmarker}{}%
\end{pgfscope}%
\begin{pgfscope}%
\pgfsys@transformshift{2.996957in}{1.828767in}%
\pgfsys@useobject{currentmarker}{}%
\end{pgfscope}%
\begin{pgfscope}%
\pgfsys@transformshift{3.017920in}{1.828767in}%
\pgfsys@useobject{currentmarker}{}%
\end{pgfscope}%
\begin{pgfscope}%
\pgfsys@transformshift{3.038884in}{1.828767in}%
\pgfsys@useobject{currentmarker}{}%
\end{pgfscope}%
\begin{pgfscope}%
\pgfsys@transformshift{3.059848in}{1.828767in}%
\pgfsys@useobject{currentmarker}{}%
\end{pgfscope}%
\begin{pgfscope}%
\pgfsys@transformshift{3.080812in}{1.828767in}%
\pgfsys@useobject{currentmarker}{}%
\end{pgfscope}%
\begin{pgfscope}%
\pgfsys@transformshift{3.101776in}{1.828767in}%
\pgfsys@useobject{currentmarker}{}%
\end{pgfscope}%
\begin{pgfscope}%
\pgfsys@transformshift{3.122740in}{1.828767in}%
\pgfsys@useobject{currentmarker}{}%
\end{pgfscope}%
\begin{pgfscope}%
\pgfsys@transformshift{3.143703in}{1.828767in}%
\pgfsys@useobject{currentmarker}{}%
\end{pgfscope}%
\begin{pgfscope}%
\pgfsys@transformshift{3.164667in}{1.828767in}%
\pgfsys@useobject{currentmarker}{}%
\end{pgfscope}%
\begin{pgfscope}%
\pgfsys@transformshift{3.185631in}{1.828767in}%
\pgfsys@useobject{currentmarker}{}%
\end{pgfscope}%
\begin{pgfscope}%
\pgfsys@transformshift{3.206595in}{1.828767in}%
\pgfsys@useobject{currentmarker}{}%
\end{pgfscope}%
\begin{pgfscope}%
\pgfsys@transformshift{3.227559in}{1.828767in}%
\pgfsys@useobject{currentmarker}{}%
\end{pgfscope}%
\begin{pgfscope}%
\pgfsys@transformshift{2.850210in}{1.837108in}%
\pgfsys@useobject{currentmarker}{}%
\end{pgfscope}%
\begin{pgfscope}%
\pgfsys@transformshift{2.871174in}{1.837108in}%
\pgfsys@useobject{currentmarker}{}%
\end{pgfscope}%
\begin{pgfscope}%
\pgfsys@transformshift{2.892137in}{1.837108in}%
\pgfsys@useobject{currentmarker}{}%
\end{pgfscope}%
\begin{pgfscope}%
\pgfsys@transformshift{2.913101in}{1.837108in}%
\pgfsys@useobject{currentmarker}{}%
\end{pgfscope}%
\begin{pgfscope}%
\pgfsys@transformshift{2.934065in}{1.837108in}%
\pgfsys@useobject{currentmarker}{}%
\end{pgfscope}%
\begin{pgfscope}%
\pgfsys@transformshift{2.955029in}{1.837108in}%
\pgfsys@useobject{currentmarker}{}%
\end{pgfscope}%
\begin{pgfscope}%
\pgfsys@transformshift{2.975993in}{1.837108in}%
\pgfsys@useobject{currentmarker}{}%
\end{pgfscope}%
\begin{pgfscope}%
\pgfsys@transformshift{2.996957in}{1.837108in}%
\pgfsys@useobject{currentmarker}{}%
\end{pgfscope}%
\begin{pgfscope}%
\pgfsys@transformshift{3.017920in}{1.837108in}%
\pgfsys@useobject{currentmarker}{}%
\end{pgfscope}%
\begin{pgfscope}%
\pgfsys@transformshift{3.038884in}{1.837108in}%
\pgfsys@useobject{currentmarker}{}%
\end{pgfscope}%
\begin{pgfscope}%
\pgfsys@transformshift{3.059848in}{1.837108in}%
\pgfsys@useobject{currentmarker}{}%
\end{pgfscope}%
\begin{pgfscope}%
\pgfsys@transformshift{3.080812in}{1.837108in}%
\pgfsys@useobject{currentmarker}{}%
\end{pgfscope}%
\begin{pgfscope}%
\pgfsys@transformshift{3.101776in}{1.837108in}%
\pgfsys@useobject{currentmarker}{}%
\end{pgfscope}%
\begin{pgfscope}%
\pgfsys@transformshift{3.122740in}{1.837108in}%
\pgfsys@useobject{currentmarker}{}%
\end{pgfscope}%
\begin{pgfscope}%
\pgfsys@transformshift{3.143703in}{1.837108in}%
\pgfsys@useobject{currentmarker}{}%
\end{pgfscope}%
\begin{pgfscope}%
\pgfsys@transformshift{3.164667in}{1.837108in}%
\pgfsys@useobject{currentmarker}{}%
\end{pgfscope}%
\begin{pgfscope}%
\pgfsys@transformshift{3.185631in}{1.837108in}%
\pgfsys@useobject{currentmarker}{}%
\end{pgfscope}%
\begin{pgfscope}%
\pgfsys@transformshift{3.206595in}{1.837108in}%
\pgfsys@useobject{currentmarker}{}%
\end{pgfscope}%
\begin{pgfscope}%
\pgfsys@transformshift{3.227559in}{1.837108in}%
\pgfsys@useobject{currentmarker}{}%
\end{pgfscope}%
\begin{pgfscope}%
\pgfsys@transformshift{2.850210in}{1.845450in}%
\pgfsys@useobject{currentmarker}{}%
\end{pgfscope}%
\begin{pgfscope}%
\pgfsys@transformshift{2.871174in}{1.845450in}%
\pgfsys@useobject{currentmarker}{}%
\end{pgfscope}%
\begin{pgfscope}%
\pgfsys@transformshift{2.892137in}{1.845450in}%
\pgfsys@useobject{currentmarker}{}%
\end{pgfscope}%
\begin{pgfscope}%
\pgfsys@transformshift{2.913101in}{1.845450in}%
\pgfsys@useobject{currentmarker}{}%
\end{pgfscope}%
\begin{pgfscope}%
\pgfsys@transformshift{2.934065in}{1.845450in}%
\pgfsys@useobject{currentmarker}{}%
\end{pgfscope}%
\begin{pgfscope}%
\pgfsys@transformshift{2.955029in}{1.845450in}%
\pgfsys@useobject{currentmarker}{}%
\end{pgfscope}%
\begin{pgfscope}%
\pgfsys@transformshift{2.975993in}{1.845450in}%
\pgfsys@useobject{currentmarker}{}%
\end{pgfscope}%
\begin{pgfscope}%
\pgfsys@transformshift{2.996957in}{1.845450in}%
\pgfsys@useobject{currentmarker}{}%
\end{pgfscope}%
\begin{pgfscope}%
\pgfsys@transformshift{3.017920in}{1.845450in}%
\pgfsys@useobject{currentmarker}{}%
\end{pgfscope}%
\begin{pgfscope}%
\pgfsys@transformshift{3.038884in}{1.845450in}%
\pgfsys@useobject{currentmarker}{}%
\end{pgfscope}%
\begin{pgfscope}%
\pgfsys@transformshift{3.059848in}{1.845450in}%
\pgfsys@useobject{currentmarker}{}%
\end{pgfscope}%
\begin{pgfscope}%
\pgfsys@transformshift{3.080812in}{1.845450in}%
\pgfsys@useobject{currentmarker}{}%
\end{pgfscope}%
\begin{pgfscope}%
\pgfsys@transformshift{3.101776in}{1.845450in}%
\pgfsys@useobject{currentmarker}{}%
\end{pgfscope}%
\begin{pgfscope}%
\pgfsys@transformshift{3.122740in}{1.845450in}%
\pgfsys@useobject{currentmarker}{}%
\end{pgfscope}%
\begin{pgfscope}%
\pgfsys@transformshift{3.143703in}{1.845450in}%
\pgfsys@useobject{currentmarker}{}%
\end{pgfscope}%
\begin{pgfscope}%
\pgfsys@transformshift{3.164667in}{1.845450in}%
\pgfsys@useobject{currentmarker}{}%
\end{pgfscope}%
\begin{pgfscope}%
\pgfsys@transformshift{3.185631in}{1.845450in}%
\pgfsys@useobject{currentmarker}{}%
\end{pgfscope}%
\begin{pgfscope}%
\pgfsys@transformshift{3.206595in}{1.845450in}%
\pgfsys@useobject{currentmarker}{}%
\end{pgfscope}%
\begin{pgfscope}%
\pgfsys@transformshift{3.227559in}{1.845450in}%
\pgfsys@useobject{currentmarker}{}%
\end{pgfscope}%
\begin{pgfscope}%
\pgfsys@transformshift{2.850210in}{1.853791in}%
\pgfsys@useobject{currentmarker}{}%
\end{pgfscope}%
\begin{pgfscope}%
\pgfsys@transformshift{2.871174in}{1.853791in}%
\pgfsys@useobject{currentmarker}{}%
\end{pgfscope}%
\begin{pgfscope}%
\pgfsys@transformshift{2.892137in}{1.853791in}%
\pgfsys@useobject{currentmarker}{}%
\end{pgfscope}%
\begin{pgfscope}%
\pgfsys@transformshift{2.913101in}{1.853791in}%
\pgfsys@useobject{currentmarker}{}%
\end{pgfscope}%
\begin{pgfscope}%
\pgfsys@transformshift{2.934065in}{1.853791in}%
\pgfsys@useobject{currentmarker}{}%
\end{pgfscope}%
\begin{pgfscope}%
\pgfsys@transformshift{2.955029in}{1.853791in}%
\pgfsys@useobject{currentmarker}{}%
\end{pgfscope}%
\begin{pgfscope}%
\pgfsys@transformshift{2.975993in}{1.853791in}%
\pgfsys@useobject{currentmarker}{}%
\end{pgfscope}%
\begin{pgfscope}%
\pgfsys@transformshift{2.996957in}{1.853791in}%
\pgfsys@useobject{currentmarker}{}%
\end{pgfscope}%
\begin{pgfscope}%
\pgfsys@transformshift{3.017920in}{1.853791in}%
\pgfsys@useobject{currentmarker}{}%
\end{pgfscope}%
\begin{pgfscope}%
\pgfsys@transformshift{3.038884in}{1.853791in}%
\pgfsys@useobject{currentmarker}{}%
\end{pgfscope}%
\begin{pgfscope}%
\pgfsys@transformshift{3.059848in}{1.853791in}%
\pgfsys@useobject{currentmarker}{}%
\end{pgfscope}%
\begin{pgfscope}%
\pgfsys@transformshift{3.080812in}{1.853791in}%
\pgfsys@useobject{currentmarker}{}%
\end{pgfscope}%
\begin{pgfscope}%
\pgfsys@transformshift{3.101776in}{1.853791in}%
\pgfsys@useobject{currentmarker}{}%
\end{pgfscope}%
\begin{pgfscope}%
\pgfsys@transformshift{3.122740in}{1.853791in}%
\pgfsys@useobject{currentmarker}{}%
\end{pgfscope}%
\begin{pgfscope}%
\pgfsys@transformshift{3.143703in}{1.853791in}%
\pgfsys@useobject{currentmarker}{}%
\end{pgfscope}%
\begin{pgfscope}%
\pgfsys@transformshift{3.164667in}{1.853791in}%
\pgfsys@useobject{currentmarker}{}%
\end{pgfscope}%
\begin{pgfscope}%
\pgfsys@transformshift{3.185631in}{1.853791in}%
\pgfsys@useobject{currentmarker}{}%
\end{pgfscope}%
\begin{pgfscope}%
\pgfsys@transformshift{3.206595in}{1.853791in}%
\pgfsys@useobject{currentmarker}{}%
\end{pgfscope}%
\begin{pgfscope}%
\pgfsys@transformshift{3.227559in}{1.853791in}%
\pgfsys@useobject{currentmarker}{}%
\end{pgfscope}%
\begin{pgfscope}%
\pgfsys@transformshift{2.850210in}{1.862132in}%
\pgfsys@useobject{currentmarker}{}%
\end{pgfscope}%
\begin{pgfscope}%
\pgfsys@transformshift{2.871174in}{1.862132in}%
\pgfsys@useobject{currentmarker}{}%
\end{pgfscope}%
\begin{pgfscope}%
\pgfsys@transformshift{2.892137in}{1.862132in}%
\pgfsys@useobject{currentmarker}{}%
\end{pgfscope}%
\begin{pgfscope}%
\pgfsys@transformshift{2.913101in}{1.862132in}%
\pgfsys@useobject{currentmarker}{}%
\end{pgfscope}%
\begin{pgfscope}%
\pgfsys@transformshift{2.934065in}{1.862132in}%
\pgfsys@useobject{currentmarker}{}%
\end{pgfscope}%
\begin{pgfscope}%
\pgfsys@transformshift{2.955029in}{1.862132in}%
\pgfsys@useobject{currentmarker}{}%
\end{pgfscope}%
\begin{pgfscope}%
\pgfsys@transformshift{2.975993in}{1.862132in}%
\pgfsys@useobject{currentmarker}{}%
\end{pgfscope}%
\begin{pgfscope}%
\pgfsys@transformshift{2.996957in}{1.862132in}%
\pgfsys@useobject{currentmarker}{}%
\end{pgfscope}%
\begin{pgfscope}%
\pgfsys@transformshift{3.017920in}{1.862132in}%
\pgfsys@useobject{currentmarker}{}%
\end{pgfscope}%
\begin{pgfscope}%
\pgfsys@transformshift{3.038884in}{1.862132in}%
\pgfsys@useobject{currentmarker}{}%
\end{pgfscope}%
\begin{pgfscope}%
\pgfsys@transformshift{3.059848in}{1.862132in}%
\pgfsys@useobject{currentmarker}{}%
\end{pgfscope}%
\begin{pgfscope}%
\pgfsys@transformshift{3.080812in}{1.862132in}%
\pgfsys@useobject{currentmarker}{}%
\end{pgfscope}%
\begin{pgfscope}%
\pgfsys@transformshift{3.101776in}{1.862132in}%
\pgfsys@useobject{currentmarker}{}%
\end{pgfscope}%
\begin{pgfscope}%
\pgfsys@transformshift{3.122740in}{1.862132in}%
\pgfsys@useobject{currentmarker}{}%
\end{pgfscope}%
\begin{pgfscope}%
\pgfsys@transformshift{3.143703in}{1.862132in}%
\pgfsys@useobject{currentmarker}{}%
\end{pgfscope}%
\begin{pgfscope}%
\pgfsys@transformshift{3.164667in}{1.862132in}%
\pgfsys@useobject{currentmarker}{}%
\end{pgfscope}%
\begin{pgfscope}%
\pgfsys@transformshift{3.185631in}{1.862132in}%
\pgfsys@useobject{currentmarker}{}%
\end{pgfscope}%
\begin{pgfscope}%
\pgfsys@transformshift{3.206595in}{1.862132in}%
\pgfsys@useobject{currentmarker}{}%
\end{pgfscope}%
\begin{pgfscope}%
\pgfsys@transformshift{3.227559in}{1.862132in}%
\pgfsys@useobject{currentmarker}{}%
\end{pgfscope}%
\begin{pgfscope}%
\pgfsys@transformshift{2.850210in}{1.870473in}%
\pgfsys@useobject{currentmarker}{}%
\end{pgfscope}%
\begin{pgfscope}%
\pgfsys@transformshift{2.871174in}{1.870473in}%
\pgfsys@useobject{currentmarker}{}%
\end{pgfscope}%
\begin{pgfscope}%
\pgfsys@transformshift{2.892137in}{1.870473in}%
\pgfsys@useobject{currentmarker}{}%
\end{pgfscope}%
\begin{pgfscope}%
\pgfsys@transformshift{2.913101in}{1.870473in}%
\pgfsys@useobject{currentmarker}{}%
\end{pgfscope}%
\begin{pgfscope}%
\pgfsys@transformshift{2.934065in}{1.870473in}%
\pgfsys@useobject{currentmarker}{}%
\end{pgfscope}%
\begin{pgfscope}%
\pgfsys@transformshift{2.955029in}{1.870473in}%
\pgfsys@useobject{currentmarker}{}%
\end{pgfscope}%
\begin{pgfscope}%
\pgfsys@transformshift{2.975993in}{1.870473in}%
\pgfsys@useobject{currentmarker}{}%
\end{pgfscope}%
\begin{pgfscope}%
\pgfsys@transformshift{2.996957in}{1.870473in}%
\pgfsys@useobject{currentmarker}{}%
\end{pgfscope}%
\begin{pgfscope}%
\pgfsys@transformshift{3.017920in}{1.870473in}%
\pgfsys@useobject{currentmarker}{}%
\end{pgfscope}%
\begin{pgfscope}%
\pgfsys@transformshift{3.038884in}{1.870473in}%
\pgfsys@useobject{currentmarker}{}%
\end{pgfscope}%
\begin{pgfscope}%
\pgfsys@transformshift{3.059848in}{1.870473in}%
\pgfsys@useobject{currentmarker}{}%
\end{pgfscope}%
\begin{pgfscope}%
\pgfsys@transformshift{3.080812in}{1.870473in}%
\pgfsys@useobject{currentmarker}{}%
\end{pgfscope}%
\begin{pgfscope}%
\pgfsys@transformshift{3.101776in}{1.870473in}%
\pgfsys@useobject{currentmarker}{}%
\end{pgfscope}%
\begin{pgfscope}%
\pgfsys@transformshift{3.122740in}{1.870473in}%
\pgfsys@useobject{currentmarker}{}%
\end{pgfscope}%
\begin{pgfscope}%
\pgfsys@transformshift{3.143703in}{1.870473in}%
\pgfsys@useobject{currentmarker}{}%
\end{pgfscope}%
\begin{pgfscope}%
\pgfsys@transformshift{3.164667in}{1.870473in}%
\pgfsys@useobject{currentmarker}{}%
\end{pgfscope}%
\begin{pgfscope}%
\pgfsys@transformshift{3.185631in}{1.870473in}%
\pgfsys@useobject{currentmarker}{}%
\end{pgfscope}%
\begin{pgfscope}%
\pgfsys@transformshift{3.206595in}{1.870473in}%
\pgfsys@useobject{currentmarker}{}%
\end{pgfscope}%
\begin{pgfscope}%
\pgfsys@transformshift{3.227559in}{1.870473in}%
\pgfsys@useobject{currentmarker}{}%
\end{pgfscope}%
\begin{pgfscope}%
\pgfsys@transformshift{2.850210in}{1.878815in}%
\pgfsys@useobject{currentmarker}{}%
\end{pgfscope}%
\begin{pgfscope}%
\pgfsys@transformshift{2.871174in}{1.878815in}%
\pgfsys@useobject{currentmarker}{}%
\end{pgfscope}%
\begin{pgfscope}%
\pgfsys@transformshift{2.892137in}{1.878815in}%
\pgfsys@useobject{currentmarker}{}%
\end{pgfscope}%
\begin{pgfscope}%
\pgfsys@transformshift{2.913101in}{1.878815in}%
\pgfsys@useobject{currentmarker}{}%
\end{pgfscope}%
\begin{pgfscope}%
\pgfsys@transformshift{2.934065in}{1.878815in}%
\pgfsys@useobject{currentmarker}{}%
\end{pgfscope}%
\begin{pgfscope}%
\pgfsys@transformshift{2.955029in}{1.878815in}%
\pgfsys@useobject{currentmarker}{}%
\end{pgfscope}%
\begin{pgfscope}%
\pgfsys@transformshift{2.975993in}{1.878815in}%
\pgfsys@useobject{currentmarker}{}%
\end{pgfscope}%
\begin{pgfscope}%
\pgfsys@transformshift{2.996957in}{1.878815in}%
\pgfsys@useobject{currentmarker}{}%
\end{pgfscope}%
\begin{pgfscope}%
\pgfsys@transformshift{3.017920in}{1.878815in}%
\pgfsys@useobject{currentmarker}{}%
\end{pgfscope}%
\begin{pgfscope}%
\pgfsys@transformshift{3.038884in}{1.878815in}%
\pgfsys@useobject{currentmarker}{}%
\end{pgfscope}%
\begin{pgfscope}%
\pgfsys@transformshift{3.059848in}{1.878815in}%
\pgfsys@useobject{currentmarker}{}%
\end{pgfscope}%
\begin{pgfscope}%
\pgfsys@transformshift{3.080812in}{1.878815in}%
\pgfsys@useobject{currentmarker}{}%
\end{pgfscope}%
\begin{pgfscope}%
\pgfsys@transformshift{3.101776in}{1.878815in}%
\pgfsys@useobject{currentmarker}{}%
\end{pgfscope}%
\begin{pgfscope}%
\pgfsys@transformshift{3.122740in}{1.878815in}%
\pgfsys@useobject{currentmarker}{}%
\end{pgfscope}%
\begin{pgfscope}%
\pgfsys@transformshift{3.143703in}{1.878815in}%
\pgfsys@useobject{currentmarker}{}%
\end{pgfscope}%
\begin{pgfscope}%
\pgfsys@transformshift{3.164667in}{1.878815in}%
\pgfsys@useobject{currentmarker}{}%
\end{pgfscope}%
\begin{pgfscope}%
\pgfsys@transformshift{3.185631in}{1.878815in}%
\pgfsys@useobject{currentmarker}{}%
\end{pgfscope}%
\begin{pgfscope}%
\pgfsys@transformshift{3.206595in}{1.878815in}%
\pgfsys@useobject{currentmarker}{}%
\end{pgfscope}%
\begin{pgfscope}%
\pgfsys@transformshift{3.227559in}{1.878815in}%
\pgfsys@useobject{currentmarker}{}%
\end{pgfscope}%
\begin{pgfscope}%
\pgfsys@transformshift{2.850210in}{1.887156in}%
\pgfsys@useobject{currentmarker}{}%
\end{pgfscope}%
\begin{pgfscope}%
\pgfsys@transformshift{2.871174in}{1.887156in}%
\pgfsys@useobject{currentmarker}{}%
\end{pgfscope}%
\begin{pgfscope}%
\pgfsys@transformshift{2.892137in}{1.887156in}%
\pgfsys@useobject{currentmarker}{}%
\end{pgfscope}%
\begin{pgfscope}%
\pgfsys@transformshift{2.913101in}{1.887156in}%
\pgfsys@useobject{currentmarker}{}%
\end{pgfscope}%
\begin{pgfscope}%
\pgfsys@transformshift{2.934065in}{1.887156in}%
\pgfsys@useobject{currentmarker}{}%
\end{pgfscope}%
\begin{pgfscope}%
\pgfsys@transformshift{2.955029in}{1.887156in}%
\pgfsys@useobject{currentmarker}{}%
\end{pgfscope}%
\begin{pgfscope}%
\pgfsys@transformshift{2.975993in}{1.887156in}%
\pgfsys@useobject{currentmarker}{}%
\end{pgfscope}%
\begin{pgfscope}%
\pgfsys@transformshift{2.996957in}{1.887156in}%
\pgfsys@useobject{currentmarker}{}%
\end{pgfscope}%
\begin{pgfscope}%
\pgfsys@transformshift{3.017920in}{1.887156in}%
\pgfsys@useobject{currentmarker}{}%
\end{pgfscope}%
\begin{pgfscope}%
\pgfsys@transformshift{3.038884in}{1.887156in}%
\pgfsys@useobject{currentmarker}{}%
\end{pgfscope}%
\begin{pgfscope}%
\pgfsys@transformshift{3.059848in}{1.887156in}%
\pgfsys@useobject{currentmarker}{}%
\end{pgfscope}%
\begin{pgfscope}%
\pgfsys@transformshift{3.080812in}{1.887156in}%
\pgfsys@useobject{currentmarker}{}%
\end{pgfscope}%
\begin{pgfscope}%
\pgfsys@transformshift{3.101776in}{1.887156in}%
\pgfsys@useobject{currentmarker}{}%
\end{pgfscope}%
\begin{pgfscope}%
\pgfsys@transformshift{3.122740in}{1.887156in}%
\pgfsys@useobject{currentmarker}{}%
\end{pgfscope}%
\begin{pgfscope}%
\pgfsys@transformshift{3.143703in}{1.887156in}%
\pgfsys@useobject{currentmarker}{}%
\end{pgfscope}%
\begin{pgfscope}%
\pgfsys@transformshift{3.164667in}{1.887156in}%
\pgfsys@useobject{currentmarker}{}%
\end{pgfscope}%
\begin{pgfscope}%
\pgfsys@transformshift{3.185631in}{1.887156in}%
\pgfsys@useobject{currentmarker}{}%
\end{pgfscope}%
\begin{pgfscope}%
\pgfsys@transformshift{3.206595in}{1.887156in}%
\pgfsys@useobject{currentmarker}{}%
\end{pgfscope}%
\begin{pgfscope}%
\pgfsys@transformshift{3.227559in}{1.887156in}%
\pgfsys@useobject{currentmarker}{}%
\end{pgfscope}%
\begin{pgfscope}%
\pgfsys@transformshift{2.850210in}{1.895497in}%
\pgfsys@useobject{currentmarker}{}%
\end{pgfscope}%
\begin{pgfscope}%
\pgfsys@transformshift{2.871174in}{1.895497in}%
\pgfsys@useobject{currentmarker}{}%
\end{pgfscope}%
\begin{pgfscope}%
\pgfsys@transformshift{2.892137in}{1.895497in}%
\pgfsys@useobject{currentmarker}{}%
\end{pgfscope}%
\begin{pgfscope}%
\pgfsys@transformshift{2.913101in}{1.895497in}%
\pgfsys@useobject{currentmarker}{}%
\end{pgfscope}%
\begin{pgfscope}%
\pgfsys@transformshift{2.934065in}{1.895497in}%
\pgfsys@useobject{currentmarker}{}%
\end{pgfscope}%
\begin{pgfscope}%
\pgfsys@transformshift{2.955029in}{1.895497in}%
\pgfsys@useobject{currentmarker}{}%
\end{pgfscope}%
\begin{pgfscope}%
\pgfsys@transformshift{2.975993in}{1.895497in}%
\pgfsys@useobject{currentmarker}{}%
\end{pgfscope}%
\begin{pgfscope}%
\pgfsys@transformshift{2.996957in}{1.895497in}%
\pgfsys@useobject{currentmarker}{}%
\end{pgfscope}%
\begin{pgfscope}%
\pgfsys@transformshift{3.017920in}{1.895497in}%
\pgfsys@useobject{currentmarker}{}%
\end{pgfscope}%
\begin{pgfscope}%
\pgfsys@transformshift{3.038884in}{1.895497in}%
\pgfsys@useobject{currentmarker}{}%
\end{pgfscope}%
\begin{pgfscope}%
\pgfsys@transformshift{3.059848in}{1.895497in}%
\pgfsys@useobject{currentmarker}{}%
\end{pgfscope}%
\begin{pgfscope}%
\pgfsys@transformshift{3.080812in}{1.895497in}%
\pgfsys@useobject{currentmarker}{}%
\end{pgfscope}%
\begin{pgfscope}%
\pgfsys@transformshift{3.101776in}{1.895497in}%
\pgfsys@useobject{currentmarker}{}%
\end{pgfscope}%
\begin{pgfscope}%
\pgfsys@transformshift{3.122740in}{1.895497in}%
\pgfsys@useobject{currentmarker}{}%
\end{pgfscope}%
\begin{pgfscope}%
\pgfsys@transformshift{3.143703in}{1.895497in}%
\pgfsys@useobject{currentmarker}{}%
\end{pgfscope}%
\begin{pgfscope}%
\pgfsys@transformshift{3.164667in}{1.895497in}%
\pgfsys@useobject{currentmarker}{}%
\end{pgfscope}%
\begin{pgfscope}%
\pgfsys@transformshift{3.185631in}{1.895497in}%
\pgfsys@useobject{currentmarker}{}%
\end{pgfscope}%
\begin{pgfscope}%
\pgfsys@transformshift{3.206595in}{1.895497in}%
\pgfsys@useobject{currentmarker}{}%
\end{pgfscope}%
\begin{pgfscope}%
\pgfsys@transformshift{3.227559in}{1.895497in}%
\pgfsys@useobject{currentmarker}{}%
\end{pgfscope}%
\begin{pgfscope}%
\pgfsys@transformshift{2.850210in}{1.903838in}%
\pgfsys@useobject{currentmarker}{}%
\end{pgfscope}%
\begin{pgfscope}%
\pgfsys@transformshift{2.871174in}{1.903838in}%
\pgfsys@useobject{currentmarker}{}%
\end{pgfscope}%
\begin{pgfscope}%
\pgfsys@transformshift{2.892137in}{1.903838in}%
\pgfsys@useobject{currentmarker}{}%
\end{pgfscope}%
\begin{pgfscope}%
\pgfsys@transformshift{2.913101in}{1.903838in}%
\pgfsys@useobject{currentmarker}{}%
\end{pgfscope}%
\begin{pgfscope}%
\pgfsys@transformshift{2.934065in}{1.903838in}%
\pgfsys@useobject{currentmarker}{}%
\end{pgfscope}%
\begin{pgfscope}%
\pgfsys@transformshift{2.955029in}{1.903838in}%
\pgfsys@useobject{currentmarker}{}%
\end{pgfscope}%
\begin{pgfscope}%
\pgfsys@transformshift{2.975993in}{1.903838in}%
\pgfsys@useobject{currentmarker}{}%
\end{pgfscope}%
\begin{pgfscope}%
\pgfsys@transformshift{2.996957in}{1.903838in}%
\pgfsys@useobject{currentmarker}{}%
\end{pgfscope}%
\begin{pgfscope}%
\pgfsys@transformshift{3.017920in}{1.903838in}%
\pgfsys@useobject{currentmarker}{}%
\end{pgfscope}%
\begin{pgfscope}%
\pgfsys@transformshift{3.038884in}{1.903838in}%
\pgfsys@useobject{currentmarker}{}%
\end{pgfscope}%
\begin{pgfscope}%
\pgfsys@transformshift{3.059848in}{1.903838in}%
\pgfsys@useobject{currentmarker}{}%
\end{pgfscope}%
\begin{pgfscope}%
\pgfsys@transformshift{3.080812in}{1.903838in}%
\pgfsys@useobject{currentmarker}{}%
\end{pgfscope}%
\begin{pgfscope}%
\pgfsys@transformshift{3.101776in}{1.903838in}%
\pgfsys@useobject{currentmarker}{}%
\end{pgfscope}%
\begin{pgfscope}%
\pgfsys@transformshift{3.122740in}{1.903838in}%
\pgfsys@useobject{currentmarker}{}%
\end{pgfscope}%
\begin{pgfscope}%
\pgfsys@transformshift{3.143703in}{1.903838in}%
\pgfsys@useobject{currentmarker}{}%
\end{pgfscope}%
\begin{pgfscope}%
\pgfsys@transformshift{3.164667in}{1.903838in}%
\pgfsys@useobject{currentmarker}{}%
\end{pgfscope}%
\begin{pgfscope}%
\pgfsys@transformshift{3.185631in}{1.903838in}%
\pgfsys@useobject{currentmarker}{}%
\end{pgfscope}%
\begin{pgfscope}%
\pgfsys@transformshift{3.206595in}{1.903838in}%
\pgfsys@useobject{currentmarker}{}%
\end{pgfscope}%
\begin{pgfscope}%
\pgfsys@transformshift{3.227559in}{1.903838in}%
\pgfsys@useobject{currentmarker}{}%
\end{pgfscope}%
\begin{pgfscope}%
\pgfsys@transformshift{2.850210in}{1.912180in}%
\pgfsys@useobject{currentmarker}{}%
\end{pgfscope}%
\begin{pgfscope}%
\pgfsys@transformshift{2.871174in}{1.912180in}%
\pgfsys@useobject{currentmarker}{}%
\end{pgfscope}%
\begin{pgfscope}%
\pgfsys@transformshift{2.892137in}{1.912180in}%
\pgfsys@useobject{currentmarker}{}%
\end{pgfscope}%
\begin{pgfscope}%
\pgfsys@transformshift{2.913101in}{1.912180in}%
\pgfsys@useobject{currentmarker}{}%
\end{pgfscope}%
\begin{pgfscope}%
\pgfsys@transformshift{2.934065in}{1.912180in}%
\pgfsys@useobject{currentmarker}{}%
\end{pgfscope}%
\begin{pgfscope}%
\pgfsys@transformshift{2.955029in}{1.912180in}%
\pgfsys@useobject{currentmarker}{}%
\end{pgfscope}%
\begin{pgfscope}%
\pgfsys@transformshift{2.975993in}{1.912180in}%
\pgfsys@useobject{currentmarker}{}%
\end{pgfscope}%
\begin{pgfscope}%
\pgfsys@transformshift{2.996957in}{1.912180in}%
\pgfsys@useobject{currentmarker}{}%
\end{pgfscope}%
\begin{pgfscope}%
\pgfsys@transformshift{3.017920in}{1.912180in}%
\pgfsys@useobject{currentmarker}{}%
\end{pgfscope}%
\begin{pgfscope}%
\pgfsys@transformshift{3.038884in}{1.912180in}%
\pgfsys@useobject{currentmarker}{}%
\end{pgfscope}%
\begin{pgfscope}%
\pgfsys@transformshift{3.059848in}{1.912180in}%
\pgfsys@useobject{currentmarker}{}%
\end{pgfscope}%
\begin{pgfscope}%
\pgfsys@transformshift{3.080812in}{1.912180in}%
\pgfsys@useobject{currentmarker}{}%
\end{pgfscope}%
\begin{pgfscope}%
\pgfsys@transformshift{3.101776in}{1.912180in}%
\pgfsys@useobject{currentmarker}{}%
\end{pgfscope}%
\begin{pgfscope}%
\pgfsys@transformshift{3.122740in}{1.912180in}%
\pgfsys@useobject{currentmarker}{}%
\end{pgfscope}%
\begin{pgfscope}%
\pgfsys@transformshift{3.143703in}{1.912180in}%
\pgfsys@useobject{currentmarker}{}%
\end{pgfscope}%
\begin{pgfscope}%
\pgfsys@transformshift{3.164667in}{1.912180in}%
\pgfsys@useobject{currentmarker}{}%
\end{pgfscope}%
\begin{pgfscope}%
\pgfsys@transformshift{3.185631in}{1.912180in}%
\pgfsys@useobject{currentmarker}{}%
\end{pgfscope}%
\begin{pgfscope}%
\pgfsys@transformshift{3.206595in}{1.912180in}%
\pgfsys@useobject{currentmarker}{}%
\end{pgfscope}%
\begin{pgfscope}%
\pgfsys@transformshift{3.227559in}{1.912180in}%
\pgfsys@useobject{currentmarker}{}%
\end{pgfscope}%
\begin{pgfscope}%
\pgfsys@transformshift{2.850210in}{1.920521in}%
\pgfsys@useobject{currentmarker}{}%
\end{pgfscope}%
\begin{pgfscope}%
\pgfsys@transformshift{2.871174in}{1.920521in}%
\pgfsys@useobject{currentmarker}{}%
\end{pgfscope}%
\begin{pgfscope}%
\pgfsys@transformshift{2.892137in}{1.920521in}%
\pgfsys@useobject{currentmarker}{}%
\end{pgfscope}%
\begin{pgfscope}%
\pgfsys@transformshift{2.913101in}{1.920521in}%
\pgfsys@useobject{currentmarker}{}%
\end{pgfscope}%
\begin{pgfscope}%
\pgfsys@transformshift{2.934065in}{1.920521in}%
\pgfsys@useobject{currentmarker}{}%
\end{pgfscope}%
\begin{pgfscope}%
\pgfsys@transformshift{2.955029in}{1.920521in}%
\pgfsys@useobject{currentmarker}{}%
\end{pgfscope}%
\begin{pgfscope}%
\pgfsys@transformshift{2.975993in}{1.920521in}%
\pgfsys@useobject{currentmarker}{}%
\end{pgfscope}%
\begin{pgfscope}%
\pgfsys@transformshift{2.996957in}{1.920521in}%
\pgfsys@useobject{currentmarker}{}%
\end{pgfscope}%
\begin{pgfscope}%
\pgfsys@transformshift{3.017920in}{1.920521in}%
\pgfsys@useobject{currentmarker}{}%
\end{pgfscope}%
\begin{pgfscope}%
\pgfsys@transformshift{3.038884in}{1.920521in}%
\pgfsys@useobject{currentmarker}{}%
\end{pgfscope}%
\begin{pgfscope}%
\pgfsys@transformshift{3.059848in}{1.920521in}%
\pgfsys@useobject{currentmarker}{}%
\end{pgfscope}%
\begin{pgfscope}%
\pgfsys@transformshift{3.080812in}{1.920521in}%
\pgfsys@useobject{currentmarker}{}%
\end{pgfscope}%
\begin{pgfscope}%
\pgfsys@transformshift{3.101776in}{1.920521in}%
\pgfsys@useobject{currentmarker}{}%
\end{pgfscope}%
\begin{pgfscope}%
\pgfsys@transformshift{3.122740in}{1.920521in}%
\pgfsys@useobject{currentmarker}{}%
\end{pgfscope}%
\begin{pgfscope}%
\pgfsys@transformshift{3.143703in}{1.920521in}%
\pgfsys@useobject{currentmarker}{}%
\end{pgfscope}%
\begin{pgfscope}%
\pgfsys@transformshift{3.164667in}{1.920521in}%
\pgfsys@useobject{currentmarker}{}%
\end{pgfscope}%
\begin{pgfscope}%
\pgfsys@transformshift{3.185631in}{1.920521in}%
\pgfsys@useobject{currentmarker}{}%
\end{pgfscope}%
\begin{pgfscope}%
\pgfsys@transformshift{3.206595in}{1.920521in}%
\pgfsys@useobject{currentmarker}{}%
\end{pgfscope}%
\begin{pgfscope}%
\pgfsys@transformshift{3.227559in}{1.920521in}%
\pgfsys@useobject{currentmarker}{}%
\end{pgfscope}%
\begin{pgfscope}%
\pgfsys@transformshift{2.850210in}{1.928862in}%
\pgfsys@useobject{currentmarker}{}%
\end{pgfscope}%
\begin{pgfscope}%
\pgfsys@transformshift{2.871174in}{1.928862in}%
\pgfsys@useobject{currentmarker}{}%
\end{pgfscope}%
\begin{pgfscope}%
\pgfsys@transformshift{2.892137in}{1.928862in}%
\pgfsys@useobject{currentmarker}{}%
\end{pgfscope}%
\begin{pgfscope}%
\pgfsys@transformshift{2.913101in}{1.928862in}%
\pgfsys@useobject{currentmarker}{}%
\end{pgfscope}%
\begin{pgfscope}%
\pgfsys@transformshift{2.934065in}{1.928862in}%
\pgfsys@useobject{currentmarker}{}%
\end{pgfscope}%
\begin{pgfscope}%
\pgfsys@transformshift{2.955029in}{1.928862in}%
\pgfsys@useobject{currentmarker}{}%
\end{pgfscope}%
\begin{pgfscope}%
\pgfsys@transformshift{2.975993in}{1.928862in}%
\pgfsys@useobject{currentmarker}{}%
\end{pgfscope}%
\begin{pgfscope}%
\pgfsys@transformshift{2.996957in}{1.928862in}%
\pgfsys@useobject{currentmarker}{}%
\end{pgfscope}%
\begin{pgfscope}%
\pgfsys@transformshift{3.017920in}{1.928862in}%
\pgfsys@useobject{currentmarker}{}%
\end{pgfscope}%
\begin{pgfscope}%
\pgfsys@transformshift{3.038884in}{1.928862in}%
\pgfsys@useobject{currentmarker}{}%
\end{pgfscope}%
\begin{pgfscope}%
\pgfsys@transformshift{3.059848in}{1.928862in}%
\pgfsys@useobject{currentmarker}{}%
\end{pgfscope}%
\begin{pgfscope}%
\pgfsys@transformshift{3.080812in}{1.928862in}%
\pgfsys@useobject{currentmarker}{}%
\end{pgfscope}%
\begin{pgfscope}%
\pgfsys@transformshift{3.101776in}{1.928862in}%
\pgfsys@useobject{currentmarker}{}%
\end{pgfscope}%
\begin{pgfscope}%
\pgfsys@transformshift{3.122740in}{1.928862in}%
\pgfsys@useobject{currentmarker}{}%
\end{pgfscope}%
\begin{pgfscope}%
\pgfsys@transformshift{3.143703in}{1.928862in}%
\pgfsys@useobject{currentmarker}{}%
\end{pgfscope}%
\begin{pgfscope}%
\pgfsys@transformshift{3.164667in}{1.928862in}%
\pgfsys@useobject{currentmarker}{}%
\end{pgfscope}%
\begin{pgfscope}%
\pgfsys@transformshift{3.185631in}{1.928862in}%
\pgfsys@useobject{currentmarker}{}%
\end{pgfscope}%
\begin{pgfscope}%
\pgfsys@transformshift{3.206595in}{1.928862in}%
\pgfsys@useobject{currentmarker}{}%
\end{pgfscope}%
\begin{pgfscope}%
\pgfsys@transformshift{3.227559in}{1.928862in}%
\pgfsys@useobject{currentmarker}{}%
\end{pgfscope}%
\begin{pgfscope}%
\pgfsys@transformshift{2.850210in}{1.937203in}%
\pgfsys@useobject{currentmarker}{}%
\end{pgfscope}%
\begin{pgfscope}%
\pgfsys@transformshift{2.871174in}{1.937203in}%
\pgfsys@useobject{currentmarker}{}%
\end{pgfscope}%
\begin{pgfscope}%
\pgfsys@transformshift{2.892137in}{1.937203in}%
\pgfsys@useobject{currentmarker}{}%
\end{pgfscope}%
\begin{pgfscope}%
\pgfsys@transformshift{2.913101in}{1.937203in}%
\pgfsys@useobject{currentmarker}{}%
\end{pgfscope}%
\begin{pgfscope}%
\pgfsys@transformshift{2.934065in}{1.937203in}%
\pgfsys@useobject{currentmarker}{}%
\end{pgfscope}%
\begin{pgfscope}%
\pgfsys@transformshift{2.955029in}{1.937203in}%
\pgfsys@useobject{currentmarker}{}%
\end{pgfscope}%
\begin{pgfscope}%
\pgfsys@transformshift{2.975993in}{1.937203in}%
\pgfsys@useobject{currentmarker}{}%
\end{pgfscope}%
\begin{pgfscope}%
\pgfsys@transformshift{2.996957in}{1.937203in}%
\pgfsys@useobject{currentmarker}{}%
\end{pgfscope}%
\begin{pgfscope}%
\pgfsys@transformshift{3.017920in}{1.937203in}%
\pgfsys@useobject{currentmarker}{}%
\end{pgfscope}%
\begin{pgfscope}%
\pgfsys@transformshift{3.038884in}{1.937203in}%
\pgfsys@useobject{currentmarker}{}%
\end{pgfscope}%
\begin{pgfscope}%
\pgfsys@transformshift{3.059848in}{1.937203in}%
\pgfsys@useobject{currentmarker}{}%
\end{pgfscope}%
\begin{pgfscope}%
\pgfsys@transformshift{3.080812in}{1.937203in}%
\pgfsys@useobject{currentmarker}{}%
\end{pgfscope}%
\begin{pgfscope}%
\pgfsys@transformshift{3.101776in}{1.937203in}%
\pgfsys@useobject{currentmarker}{}%
\end{pgfscope}%
\begin{pgfscope}%
\pgfsys@transformshift{3.122740in}{1.937203in}%
\pgfsys@useobject{currentmarker}{}%
\end{pgfscope}%
\begin{pgfscope}%
\pgfsys@transformshift{3.143703in}{1.937203in}%
\pgfsys@useobject{currentmarker}{}%
\end{pgfscope}%
\begin{pgfscope}%
\pgfsys@transformshift{3.164667in}{1.937203in}%
\pgfsys@useobject{currentmarker}{}%
\end{pgfscope}%
\begin{pgfscope}%
\pgfsys@transformshift{3.185631in}{1.937203in}%
\pgfsys@useobject{currentmarker}{}%
\end{pgfscope}%
\begin{pgfscope}%
\pgfsys@transformshift{3.206595in}{1.937203in}%
\pgfsys@useobject{currentmarker}{}%
\end{pgfscope}%
\begin{pgfscope}%
\pgfsys@transformshift{3.227559in}{1.937203in}%
\pgfsys@useobject{currentmarker}{}%
\end{pgfscope}%
\begin{pgfscope}%
\pgfsys@transformshift{2.850210in}{1.945545in}%
\pgfsys@useobject{currentmarker}{}%
\end{pgfscope}%
\begin{pgfscope}%
\pgfsys@transformshift{2.871174in}{1.945545in}%
\pgfsys@useobject{currentmarker}{}%
\end{pgfscope}%
\begin{pgfscope}%
\pgfsys@transformshift{2.892137in}{1.945545in}%
\pgfsys@useobject{currentmarker}{}%
\end{pgfscope}%
\begin{pgfscope}%
\pgfsys@transformshift{2.913101in}{1.945545in}%
\pgfsys@useobject{currentmarker}{}%
\end{pgfscope}%
\begin{pgfscope}%
\pgfsys@transformshift{2.934065in}{1.945545in}%
\pgfsys@useobject{currentmarker}{}%
\end{pgfscope}%
\begin{pgfscope}%
\pgfsys@transformshift{2.955029in}{1.945545in}%
\pgfsys@useobject{currentmarker}{}%
\end{pgfscope}%
\begin{pgfscope}%
\pgfsys@transformshift{2.975993in}{1.945545in}%
\pgfsys@useobject{currentmarker}{}%
\end{pgfscope}%
\begin{pgfscope}%
\pgfsys@transformshift{2.996957in}{1.945545in}%
\pgfsys@useobject{currentmarker}{}%
\end{pgfscope}%
\begin{pgfscope}%
\pgfsys@transformshift{3.017920in}{1.945545in}%
\pgfsys@useobject{currentmarker}{}%
\end{pgfscope}%
\begin{pgfscope}%
\pgfsys@transformshift{3.038884in}{1.945545in}%
\pgfsys@useobject{currentmarker}{}%
\end{pgfscope}%
\begin{pgfscope}%
\pgfsys@transformshift{3.059848in}{1.945545in}%
\pgfsys@useobject{currentmarker}{}%
\end{pgfscope}%
\begin{pgfscope}%
\pgfsys@transformshift{3.080812in}{1.945545in}%
\pgfsys@useobject{currentmarker}{}%
\end{pgfscope}%
\begin{pgfscope}%
\pgfsys@transformshift{3.101776in}{1.945545in}%
\pgfsys@useobject{currentmarker}{}%
\end{pgfscope}%
\begin{pgfscope}%
\pgfsys@transformshift{3.122740in}{1.945545in}%
\pgfsys@useobject{currentmarker}{}%
\end{pgfscope}%
\begin{pgfscope}%
\pgfsys@transformshift{3.143703in}{1.945545in}%
\pgfsys@useobject{currentmarker}{}%
\end{pgfscope}%
\begin{pgfscope}%
\pgfsys@transformshift{3.164667in}{1.945545in}%
\pgfsys@useobject{currentmarker}{}%
\end{pgfscope}%
\begin{pgfscope}%
\pgfsys@transformshift{3.185631in}{1.945545in}%
\pgfsys@useobject{currentmarker}{}%
\end{pgfscope}%
\begin{pgfscope}%
\pgfsys@transformshift{3.206595in}{1.945545in}%
\pgfsys@useobject{currentmarker}{}%
\end{pgfscope}%
\begin{pgfscope}%
\pgfsys@transformshift{3.227559in}{1.945545in}%
\pgfsys@useobject{currentmarker}{}%
\end{pgfscope}%
\begin{pgfscope}%
\pgfsys@transformshift{2.850210in}{1.953886in}%
\pgfsys@useobject{currentmarker}{}%
\end{pgfscope}%
\begin{pgfscope}%
\pgfsys@transformshift{2.871174in}{1.953886in}%
\pgfsys@useobject{currentmarker}{}%
\end{pgfscope}%
\begin{pgfscope}%
\pgfsys@transformshift{2.892137in}{1.953886in}%
\pgfsys@useobject{currentmarker}{}%
\end{pgfscope}%
\begin{pgfscope}%
\pgfsys@transformshift{2.913101in}{1.953886in}%
\pgfsys@useobject{currentmarker}{}%
\end{pgfscope}%
\begin{pgfscope}%
\pgfsys@transformshift{2.934065in}{1.953886in}%
\pgfsys@useobject{currentmarker}{}%
\end{pgfscope}%
\begin{pgfscope}%
\pgfsys@transformshift{2.955029in}{1.953886in}%
\pgfsys@useobject{currentmarker}{}%
\end{pgfscope}%
\begin{pgfscope}%
\pgfsys@transformshift{2.975993in}{1.953886in}%
\pgfsys@useobject{currentmarker}{}%
\end{pgfscope}%
\begin{pgfscope}%
\pgfsys@transformshift{2.850210in}{1.962227in}%
\pgfsys@useobject{currentmarker}{}%
\end{pgfscope}%
\begin{pgfscope}%
\pgfsys@transformshift{2.871174in}{1.962227in}%
\pgfsys@useobject{currentmarker}{}%
\end{pgfscope}%
\begin{pgfscope}%
\pgfsys@transformshift{2.892137in}{1.962227in}%
\pgfsys@useobject{currentmarker}{}%
\end{pgfscope}%
\begin{pgfscope}%
\pgfsys@transformshift{2.913101in}{1.962227in}%
\pgfsys@useobject{currentmarker}{}%
\end{pgfscope}%
\begin{pgfscope}%
\pgfsys@transformshift{2.934065in}{1.962227in}%
\pgfsys@useobject{currentmarker}{}%
\end{pgfscope}%
\begin{pgfscope}%
\pgfsys@transformshift{2.955029in}{1.962227in}%
\pgfsys@useobject{currentmarker}{}%
\end{pgfscope}%
\begin{pgfscope}%
\pgfsys@transformshift{2.975993in}{1.962227in}%
\pgfsys@useobject{currentmarker}{}%
\end{pgfscope}%
\begin{pgfscope}%
\pgfsys@transformshift{2.850210in}{1.970568in}%
\pgfsys@useobject{currentmarker}{}%
\end{pgfscope}%
\begin{pgfscope}%
\pgfsys@transformshift{2.871174in}{1.970568in}%
\pgfsys@useobject{currentmarker}{}%
\end{pgfscope}%
\begin{pgfscope}%
\pgfsys@transformshift{2.892137in}{1.970568in}%
\pgfsys@useobject{currentmarker}{}%
\end{pgfscope}%
\begin{pgfscope}%
\pgfsys@transformshift{2.913101in}{1.970568in}%
\pgfsys@useobject{currentmarker}{}%
\end{pgfscope}%
\begin{pgfscope}%
\pgfsys@transformshift{2.934065in}{1.970568in}%
\pgfsys@useobject{currentmarker}{}%
\end{pgfscope}%
\begin{pgfscope}%
\pgfsys@transformshift{2.955029in}{1.970568in}%
\pgfsys@useobject{currentmarker}{}%
\end{pgfscope}%
\begin{pgfscope}%
\pgfsys@transformshift{2.975993in}{1.970568in}%
\pgfsys@useobject{currentmarker}{}%
\end{pgfscope}%
\begin{pgfscope}%
\pgfsys@transformshift{2.850210in}{1.978910in}%
\pgfsys@useobject{currentmarker}{}%
\end{pgfscope}%
\begin{pgfscope}%
\pgfsys@transformshift{2.871174in}{1.978910in}%
\pgfsys@useobject{currentmarker}{}%
\end{pgfscope}%
\begin{pgfscope}%
\pgfsys@transformshift{2.892137in}{1.978910in}%
\pgfsys@useobject{currentmarker}{}%
\end{pgfscope}%
\begin{pgfscope}%
\pgfsys@transformshift{2.913101in}{1.978910in}%
\pgfsys@useobject{currentmarker}{}%
\end{pgfscope}%
\begin{pgfscope}%
\pgfsys@transformshift{2.934065in}{1.978910in}%
\pgfsys@useobject{currentmarker}{}%
\end{pgfscope}%
\begin{pgfscope}%
\pgfsys@transformshift{2.955029in}{1.978910in}%
\pgfsys@useobject{currentmarker}{}%
\end{pgfscope}%
\begin{pgfscope}%
\pgfsys@transformshift{2.975993in}{1.978910in}%
\pgfsys@useobject{currentmarker}{}%
\end{pgfscope}%
\begin{pgfscope}%
\pgfsys@transformshift{2.850210in}{1.987251in}%
\pgfsys@useobject{currentmarker}{}%
\end{pgfscope}%
\begin{pgfscope}%
\pgfsys@transformshift{2.871174in}{1.987251in}%
\pgfsys@useobject{currentmarker}{}%
\end{pgfscope}%
\begin{pgfscope}%
\pgfsys@transformshift{2.892137in}{1.987251in}%
\pgfsys@useobject{currentmarker}{}%
\end{pgfscope}%
\begin{pgfscope}%
\pgfsys@transformshift{2.913101in}{1.987251in}%
\pgfsys@useobject{currentmarker}{}%
\end{pgfscope}%
\begin{pgfscope}%
\pgfsys@transformshift{2.934065in}{1.987251in}%
\pgfsys@useobject{currentmarker}{}%
\end{pgfscope}%
\begin{pgfscope}%
\pgfsys@transformshift{2.955029in}{1.987251in}%
\pgfsys@useobject{currentmarker}{}%
\end{pgfscope}%
\begin{pgfscope}%
\pgfsys@transformshift{2.975993in}{1.987251in}%
\pgfsys@useobject{currentmarker}{}%
\end{pgfscope}%
\begin{pgfscope}%
\pgfsys@transformshift{2.850210in}{1.995592in}%
\pgfsys@useobject{currentmarker}{}%
\end{pgfscope}%
\begin{pgfscope}%
\pgfsys@transformshift{2.871174in}{1.995592in}%
\pgfsys@useobject{currentmarker}{}%
\end{pgfscope}%
\begin{pgfscope}%
\pgfsys@transformshift{2.892137in}{1.995592in}%
\pgfsys@useobject{currentmarker}{}%
\end{pgfscope}%
\begin{pgfscope}%
\pgfsys@transformshift{2.913101in}{1.995592in}%
\pgfsys@useobject{currentmarker}{}%
\end{pgfscope}%
\begin{pgfscope}%
\pgfsys@transformshift{2.934065in}{1.995592in}%
\pgfsys@useobject{currentmarker}{}%
\end{pgfscope}%
\begin{pgfscope}%
\pgfsys@transformshift{2.955029in}{1.995592in}%
\pgfsys@useobject{currentmarker}{}%
\end{pgfscope}%
\begin{pgfscope}%
\pgfsys@transformshift{2.975993in}{1.995592in}%
\pgfsys@useobject{currentmarker}{}%
\end{pgfscope}%
\begin{pgfscope}%
\pgfsys@transformshift{2.850210in}{2.003933in}%
\pgfsys@useobject{currentmarker}{}%
\end{pgfscope}%
\begin{pgfscope}%
\pgfsys@transformshift{2.871174in}{2.003933in}%
\pgfsys@useobject{currentmarker}{}%
\end{pgfscope}%
\begin{pgfscope}%
\pgfsys@transformshift{2.892137in}{2.003933in}%
\pgfsys@useobject{currentmarker}{}%
\end{pgfscope}%
\begin{pgfscope}%
\pgfsys@transformshift{2.913101in}{2.003933in}%
\pgfsys@useobject{currentmarker}{}%
\end{pgfscope}%
\begin{pgfscope}%
\pgfsys@transformshift{2.934065in}{2.003933in}%
\pgfsys@useobject{currentmarker}{}%
\end{pgfscope}%
\begin{pgfscope}%
\pgfsys@transformshift{2.955029in}{2.003933in}%
\pgfsys@useobject{currentmarker}{}%
\end{pgfscope}%
\begin{pgfscope}%
\pgfsys@transformshift{2.975993in}{2.003933in}%
\pgfsys@useobject{currentmarker}{}%
\end{pgfscope}%
\begin{pgfscope}%
\pgfsys@transformshift{2.850210in}{2.012275in}%
\pgfsys@useobject{currentmarker}{}%
\end{pgfscope}%
\begin{pgfscope}%
\pgfsys@transformshift{2.871174in}{2.012275in}%
\pgfsys@useobject{currentmarker}{}%
\end{pgfscope}%
\begin{pgfscope}%
\pgfsys@transformshift{2.892137in}{2.012275in}%
\pgfsys@useobject{currentmarker}{}%
\end{pgfscope}%
\begin{pgfscope}%
\pgfsys@transformshift{2.913101in}{2.012275in}%
\pgfsys@useobject{currentmarker}{}%
\end{pgfscope}%
\begin{pgfscope}%
\pgfsys@transformshift{2.934065in}{2.012275in}%
\pgfsys@useobject{currentmarker}{}%
\end{pgfscope}%
\begin{pgfscope}%
\pgfsys@transformshift{2.955029in}{2.012275in}%
\pgfsys@useobject{currentmarker}{}%
\end{pgfscope}%
\begin{pgfscope}%
\pgfsys@transformshift{2.975993in}{2.012275in}%
\pgfsys@useobject{currentmarker}{}%
\end{pgfscope}%
\begin{pgfscope}%
\pgfsys@transformshift{2.850210in}{2.020616in}%
\pgfsys@useobject{currentmarker}{}%
\end{pgfscope}%
\begin{pgfscope}%
\pgfsys@transformshift{2.871174in}{2.020616in}%
\pgfsys@useobject{currentmarker}{}%
\end{pgfscope}%
\begin{pgfscope}%
\pgfsys@transformshift{2.892137in}{2.020616in}%
\pgfsys@useobject{currentmarker}{}%
\end{pgfscope}%
\begin{pgfscope}%
\pgfsys@transformshift{2.913101in}{2.020616in}%
\pgfsys@useobject{currentmarker}{}%
\end{pgfscope}%
\begin{pgfscope}%
\pgfsys@transformshift{2.934065in}{2.020616in}%
\pgfsys@useobject{currentmarker}{}%
\end{pgfscope}%
\begin{pgfscope}%
\pgfsys@transformshift{2.955029in}{2.020616in}%
\pgfsys@useobject{currentmarker}{}%
\end{pgfscope}%
\begin{pgfscope}%
\pgfsys@transformshift{2.975993in}{2.020616in}%
\pgfsys@useobject{currentmarker}{}%
\end{pgfscope}%
\begin{pgfscope}%
\pgfsys@transformshift{2.850210in}{2.028957in}%
\pgfsys@useobject{currentmarker}{}%
\end{pgfscope}%
\begin{pgfscope}%
\pgfsys@transformshift{2.871174in}{2.028957in}%
\pgfsys@useobject{currentmarker}{}%
\end{pgfscope}%
\begin{pgfscope}%
\pgfsys@transformshift{2.892137in}{2.028957in}%
\pgfsys@useobject{currentmarker}{}%
\end{pgfscope}%
\begin{pgfscope}%
\pgfsys@transformshift{2.913101in}{2.028957in}%
\pgfsys@useobject{currentmarker}{}%
\end{pgfscope}%
\begin{pgfscope}%
\pgfsys@transformshift{2.934065in}{2.028957in}%
\pgfsys@useobject{currentmarker}{}%
\end{pgfscope}%
\begin{pgfscope}%
\pgfsys@transformshift{2.955029in}{2.028957in}%
\pgfsys@useobject{currentmarker}{}%
\end{pgfscope}%
\begin{pgfscope}%
\pgfsys@transformshift{2.975993in}{2.028957in}%
\pgfsys@useobject{currentmarker}{}%
\end{pgfscope}%
\begin{pgfscope}%
\pgfsys@transformshift{2.850210in}{2.037298in}%
\pgfsys@useobject{currentmarker}{}%
\end{pgfscope}%
\begin{pgfscope}%
\pgfsys@transformshift{2.871174in}{2.037298in}%
\pgfsys@useobject{currentmarker}{}%
\end{pgfscope}%
\begin{pgfscope}%
\pgfsys@transformshift{2.892137in}{2.037298in}%
\pgfsys@useobject{currentmarker}{}%
\end{pgfscope}%
\begin{pgfscope}%
\pgfsys@transformshift{2.913101in}{2.037298in}%
\pgfsys@useobject{currentmarker}{}%
\end{pgfscope}%
\begin{pgfscope}%
\pgfsys@transformshift{2.934065in}{2.037298in}%
\pgfsys@useobject{currentmarker}{}%
\end{pgfscope}%
\begin{pgfscope}%
\pgfsys@transformshift{2.955029in}{2.037298in}%
\pgfsys@useobject{currentmarker}{}%
\end{pgfscope}%
\begin{pgfscope}%
\pgfsys@transformshift{2.975993in}{2.037298in}%
\pgfsys@useobject{currentmarker}{}%
\end{pgfscope}%
\begin{pgfscope}%
\pgfsys@transformshift{2.850210in}{2.045640in}%
\pgfsys@useobject{currentmarker}{}%
\end{pgfscope}%
\begin{pgfscope}%
\pgfsys@transformshift{2.871174in}{2.045640in}%
\pgfsys@useobject{currentmarker}{}%
\end{pgfscope}%
\begin{pgfscope}%
\pgfsys@transformshift{2.892137in}{2.045640in}%
\pgfsys@useobject{currentmarker}{}%
\end{pgfscope}%
\begin{pgfscope}%
\pgfsys@transformshift{2.913101in}{2.045640in}%
\pgfsys@useobject{currentmarker}{}%
\end{pgfscope}%
\begin{pgfscope}%
\pgfsys@transformshift{2.934065in}{2.045640in}%
\pgfsys@useobject{currentmarker}{}%
\end{pgfscope}%
\begin{pgfscope}%
\pgfsys@transformshift{2.955029in}{2.045640in}%
\pgfsys@useobject{currentmarker}{}%
\end{pgfscope}%
\begin{pgfscope}%
\pgfsys@transformshift{2.975993in}{2.045640in}%
\pgfsys@useobject{currentmarker}{}%
\end{pgfscope}%
\begin{pgfscope}%
\pgfsys@transformshift{2.850210in}{2.053981in}%
\pgfsys@useobject{currentmarker}{}%
\end{pgfscope}%
\begin{pgfscope}%
\pgfsys@transformshift{2.871174in}{2.053981in}%
\pgfsys@useobject{currentmarker}{}%
\end{pgfscope}%
\begin{pgfscope}%
\pgfsys@transformshift{2.892137in}{2.053981in}%
\pgfsys@useobject{currentmarker}{}%
\end{pgfscope}%
\begin{pgfscope}%
\pgfsys@transformshift{2.913101in}{2.053981in}%
\pgfsys@useobject{currentmarker}{}%
\end{pgfscope}%
\begin{pgfscope}%
\pgfsys@transformshift{2.934065in}{2.053981in}%
\pgfsys@useobject{currentmarker}{}%
\end{pgfscope}%
\begin{pgfscope}%
\pgfsys@transformshift{2.955029in}{2.053981in}%
\pgfsys@useobject{currentmarker}{}%
\end{pgfscope}%
\begin{pgfscope}%
\pgfsys@transformshift{2.975993in}{2.053981in}%
\pgfsys@useobject{currentmarker}{}%
\end{pgfscope}%
\begin{pgfscope}%
\pgfsys@transformshift{2.850210in}{2.062322in}%
\pgfsys@useobject{currentmarker}{}%
\end{pgfscope}%
\begin{pgfscope}%
\pgfsys@transformshift{2.871174in}{2.062322in}%
\pgfsys@useobject{currentmarker}{}%
\end{pgfscope}%
\begin{pgfscope}%
\pgfsys@transformshift{2.892137in}{2.062322in}%
\pgfsys@useobject{currentmarker}{}%
\end{pgfscope}%
\begin{pgfscope}%
\pgfsys@transformshift{2.913101in}{2.062322in}%
\pgfsys@useobject{currentmarker}{}%
\end{pgfscope}%
\begin{pgfscope}%
\pgfsys@transformshift{2.934065in}{2.062322in}%
\pgfsys@useobject{currentmarker}{}%
\end{pgfscope}%
\begin{pgfscope}%
\pgfsys@transformshift{2.955029in}{2.062322in}%
\pgfsys@useobject{currentmarker}{}%
\end{pgfscope}%
\begin{pgfscope}%
\pgfsys@transformshift{2.975993in}{2.062322in}%
\pgfsys@useobject{currentmarker}{}%
\end{pgfscope}%
\begin{pgfscope}%
\pgfsys@transformshift{2.850210in}{2.070663in}%
\pgfsys@useobject{currentmarker}{}%
\end{pgfscope}%
\begin{pgfscope}%
\pgfsys@transformshift{2.871174in}{2.070663in}%
\pgfsys@useobject{currentmarker}{}%
\end{pgfscope}%
\begin{pgfscope}%
\pgfsys@transformshift{2.892137in}{2.070663in}%
\pgfsys@useobject{currentmarker}{}%
\end{pgfscope}%
\begin{pgfscope}%
\pgfsys@transformshift{2.913101in}{2.070663in}%
\pgfsys@useobject{currentmarker}{}%
\end{pgfscope}%
\begin{pgfscope}%
\pgfsys@transformshift{2.934065in}{2.070663in}%
\pgfsys@useobject{currentmarker}{}%
\end{pgfscope}%
\begin{pgfscope}%
\pgfsys@transformshift{2.955029in}{2.070663in}%
\pgfsys@useobject{currentmarker}{}%
\end{pgfscope}%
\begin{pgfscope}%
\pgfsys@transformshift{2.975993in}{2.070663in}%
\pgfsys@useobject{currentmarker}{}%
\end{pgfscope}%
\begin{pgfscope}%
\pgfsys@transformshift{2.850210in}{2.079004in}%
\pgfsys@useobject{currentmarker}{}%
\end{pgfscope}%
\begin{pgfscope}%
\pgfsys@transformshift{2.871174in}{2.079004in}%
\pgfsys@useobject{currentmarker}{}%
\end{pgfscope}%
\begin{pgfscope}%
\pgfsys@transformshift{2.892137in}{2.079004in}%
\pgfsys@useobject{currentmarker}{}%
\end{pgfscope}%
\begin{pgfscope}%
\pgfsys@transformshift{2.913101in}{2.079004in}%
\pgfsys@useobject{currentmarker}{}%
\end{pgfscope}%
\begin{pgfscope}%
\pgfsys@transformshift{2.934065in}{2.079004in}%
\pgfsys@useobject{currentmarker}{}%
\end{pgfscope}%
\begin{pgfscope}%
\pgfsys@transformshift{2.955029in}{2.079004in}%
\pgfsys@useobject{currentmarker}{}%
\end{pgfscope}%
\begin{pgfscope}%
\pgfsys@transformshift{2.975993in}{2.079004in}%
\pgfsys@useobject{currentmarker}{}%
\end{pgfscope}%
\begin{pgfscope}%
\pgfsys@transformshift{2.850210in}{2.087346in}%
\pgfsys@useobject{currentmarker}{}%
\end{pgfscope}%
\begin{pgfscope}%
\pgfsys@transformshift{2.871174in}{2.087346in}%
\pgfsys@useobject{currentmarker}{}%
\end{pgfscope}%
\begin{pgfscope}%
\pgfsys@transformshift{2.892137in}{2.087346in}%
\pgfsys@useobject{currentmarker}{}%
\end{pgfscope}%
\begin{pgfscope}%
\pgfsys@transformshift{2.913101in}{2.087346in}%
\pgfsys@useobject{currentmarker}{}%
\end{pgfscope}%
\begin{pgfscope}%
\pgfsys@transformshift{2.934065in}{2.087346in}%
\pgfsys@useobject{currentmarker}{}%
\end{pgfscope}%
\begin{pgfscope}%
\pgfsys@transformshift{2.955029in}{2.087346in}%
\pgfsys@useobject{currentmarker}{}%
\end{pgfscope}%
\begin{pgfscope}%
\pgfsys@transformshift{2.975993in}{2.087346in}%
\pgfsys@useobject{currentmarker}{}%
\end{pgfscope}%
\begin{pgfscope}%
\pgfsys@transformshift{2.850210in}{2.095687in}%
\pgfsys@useobject{currentmarker}{}%
\end{pgfscope}%
\begin{pgfscope}%
\pgfsys@transformshift{2.871174in}{2.095687in}%
\pgfsys@useobject{currentmarker}{}%
\end{pgfscope}%
\begin{pgfscope}%
\pgfsys@transformshift{2.892137in}{2.095687in}%
\pgfsys@useobject{currentmarker}{}%
\end{pgfscope}%
\begin{pgfscope}%
\pgfsys@transformshift{2.913101in}{2.095687in}%
\pgfsys@useobject{currentmarker}{}%
\end{pgfscope}%
\begin{pgfscope}%
\pgfsys@transformshift{2.934065in}{2.095687in}%
\pgfsys@useobject{currentmarker}{}%
\end{pgfscope}%
\begin{pgfscope}%
\pgfsys@transformshift{2.955029in}{2.095687in}%
\pgfsys@useobject{currentmarker}{}%
\end{pgfscope}%
\begin{pgfscope}%
\pgfsys@transformshift{2.975993in}{2.095687in}%
\pgfsys@useobject{currentmarker}{}%
\end{pgfscope}%
\begin{pgfscope}%
\pgfsys@transformshift{2.850210in}{2.104028in}%
\pgfsys@useobject{currentmarker}{}%
\end{pgfscope}%
\begin{pgfscope}%
\pgfsys@transformshift{2.871174in}{2.104028in}%
\pgfsys@useobject{currentmarker}{}%
\end{pgfscope}%
\begin{pgfscope}%
\pgfsys@transformshift{2.892137in}{2.104028in}%
\pgfsys@useobject{currentmarker}{}%
\end{pgfscope}%
\begin{pgfscope}%
\pgfsys@transformshift{2.913101in}{2.104028in}%
\pgfsys@useobject{currentmarker}{}%
\end{pgfscope}%
\begin{pgfscope}%
\pgfsys@transformshift{2.934065in}{2.104028in}%
\pgfsys@useobject{currentmarker}{}%
\end{pgfscope}%
\begin{pgfscope}%
\pgfsys@transformshift{2.955029in}{2.104028in}%
\pgfsys@useobject{currentmarker}{}%
\end{pgfscope}%
\begin{pgfscope}%
\pgfsys@transformshift{2.975993in}{2.104028in}%
\pgfsys@useobject{currentmarker}{}%
\end{pgfscope}%
\begin{pgfscope}%
\pgfsys@transformshift{2.850210in}{2.112369in}%
\pgfsys@useobject{currentmarker}{}%
\end{pgfscope}%
\begin{pgfscope}%
\pgfsys@transformshift{2.871174in}{2.112369in}%
\pgfsys@useobject{currentmarker}{}%
\end{pgfscope}%
\begin{pgfscope}%
\pgfsys@transformshift{2.892137in}{2.112369in}%
\pgfsys@useobject{currentmarker}{}%
\end{pgfscope}%
\begin{pgfscope}%
\pgfsys@transformshift{2.913101in}{2.112369in}%
\pgfsys@useobject{currentmarker}{}%
\end{pgfscope}%
\begin{pgfscope}%
\pgfsys@transformshift{2.934065in}{2.112369in}%
\pgfsys@useobject{currentmarker}{}%
\end{pgfscope}%
\begin{pgfscope}%
\pgfsys@transformshift{2.955029in}{2.112369in}%
\pgfsys@useobject{currentmarker}{}%
\end{pgfscope}%
\begin{pgfscope}%
\pgfsys@transformshift{2.975993in}{2.112369in}%
\pgfsys@useobject{currentmarker}{}%
\end{pgfscope}%
\begin{pgfscope}%
\pgfsys@transformshift{2.850210in}{2.120711in}%
\pgfsys@useobject{currentmarker}{}%
\end{pgfscope}%
\begin{pgfscope}%
\pgfsys@transformshift{2.871174in}{2.120711in}%
\pgfsys@useobject{currentmarker}{}%
\end{pgfscope}%
\begin{pgfscope}%
\pgfsys@transformshift{2.892137in}{2.120711in}%
\pgfsys@useobject{currentmarker}{}%
\end{pgfscope}%
\begin{pgfscope}%
\pgfsys@transformshift{2.913101in}{2.120711in}%
\pgfsys@useobject{currentmarker}{}%
\end{pgfscope}%
\begin{pgfscope}%
\pgfsys@transformshift{2.934065in}{2.120711in}%
\pgfsys@useobject{currentmarker}{}%
\end{pgfscope}%
\begin{pgfscope}%
\pgfsys@transformshift{2.955029in}{2.120711in}%
\pgfsys@useobject{currentmarker}{}%
\end{pgfscope}%
\begin{pgfscope}%
\pgfsys@transformshift{2.975993in}{2.120711in}%
\pgfsys@useobject{currentmarker}{}%
\end{pgfscope}%
\begin{pgfscope}%
\pgfsys@transformshift{2.850210in}{2.129052in}%
\pgfsys@useobject{currentmarker}{}%
\end{pgfscope}%
\begin{pgfscope}%
\pgfsys@transformshift{2.871174in}{2.129052in}%
\pgfsys@useobject{currentmarker}{}%
\end{pgfscope}%
\begin{pgfscope}%
\pgfsys@transformshift{2.892137in}{2.129052in}%
\pgfsys@useobject{currentmarker}{}%
\end{pgfscope}%
\begin{pgfscope}%
\pgfsys@transformshift{2.913101in}{2.129052in}%
\pgfsys@useobject{currentmarker}{}%
\end{pgfscope}%
\begin{pgfscope}%
\pgfsys@transformshift{2.934065in}{2.129052in}%
\pgfsys@useobject{currentmarker}{}%
\end{pgfscope}%
\begin{pgfscope}%
\pgfsys@transformshift{2.955029in}{2.129052in}%
\pgfsys@useobject{currentmarker}{}%
\end{pgfscope}%
\begin{pgfscope}%
\pgfsys@transformshift{2.975993in}{2.129052in}%
\pgfsys@useobject{currentmarker}{}%
\end{pgfscope}%
\begin{pgfscope}%
\pgfsys@transformshift{2.850210in}{2.137393in}%
\pgfsys@useobject{currentmarker}{}%
\end{pgfscope}%
\begin{pgfscope}%
\pgfsys@transformshift{2.871174in}{2.137393in}%
\pgfsys@useobject{currentmarker}{}%
\end{pgfscope}%
\begin{pgfscope}%
\pgfsys@transformshift{2.892137in}{2.137393in}%
\pgfsys@useobject{currentmarker}{}%
\end{pgfscope}%
\begin{pgfscope}%
\pgfsys@transformshift{2.913101in}{2.137393in}%
\pgfsys@useobject{currentmarker}{}%
\end{pgfscope}%
\begin{pgfscope}%
\pgfsys@transformshift{2.934065in}{2.137393in}%
\pgfsys@useobject{currentmarker}{}%
\end{pgfscope}%
\begin{pgfscope}%
\pgfsys@transformshift{2.955029in}{2.137393in}%
\pgfsys@useobject{currentmarker}{}%
\end{pgfscope}%
\begin{pgfscope}%
\pgfsys@transformshift{2.975993in}{2.137393in}%
\pgfsys@useobject{currentmarker}{}%
\end{pgfscope}%
\begin{pgfscope}%
\pgfsys@transformshift{2.850210in}{2.145734in}%
\pgfsys@useobject{currentmarker}{}%
\end{pgfscope}%
\begin{pgfscope}%
\pgfsys@transformshift{2.871174in}{2.145734in}%
\pgfsys@useobject{currentmarker}{}%
\end{pgfscope}%
\begin{pgfscope}%
\pgfsys@transformshift{2.892137in}{2.145734in}%
\pgfsys@useobject{currentmarker}{}%
\end{pgfscope}%
\begin{pgfscope}%
\pgfsys@transformshift{2.913101in}{2.145734in}%
\pgfsys@useobject{currentmarker}{}%
\end{pgfscope}%
\begin{pgfscope}%
\pgfsys@transformshift{2.934065in}{2.145734in}%
\pgfsys@useobject{currentmarker}{}%
\end{pgfscope}%
\begin{pgfscope}%
\pgfsys@transformshift{2.955029in}{2.145734in}%
\pgfsys@useobject{currentmarker}{}%
\end{pgfscope}%
\begin{pgfscope}%
\pgfsys@transformshift{2.975993in}{2.145734in}%
\pgfsys@useobject{currentmarker}{}%
\end{pgfscope}%
\begin{pgfscope}%
\pgfsys@transformshift{2.850210in}{2.154076in}%
\pgfsys@useobject{currentmarker}{}%
\end{pgfscope}%
\begin{pgfscope}%
\pgfsys@transformshift{2.871174in}{2.154076in}%
\pgfsys@useobject{currentmarker}{}%
\end{pgfscope}%
\begin{pgfscope}%
\pgfsys@transformshift{2.892137in}{2.154076in}%
\pgfsys@useobject{currentmarker}{}%
\end{pgfscope}%
\begin{pgfscope}%
\pgfsys@transformshift{2.913101in}{2.154076in}%
\pgfsys@useobject{currentmarker}{}%
\end{pgfscope}%
\begin{pgfscope}%
\pgfsys@transformshift{2.934065in}{2.154076in}%
\pgfsys@useobject{currentmarker}{}%
\end{pgfscope}%
\begin{pgfscope}%
\pgfsys@transformshift{2.955029in}{2.154076in}%
\pgfsys@useobject{currentmarker}{}%
\end{pgfscope}%
\begin{pgfscope}%
\pgfsys@transformshift{2.975993in}{2.154076in}%
\pgfsys@useobject{currentmarker}{}%
\end{pgfscope}%
\begin{pgfscope}%
\pgfsys@transformshift{2.850210in}{2.162417in}%
\pgfsys@useobject{currentmarker}{}%
\end{pgfscope}%
\begin{pgfscope}%
\pgfsys@transformshift{2.871174in}{2.162417in}%
\pgfsys@useobject{currentmarker}{}%
\end{pgfscope}%
\begin{pgfscope}%
\pgfsys@transformshift{2.892137in}{2.162417in}%
\pgfsys@useobject{currentmarker}{}%
\end{pgfscope}%
\begin{pgfscope}%
\pgfsys@transformshift{2.913101in}{2.162417in}%
\pgfsys@useobject{currentmarker}{}%
\end{pgfscope}%
\begin{pgfscope}%
\pgfsys@transformshift{2.934065in}{2.162417in}%
\pgfsys@useobject{currentmarker}{}%
\end{pgfscope}%
\begin{pgfscope}%
\pgfsys@transformshift{2.955029in}{2.162417in}%
\pgfsys@useobject{currentmarker}{}%
\end{pgfscope}%
\begin{pgfscope}%
\pgfsys@transformshift{2.975993in}{2.162417in}%
\pgfsys@useobject{currentmarker}{}%
\end{pgfscope}%
\begin{pgfscope}%
\pgfsys@transformshift{2.850210in}{2.170758in}%
\pgfsys@useobject{currentmarker}{}%
\end{pgfscope}%
\begin{pgfscope}%
\pgfsys@transformshift{2.871174in}{2.170758in}%
\pgfsys@useobject{currentmarker}{}%
\end{pgfscope}%
\begin{pgfscope}%
\pgfsys@transformshift{2.892137in}{2.170758in}%
\pgfsys@useobject{currentmarker}{}%
\end{pgfscope}%
\begin{pgfscope}%
\pgfsys@transformshift{2.913101in}{2.170758in}%
\pgfsys@useobject{currentmarker}{}%
\end{pgfscope}%
\begin{pgfscope}%
\pgfsys@transformshift{2.934065in}{2.170758in}%
\pgfsys@useobject{currentmarker}{}%
\end{pgfscope}%
\begin{pgfscope}%
\pgfsys@transformshift{2.955029in}{2.170758in}%
\pgfsys@useobject{currentmarker}{}%
\end{pgfscope}%
\begin{pgfscope}%
\pgfsys@transformshift{2.975993in}{2.170758in}%
\pgfsys@useobject{currentmarker}{}%
\end{pgfscope}%
\begin{pgfscope}%
\pgfsys@transformshift{2.850210in}{2.179099in}%
\pgfsys@useobject{currentmarker}{}%
\end{pgfscope}%
\begin{pgfscope}%
\pgfsys@transformshift{2.871174in}{2.179099in}%
\pgfsys@useobject{currentmarker}{}%
\end{pgfscope}%
\begin{pgfscope}%
\pgfsys@transformshift{2.892137in}{2.179099in}%
\pgfsys@useobject{currentmarker}{}%
\end{pgfscope}%
\begin{pgfscope}%
\pgfsys@transformshift{2.913101in}{2.179099in}%
\pgfsys@useobject{currentmarker}{}%
\end{pgfscope}%
\begin{pgfscope}%
\pgfsys@transformshift{2.934065in}{2.179099in}%
\pgfsys@useobject{currentmarker}{}%
\end{pgfscope}%
\begin{pgfscope}%
\pgfsys@transformshift{2.955029in}{2.179099in}%
\pgfsys@useobject{currentmarker}{}%
\end{pgfscope}%
\begin{pgfscope}%
\pgfsys@transformshift{2.975993in}{2.179099in}%
\pgfsys@useobject{currentmarker}{}%
\end{pgfscope}%
\begin{pgfscope}%
\pgfsys@transformshift{2.850210in}{2.187441in}%
\pgfsys@useobject{currentmarker}{}%
\end{pgfscope}%
\begin{pgfscope}%
\pgfsys@transformshift{2.871174in}{2.187441in}%
\pgfsys@useobject{currentmarker}{}%
\end{pgfscope}%
\begin{pgfscope}%
\pgfsys@transformshift{2.892137in}{2.187441in}%
\pgfsys@useobject{currentmarker}{}%
\end{pgfscope}%
\begin{pgfscope}%
\pgfsys@transformshift{2.913101in}{2.187441in}%
\pgfsys@useobject{currentmarker}{}%
\end{pgfscope}%
\begin{pgfscope}%
\pgfsys@transformshift{2.934065in}{2.187441in}%
\pgfsys@useobject{currentmarker}{}%
\end{pgfscope}%
\begin{pgfscope}%
\pgfsys@transformshift{2.955029in}{2.187441in}%
\pgfsys@useobject{currentmarker}{}%
\end{pgfscope}%
\begin{pgfscope}%
\pgfsys@transformshift{2.975993in}{2.187441in}%
\pgfsys@useobject{currentmarker}{}%
\end{pgfscope}%
\begin{pgfscope}%
\pgfsys@transformshift{2.850210in}{2.195782in}%
\pgfsys@useobject{currentmarker}{}%
\end{pgfscope}%
\begin{pgfscope}%
\pgfsys@transformshift{2.871174in}{2.195782in}%
\pgfsys@useobject{currentmarker}{}%
\end{pgfscope}%
\begin{pgfscope}%
\pgfsys@transformshift{2.892137in}{2.195782in}%
\pgfsys@useobject{currentmarker}{}%
\end{pgfscope}%
\begin{pgfscope}%
\pgfsys@transformshift{2.913101in}{2.195782in}%
\pgfsys@useobject{currentmarker}{}%
\end{pgfscope}%
\begin{pgfscope}%
\pgfsys@transformshift{2.934065in}{2.195782in}%
\pgfsys@useobject{currentmarker}{}%
\end{pgfscope}%
\begin{pgfscope}%
\pgfsys@transformshift{2.955029in}{2.195782in}%
\pgfsys@useobject{currentmarker}{}%
\end{pgfscope}%
\begin{pgfscope}%
\pgfsys@transformshift{2.975993in}{2.195782in}%
\pgfsys@useobject{currentmarker}{}%
\end{pgfscope}%
\begin{pgfscope}%
\pgfsys@transformshift{2.850210in}{2.204123in}%
\pgfsys@useobject{currentmarker}{}%
\end{pgfscope}%
\begin{pgfscope}%
\pgfsys@transformshift{2.871174in}{2.204123in}%
\pgfsys@useobject{currentmarker}{}%
\end{pgfscope}%
\begin{pgfscope}%
\pgfsys@transformshift{2.892137in}{2.204123in}%
\pgfsys@useobject{currentmarker}{}%
\end{pgfscope}%
\begin{pgfscope}%
\pgfsys@transformshift{2.913101in}{2.204123in}%
\pgfsys@useobject{currentmarker}{}%
\end{pgfscope}%
\begin{pgfscope}%
\pgfsys@transformshift{2.934065in}{2.204123in}%
\pgfsys@useobject{currentmarker}{}%
\end{pgfscope}%
\begin{pgfscope}%
\pgfsys@transformshift{2.955029in}{2.204123in}%
\pgfsys@useobject{currentmarker}{}%
\end{pgfscope}%
\begin{pgfscope}%
\pgfsys@transformshift{2.975993in}{2.204123in}%
\pgfsys@useobject{currentmarker}{}%
\end{pgfscope}%
\begin{pgfscope}%
\pgfsys@transformshift{2.850210in}{2.212464in}%
\pgfsys@useobject{currentmarker}{}%
\end{pgfscope}%
\begin{pgfscope}%
\pgfsys@transformshift{2.871174in}{2.212464in}%
\pgfsys@useobject{currentmarker}{}%
\end{pgfscope}%
\begin{pgfscope}%
\pgfsys@transformshift{2.892137in}{2.212464in}%
\pgfsys@useobject{currentmarker}{}%
\end{pgfscope}%
\begin{pgfscope}%
\pgfsys@transformshift{2.913101in}{2.212464in}%
\pgfsys@useobject{currentmarker}{}%
\end{pgfscope}%
\begin{pgfscope}%
\pgfsys@transformshift{2.934065in}{2.212464in}%
\pgfsys@useobject{currentmarker}{}%
\end{pgfscope}%
\begin{pgfscope}%
\pgfsys@transformshift{2.955029in}{2.212464in}%
\pgfsys@useobject{currentmarker}{}%
\end{pgfscope}%
\begin{pgfscope}%
\pgfsys@transformshift{2.975993in}{2.212464in}%
\pgfsys@useobject{currentmarker}{}%
\end{pgfscope}%
\begin{pgfscope}%
\pgfsys@transformshift{2.850210in}{2.220806in}%
\pgfsys@useobject{currentmarker}{}%
\end{pgfscope}%
\begin{pgfscope}%
\pgfsys@transformshift{2.871174in}{2.220806in}%
\pgfsys@useobject{currentmarker}{}%
\end{pgfscope}%
\begin{pgfscope}%
\pgfsys@transformshift{2.892137in}{2.220806in}%
\pgfsys@useobject{currentmarker}{}%
\end{pgfscope}%
\begin{pgfscope}%
\pgfsys@transformshift{2.913101in}{2.220806in}%
\pgfsys@useobject{currentmarker}{}%
\end{pgfscope}%
\begin{pgfscope}%
\pgfsys@transformshift{2.934065in}{2.220806in}%
\pgfsys@useobject{currentmarker}{}%
\end{pgfscope}%
\begin{pgfscope}%
\pgfsys@transformshift{2.955029in}{2.220806in}%
\pgfsys@useobject{currentmarker}{}%
\end{pgfscope}%
\begin{pgfscope}%
\pgfsys@transformshift{2.975993in}{2.220806in}%
\pgfsys@useobject{currentmarker}{}%
\end{pgfscope}%
\begin{pgfscope}%
\pgfsys@transformshift{2.850210in}{2.229147in}%
\pgfsys@useobject{currentmarker}{}%
\end{pgfscope}%
\begin{pgfscope}%
\pgfsys@transformshift{2.871174in}{2.229147in}%
\pgfsys@useobject{currentmarker}{}%
\end{pgfscope}%
\begin{pgfscope}%
\pgfsys@transformshift{2.892137in}{2.229147in}%
\pgfsys@useobject{currentmarker}{}%
\end{pgfscope}%
\begin{pgfscope}%
\pgfsys@transformshift{2.913101in}{2.229147in}%
\pgfsys@useobject{currentmarker}{}%
\end{pgfscope}%
\begin{pgfscope}%
\pgfsys@transformshift{2.934065in}{2.229147in}%
\pgfsys@useobject{currentmarker}{}%
\end{pgfscope}%
\begin{pgfscope}%
\pgfsys@transformshift{2.955029in}{2.229147in}%
\pgfsys@useobject{currentmarker}{}%
\end{pgfscope}%
\begin{pgfscope}%
\pgfsys@transformshift{2.975993in}{2.229147in}%
\pgfsys@useobject{currentmarker}{}%
\end{pgfscope}%
\begin{pgfscope}%
\pgfsys@transformshift{2.850210in}{2.237488in}%
\pgfsys@useobject{currentmarker}{}%
\end{pgfscope}%
\begin{pgfscope}%
\pgfsys@transformshift{2.871174in}{2.237488in}%
\pgfsys@useobject{currentmarker}{}%
\end{pgfscope}%
\begin{pgfscope}%
\pgfsys@transformshift{2.892137in}{2.237488in}%
\pgfsys@useobject{currentmarker}{}%
\end{pgfscope}%
\begin{pgfscope}%
\pgfsys@transformshift{2.913101in}{2.237488in}%
\pgfsys@useobject{currentmarker}{}%
\end{pgfscope}%
\begin{pgfscope}%
\pgfsys@transformshift{2.934065in}{2.237488in}%
\pgfsys@useobject{currentmarker}{}%
\end{pgfscope}%
\begin{pgfscope}%
\pgfsys@transformshift{2.955029in}{2.237488in}%
\pgfsys@useobject{currentmarker}{}%
\end{pgfscope}%
\begin{pgfscope}%
\pgfsys@transformshift{2.975993in}{2.237488in}%
\pgfsys@useobject{currentmarker}{}%
\end{pgfscope}%
\begin{pgfscope}%
\pgfsys@transformshift{2.850210in}{2.245829in}%
\pgfsys@useobject{currentmarker}{}%
\end{pgfscope}%
\begin{pgfscope}%
\pgfsys@transformshift{2.871174in}{2.245829in}%
\pgfsys@useobject{currentmarker}{}%
\end{pgfscope}%
\begin{pgfscope}%
\pgfsys@transformshift{2.892137in}{2.245829in}%
\pgfsys@useobject{currentmarker}{}%
\end{pgfscope}%
\begin{pgfscope}%
\pgfsys@transformshift{2.913101in}{2.245829in}%
\pgfsys@useobject{currentmarker}{}%
\end{pgfscope}%
\begin{pgfscope}%
\pgfsys@transformshift{2.934065in}{2.245829in}%
\pgfsys@useobject{currentmarker}{}%
\end{pgfscope}%
\begin{pgfscope}%
\pgfsys@transformshift{2.955029in}{2.245829in}%
\pgfsys@useobject{currentmarker}{}%
\end{pgfscope}%
\begin{pgfscope}%
\pgfsys@transformshift{2.975993in}{2.245829in}%
\pgfsys@useobject{currentmarker}{}%
\end{pgfscope}%
\begin{pgfscope}%
\pgfsys@transformshift{2.850210in}{2.254171in}%
\pgfsys@useobject{currentmarker}{}%
\end{pgfscope}%
\begin{pgfscope}%
\pgfsys@transformshift{2.871174in}{2.254171in}%
\pgfsys@useobject{currentmarker}{}%
\end{pgfscope}%
\begin{pgfscope}%
\pgfsys@transformshift{2.892137in}{2.254171in}%
\pgfsys@useobject{currentmarker}{}%
\end{pgfscope}%
\begin{pgfscope}%
\pgfsys@transformshift{2.913101in}{2.254171in}%
\pgfsys@useobject{currentmarker}{}%
\end{pgfscope}%
\begin{pgfscope}%
\pgfsys@transformshift{2.934065in}{2.254171in}%
\pgfsys@useobject{currentmarker}{}%
\end{pgfscope}%
\begin{pgfscope}%
\pgfsys@transformshift{2.955029in}{2.254171in}%
\pgfsys@useobject{currentmarker}{}%
\end{pgfscope}%
\begin{pgfscope}%
\pgfsys@transformshift{2.975993in}{2.254171in}%
\pgfsys@useobject{currentmarker}{}%
\end{pgfscope}%
\begin{pgfscope}%
\pgfsys@transformshift{2.850210in}{2.262512in}%
\pgfsys@useobject{currentmarker}{}%
\end{pgfscope}%
\begin{pgfscope}%
\pgfsys@transformshift{2.871174in}{2.262512in}%
\pgfsys@useobject{currentmarker}{}%
\end{pgfscope}%
\begin{pgfscope}%
\pgfsys@transformshift{2.892137in}{2.262512in}%
\pgfsys@useobject{currentmarker}{}%
\end{pgfscope}%
\begin{pgfscope}%
\pgfsys@transformshift{2.913101in}{2.262512in}%
\pgfsys@useobject{currentmarker}{}%
\end{pgfscope}%
\begin{pgfscope}%
\pgfsys@transformshift{2.934065in}{2.262512in}%
\pgfsys@useobject{currentmarker}{}%
\end{pgfscope}%
\begin{pgfscope}%
\pgfsys@transformshift{2.955029in}{2.262512in}%
\pgfsys@useobject{currentmarker}{}%
\end{pgfscope}%
\begin{pgfscope}%
\pgfsys@transformshift{2.975993in}{2.262512in}%
\pgfsys@useobject{currentmarker}{}%
\end{pgfscope}%
\begin{pgfscope}%
\pgfsys@transformshift{2.850210in}{2.270853in}%
\pgfsys@useobject{currentmarker}{}%
\end{pgfscope}%
\begin{pgfscope}%
\pgfsys@transformshift{2.871174in}{2.270853in}%
\pgfsys@useobject{currentmarker}{}%
\end{pgfscope}%
\begin{pgfscope}%
\pgfsys@transformshift{2.892137in}{2.270853in}%
\pgfsys@useobject{currentmarker}{}%
\end{pgfscope}%
\begin{pgfscope}%
\pgfsys@transformshift{2.913101in}{2.270853in}%
\pgfsys@useobject{currentmarker}{}%
\end{pgfscope}%
\begin{pgfscope}%
\pgfsys@transformshift{2.934065in}{2.270853in}%
\pgfsys@useobject{currentmarker}{}%
\end{pgfscope}%
\begin{pgfscope}%
\pgfsys@transformshift{2.955029in}{2.270853in}%
\pgfsys@useobject{currentmarker}{}%
\end{pgfscope}%
\begin{pgfscope}%
\pgfsys@transformshift{2.975993in}{2.270853in}%
\pgfsys@useobject{currentmarker}{}%
\end{pgfscope}%
\begin{pgfscope}%
\pgfsys@transformshift{2.850210in}{2.279194in}%
\pgfsys@useobject{currentmarker}{}%
\end{pgfscope}%
\begin{pgfscope}%
\pgfsys@transformshift{2.871174in}{2.279194in}%
\pgfsys@useobject{currentmarker}{}%
\end{pgfscope}%
\begin{pgfscope}%
\pgfsys@transformshift{2.892137in}{2.279194in}%
\pgfsys@useobject{currentmarker}{}%
\end{pgfscope}%
\begin{pgfscope}%
\pgfsys@transformshift{2.913101in}{2.279194in}%
\pgfsys@useobject{currentmarker}{}%
\end{pgfscope}%
\begin{pgfscope}%
\pgfsys@transformshift{2.934065in}{2.279194in}%
\pgfsys@useobject{currentmarker}{}%
\end{pgfscope}%
\begin{pgfscope}%
\pgfsys@transformshift{2.955029in}{2.279194in}%
\pgfsys@useobject{currentmarker}{}%
\end{pgfscope}%
\begin{pgfscope}%
\pgfsys@transformshift{2.975993in}{2.279194in}%
\pgfsys@useobject{currentmarker}{}%
\end{pgfscope}%
\begin{pgfscope}%
\pgfsys@transformshift{2.850210in}{2.287536in}%
\pgfsys@useobject{currentmarker}{}%
\end{pgfscope}%
\begin{pgfscope}%
\pgfsys@transformshift{2.871174in}{2.287536in}%
\pgfsys@useobject{currentmarker}{}%
\end{pgfscope}%
\begin{pgfscope}%
\pgfsys@transformshift{2.892137in}{2.287536in}%
\pgfsys@useobject{currentmarker}{}%
\end{pgfscope}%
\begin{pgfscope}%
\pgfsys@transformshift{2.913101in}{2.287536in}%
\pgfsys@useobject{currentmarker}{}%
\end{pgfscope}%
\begin{pgfscope}%
\pgfsys@transformshift{2.934065in}{2.287536in}%
\pgfsys@useobject{currentmarker}{}%
\end{pgfscope}%
\begin{pgfscope}%
\pgfsys@transformshift{2.955029in}{2.287536in}%
\pgfsys@useobject{currentmarker}{}%
\end{pgfscope}%
\begin{pgfscope}%
\pgfsys@transformshift{2.975993in}{2.287536in}%
\pgfsys@useobject{currentmarker}{}%
\end{pgfscope}%
\begin{pgfscope}%
\pgfsys@transformshift{2.850210in}{2.295877in}%
\pgfsys@useobject{currentmarker}{}%
\end{pgfscope}%
\begin{pgfscope}%
\pgfsys@transformshift{2.871174in}{2.295877in}%
\pgfsys@useobject{currentmarker}{}%
\end{pgfscope}%
\begin{pgfscope}%
\pgfsys@transformshift{2.892137in}{2.295877in}%
\pgfsys@useobject{currentmarker}{}%
\end{pgfscope}%
\begin{pgfscope}%
\pgfsys@transformshift{2.913101in}{2.295877in}%
\pgfsys@useobject{currentmarker}{}%
\end{pgfscope}%
\begin{pgfscope}%
\pgfsys@transformshift{2.934065in}{2.295877in}%
\pgfsys@useobject{currentmarker}{}%
\end{pgfscope}%
\begin{pgfscope}%
\pgfsys@transformshift{2.955029in}{2.295877in}%
\pgfsys@useobject{currentmarker}{}%
\end{pgfscope}%
\begin{pgfscope}%
\pgfsys@transformshift{2.975993in}{2.295877in}%
\pgfsys@useobject{currentmarker}{}%
\end{pgfscope}%
\begin{pgfscope}%
\pgfsys@transformshift{2.850210in}{2.304218in}%
\pgfsys@useobject{currentmarker}{}%
\end{pgfscope}%
\begin{pgfscope}%
\pgfsys@transformshift{2.871174in}{2.304218in}%
\pgfsys@useobject{currentmarker}{}%
\end{pgfscope}%
\begin{pgfscope}%
\pgfsys@transformshift{2.892137in}{2.304218in}%
\pgfsys@useobject{currentmarker}{}%
\end{pgfscope}%
\begin{pgfscope}%
\pgfsys@transformshift{2.913101in}{2.304218in}%
\pgfsys@useobject{currentmarker}{}%
\end{pgfscope}%
\begin{pgfscope}%
\pgfsys@transformshift{2.934065in}{2.304218in}%
\pgfsys@useobject{currentmarker}{}%
\end{pgfscope}%
\begin{pgfscope}%
\pgfsys@transformshift{2.955029in}{2.304218in}%
\pgfsys@useobject{currentmarker}{}%
\end{pgfscope}%
\begin{pgfscope}%
\pgfsys@transformshift{2.975993in}{2.304218in}%
\pgfsys@useobject{currentmarker}{}%
\end{pgfscope}%
\begin{pgfscope}%
\pgfsys@transformshift{2.850210in}{2.312559in}%
\pgfsys@useobject{currentmarker}{}%
\end{pgfscope}%
\begin{pgfscope}%
\pgfsys@transformshift{2.871174in}{2.312559in}%
\pgfsys@useobject{currentmarker}{}%
\end{pgfscope}%
\begin{pgfscope}%
\pgfsys@transformshift{2.892137in}{2.312559in}%
\pgfsys@useobject{currentmarker}{}%
\end{pgfscope}%
\begin{pgfscope}%
\pgfsys@transformshift{2.913101in}{2.312559in}%
\pgfsys@useobject{currentmarker}{}%
\end{pgfscope}%
\begin{pgfscope}%
\pgfsys@transformshift{2.934065in}{2.312559in}%
\pgfsys@useobject{currentmarker}{}%
\end{pgfscope}%
\begin{pgfscope}%
\pgfsys@transformshift{2.955029in}{2.312559in}%
\pgfsys@useobject{currentmarker}{}%
\end{pgfscope}%
\begin{pgfscope}%
\pgfsys@transformshift{2.975993in}{2.312559in}%
\pgfsys@useobject{currentmarker}{}%
\end{pgfscope}%
\begin{pgfscope}%
\pgfsys@transformshift{2.850210in}{2.320901in}%
\pgfsys@useobject{currentmarker}{}%
\end{pgfscope}%
\begin{pgfscope}%
\pgfsys@transformshift{2.871174in}{2.320901in}%
\pgfsys@useobject{currentmarker}{}%
\end{pgfscope}%
\begin{pgfscope}%
\pgfsys@transformshift{2.892137in}{2.320901in}%
\pgfsys@useobject{currentmarker}{}%
\end{pgfscope}%
\begin{pgfscope}%
\pgfsys@transformshift{2.913101in}{2.320901in}%
\pgfsys@useobject{currentmarker}{}%
\end{pgfscope}%
\begin{pgfscope}%
\pgfsys@transformshift{2.934065in}{2.320901in}%
\pgfsys@useobject{currentmarker}{}%
\end{pgfscope}%
\begin{pgfscope}%
\pgfsys@transformshift{2.955029in}{2.320901in}%
\pgfsys@useobject{currentmarker}{}%
\end{pgfscope}%
\begin{pgfscope}%
\pgfsys@transformshift{2.975993in}{2.320901in}%
\pgfsys@useobject{currentmarker}{}%
\end{pgfscope}%
\begin{pgfscope}%
\pgfsys@transformshift{2.850210in}{2.329242in}%
\pgfsys@useobject{currentmarker}{}%
\end{pgfscope}%
\begin{pgfscope}%
\pgfsys@transformshift{2.871174in}{2.329242in}%
\pgfsys@useobject{currentmarker}{}%
\end{pgfscope}%
\begin{pgfscope}%
\pgfsys@transformshift{2.892137in}{2.329242in}%
\pgfsys@useobject{currentmarker}{}%
\end{pgfscope}%
\begin{pgfscope}%
\pgfsys@transformshift{2.913101in}{2.329242in}%
\pgfsys@useobject{currentmarker}{}%
\end{pgfscope}%
\begin{pgfscope}%
\pgfsys@transformshift{2.934065in}{2.329242in}%
\pgfsys@useobject{currentmarker}{}%
\end{pgfscope}%
\begin{pgfscope}%
\pgfsys@transformshift{2.955029in}{2.329242in}%
\pgfsys@useobject{currentmarker}{}%
\end{pgfscope}%
\begin{pgfscope}%
\pgfsys@transformshift{2.975993in}{2.329242in}%
\pgfsys@useobject{currentmarker}{}%
\end{pgfscope}%
\begin{pgfscope}%
\pgfsys@transformshift{2.850210in}{2.337583in}%
\pgfsys@useobject{currentmarker}{}%
\end{pgfscope}%
\begin{pgfscope}%
\pgfsys@transformshift{2.871174in}{2.337583in}%
\pgfsys@useobject{currentmarker}{}%
\end{pgfscope}%
\begin{pgfscope}%
\pgfsys@transformshift{2.892137in}{2.337583in}%
\pgfsys@useobject{currentmarker}{}%
\end{pgfscope}%
\begin{pgfscope}%
\pgfsys@transformshift{2.913101in}{2.337583in}%
\pgfsys@useobject{currentmarker}{}%
\end{pgfscope}%
\begin{pgfscope}%
\pgfsys@transformshift{2.934065in}{2.337583in}%
\pgfsys@useobject{currentmarker}{}%
\end{pgfscope}%
\begin{pgfscope}%
\pgfsys@transformshift{2.955029in}{2.337583in}%
\pgfsys@useobject{currentmarker}{}%
\end{pgfscope}%
\begin{pgfscope}%
\pgfsys@transformshift{2.975993in}{2.337583in}%
\pgfsys@useobject{currentmarker}{}%
\end{pgfscope}%
\begin{pgfscope}%
\pgfsys@transformshift{2.850210in}{2.345924in}%
\pgfsys@useobject{currentmarker}{}%
\end{pgfscope}%
\begin{pgfscope}%
\pgfsys@transformshift{2.871174in}{2.345924in}%
\pgfsys@useobject{currentmarker}{}%
\end{pgfscope}%
\begin{pgfscope}%
\pgfsys@transformshift{2.892137in}{2.345924in}%
\pgfsys@useobject{currentmarker}{}%
\end{pgfscope}%
\begin{pgfscope}%
\pgfsys@transformshift{2.913101in}{2.345924in}%
\pgfsys@useobject{currentmarker}{}%
\end{pgfscope}%
\begin{pgfscope}%
\pgfsys@transformshift{2.934065in}{2.345924in}%
\pgfsys@useobject{currentmarker}{}%
\end{pgfscope}%
\begin{pgfscope}%
\pgfsys@transformshift{2.955029in}{2.345924in}%
\pgfsys@useobject{currentmarker}{}%
\end{pgfscope}%
\begin{pgfscope}%
\pgfsys@transformshift{2.975993in}{2.345924in}%
\pgfsys@useobject{currentmarker}{}%
\end{pgfscope}%
\begin{pgfscope}%
\pgfsys@transformshift{2.850210in}{2.354266in}%
\pgfsys@useobject{currentmarker}{}%
\end{pgfscope}%
\begin{pgfscope}%
\pgfsys@transformshift{2.871174in}{2.354266in}%
\pgfsys@useobject{currentmarker}{}%
\end{pgfscope}%
\begin{pgfscope}%
\pgfsys@transformshift{2.892137in}{2.354266in}%
\pgfsys@useobject{currentmarker}{}%
\end{pgfscope}%
\begin{pgfscope}%
\pgfsys@transformshift{2.913101in}{2.354266in}%
\pgfsys@useobject{currentmarker}{}%
\end{pgfscope}%
\begin{pgfscope}%
\pgfsys@transformshift{2.934065in}{2.354266in}%
\pgfsys@useobject{currentmarker}{}%
\end{pgfscope}%
\begin{pgfscope}%
\pgfsys@transformshift{2.955029in}{2.354266in}%
\pgfsys@useobject{currentmarker}{}%
\end{pgfscope}%
\begin{pgfscope}%
\pgfsys@transformshift{2.975993in}{2.354266in}%
\pgfsys@useobject{currentmarker}{}%
\end{pgfscope}%
\begin{pgfscope}%
\pgfsys@transformshift{2.850210in}{2.362607in}%
\pgfsys@useobject{currentmarker}{}%
\end{pgfscope}%
\begin{pgfscope}%
\pgfsys@transformshift{2.871174in}{2.362607in}%
\pgfsys@useobject{currentmarker}{}%
\end{pgfscope}%
\begin{pgfscope}%
\pgfsys@transformshift{2.892137in}{2.362607in}%
\pgfsys@useobject{currentmarker}{}%
\end{pgfscope}%
\begin{pgfscope}%
\pgfsys@transformshift{2.913101in}{2.362607in}%
\pgfsys@useobject{currentmarker}{}%
\end{pgfscope}%
\begin{pgfscope}%
\pgfsys@transformshift{2.934065in}{2.362607in}%
\pgfsys@useobject{currentmarker}{}%
\end{pgfscope}%
\begin{pgfscope}%
\pgfsys@transformshift{2.955029in}{2.362607in}%
\pgfsys@useobject{currentmarker}{}%
\end{pgfscope}%
\begin{pgfscope}%
\pgfsys@transformshift{2.975993in}{2.362607in}%
\pgfsys@useobject{currentmarker}{}%
\end{pgfscope}%
\begin{pgfscope}%
\pgfsys@transformshift{2.850210in}{2.370948in}%
\pgfsys@useobject{currentmarker}{}%
\end{pgfscope}%
\begin{pgfscope}%
\pgfsys@transformshift{2.871174in}{2.370948in}%
\pgfsys@useobject{currentmarker}{}%
\end{pgfscope}%
\begin{pgfscope}%
\pgfsys@transformshift{2.892137in}{2.370948in}%
\pgfsys@useobject{currentmarker}{}%
\end{pgfscope}%
\begin{pgfscope}%
\pgfsys@transformshift{2.913101in}{2.370948in}%
\pgfsys@useobject{currentmarker}{}%
\end{pgfscope}%
\begin{pgfscope}%
\pgfsys@transformshift{2.934065in}{2.370948in}%
\pgfsys@useobject{currentmarker}{}%
\end{pgfscope}%
\begin{pgfscope}%
\pgfsys@transformshift{2.955029in}{2.370948in}%
\pgfsys@useobject{currentmarker}{}%
\end{pgfscope}%
\begin{pgfscope}%
\pgfsys@transformshift{2.975993in}{2.370948in}%
\pgfsys@useobject{currentmarker}{}%
\end{pgfscope}%
\begin{pgfscope}%
\pgfsys@transformshift{2.850210in}{2.379289in}%
\pgfsys@useobject{currentmarker}{}%
\end{pgfscope}%
\begin{pgfscope}%
\pgfsys@transformshift{2.871174in}{2.379289in}%
\pgfsys@useobject{currentmarker}{}%
\end{pgfscope}%
\begin{pgfscope}%
\pgfsys@transformshift{2.892137in}{2.379289in}%
\pgfsys@useobject{currentmarker}{}%
\end{pgfscope}%
\begin{pgfscope}%
\pgfsys@transformshift{2.913101in}{2.379289in}%
\pgfsys@useobject{currentmarker}{}%
\end{pgfscope}%
\begin{pgfscope}%
\pgfsys@transformshift{2.934065in}{2.379289in}%
\pgfsys@useobject{currentmarker}{}%
\end{pgfscope}%
\begin{pgfscope}%
\pgfsys@transformshift{2.955029in}{2.379289in}%
\pgfsys@useobject{currentmarker}{}%
\end{pgfscope}%
\begin{pgfscope}%
\pgfsys@transformshift{2.975993in}{2.379289in}%
\pgfsys@useobject{currentmarker}{}%
\end{pgfscope}%
\begin{pgfscope}%
\pgfsys@transformshift{2.850210in}{2.387631in}%
\pgfsys@useobject{currentmarker}{}%
\end{pgfscope}%
\begin{pgfscope}%
\pgfsys@transformshift{2.871174in}{2.387631in}%
\pgfsys@useobject{currentmarker}{}%
\end{pgfscope}%
\begin{pgfscope}%
\pgfsys@transformshift{2.892137in}{2.387631in}%
\pgfsys@useobject{currentmarker}{}%
\end{pgfscope}%
\begin{pgfscope}%
\pgfsys@transformshift{2.913101in}{2.387631in}%
\pgfsys@useobject{currentmarker}{}%
\end{pgfscope}%
\begin{pgfscope}%
\pgfsys@transformshift{2.934065in}{2.387631in}%
\pgfsys@useobject{currentmarker}{}%
\end{pgfscope}%
\begin{pgfscope}%
\pgfsys@transformshift{2.955029in}{2.387631in}%
\pgfsys@useobject{currentmarker}{}%
\end{pgfscope}%
\begin{pgfscope}%
\pgfsys@transformshift{2.975993in}{2.387631in}%
\pgfsys@useobject{currentmarker}{}%
\end{pgfscope}%
\begin{pgfscope}%
\pgfsys@transformshift{2.850210in}{2.395972in}%
\pgfsys@useobject{currentmarker}{}%
\end{pgfscope}%
\begin{pgfscope}%
\pgfsys@transformshift{2.871174in}{2.395972in}%
\pgfsys@useobject{currentmarker}{}%
\end{pgfscope}%
\begin{pgfscope}%
\pgfsys@transformshift{2.892137in}{2.395972in}%
\pgfsys@useobject{currentmarker}{}%
\end{pgfscope}%
\begin{pgfscope}%
\pgfsys@transformshift{2.913101in}{2.395972in}%
\pgfsys@useobject{currentmarker}{}%
\end{pgfscope}%
\begin{pgfscope}%
\pgfsys@transformshift{2.934065in}{2.395972in}%
\pgfsys@useobject{currentmarker}{}%
\end{pgfscope}%
\begin{pgfscope}%
\pgfsys@transformshift{2.955029in}{2.395972in}%
\pgfsys@useobject{currentmarker}{}%
\end{pgfscope}%
\begin{pgfscope}%
\pgfsys@transformshift{2.975993in}{2.395972in}%
\pgfsys@useobject{currentmarker}{}%
\end{pgfscope}%
\begin{pgfscope}%
\pgfsys@transformshift{2.850210in}{2.404313in}%
\pgfsys@useobject{currentmarker}{}%
\end{pgfscope}%
\begin{pgfscope}%
\pgfsys@transformshift{2.871174in}{2.404313in}%
\pgfsys@useobject{currentmarker}{}%
\end{pgfscope}%
\begin{pgfscope}%
\pgfsys@transformshift{2.892137in}{2.404313in}%
\pgfsys@useobject{currentmarker}{}%
\end{pgfscope}%
\begin{pgfscope}%
\pgfsys@transformshift{2.913101in}{2.404313in}%
\pgfsys@useobject{currentmarker}{}%
\end{pgfscope}%
\begin{pgfscope}%
\pgfsys@transformshift{2.934065in}{2.404313in}%
\pgfsys@useobject{currentmarker}{}%
\end{pgfscope}%
\begin{pgfscope}%
\pgfsys@transformshift{2.955029in}{2.404313in}%
\pgfsys@useobject{currentmarker}{}%
\end{pgfscope}%
\begin{pgfscope}%
\pgfsys@transformshift{2.975993in}{2.404313in}%
\pgfsys@useobject{currentmarker}{}%
\end{pgfscope}%
\begin{pgfscope}%
\pgfsys@transformshift{2.850210in}{2.412654in}%
\pgfsys@useobject{currentmarker}{}%
\end{pgfscope}%
\begin{pgfscope}%
\pgfsys@transformshift{2.871174in}{2.412654in}%
\pgfsys@useobject{currentmarker}{}%
\end{pgfscope}%
\begin{pgfscope}%
\pgfsys@transformshift{2.892137in}{2.412654in}%
\pgfsys@useobject{currentmarker}{}%
\end{pgfscope}%
\begin{pgfscope}%
\pgfsys@transformshift{2.913101in}{2.412654in}%
\pgfsys@useobject{currentmarker}{}%
\end{pgfscope}%
\begin{pgfscope}%
\pgfsys@transformshift{2.934065in}{2.412654in}%
\pgfsys@useobject{currentmarker}{}%
\end{pgfscope}%
\begin{pgfscope}%
\pgfsys@transformshift{2.955029in}{2.412654in}%
\pgfsys@useobject{currentmarker}{}%
\end{pgfscope}%
\begin{pgfscope}%
\pgfsys@transformshift{2.975993in}{2.412654in}%
\pgfsys@useobject{currentmarker}{}%
\end{pgfscope}%
\end{pgfscope}%
\begin{pgfscope}%
\pgfsetrectcap%
\pgfsetmiterjoin%
\pgfsetlinewidth{0.803000pt}%
\definecolor{currentstroke}{rgb}{0.000000,0.000000,0.000000}%
\pgfsetstrokecolor{currentstroke}%
\pgfsetdash{}{0pt}%
\pgfpathmoveto{\pgfqpoint{0.495251in}{1.287254in}}%
\pgfpathlineto{\pgfqpoint{0.495251in}{3.212746in}}%
\pgfusepath{stroke}%
\end{pgfscope}%
\begin{pgfscope}%
\pgfsetrectcap%
\pgfsetmiterjoin%
\pgfsetlinewidth{0.803000pt}%
\definecolor{currentstroke}{rgb}{0.000000,0.000000,0.000000}%
\pgfsetstrokecolor{currentstroke}%
\pgfsetdash{}{0pt}%
\pgfpathmoveto{\pgfqpoint{4.932639in}{1.287254in}}%
\pgfpathlineto{\pgfqpoint{4.932639in}{3.212746in}}%
\pgfusepath{stroke}%
\end{pgfscope}%
\begin{pgfscope}%
\pgfsetrectcap%
\pgfsetmiterjoin%
\pgfsetlinewidth{0.803000pt}%
\definecolor{currentstroke}{rgb}{0.000000,0.000000,0.000000}%
\pgfsetstrokecolor{currentstroke}%
\pgfsetdash{}{0pt}%
\pgfpathmoveto{\pgfqpoint{0.495251in}{1.287254in}}%
\pgfpathlineto{\pgfqpoint{4.932639in}{1.287254in}}%
\pgfusepath{stroke}%
\end{pgfscope}%
\begin{pgfscope}%
\pgfsetrectcap%
\pgfsetmiterjoin%
\pgfsetlinewidth{0.803000pt}%
\definecolor{currentstroke}{rgb}{0.000000,0.000000,0.000000}%
\pgfsetstrokecolor{currentstroke}%
\pgfsetdash{}{0pt}%
\pgfpathmoveto{\pgfqpoint{0.495251in}{3.212746in}}%
\pgfpathlineto{\pgfqpoint{4.932639in}{3.212746in}}%
\pgfusepath{stroke}%
\end{pgfscope}%
\begin{pgfscope}%
\definecolor{textcolor}{rgb}{0.000000,0.000000,0.000000}%
\pgfsetstrokecolor{textcolor}%
\pgfsetfillcolor{textcolor}%
\pgftext[x=2.713945in,y=3.296080in,,base]{\color{textcolor}{\rmfamily\fontsize{12.000000}{14.400000}\selectfont\catcode`\^=\active\def^{\ifmmode\sp\else\^{}\fi}\catcode`\%=\active\def%{\%}z-plane: $z = x + iy$}}%
\end{pgfscope}%
\begin{pgfscope}%
\pgfsetbuttcap%
\pgfsetmiterjoin%
\definecolor{currentfill}{rgb}{1.000000,1.000000,1.000000}%
\pgfsetfillcolor{currentfill}%
\pgfsetlinewidth{0.000000pt}%
\definecolor{currentstroke}{rgb}{0.000000,0.000000,0.000000}%
\pgfsetstrokecolor{currentstroke}%
\pgfsetstrokeopacity{0.000000}%
\pgfsetdash{}{0pt}%
\pgfpathmoveto{\pgfqpoint{5.486220in}{0.585980in}}%
\pgfpathlineto{\pgfqpoint{9.923608in}{0.585980in}}%
\pgfpathlineto{\pgfqpoint{9.923608in}{3.914020in}}%
\pgfpathlineto{\pgfqpoint{5.486220in}{3.914020in}}%
\pgfpathlineto{\pgfqpoint{5.486220in}{0.585980in}}%
\pgfpathclose%
\pgfusepath{fill}%
\end{pgfscope}%
\begin{pgfscope}%
\pgfpathrectangle{\pgfqpoint{5.486220in}{0.585980in}}{\pgfqpoint{4.437388in}{3.328041in}}%
\pgfusepath{clip}%
\pgfsetrectcap%
\pgfsetroundjoin%
\pgfsetlinewidth{0.803000pt}%
\definecolor{currentstroke}{rgb}{0.690196,0.690196,0.690196}%
\pgfsetstrokecolor{currentstroke}%
\pgfsetstrokeopacity{0.300000}%
\pgfsetdash{}{0pt}%
\pgfpathmoveto{\pgfqpoint{5.486220in}{0.585980in}}%
\pgfpathlineto{\pgfqpoint{5.486220in}{3.914020in}}%
\pgfusepath{stroke}%
\end{pgfscope}%
\begin{pgfscope}%
\pgfsetbuttcap%
\pgfsetroundjoin%
\definecolor{currentfill}{rgb}{0.000000,0.000000,0.000000}%
\pgfsetfillcolor{currentfill}%
\pgfsetlinewidth{0.803000pt}%
\definecolor{currentstroke}{rgb}{0.000000,0.000000,0.000000}%
\pgfsetstrokecolor{currentstroke}%
\pgfsetdash{}{0pt}%
\pgfsys@defobject{currentmarker}{\pgfqpoint{0.000000in}{-0.048611in}}{\pgfqpoint{0.000000in}{0.000000in}}{%
\pgfpathmoveto{\pgfqpoint{0.000000in}{0.000000in}}%
\pgfpathlineto{\pgfqpoint{0.000000in}{-0.048611in}}%
\pgfusepath{stroke,fill}%
}%
\begin{pgfscope}%
\pgfsys@transformshift{5.486220in}{0.585980in}%
\pgfsys@useobject{currentmarker}{}%
\end{pgfscope}%
\end{pgfscope}%
\begin{pgfscope}%
\definecolor{textcolor}{rgb}{0.000000,0.000000,0.000000}%
\pgfsetstrokecolor{textcolor}%
\pgfsetfillcolor{textcolor}%
\pgftext[x=5.486220in,y=0.488757in,,top]{\color{textcolor}{\rmfamily\fontsize{10.000000}{12.000000}\selectfont\catcode`\^=\active\def^{\ifmmode\sp\else\^{}\fi}\catcode`\%=\active\def%{\%}$\mathdefault{\ensuremath{-}4}$}}%
\end{pgfscope}%
\begin{pgfscope}%
\pgfpathrectangle{\pgfqpoint{5.486220in}{0.585980in}}{\pgfqpoint{4.437388in}{3.328041in}}%
\pgfusepath{clip}%
\pgfsetrectcap%
\pgfsetroundjoin%
\pgfsetlinewidth{0.803000pt}%
\definecolor{currentstroke}{rgb}{0.690196,0.690196,0.690196}%
\pgfsetstrokecolor{currentstroke}%
\pgfsetstrokeopacity{0.300000}%
\pgfsetdash{}{0pt}%
\pgfpathmoveto{\pgfqpoint{6.040893in}{0.585980in}}%
\pgfpathlineto{\pgfqpoint{6.040893in}{3.914020in}}%
\pgfusepath{stroke}%
\end{pgfscope}%
\begin{pgfscope}%
\pgfsetbuttcap%
\pgfsetroundjoin%
\definecolor{currentfill}{rgb}{0.000000,0.000000,0.000000}%
\pgfsetfillcolor{currentfill}%
\pgfsetlinewidth{0.803000pt}%
\definecolor{currentstroke}{rgb}{0.000000,0.000000,0.000000}%
\pgfsetstrokecolor{currentstroke}%
\pgfsetdash{}{0pt}%
\pgfsys@defobject{currentmarker}{\pgfqpoint{0.000000in}{-0.048611in}}{\pgfqpoint{0.000000in}{0.000000in}}{%
\pgfpathmoveto{\pgfqpoint{0.000000in}{0.000000in}}%
\pgfpathlineto{\pgfqpoint{0.000000in}{-0.048611in}}%
\pgfusepath{stroke,fill}%
}%
\begin{pgfscope}%
\pgfsys@transformshift{6.040893in}{0.585980in}%
\pgfsys@useobject{currentmarker}{}%
\end{pgfscope}%
\end{pgfscope}%
\begin{pgfscope}%
\definecolor{textcolor}{rgb}{0.000000,0.000000,0.000000}%
\pgfsetstrokecolor{textcolor}%
\pgfsetfillcolor{textcolor}%
\pgftext[x=6.040893in,y=0.488757in,,top]{\color{textcolor}{\rmfamily\fontsize{10.000000}{12.000000}\selectfont\catcode`\^=\active\def^{\ifmmode\sp\else\^{}\fi}\catcode`\%=\active\def%{\%}$\mathdefault{\ensuremath{-}3}$}}%
\end{pgfscope}%
\begin{pgfscope}%
\pgfpathrectangle{\pgfqpoint{5.486220in}{0.585980in}}{\pgfqpoint{4.437388in}{3.328041in}}%
\pgfusepath{clip}%
\pgfsetrectcap%
\pgfsetroundjoin%
\pgfsetlinewidth{0.803000pt}%
\definecolor{currentstroke}{rgb}{0.690196,0.690196,0.690196}%
\pgfsetstrokecolor{currentstroke}%
\pgfsetstrokeopacity{0.300000}%
\pgfsetdash{}{0pt}%
\pgfpathmoveto{\pgfqpoint{6.595567in}{0.585980in}}%
\pgfpathlineto{\pgfqpoint{6.595567in}{3.914020in}}%
\pgfusepath{stroke}%
\end{pgfscope}%
\begin{pgfscope}%
\pgfsetbuttcap%
\pgfsetroundjoin%
\definecolor{currentfill}{rgb}{0.000000,0.000000,0.000000}%
\pgfsetfillcolor{currentfill}%
\pgfsetlinewidth{0.803000pt}%
\definecolor{currentstroke}{rgb}{0.000000,0.000000,0.000000}%
\pgfsetstrokecolor{currentstroke}%
\pgfsetdash{}{0pt}%
\pgfsys@defobject{currentmarker}{\pgfqpoint{0.000000in}{-0.048611in}}{\pgfqpoint{0.000000in}{0.000000in}}{%
\pgfpathmoveto{\pgfqpoint{0.000000in}{0.000000in}}%
\pgfpathlineto{\pgfqpoint{0.000000in}{-0.048611in}}%
\pgfusepath{stroke,fill}%
}%
\begin{pgfscope}%
\pgfsys@transformshift{6.595567in}{0.585980in}%
\pgfsys@useobject{currentmarker}{}%
\end{pgfscope}%
\end{pgfscope}%
\begin{pgfscope}%
\definecolor{textcolor}{rgb}{0.000000,0.000000,0.000000}%
\pgfsetstrokecolor{textcolor}%
\pgfsetfillcolor{textcolor}%
\pgftext[x=6.595567in,y=0.488757in,,top]{\color{textcolor}{\rmfamily\fontsize{10.000000}{12.000000}\selectfont\catcode`\^=\active\def^{\ifmmode\sp\else\^{}\fi}\catcode`\%=\active\def%{\%}$\mathdefault{\ensuremath{-}2}$}}%
\end{pgfscope}%
\begin{pgfscope}%
\pgfpathrectangle{\pgfqpoint{5.486220in}{0.585980in}}{\pgfqpoint{4.437388in}{3.328041in}}%
\pgfusepath{clip}%
\pgfsetrectcap%
\pgfsetroundjoin%
\pgfsetlinewidth{0.803000pt}%
\definecolor{currentstroke}{rgb}{0.690196,0.690196,0.690196}%
\pgfsetstrokecolor{currentstroke}%
\pgfsetstrokeopacity{0.300000}%
\pgfsetdash{}{0pt}%
\pgfpathmoveto{\pgfqpoint{7.150240in}{0.585980in}}%
\pgfpathlineto{\pgfqpoint{7.150240in}{3.914020in}}%
\pgfusepath{stroke}%
\end{pgfscope}%
\begin{pgfscope}%
\pgfsetbuttcap%
\pgfsetroundjoin%
\definecolor{currentfill}{rgb}{0.000000,0.000000,0.000000}%
\pgfsetfillcolor{currentfill}%
\pgfsetlinewidth{0.803000pt}%
\definecolor{currentstroke}{rgb}{0.000000,0.000000,0.000000}%
\pgfsetstrokecolor{currentstroke}%
\pgfsetdash{}{0pt}%
\pgfsys@defobject{currentmarker}{\pgfqpoint{0.000000in}{-0.048611in}}{\pgfqpoint{0.000000in}{0.000000in}}{%
\pgfpathmoveto{\pgfqpoint{0.000000in}{0.000000in}}%
\pgfpathlineto{\pgfqpoint{0.000000in}{-0.048611in}}%
\pgfusepath{stroke,fill}%
}%
\begin{pgfscope}%
\pgfsys@transformshift{7.150240in}{0.585980in}%
\pgfsys@useobject{currentmarker}{}%
\end{pgfscope}%
\end{pgfscope}%
\begin{pgfscope}%
\definecolor{textcolor}{rgb}{0.000000,0.000000,0.000000}%
\pgfsetstrokecolor{textcolor}%
\pgfsetfillcolor{textcolor}%
\pgftext[x=7.150240in,y=0.488757in,,top]{\color{textcolor}{\rmfamily\fontsize{10.000000}{12.000000}\selectfont\catcode`\^=\active\def^{\ifmmode\sp\else\^{}\fi}\catcode`\%=\active\def%{\%}$\mathdefault{\ensuremath{-}1}$}}%
\end{pgfscope}%
\begin{pgfscope}%
\pgfpathrectangle{\pgfqpoint{5.486220in}{0.585980in}}{\pgfqpoint{4.437388in}{3.328041in}}%
\pgfusepath{clip}%
\pgfsetrectcap%
\pgfsetroundjoin%
\pgfsetlinewidth{0.803000pt}%
\definecolor{currentstroke}{rgb}{0.690196,0.690196,0.690196}%
\pgfsetstrokecolor{currentstroke}%
\pgfsetstrokeopacity{0.300000}%
\pgfsetdash{}{0pt}%
\pgfpathmoveto{\pgfqpoint{7.704914in}{0.585980in}}%
\pgfpathlineto{\pgfqpoint{7.704914in}{3.914020in}}%
\pgfusepath{stroke}%
\end{pgfscope}%
\begin{pgfscope}%
\pgfsetbuttcap%
\pgfsetroundjoin%
\definecolor{currentfill}{rgb}{0.000000,0.000000,0.000000}%
\pgfsetfillcolor{currentfill}%
\pgfsetlinewidth{0.803000pt}%
\definecolor{currentstroke}{rgb}{0.000000,0.000000,0.000000}%
\pgfsetstrokecolor{currentstroke}%
\pgfsetdash{}{0pt}%
\pgfsys@defobject{currentmarker}{\pgfqpoint{0.000000in}{-0.048611in}}{\pgfqpoint{0.000000in}{0.000000in}}{%
\pgfpathmoveto{\pgfqpoint{0.000000in}{0.000000in}}%
\pgfpathlineto{\pgfqpoint{0.000000in}{-0.048611in}}%
\pgfusepath{stroke,fill}%
}%
\begin{pgfscope}%
\pgfsys@transformshift{7.704914in}{0.585980in}%
\pgfsys@useobject{currentmarker}{}%
\end{pgfscope}%
\end{pgfscope}%
\begin{pgfscope}%
\definecolor{textcolor}{rgb}{0.000000,0.000000,0.000000}%
\pgfsetstrokecolor{textcolor}%
\pgfsetfillcolor{textcolor}%
\pgftext[x=7.704914in,y=0.488757in,,top]{\color{textcolor}{\rmfamily\fontsize{10.000000}{12.000000}\selectfont\catcode`\^=\active\def^{\ifmmode\sp\else\^{}\fi}\catcode`\%=\active\def%{\%}$\mathdefault{0}$}}%
\end{pgfscope}%
\begin{pgfscope}%
\pgfpathrectangle{\pgfqpoint{5.486220in}{0.585980in}}{\pgfqpoint{4.437388in}{3.328041in}}%
\pgfusepath{clip}%
\pgfsetrectcap%
\pgfsetroundjoin%
\pgfsetlinewidth{0.803000pt}%
\definecolor{currentstroke}{rgb}{0.690196,0.690196,0.690196}%
\pgfsetstrokecolor{currentstroke}%
\pgfsetstrokeopacity{0.300000}%
\pgfsetdash{}{0pt}%
\pgfpathmoveto{\pgfqpoint{8.259587in}{0.585980in}}%
\pgfpathlineto{\pgfqpoint{8.259587in}{3.914020in}}%
\pgfusepath{stroke}%
\end{pgfscope}%
\begin{pgfscope}%
\pgfsetbuttcap%
\pgfsetroundjoin%
\definecolor{currentfill}{rgb}{0.000000,0.000000,0.000000}%
\pgfsetfillcolor{currentfill}%
\pgfsetlinewidth{0.803000pt}%
\definecolor{currentstroke}{rgb}{0.000000,0.000000,0.000000}%
\pgfsetstrokecolor{currentstroke}%
\pgfsetdash{}{0pt}%
\pgfsys@defobject{currentmarker}{\pgfqpoint{0.000000in}{-0.048611in}}{\pgfqpoint{0.000000in}{0.000000in}}{%
\pgfpathmoveto{\pgfqpoint{0.000000in}{0.000000in}}%
\pgfpathlineto{\pgfqpoint{0.000000in}{-0.048611in}}%
\pgfusepath{stroke,fill}%
}%
\begin{pgfscope}%
\pgfsys@transformshift{8.259587in}{0.585980in}%
\pgfsys@useobject{currentmarker}{}%
\end{pgfscope}%
\end{pgfscope}%
\begin{pgfscope}%
\definecolor{textcolor}{rgb}{0.000000,0.000000,0.000000}%
\pgfsetstrokecolor{textcolor}%
\pgfsetfillcolor{textcolor}%
\pgftext[x=8.259587in,y=0.488757in,,top]{\color{textcolor}{\rmfamily\fontsize{10.000000}{12.000000}\selectfont\catcode`\^=\active\def^{\ifmmode\sp\else\^{}\fi}\catcode`\%=\active\def%{\%}$\mathdefault{1}$}}%
\end{pgfscope}%
\begin{pgfscope}%
\pgfpathrectangle{\pgfqpoint{5.486220in}{0.585980in}}{\pgfqpoint{4.437388in}{3.328041in}}%
\pgfusepath{clip}%
\pgfsetrectcap%
\pgfsetroundjoin%
\pgfsetlinewidth{0.803000pt}%
\definecolor{currentstroke}{rgb}{0.690196,0.690196,0.690196}%
\pgfsetstrokecolor{currentstroke}%
\pgfsetstrokeopacity{0.300000}%
\pgfsetdash{}{0pt}%
\pgfpathmoveto{\pgfqpoint{8.814261in}{0.585980in}}%
\pgfpathlineto{\pgfqpoint{8.814261in}{3.914020in}}%
\pgfusepath{stroke}%
\end{pgfscope}%
\begin{pgfscope}%
\pgfsetbuttcap%
\pgfsetroundjoin%
\definecolor{currentfill}{rgb}{0.000000,0.000000,0.000000}%
\pgfsetfillcolor{currentfill}%
\pgfsetlinewidth{0.803000pt}%
\definecolor{currentstroke}{rgb}{0.000000,0.000000,0.000000}%
\pgfsetstrokecolor{currentstroke}%
\pgfsetdash{}{0pt}%
\pgfsys@defobject{currentmarker}{\pgfqpoint{0.000000in}{-0.048611in}}{\pgfqpoint{0.000000in}{0.000000in}}{%
\pgfpathmoveto{\pgfqpoint{0.000000in}{0.000000in}}%
\pgfpathlineto{\pgfqpoint{0.000000in}{-0.048611in}}%
\pgfusepath{stroke,fill}%
}%
\begin{pgfscope}%
\pgfsys@transformshift{8.814261in}{0.585980in}%
\pgfsys@useobject{currentmarker}{}%
\end{pgfscope}%
\end{pgfscope}%
\begin{pgfscope}%
\definecolor{textcolor}{rgb}{0.000000,0.000000,0.000000}%
\pgfsetstrokecolor{textcolor}%
\pgfsetfillcolor{textcolor}%
\pgftext[x=8.814261in,y=0.488757in,,top]{\color{textcolor}{\rmfamily\fontsize{10.000000}{12.000000}\selectfont\catcode`\^=\active\def^{\ifmmode\sp\else\^{}\fi}\catcode`\%=\active\def%{\%}$\mathdefault{2}$}}%
\end{pgfscope}%
\begin{pgfscope}%
\pgfpathrectangle{\pgfqpoint{5.486220in}{0.585980in}}{\pgfqpoint{4.437388in}{3.328041in}}%
\pgfusepath{clip}%
\pgfsetrectcap%
\pgfsetroundjoin%
\pgfsetlinewidth{0.803000pt}%
\definecolor{currentstroke}{rgb}{0.690196,0.690196,0.690196}%
\pgfsetstrokecolor{currentstroke}%
\pgfsetstrokeopacity{0.300000}%
\pgfsetdash{}{0pt}%
\pgfpathmoveto{\pgfqpoint{9.368934in}{0.585980in}}%
\pgfpathlineto{\pgfqpoint{9.368934in}{3.914020in}}%
\pgfusepath{stroke}%
\end{pgfscope}%
\begin{pgfscope}%
\pgfsetbuttcap%
\pgfsetroundjoin%
\definecolor{currentfill}{rgb}{0.000000,0.000000,0.000000}%
\pgfsetfillcolor{currentfill}%
\pgfsetlinewidth{0.803000pt}%
\definecolor{currentstroke}{rgb}{0.000000,0.000000,0.000000}%
\pgfsetstrokecolor{currentstroke}%
\pgfsetdash{}{0pt}%
\pgfsys@defobject{currentmarker}{\pgfqpoint{0.000000in}{-0.048611in}}{\pgfqpoint{0.000000in}{0.000000in}}{%
\pgfpathmoveto{\pgfqpoint{0.000000in}{0.000000in}}%
\pgfpathlineto{\pgfqpoint{0.000000in}{-0.048611in}}%
\pgfusepath{stroke,fill}%
}%
\begin{pgfscope}%
\pgfsys@transformshift{9.368934in}{0.585980in}%
\pgfsys@useobject{currentmarker}{}%
\end{pgfscope}%
\end{pgfscope}%
\begin{pgfscope}%
\definecolor{textcolor}{rgb}{0.000000,0.000000,0.000000}%
\pgfsetstrokecolor{textcolor}%
\pgfsetfillcolor{textcolor}%
\pgftext[x=9.368934in,y=0.488757in,,top]{\color{textcolor}{\rmfamily\fontsize{10.000000}{12.000000}\selectfont\catcode`\^=\active\def^{\ifmmode\sp\else\^{}\fi}\catcode`\%=\active\def%{\%}$\mathdefault{3}$}}%
\end{pgfscope}%
\begin{pgfscope}%
\pgfpathrectangle{\pgfqpoint{5.486220in}{0.585980in}}{\pgfqpoint{4.437388in}{3.328041in}}%
\pgfusepath{clip}%
\pgfsetrectcap%
\pgfsetroundjoin%
\pgfsetlinewidth{0.803000pt}%
\definecolor{currentstroke}{rgb}{0.690196,0.690196,0.690196}%
\pgfsetstrokecolor{currentstroke}%
\pgfsetstrokeopacity{0.300000}%
\pgfsetdash{}{0pt}%
\pgfpathmoveto{\pgfqpoint{9.923608in}{0.585980in}}%
\pgfpathlineto{\pgfqpoint{9.923608in}{3.914020in}}%
\pgfusepath{stroke}%
\end{pgfscope}%
\begin{pgfscope}%
\pgfsetbuttcap%
\pgfsetroundjoin%
\definecolor{currentfill}{rgb}{0.000000,0.000000,0.000000}%
\pgfsetfillcolor{currentfill}%
\pgfsetlinewidth{0.803000pt}%
\definecolor{currentstroke}{rgb}{0.000000,0.000000,0.000000}%
\pgfsetstrokecolor{currentstroke}%
\pgfsetdash{}{0pt}%
\pgfsys@defobject{currentmarker}{\pgfqpoint{0.000000in}{-0.048611in}}{\pgfqpoint{0.000000in}{0.000000in}}{%
\pgfpathmoveto{\pgfqpoint{0.000000in}{0.000000in}}%
\pgfpathlineto{\pgfqpoint{0.000000in}{-0.048611in}}%
\pgfusepath{stroke,fill}%
}%
\begin{pgfscope}%
\pgfsys@transformshift{9.923608in}{0.585980in}%
\pgfsys@useobject{currentmarker}{}%
\end{pgfscope}%
\end{pgfscope}%
\begin{pgfscope}%
\definecolor{textcolor}{rgb}{0.000000,0.000000,0.000000}%
\pgfsetstrokecolor{textcolor}%
\pgfsetfillcolor{textcolor}%
\pgftext[x=9.923608in,y=0.488757in,,top]{\color{textcolor}{\rmfamily\fontsize{10.000000}{12.000000}\selectfont\catcode`\^=\active\def^{\ifmmode\sp\else\^{}\fi}\catcode`\%=\active\def%{\%}$\mathdefault{4}$}}%
\end{pgfscope}%
\begin{pgfscope}%
\definecolor{textcolor}{rgb}{0.000000,0.000000,0.000000}%
\pgfsetstrokecolor{textcolor}%
\pgfsetfillcolor{textcolor}%
\pgftext[x=7.704914in,y=0.309868in,,top]{\color{textcolor}{\rmfamily\fontsize{10.000000}{12.000000}\selectfont\catcode`\^=\active\def^{\ifmmode\sp\else\^{}\fi}\catcode`\%=\active\def%{\%}$u$}}%
\end{pgfscope}%
\begin{pgfscope}%
\pgfpathrectangle{\pgfqpoint{5.486220in}{0.585980in}}{\pgfqpoint{4.437388in}{3.328041in}}%
\pgfusepath{clip}%
\pgfsetrectcap%
\pgfsetroundjoin%
\pgfsetlinewidth{0.803000pt}%
\definecolor{currentstroke}{rgb}{0.690196,0.690196,0.690196}%
\pgfsetstrokecolor{currentstroke}%
\pgfsetstrokeopacity{0.300000}%
\pgfsetdash{}{0pt}%
\pgfpathmoveto{\pgfqpoint{5.486220in}{0.585980in}}%
\pgfpathlineto{\pgfqpoint{9.923608in}{0.585980in}}%
\pgfusepath{stroke}%
\end{pgfscope}%
\begin{pgfscope}%
\pgfsetbuttcap%
\pgfsetroundjoin%
\definecolor{currentfill}{rgb}{0.000000,0.000000,0.000000}%
\pgfsetfillcolor{currentfill}%
\pgfsetlinewidth{0.803000pt}%
\definecolor{currentstroke}{rgb}{0.000000,0.000000,0.000000}%
\pgfsetstrokecolor{currentstroke}%
\pgfsetdash{}{0pt}%
\pgfsys@defobject{currentmarker}{\pgfqpoint{-0.048611in}{0.000000in}}{\pgfqpoint{-0.000000in}{0.000000in}}{%
\pgfpathmoveto{\pgfqpoint{-0.000000in}{0.000000in}}%
\pgfpathlineto{\pgfqpoint{-0.048611in}{0.000000in}}%
\pgfusepath{stroke,fill}%
}%
\begin{pgfscope}%
\pgfsys@transformshift{5.486220in}{0.585980in}%
\pgfsys@useobject{currentmarker}{}%
\end{pgfscope}%
\end{pgfscope}%
\begin{pgfscope}%
\definecolor{textcolor}{rgb}{0.000000,0.000000,0.000000}%
\pgfsetstrokecolor{textcolor}%
\pgfsetfillcolor{textcolor}%
\pgftext[x=5.211528in, y=0.537785in, left, base]{\color{textcolor}{\rmfamily\fontsize{10.000000}{12.000000}\selectfont\catcode`\^=\active\def^{\ifmmode\sp\else\^{}\fi}\catcode`\%=\active\def%{\%}$\mathdefault{\ensuremath{-}3}$}}%
\end{pgfscope}%
\begin{pgfscope}%
\pgfpathrectangle{\pgfqpoint{5.486220in}{0.585980in}}{\pgfqpoint{4.437388in}{3.328041in}}%
\pgfusepath{clip}%
\pgfsetrectcap%
\pgfsetroundjoin%
\pgfsetlinewidth{0.803000pt}%
\definecolor{currentstroke}{rgb}{0.690196,0.690196,0.690196}%
\pgfsetstrokecolor{currentstroke}%
\pgfsetstrokeopacity{0.300000}%
\pgfsetdash{}{0pt}%
\pgfpathmoveto{\pgfqpoint{5.486220in}{1.140653in}}%
\pgfpathlineto{\pgfqpoint{9.923608in}{1.140653in}}%
\pgfusepath{stroke}%
\end{pgfscope}%
\begin{pgfscope}%
\pgfsetbuttcap%
\pgfsetroundjoin%
\definecolor{currentfill}{rgb}{0.000000,0.000000,0.000000}%
\pgfsetfillcolor{currentfill}%
\pgfsetlinewidth{0.803000pt}%
\definecolor{currentstroke}{rgb}{0.000000,0.000000,0.000000}%
\pgfsetstrokecolor{currentstroke}%
\pgfsetdash{}{0pt}%
\pgfsys@defobject{currentmarker}{\pgfqpoint{-0.048611in}{0.000000in}}{\pgfqpoint{-0.000000in}{0.000000in}}{%
\pgfpathmoveto{\pgfqpoint{-0.000000in}{0.000000in}}%
\pgfpathlineto{\pgfqpoint{-0.048611in}{0.000000in}}%
\pgfusepath{stroke,fill}%
}%
\begin{pgfscope}%
\pgfsys@transformshift{5.486220in}{1.140653in}%
\pgfsys@useobject{currentmarker}{}%
\end{pgfscope}%
\end{pgfscope}%
\begin{pgfscope}%
\definecolor{textcolor}{rgb}{0.000000,0.000000,0.000000}%
\pgfsetstrokecolor{textcolor}%
\pgfsetfillcolor{textcolor}%
\pgftext[x=5.211528in, y=1.092459in, left, base]{\color{textcolor}{\rmfamily\fontsize{10.000000}{12.000000}\selectfont\catcode`\^=\active\def^{\ifmmode\sp\else\^{}\fi}\catcode`\%=\active\def%{\%}$\mathdefault{\ensuremath{-}2}$}}%
\end{pgfscope}%
\begin{pgfscope}%
\pgfpathrectangle{\pgfqpoint{5.486220in}{0.585980in}}{\pgfqpoint{4.437388in}{3.328041in}}%
\pgfusepath{clip}%
\pgfsetrectcap%
\pgfsetroundjoin%
\pgfsetlinewidth{0.803000pt}%
\definecolor{currentstroke}{rgb}{0.690196,0.690196,0.690196}%
\pgfsetstrokecolor{currentstroke}%
\pgfsetstrokeopacity{0.300000}%
\pgfsetdash{}{0pt}%
\pgfpathmoveto{\pgfqpoint{5.486220in}{1.695327in}}%
\pgfpathlineto{\pgfqpoint{9.923608in}{1.695327in}}%
\pgfusepath{stroke}%
\end{pgfscope}%
\begin{pgfscope}%
\pgfsetbuttcap%
\pgfsetroundjoin%
\definecolor{currentfill}{rgb}{0.000000,0.000000,0.000000}%
\pgfsetfillcolor{currentfill}%
\pgfsetlinewidth{0.803000pt}%
\definecolor{currentstroke}{rgb}{0.000000,0.000000,0.000000}%
\pgfsetstrokecolor{currentstroke}%
\pgfsetdash{}{0pt}%
\pgfsys@defobject{currentmarker}{\pgfqpoint{-0.048611in}{0.000000in}}{\pgfqpoint{-0.000000in}{0.000000in}}{%
\pgfpathmoveto{\pgfqpoint{-0.000000in}{0.000000in}}%
\pgfpathlineto{\pgfqpoint{-0.048611in}{0.000000in}}%
\pgfusepath{stroke,fill}%
}%
\begin{pgfscope}%
\pgfsys@transformshift{5.486220in}{1.695327in}%
\pgfsys@useobject{currentmarker}{}%
\end{pgfscope}%
\end{pgfscope}%
\begin{pgfscope}%
\definecolor{textcolor}{rgb}{0.000000,0.000000,0.000000}%
\pgfsetstrokecolor{textcolor}%
\pgfsetfillcolor{textcolor}%
\pgftext[x=5.211528in, y=1.647132in, left, base]{\color{textcolor}{\rmfamily\fontsize{10.000000}{12.000000}\selectfont\catcode`\^=\active\def^{\ifmmode\sp\else\^{}\fi}\catcode`\%=\active\def%{\%}$\mathdefault{\ensuremath{-}1}$}}%
\end{pgfscope}%
\begin{pgfscope}%
\pgfpathrectangle{\pgfqpoint{5.486220in}{0.585980in}}{\pgfqpoint{4.437388in}{3.328041in}}%
\pgfusepath{clip}%
\pgfsetrectcap%
\pgfsetroundjoin%
\pgfsetlinewidth{0.803000pt}%
\definecolor{currentstroke}{rgb}{0.690196,0.690196,0.690196}%
\pgfsetstrokecolor{currentstroke}%
\pgfsetstrokeopacity{0.300000}%
\pgfsetdash{}{0pt}%
\pgfpathmoveto{\pgfqpoint{5.486220in}{2.250000in}}%
\pgfpathlineto{\pgfqpoint{9.923608in}{2.250000in}}%
\pgfusepath{stroke}%
\end{pgfscope}%
\begin{pgfscope}%
\pgfsetbuttcap%
\pgfsetroundjoin%
\definecolor{currentfill}{rgb}{0.000000,0.000000,0.000000}%
\pgfsetfillcolor{currentfill}%
\pgfsetlinewidth{0.803000pt}%
\definecolor{currentstroke}{rgb}{0.000000,0.000000,0.000000}%
\pgfsetstrokecolor{currentstroke}%
\pgfsetdash{}{0pt}%
\pgfsys@defobject{currentmarker}{\pgfqpoint{-0.048611in}{0.000000in}}{\pgfqpoint{-0.000000in}{0.000000in}}{%
\pgfpathmoveto{\pgfqpoint{-0.000000in}{0.000000in}}%
\pgfpathlineto{\pgfqpoint{-0.048611in}{0.000000in}}%
\pgfusepath{stroke,fill}%
}%
\begin{pgfscope}%
\pgfsys@transformshift{5.486220in}{2.250000in}%
\pgfsys@useobject{currentmarker}{}%
\end{pgfscope}%
\end{pgfscope}%
\begin{pgfscope}%
\definecolor{textcolor}{rgb}{0.000000,0.000000,0.000000}%
\pgfsetstrokecolor{textcolor}%
\pgfsetfillcolor{textcolor}%
\pgftext[x=5.319553in, y=2.201806in, left, base]{\color{textcolor}{\rmfamily\fontsize{10.000000}{12.000000}\selectfont\catcode`\^=\active\def^{\ifmmode\sp\else\^{}\fi}\catcode`\%=\active\def%{\%}$\mathdefault{0}$}}%
\end{pgfscope}%
\begin{pgfscope}%
\pgfpathrectangle{\pgfqpoint{5.486220in}{0.585980in}}{\pgfqpoint{4.437388in}{3.328041in}}%
\pgfusepath{clip}%
\pgfsetrectcap%
\pgfsetroundjoin%
\pgfsetlinewidth{0.803000pt}%
\definecolor{currentstroke}{rgb}{0.690196,0.690196,0.690196}%
\pgfsetstrokecolor{currentstroke}%
\pgfsetstrokeopacity{0.300000}%
\pgfsetdash{}{0pt}%
\pgfpathmoveto{\pgfqpoint{5.486220in}{2.804673in}}%
\pgfpathlineto{\pgfqpoint{9.923608in}{2.804673in}}%
\pgfusepath{stroke}%
\end{pgfscope}%
\begin{pgfscope}%
\pgfsetbuttcap%
\pgfsetroundjoin%
\definecolor{currentfill}{rgb}{0.000000,0.000000,0.000000}%
\pgfsetfillcolor{currentfill}%
\pgfsetlinewidth{0.803000pt}%
\definecolor{currentstroke}{rgb}{0.000000,0.000000,0.000000}%
\pgfsetstrokecolor{currentstroke}%
\pgfsetdash{}{0pt}%
\pgfsys@defobject{currentmarker}{\pgfqpoint{-0.048611in}{0.000000in}}{\pgfqpoint{-0.000000in}{0.000000in}}{%
\pgfpathmoveto{\pgfqpoint{-0.000000in}{0.000000in}}%
\pgfpathlineto{\pgfqpoint{-0.048611in}{0.000000in}}%
\pgfusepath{stroke,fill}%
}%
\begin{pgfscope}%
\pgfsys@transformshift{5.486220in}{2.804673in}%
\pgfsys@useobject{currentmarker}{}%
\end{pgfscope}%
\end{pgfscope}%
\begin{pgfscope}%
\definecolor{textcolor}{rgb}{0.000000,0.000000,0.000000}%
\pgfsetstrokecolor{textcolor}%
\pgfsetfillcolor{textcolor}%
\pgftext[x=5.319553in, y=2.756479in, left, base]{\color{textcolor}{\rmfamily\fontsize{10.000000}{12.000000}\selectfont\catcode`\^=\active\def^{\ifmmode\sp\else\^{}\fi}\catcode`\%=\active\def%{\%}$\mathdefault{1}$}}%
\end{pgfscope}%
\begin{pgfscope}%
\pgfpathrectangle{\pgfqpoint{5.486220in}{0.585980in}}{\pgfqpoint{4.437388in}{3.328041in}}%
\pgfusepath{clip}%
\pgfsetrectcap%
\pgfsetroundjoin%
\pgfsetlinewidth{0.803000pt}%
\definecolor{currentstroke}{rgb}{0.690196,0.690196,0.690196}%
\pgfsetstrokecolor{currentstroke}%
\pgfsetstrokeopacity{0.300000}%
\pgfsetdash{}{0pt}%
\pgfpathmoveto{\pgfqpoint{5.486220in}{3.359347in}}%
\pgfpathlineto{\pgfqpoint{9.923608in}{3.359347in}}%
\pgfusepath{stroke}%
\end{pgfscope}%
\begin{pgfscope}%
\pgfsetbuttcap%
\pgfsetroundjoin%
\definecolor{currentfill}{rgb}{0.000000,0.000000,0.000000}%
\pgfsetfillcolor{currentfill}%
\pgfsetlinewidth{0.803000pt}%
\definecolor{currentstroke}{rgb}{0.000000,0.000000,0.000000}%
\pgfsetstrokecolor{currentstroke}%
\pgfsetdash{}{0pt}%
\pgfsys@defobject{currentmarker}{\pgfqpoint{-0.048611in}{0.000000in}}{\pgfqpoint{-0.000000in}{0.000000in}}{%
\pgfpathmoveto{\pgfqpoint{-0.000000in}{0.000000in}}%
\pgfpathlineto{\pgfqpoint{-0.048611in}{0.000000in}}%
\pgfusepath{stroke,fill}%
}%
\begin{pgfscope}%
\pgfsys@transformshift{5.486220in}{3.359347in}%
\pgfsys@useobject{currentmarker}{}%
\end{pgfscope}%
\end{pgfscope}%
\begin{pgfscope}%
\definecolor{textcolor}{rgb}{0.000000,0.000000,0.000000}%
\pgfsetstrokecolor{textcolor}%
\pgfsetfillcolor{textcolor}%
\pgftext[x=5.319553in, y=3.311153in, left, base]{\color{textcolor}{\rmfamily\fontsize{10.000000}{12.000000}\selectfont\catcode`\^=\active\def^{\ifmmode\sp\else\^{}\fi}\catcode`\%=\active\def%{\%}$\mathdefault{2}$}}%
\end{pgfscope}%
\begin{pgfscope}%
\pgfpathrectangle{\pgfqpoint{5.486220in}{0.585980in}}{\pgfqpoint{4.437388in}{3.328041in}}%
\pgfusepath{clip}%
\pgfsetrectcap%
\pgfsetroundjoin%
\pgfsetlinewidth{0.803000pt}%
\definecolor{currentstroke}{rgb}{0.690196,0.690196,0.690196}%
\pgfsetstrokecolor{currentstroke}%
\pgfsetstrokeopacity{0.300000}%
\pgfsetdash{}{0pt}%
\pgfpathmoveto{\pgfqpoint{5.486220in}{3.914020in}}%
\pgfpathlineto{\pgfqpoint{9.923608in}{3.914020in}}%
\pgfusepath{stroke}%
\end{pgfscope}%
\begin{pgfscope}%
\pgfsetbuttcap%
\pgfsetroundjoin%
\definecolor{currentfill}{rgb}{0.000000,0.000000,0.000000}%
\pgfsetfillcolor{currentfill}%
\pgfsetlinewidth{0.803000pt}%
\definecolor{currentstroke}{rgb}{0.000000,0.000000,0.000000}%
\pgfsetstrokecolor{currentstroke}%
\pgfsetdash{}{0pt}%
\pgfsys@defobject{currentmarker}{\pgfqpoint{-0.048611in}{0.000000in}}{\pgfqpoint{-0.000000in}{0.000000in}}{%
\pgfpathmoveto{\pgfqpoint{-0.000000in}{0.000000in}}%
\pgfpathlineto{\pgfqpoint{-0.048611in}{0.000000in}}%
\pgfusepath{stroke,fill}%
}%
\begin{pgfscope}%
\pgfsys@transformshift{5.486220in}{3.914020in}%
\pgfsys@useobject{currentmarker}{}%
\end{pgfscope}%
\end{pgfscope}%
\begin{pgfscope}%
\definecolor{textcolor}{rgb}{0.000000,0.000000,0.000000}%
\pgfsetstrokecolor{textcolor}%
\pgfsetfillcolor{textcolor}%
\pgftext[x=5.319553in, y=3.865826in, left, base]{\color{textcolor}{\rmfamily\fontsize{10.000000}{12.000000}\selectfont\catcode`\^=\active\def^{\ifmmode\sp\else\^{}\fi}\catcode`\%=\active\def%{\%}$\mathdefault{3}$}}%
\end{pgfscope}%
\begin{pgfscope}%
\definecolor{textcolor}{rgb}{0.000000,0.000000,0.000000}%
\pgfsetstrokecolor{textcolor}%
\pgfsetfillcolor{textcolor}%
\pgftext[x=5.155972in,y=2.250000in,,bottom,rotate=90.000000]{\color{textcolor}{\rmfamily\fontsize{10.000000}{12.000000}\selectfont\catcode`\^=\active\def^{\ifmmode\sp\else\^{}\fi}\catcode`\%=\active\def%{\%}$v$}}%
\end{pgfscope}%
\begin{pgfscope}%
\pgfpathrectangle{\pgfqpoint{5.486220in}{0.585980in}}{\pgfqpoint{4.437388in}{3.328041in}}%
\pgfusepath{clip}%
\pgfsetrectcap%
\pgfsetroundjoin%
\pgfsetlinewidth{1.505625pt}%
\definecolor{currentstroke}{rgb}{0.267004,0.004874,0.329415}%
\pgfsetstrokecolor{currentstroke}%
\pgfsetstrokeopacity{0.800000}%
\pgfsetdash{}{0pt}%
\pgfpathmoveto{\pgfqpoint{7.150240in}{2.250000in}}%
\pgfpathlineto{\pgfqpoint{5.476220in}{2.250000in}}%
\pgfpathlineto{\pgfqpoint{5.476220in}{2.250000in}}%
\pgfusepath{stroke}%
\end{pgfscope}%
\begin{pgfscope}%
\pgfpathrectangle{\pgfqpoint{5.486220in}{0.585980in}}{\pgfqpoint{4.437388in}{3.328041in}}%
\pgfusepath{clip}%
\pgfsetrectcap%
\pgfsetroundjoin%
\pgfsetlinewidth{1.505625pt}%
\definecolor{currentstroke}{rgb}{0.283229,0.120777,0.440584}%
\pgfsetstrokecolor{currentstroke}%
\pgfsetstrokeopacity{0.800000}%
\pgfsetdash{}{0pt}%
\pgfpathmoveto{\pgfqpoint{7.224552in}{2.250000in}}%
\pgfpathlineto{\pgfqpoint{7.223604in}{2.267432in}}%
\pgfpathlineto{\pgfqpoint{7.220757in}{2.284933in}}%
\pgfpathlineto{\pgfqpoint{7.214816in}{2.306122in}}%
\pgfpathlineto{\pgfqpoint{7.206089in}{2.327630in}}%
\pgfpathlineto{\pgfqpoint{7.194526in}{2.349580in}}%
\pgfpathlineto{\pgfqpoint{7.180062in}{2.372096in}}%
\pgfpathlineto{\pgfqpoint{7.162615in}{2.395305in}}%
\pgfpathlineto{\pgfqpoint{7.138356in}{2.423436in}}%
\pgfpathlineto{\pgfqpoint{7.109713in}{2.452910in}}%
\pgfpathlineto{\pgfqpoint{7.076463in}{2.483953in}}%
\pgfpathlineto{\pgfqpoint{7.032493in}{2.521663in}}%
\pgfpathlineto{\pgfqpoint{6.981724in}{2.562120in}}%
\pgfpathlineto{\pgfqpoint{6.915843in}{2.611426in}}%
\pgfpathlineto{\pgfqpoint{6.839863in}{2.665358in}}%
\pgfpathlineto{\pgfqpoint{6.742410in}{2.731548in}}%
\pgfpathlineto{\pgfqpoint{6.629747in}{2.805349in}}%
\pgfpathlineto{\pgfqpoint{6.486116in}{2.896715in}}%
\pgfpathlineto{\pgfqpoint{6.302734in}{3.010561in}}%
\pgfpathlineto{\pgfqpoint{6.067893in}{3.153528in}}%
\pgfpathlineto{\pgfqpoint{5.765844in}{3.334627in}}%
\pgfpathlineto{\pgfqpoint{5.476220in}{3.506493in}}%
\pgfpathlineto{\pgfqpoint{5.476220in}{3.506493in}}%
\pgfusepath{stroke}%
\end{pgfscope}%
\begin{pgfscope}%
\pgfpathrectangle{\pgfqpoint{5.486220in}{0.585980in}}{\pgfqpoint{4.437388in}{3.328041in}}%
\pgfusepath{clip}%
\pgfsetrectcap%
\pgfsetroundjoin%
\pgfsetlinewidth{1.505625pt}%
\definecolor{currentstroke}{rgb}{0.267968,0.223549,0.512008}%
\pgfsetstrokecolor{currentstroke}%
\pgfsetstrokeopacity{0.800000}%
\pgfsetdash{}{0pt}%
\pgfpathmoveto{\pgfqpoint{7.427577in}{2.250000in}}%
\pgfpathlineto{\pgfqpoint{7.426504in}{2.292297in}}%
\pgfpathlineto{\pgfqpoint{7.423276in}{2.334922in}}%
\pgfpathlineto{\pgfqpoint{7.417869in}{2.378204in}}%
\pgfpathlineto{\pgfqpoint{7.410241in}{2.422478in}}%
\pgfpathlineto{\pgfqpoint{7.400333in}{2.468086in}}%
\pgfpathlineto{\pgfqpoint{7.386117in}{2.522298in}}%
\pgfpathlineto{\pgfqpoint{7.368679in}{2.579262in}}%
\pgfpathlineto{\pgfqpoint{7.347842in}{2.639555in}}%
\pgfpathlineto{\pgfqpoint{7.323395in}{2.703786in}}%
\pgfpathlineto{\pgfqpoint{7.291268in}{2.781564in}}%
\pgfpathlineto{\pgfqpoint{7.253848in}{2.866144in}}%
\pgfpathlineto{\pgfqpoint{7.205476in}{2.969420in}}%
\pgfpathlineto{\pgfqpoint{7.149212in}{3.084066in}}%
\pgfpathlineto{\pgfqpoint{7.077140in}{3.225475in}}%
\pgfpathlineto{\pgfqpoint{6.993060in}{3.385543in}}%
\pgfpathlineto{\pgfqpoint{6.885746in}{3.585050in}}%
\pgfpathlineto{\pgfqpoint{6.748354in}{3.835646in}}%
\pgfpathlineto{\pgfqpoint{6.699413in}{3.924020in}}%
\pgfpathlineto{\pgfqpoint{6.699413in}{3.924020in}}%
\pgfusepath{stroke}%
\end{pgfscope}%
\begin{pgfscope}%
\pgfpathrectangle{\pgfqpoint{5.486220in}{0.585980in}}{\pgfqpoint{4.437388in}{3.328041in}}%
\pgfusepath{clip}%
\pgfsetrectcap%
\pgfsetroundjoin%
\pgfsetlinewidth{1.505625pt}%
\definecolor{currentstroke}{rgb}{0.229739,0.322361,0.545706}%
\pgfsetstrokecolor{currentstroke}%
\pgfsetstrokeopacity{0.800000}%
\pgfsetdash{}{0pt}%
\pgfpathmoveto{\pgfqpoint{7.704914in}{2.250000in}}%
\pgfpathlineto{\pgfqpoint{7.704914in}{3.924020in}}%
\pgfpathlineto{\pgfqpoint{7.704914in}{3.924020in}}%
\pgfusepath{stroke}%
\end{pgfscope}%
\begin{pgfscope}%
\pgfpathrectangle{\pgfqpoint{5.486220in}{0.585980in}}{\pgfqpoint{4.437388in}{3.328041in}}%
\pgfusepath{clip}%
\pgfsetrectcap%
\pgfsetroundjoin%
\pgfsetlinewidth{1.505625pt}%
\definecolor{currentstroke}{rgb}{0.190631,0.407061,0.556089}%
\pgfsetstrokecolor{currentstroke}%
\pgfsetstrokeopacity{0.800000}%
\pgfsetdash{}{0pt}%
\pgfpathmoveto{\pgfqpoint{7.982250in}{2.250000in}}%
\pgfpathlineto{\pgfqpoint{7.983323in}{2.292297in}}%
\pgfpathlineto{\pgfqpoint{7.986551in}{2.334922in}}%
\pgfpathlineto{\pgfqpoint{7.991958in}{2.378204in}}%
\pgfpathlineto{\pgfqpoint{7.999586in}{2.422478in}}%
\pgfpathlineto{\pgfqpoint{8.009495in}{2.468086in}}%
\pgfpathlineto{\pgfqpoint{8.023710in}{2.522298in}}%
\pgfpathlineto{\pgfqpoint{8.041148in}{2.579262in}}%
\pgfpathlineto{\pgfqpoint{8.061985in}{2.639555in}}%
\pgfpathlineto{\pgfqpoint{8.086432in}{2.703786in}}%
\pgfpathlineto{\pgfqpoint{8.118559in}{2.781564in}}%
\pgfpathlineto{\pgfqpoint{8.155979in}{2.866144in}}%
\pgfpathlineto{\pgfqpoint{8.204351in}{2.969420in}}%
\pgfpathlineto{\pgfqpoint{8.260615in}{3.084066in}}%
\pgfpathlineto{\pgfqpoint{8.332687in}{3.225475in}}%
\pgfpathlineto{\pgfqpoint{8.416767in}{3.385543in}}%
\pgfpathlineto{\pgfqpoint{8.524081in}{3.585050in}}%
\pgfpathlineto{\pgfqpoint{8.661473in}{3.835646in}}%
\pgfpathlineto{\pgfqpoint{8.710415in}{3.924020in}}%
\pgfpathlineto{\pgfqpoint{8.710415in}{3.924020in}}%
\pgfusepath{stroke}%
\end{pgfscope}%
\begin{pgfscope}%
\pgfpathrectangle{\pgfqpoint{5.486220in}{0.585980in}}{\pgfqpoint{4.437388in}{3.328041in}}%
\pgfusepath{clip}%
\pgfsetrectcap%
\pgfsetroundjoin%
\pgfsetlinewidth{1.505625pt}%
\definecolor{currentstroke}{rgb}{0.157729,0.485932,0.558013}%
\pgfsetstrokecolor{currentstroke}%
\pgfsetstrokeopacity{0.800000}%
\pgfsetdash{}{0pt}%
\pgfpathmoveto{\pgfqpoint{8.185275in}{2.250000in}}%
\pgfpathlineto{\pgfqpoint{8.186223in}{2.267432in}}%
\pgfpathlineto{\pgfqpoint{8.189071in}{2.284933in}}%
\pgfpathlineto{\pgfqpoint{8.195012in}{2.306122in}}%
\pgfpathlineto{\pgfqpoint{8.203739in}{2.327630in}}%
\pgfpathlineto{\pgfqpoint{8.215301in}{2.349580in}}%
\pgfpathlineto{\pgfqpoint{8.229765in}{2.372096in}}%
\pgfpathlineto{\pgfqpoint{8.247212in}{2.395305in}}%
\pgfpathlineto{\pgfqpoint{8.271471in}{2.423436in}}%
\pgfpathlineto{\pgfqpoint{8.300115in}{2.452910in}}%
\pgfpathlineto{\pgfqpoint{8.333364in}{2.483953in}}%
\pgfpathlineto{\pgfqpoint{8.377335in}{2.521663in}}%
\pgfpathlineto{\pgfqpoint{8.428103in}{2.562120in}}%
\pgfpathlineto{\pgfqpoint{8.493985in}{2.611426in}}%
\pgfpathlineto{\pgfqpoint{8.569964in}{2.665358in}}%
\pgfpathlineto{\pgfqpoint{8.667417in}{2.731548in}}%
\pgfpathlineto{\pgfqpoint{8.780080in}{2.805349in}}%
\pgfpathlineto{\pgfqpoint{8.923711in}{2.896715in}}%
\pgfpathlineto{\pgfqpoint{9.107093in}{3.010561in}}%
\pgfpathlineto{\pgfqpoint{9.341934in}{3.153528in}}%
\pgfpathlineto{\pgfqpoint{9.643983in}{3.334627in}}%
\pgfpathlineto{\pgfqpoint{9.933608in}{3.506493in}}%
\pgfpathlineto{\pgfqpoint{9.933608in}{3.506493in}}%
\pgfusepath{stroke}%
\end{pgfscope}%
\begin{pgfscope}%
\pgfpathrectangle{\pgfqpoint{5.486220in}{0.585980in}}{\pgfqpoint{4.437388in}{3.328041in}}%
\pgfusepath{clip}%
\pgfsetrectcap%
\pgfsetroundjoin%
\pgfsetlinewidth{1.505625pt}%
\definecolor{currentstroke}{rgb}{0.127568,0.566949,0.550556}%
\pgfsetstrokecolor{currentstroke}%
\pgfsetstrokeopacity{0.800000}%
\pgfsetdash{}{0pt}%
\pgfpathmoveto{\pgfqpoint{8.259587in}{2.250000in}}%
\pgfpathlineto{\pgfqpoint{9.933608in}{2.250000in}}%
\pgfpathlineto{\pgfqpoint{9.933608in}{2.250000in}}%
\pgfusepath{stroke}%
\end{pgfscope}%
\begin{pgfscope}%
\pgfpathrectangle{\pgfqpoint{5.486220in}{0.585980in}}{\pgfqpoint{4.437388in}{3.328041in}}%
\pgfusepath{clip}%
\pgfsetrectcap%
\pgfsetroundjoin%
\pgfsetlinewidth{1.505625pt}%
\definecolor{currentstroke}{rgb}{0.126326,0.644107,0.525311}%
\pgfsetstrokecolor{currentstroke}%
\pgfsetstrokeopacity{0.800000}%
\pgfsetdash{}{0pt}%
\pgfpathmoveto{\pgfqpoint{8.185275in}{2.250000in}}%
\pgfpathlineto{\pgfqpoint{8.186223in}{2.232568in}}%
\pgfpathlineto{\pgfqpoint{8.189071in}{2.215067in}}%
\pgfpathlineto{\pgfqpoint{8.195012in}{2.193878in}}%
\pgfpathlineto{\pgfqpoint{8.203739in}{2.172370in}}%
\pgfpathlineto{\pgfqpoint{8.215301in}{2.150420in}}%
\pgfpathlineto{\pgfqpoint{8.229765in}{2.127904in}}%
\pgfpathlineto{\pgfqpoint{8.247212in}{2.104695in}}%
\pgfpathlineto{\pgfqpoint{8.271471in}{2.076564in}}%
\pgfpathlineto{\pgfqpoint{8.300115in}{2.047090in}}%
\pgfpathlineto{\pgfqpoint{8.333364in}{2.016047in}}%
\pgfpathlineto{\pgfqpoint{8.377335in}{1.978337in}}%
\pgfpathlineto{\pgfqpoint{8.428103in}{1.937880in}}%
\pgfpathlineto{\pgfqpoint{8.493985in}{1.888574in}}%
\pgfpathlineto{\pgfqpoint{8.569964in}{1.834642in}}%
\pgfpathlineto{\pgfqpoint{8.667417in}{1.768452in}}%
\pgfpathlineto{\pgfqpoint{8.780080in}{1.694651in}}%
\pgfpathlineto{\pgfqpoint{8.923711in}{1.603285in}}%
\pgfpathlineto{\pgfqpoint{9.107093in}{1.489439in}}%
\pgfpathlineto{\pgfqpoint{9.341934in}{1.346472in}}%
\pgfpathlineto{\pgfqpoint{9.643983in}{1.165373in}}%
\pgfpathlineto{\pgfqpoint{9.933608in}{0.993507in}}%
\pgfpathlineto{\pgfqpoint{9.933608in}{0.993507in}}%
\pgfusepath{stroke}%
\end{pgfscope}%
\begin{pgfscope}%
\pgfpathrectangle{\pgfqpoint{5.486220in}{0.585980in}}{\pgfqpoint{4.437388in}{3.328041in}}%
\pgfusepath{clip}%
\pgfsetrectcap%
\pgfsetroundjoin%
\pgfsetlinewidth{1.505625pt}%
\definecolor{currentstroke}{rgb}{0.208030,0.718701,0.472873}%
\pgfsetstrokecolor{currentstroke}%
\pgfsetstrokeopacity{0.800000}%
\pgfsetdash{}{0pt}%
\pgfpathmoveto{\pgfqpoint{7.982250in}{2.250000in}}%
\pgfpathlineto{\pgfqpoint{7.983323in}{2.207703in}}%
\pgfpathlineto{\pgfqpoint{7.986551in}{2.165078in}}%
\pgfpathlineto{\pgfqpoint{7.991958in}{2.121796in}}%
\pgfpathlineto{\pgfqpoint{7.999586in}{2.077522in}}%
\pgfpathlineto{\pgfqpoint{8.009495in}{2.031914in}}%
\pgfpathlineto{\pgfqpoint{8.023710in}{1.977702in}}%
\pgfpathlineto{\pgfqpoint{8.041148in}{1.920738in}}%
\pgfpathlineto{\pgfqpoint{8.061985in}{1.860445in}}%
\pgfpathlineto{\pgfqpoint{8.086432in}{1.796214in}}%
\pgfpathlineto{\pgfqpoint{8.118559in}{1.718436in}}%
\pgfpathlineto{\pgfqpoint{8.155979in}{1.633856in}}%
\pgfpathlineto{\pgfqpoint{8.204351in}{1.530580in}}%
\pgfpathlineto{\pgfqpoint{8.260615in}{1.415934in}}%
\pgfpathlineto{\pgfqpoint{8.332687in}{1.274525in}}%
\pgfpathlineto{\pgfqpoint{8.416767in}{1.114457in}}%
\pgfpathlineto{\pgfqpoint{8.524081in}{0.914950in}}%
\pgfpathlineto{\pgfqpoint{8.661473in}{0.664354in}}%
\pgfpathlineto{\pgfqpoint{8.710415in}{0.575980in}}%
\pgfpathlineto{\pgfqpoint{8.710415in}{0.575980in}}%
\pgfusepath{stroke}%
\end{pgfscope}%
\begin{pgfscope}%
\pgfpathrectangle{\pgfqpoint{5.486220in}{0.585980in}}{\pgfqpoint{4.437388in}{3.328041in}}%
\pgfusepath{clip}%
\pgfsetrectcap%
\pgfsetroundjoin%
\pgfsetlinewidth{1.505625pt}%
\definecolor{currentstroke}{rgb}{0.369214,0.788888,0.382914}%
\pgfsetstrokecolor{currentstroke}%
\pgfsetstrokeopacity{0.800000}%
\pgfsetdash{}{0pt}%
\pgfpathmoveto{\pgfqpoint{7.704914in}{2.250000in}}%
\pgfpathlineto{\pgfqpoint{7.704914in}{0.575980in}}%
\pgfpathlineto{\pgfqpoint{7.704914in}{0.575980in}}%
\pgfusepath{stroke}%
\end{pgfscope}%
\begin{pgfscope}%
\pgfpathrectangle{\pgfqpoint{5.486220in}{0.585980in}}{\pgfqpoint{4.437388in}{3.328041in}}%
\pgfusepath{clip}%
\pgfsetrectcap%
\pgfsetroundjoin%
\pgfsetlinewidth{1.505625pt}%
\definecolor{currentstroke}{rgb}{0.565498,0.842430,0.262877}%
\pgfsetstrokecolor{currentstroke}%
\pgfsetstrokeopacity{0.800000}%
\pgfsetdash{}{0pt}%
\pgfpathmoveto{\pgfqpoint{7.427577in}{2.250000in}}%
\pgfpathlineto{\pgfqpoint{7.426504in}{2.207703in}}%
\pgfpathlineto{\pgfqpoint{7.423276in}{2.165078in}}%
\pgfpathlineto{\pgfqpoint{7.417869in}{2.121796in}}%
\pgfpathlineto{\pgfqpoint{7.410241in}{2.077522in}}%
\pgfpathlineto{\pgfqpoint{7.400333in}{2.031914in}}%
\pgfpathlineto{\pgfqpoint{7.386117in}{1.977702in}}%
\pgfpathlineto{\pgfqpoint{7.368679in}{1.920738in}}%
\pgfpathlineto{\pgfqpoint{7.347842in}{1.860445in}}%
\pgfpathlineto{\pgfqpoint{7.323395in}{1.796214in}}%
\pgfpathlineto{\pgfqpoint{7.291268in}{1.718436in}}%
\pgfpathlineto{\pgfqpoint{7.253848in}{1.633856in}}%
\pgfpathlineto{\pgfqpoint{7.205476in}{1.530580in}}%
\pgfpathlineto{\pgfqpoint{7.149212in}{1.415934in}}%
\pgfpathlineto{\pgfqpoint{7.077140in}{1.274525in}}%
\pgfpathlineto{\pgfqpoint{6.993060in}{1.114457in}}%
\pgfpathlineto{\pgfqpoint{6.885746in}{0.914950in}}%
\pgfpathlineto{\pgfqpoint{6.748354in}{0.664354in}}%
\pgfpathlineto{\pgfqpoint{6.699413in}{0.575980in}}%
\pgfpathlineto{\pgfqpoint{6.699413in}{0.575980in}}%
\pgfusepath{stroke}%
\end{pgfscope}%
\begin{pgfscope}%
\pgfpathrectangle{\pgfqpoint{5.486220in}{0.585980in}}{\pgfqpoint{4.437388in}{3.328041in}}%
\pgfusepath{clip}%
\pgfsetrectcap%
\pgfsetroundjoin%
\pgfsetlinewidth{1.505625pt}%
\definecolor{currentstroke}{rgb}{0.783315,0.879285,0.125405}%
\pgfsetstrokecolor{currentstroke}%
\pgfsetstrokeopacity{0.800000}%
\pgfsetdash{}{0pt}%
\pgfpathmoveto{\pgfqpoint{7.224552in}{2.250000in}}%
\pgfpathlineto{\pgfqpoint{7.223604in}{2.232568in}}%
\pgfpathlineto{\pgfqpoint{7.220757in}{2.215067in}}%
\pgfpathlineto{\pgfqpoint{7.214816in}{2.193878in}}%
\pgfpathlineto{\pgfqpoint{7.206089in}{2.172370in}}%
\pgfpathlineto{\pgfqpoint{7.194526in}{2.150420in}}%
\pgfpathlineto{\pgfqpoint{7.180062in}{2.127904in}}%
\pgfpathlineto{\pgfqpoint{7.162615in}{2.104695in}}%
\pgfpathlineto{\pgfqpoint{7.138356in}{2.076564in}}%
\pgfpathlineto{\pgfqpoint{7.109713in}{2.047090in}}%
\pgfpathlineto{\pgfqpoint{7.076463in}{2.016047in}}%
\pgfpathlineto{\pgfqpoint{7.032493in}{1.978337in}}%
\pgfpathlineto{\pgfqpoint{6.981724in}{1.937880in}}%
\pgfpathlineto{\pgfqpoint{6.915843in}{1.888574in}}%
\pgfpathlineto{\pgfqpoint{6.839863in}{1.834642in}}%
\pgfpathlineto{\pgfqpoint{6.742410in}{1.768452in}}%
\pgfpathlineto{\pgfqpoint{6.629747in}{1.694651in}}%
\pgfpathlineto{\pgfqpoint{6.486116in}{1.603285in}}%
\pgfpathlineto{\pgfqpoint{6.302734in}{1.489439in}}%
\pgfpathlineto{\pgfqpoint{6.067893in}{1.346472in}}%
\pgfpathlineto{\pgfqpoint{5.765844in}{1.165373in}}%
\pgfpathlineto{\pgfqpoint{5.476220in}{0.993507in}}%
\pgfpathlineto{\pgfqpoint{5.476220in}{0.993507in}}%
\pgfusepath{stroke}%
\end{pgfscope}%
\begin{pgfscope}%
\pgfpathrectangle{\pgfqpoint{5.486220in}{0.585980in}}{\pgfqpoint{4.437388in}{3.328041in}}%
\pgfusepath{clip}%
\pgfsetrectcap%
\pgfsetroundjoin%
\pgfsetlinewidth{1.505625pt}%
\definecolor{currentstroke}{rgb}{0.993248,0.906157,0.143936}%
\pgfsetstrokecolor{currentstroke}%
\pgfsetstrokeopacity{0.800000}%
\pgfsetdash{}{0pt}%
\pgfpathmoveto{\pgfqpoint{7.150240in}{2.250000in}}%
\pgfpathlineto{\pgfqpoint{5.476220in}{2.250000in}}%
\pgfpathlineto{\pgfqpoint{5.476220in}{2.250000in}}%
\pgfusepath{stroke}%
\end{pgfscope}%
\begin{pgfscope}%
\pgfpathrectangle{\pgfqpoint{5.486220in}{0.585980in}}{\pgfqpoint{4.437388in}{3.328041in}}%
\pgfusepath{clip}%
\pgfsetbuttcap%
\pgfsetroundjoin%
\pgfsetlinewidth{1.505625pt}%
\definecolor{currentstroke}{rgb}{0.050383,0.029803,0.527975}%
\pgfsetstrokecolor{currentstroke}%
\pgfsetstrokeopacity{0.800000}%
\pgfsetdash{{5.550000pt}{2.400000pt}}{0.000000pt}%
\pgfpathmoveto{\pgfqpoint{7.150240in}{2.250000in}}%
\pgfpathlineto{\pgfqpoint{8.259518in}{2.250000in}}%
\pgfpathlineto{\pgfqpoint{7.150240in}{2.250000in}}%
\pgfusepath{stroke}%
\end{pgfscope}%
\begin{pgfscope}%
\pgfpathrectangle{\pgfqpoint{5.486220in}{0.585980in}}{\pgfqpoint{4.437388in}{3.328041in}}%
\pgfusepath{clip}%
\pgfsetbuttcap%
\pgfsetroundjoin%
\pgfsetlinewidth{1.505625pt}%
\definecolor{currentstroke}{rgb}{0.299855,0.009561,0.631624}%
\pgfsetstrokecolor{currentstroke}%
\pgfsetstrokeopacity{0.800000}%
\pgfsetdash{{5.550000pt}{2.400000pt}}{0.000000pt}%
\pgfpathmoveto{\pgfqpoint{7.122935in}{2.250000in}}%
\pgfpathlineto{\pgfqpoint{7.124095in}{2.261117in}}%
\pgfpathlineto{\pgfqpoint{7.127571in}{2.272190in}}%
\pgfpathlineto{\pgfqpoint{7.133347in}{2.283175in}}%
\pgfpathlineto{\pgfqpoint{7.141403in}{2.294027in}}%
\pgfpathlineto{\pgfqpoint{7.151704in}{2.304704in}}%
\pgfpathlineto{\pgfqpoint{7.164211in}{2.315163in}}%
\pgfpathlineto{\pgfqpoint{7.186994in}{2.330351in}}%
\pgfpathlineto{\pgfqpoint{7.214421in}{2.344819in}}%
\pgfpathlineto{\pgfqpoint{7.246246in}{2.358437in}}%
\pgfpathlineto{\pgfqpoint{7.282182in}{2.371083in}}%
\pgfpathlineto{\pgfqpoint{7.321909in}{2.382643in}}%
\pgfpathlineto{\pgfqpoint{7.365069in}{2.393014in}}%
\pgfpathlineto{\pgfqpoint{7.411276in}{2.402102in}}%
\pgfpathlineto{\pgfqpoint{7.460116in}{2.409828in}}%
\pgfpathlineto{\pgfqpoint{7.511151in}{2.416120in}}%
\pgfpathlineto{\pgfqpoint{7.581817in}{2.422185in}}%
\pgfpathlineto{\pgfqpoint{7.654445in}{2.425507in}}%
\pgfpathlineto{\pgfqpoint{7.727877in}{2.426033in}}%
\pgfpathlineto{\pgfqpoint{7.800943in}{2.423756in}}%
\pgfpathlineto{\pgfqpoint{7.872479in}{2.418710in}}%
\pgfpathlineto{\pgfqpoint{7.941346in}{2.410977in}}%
\pgfpathlineto{\pgfqpoint{7.990582in}{2.403487in}}%
\pgfpathlineto{\pgfqpoint{8.037258in}{2.394620in}}%
\pgfpathlineto{\pgfqpoint{8.080954in}{2.384457in}}%
\pgfpathlineto{\pgfqpoint{8.121278in}{2.373088in}}%
\pgfpathlineto{\pgfqpoint{8.157870in}{2.360615in}}%
\pgfpathlineto{\pgfqpoint{8.190400in}{2.347151in}}%
\pgfpathlineto{\pgfqpoint{8.218578in}{2.332816in}}%
\pgfpathlineto{\pgfqpoint{8.242151in}{2.317739in}}%
\pgfpathlineto{\pgfqpoint{8.255202in}{2.307341in}}%
\pgfpathlineto{\pgfqpoint{8.266059in}{2.296715in}}%
\pgfpathlineto{\pgfqpoint{8.274679in}{2.285902in}}%
\pgfpathlineto{\pgfqpoint{8.281028in}{2.274947in}}%
\pgfpathlineto{\pgfqpoint{8.285080in}{2.263891in}}%
\pgfpathlineto{\pgfqpoint{8.286819in}{2.252781in}}%
\pgfpathlineto{\pgfqpoint{8.286239in}{2.241660in}}%
\pgfpathlineto{\pgfqpoint{8.283342in}{2.230571in}}%
\pgfpathlineto{\pgfqpoint{8.278139in}{2.219560in}}%
\pgfpathlineto{\pgfqpoint{8.270651in}{2.208671in}}%
\pgfpathlineto{\pgfqpoint{8.260907in}{2.197946in}}%
\pgfpathlineto{\pgfqpoint{8.248947in}{2.187429in}}%
\pgfpathlineto{\pgfqpoint{8.234819in}{2.177161in}}%
\pgfpathlineto{\pgfqpoint{8.209686in}{2.162316in}}%
\pgfpathlineto{\pgfqpoint{8.180027in}{2.148258in}}%
\pgfpathlineto{\pgfqpoint{8.146108in}{2.135111in}}%
\pgfpathlineto{\pgfqpoint{8.108234in}{2.122995in}}%
\pgfpathlineto{\pgfqpoint{8.066744in}{2.112017in}}%
\pgfpathlineto{\pgfqpoint{8.022010in}{2.102276in}}%
\pgfpathlineto{\pgfqpoint{7.974433in}{2.093860in}}%
\pgfpathlineto{\pgfqpoint{7.924440in}{2.086843in}}%
\pgfpathlineto{\pgfqpoint{7.872479in}{2.081290in}}%
\pgfpathlineto{\pgfqpoint{7.800943in}{2.076244in}}%
\pgfpathlineto{\pgfqpoint{7.727877in}{2.073967in}}%
\pgfpathlineto{\pgfqpoint{7.654445in}{2.074493in}}%
\pgfpathlineto{\pgfqpoint{7.581817in}{2.077815in}}%
\pgfpathlineto{\pgfqpoint{7.511151in}{2.083880in}}%
\pgfpathlineto{\pgfqpoint{7.460116in}{2.090172in}}%
\pgfpathlineto{\pgfqpoint{7.411276in}{2.097898in}}%
\pgfpathlineto{\pgfqpoint{7.365069in}{2.106986in}}%
\pgfpathlineto{\pgfqpoint{7.321909in}{2.117357in}}%
\pgfpathlineto{\pgfqpoint{7.282182in}{2.128917in}}%
\pgfpathlineto{\pgfqpoint{7.246246in}{2.141563in}}%
\pgfpathlineto{\pgfqpoint{7.214421in}{2.155181in}}%
\pgfpathlineto{\pgfqpoint{7.186994in}{2.169649in}}%
\pgfpathlineto{\pgfqpoint{7.164211in}{2.184837in}}%
\pgfpathlineto{\pgfqpoint{7.151704in}{2.195296in}}%
\pgfpathlineto{\pgfqpoint{7.141403in}{2.205973in}}%
\pgfpathlineto{\pgfqpoint{7.133347in}{2.216825in}}%
\pgfpathlineto{\pgfqpoint{7.127571in}{2.227810in}}%
\pgfpathlineto{\pgfqpoint{7.124095in}{2.238883in}}%
\pgfpathlineto{\pgfqpoint{7.122935in}{2.250000in}}%
\pgfpathlineto{\pgfqpoint{7.122935in}{2.250000in}}%
\pgfusepath{stroke}%
\end{pgfscope}%
\begin{pgfscope}%
\pgfpathrectangle{\pgfqpoint{5.486220in}{0.585980in}}{\pgfqpoint{4.437388in}{3.328041in}}%
\pgfusepath{clip}%
\pgfsetbuttcap%
\pgfsetroundjoin%
\pgfsetlinewidth{1.505625pt}%
\definecolor{currentstroke}{rgb}{0.494877,0.011990,0.657865}%
\pgfsetstrokecolor{currentstroke}%
\pgfsetstrokeopacity{0.800000}%
\pgfsetdash{{5.550000pt}{2.400000pt}}{0.000000pt}%
\pgfpathmoveto{\pgfqpoint{7.038333in}{2.250000in}}%
\pgfpathlineto{\pgfqpoint{7.039661in}{2.273329in}}%
\pgfpathlineto{\pgfqpoint{7.043642in}{2.296565in}}%
\pgfpathlineto{\pgfqpoint{7.050258in}{2.319616in}}%
\pgfpathlineto{\pgfqpoint{7.059484in}{2.342389in}}%
\pgfpathlineto{\pgfqpoint{7.071283in}{2.364794in}}%
\pgfpathlineto{\pgfqpoint{7.085608in}{2.386741in}}%
\pgfpathlineto{\pgfqpoint{7.102402in}{2.408143in}}%
\pgfpathlineto{\pgfqpoint{7.121597in}{2.428915in}}%
\pgfpathlineto{\pgfqpoint{7.143118in}{2.448974in}}%
\pgfpathlineto{\pgfqpoint{7.166878in}{2.468239in}}%
\pgfpathlineto{\pgfqpoint{7.192783in}{2.486634in}}%
\pgfpathlineto{\pgfqpoint{7.220729in}{2.504086in}}%
\pgfpathlineto{\pgfqpoint{7.250606in}{2.520526in}}%
\pgfpathlineto{\pgfqpoint{7.282293in}{2.535886in}}%
\pgfpathlineto{\pgfqpoint{7.315666in}{2.550108in}}%
\pgfpathlineto{\pgfqpoint{7.350589in}{2.563133in}}%
\pgfpathlineto{\pgfqpoint{7.386926in}{2.574909in}}%
\pgfpathlineto{\pgfqpoint{7.424530in}{2.585391in}}%
\pgfpathlineto{\pgfqpoint{7.463251in}{2.594535in}}%
\pgfpathlineto{\pgfqpoint{7.502936in}{2.602306in}}%
\pgfpathlineto{\pgfqpoint{7.543426in}{2.608673in}}%
\pgfpathlineto{\pgfqpoint{7.584560in}{2.613610in}}%
\pgfpathlineto{\pgfqpoint{7.626174in}{2.617097in}}%
\pgfpathlineto{\pgfqpoint{7.689130in}{2.619582in}}%
\pgfpathlineto{\pgfqpoint{7.752228in}{2.618753in}}%
\pgfpathlineto{\pgfqpoint{7.814902in}{2.614618in}}%
\pgfpathlineto{\pgfqpoint{7.856172in}{2.610042in}}%
\pgfpathlineto{\pgfqpoint{7.896838in}{2.604031in}}%
\pgfpathlineto{\pgfqpoint{7.936739in}{2.596608in}}%
\pgfpathlineto{\pgfqpoint{7.975716in}{2.587804in}}%
\pgfpathlineto{\pgfqpoint{8.013614in}{2.577653in}}%
\pgfpathlineto{\pgfqpoint{8.050281in}{2.566196in}}%
\pgfpathlineto{\pgfqpoint{8.085571in}{2.553478in}}%
\pgfpathlineto{\pgfqpoint{8.119344in}{2.539551in}}%
\pgfpathlineto{\pgfqpoint{8.151464in}{2.524469in}}%
\pgfpathlineto{\pgfqpoint{8.181805in}{2.508294in}}%
\pgfpathlineto{\pgfqpoint{8.210245in}{2.491089in}}%
\pgfpathlineto{\pgfqpoint{8.236670in}{2.472922in}}%
\pgfpathlineto{\pgfqpoint{8.260976in}{2.453867in}}%
\pgfpathlineto{\pgfqpoint{8.283065in}{2.434000in}}%
\pgfpathlineto{\pgfqpoint{8.302849in}{2.413399in}}%
\pgfpathlineto{\pgfqpoint{8.320250in}{2.392146in}}%
\pgfpathlineto{\pgfqpoint{8.335197in}{2.370327in}}%
\pgfpathlineto{\pgfqpoint{8.347633in}{2.348028in}}%
\pgfpathlineto{\pgfqpoint{8.357506in}{2.325339in}}%
\pgfpathlineto{\pgfqpoint{8.364778in}{2.302349in}}%
\pgfpathlineto{\pgfqpoint{8.369419in}{2.279151in}}%
\pgfpathlineto{\pgfqpoint{8.371412in}{2.255836in}}%
\pgfpathlineto{\pgfqpoint{8.370747in}{2.232498in}}%
\pgfpathlineto{\pgfqpoint{8.367429in}{2.209230in}}%
\pgfpathlineto{\pgfqpoint{8.361469in}{2.186124in}}%
\pgfpathlineto{\pgfqpoint{8.352892in}{2.163273in}}%
\pgfpathlineto{\pgfqpoint{8.341732in}{2.140768in}}%
\pgfpathlineto{\pgfqpoint{8.328034in}{2.118698in}}%
\pgfpathlineto{\pgfqpoint{8.311852in}{2.097151in}}%
\pgfpathlineto{\pgfqpoint{8.293250in}{2.076214in}}%
\pgfpathlineto{\pgfqpoint{8.272303in}{2.055970in}}%
\pgfpathlineto{\pgfqpoint{8.249094in}{2.036499in}}%
\pgfpathlineto{\pgfqpoint{8.223716in}{2.017879in}}%
\pgfpathlineto{\pgfqpoint{8.196270in}{2.000184in}}%
\pgfpathlineto{\pgfqpoint{8.166865in}{1.983486in}}%
\pgfpathlineto{\pgfqpoint{8.135619in}{1.967849in}}%
\pgfpathlineto{\pgfqpoint{8.102655in}{1.953338in}}%
\pgfpathlineto{\pgfqpoint{8.068107in}{1.940009in}}%
\pgfpathlineto{\pgfqpoint{8.032110in}{1.927915in}}%
\pgfpathlineto{\pgfqpoint{7.994809in}{1.917106in}}%
\pgfpathlineto{\pgfqpoint{7.956353in}{1.907624in}}%
\pgfpathlineto{\pgfqpoint{7.916894in}{1.899506in}}%
\pgfpathlineto{\pgfqpoint{7.876590in}{1.892786in}}%
\pgfpathlineto{\pgfqpoint{7.835602in}{1.887489in}}%
\pgfpathlineto{\pgfqpoint{7.794093in}{1.883638in}}%
\pgfpathlineto{\pgfqpoint{7.731215in}{1.880602in}}%
\pgfpathlineto{\pgfqpoint{7.668101in}{1.880879in}}%
\pgfpathlineto{\pgfqpoint{7.605317in}{1.884464in}}%
\pgfpathlineto{\pgfqpoint{7.563923in}{1.888679in}}%
\pgfpathlineto{\pgfqpoint{7.523091in}{1.894333in}}%
\pgfpathlineto{\pgfqpoint{7.482983in}{1.901406in}}%
\pgfpathlineto{\pgfqpoint{7.443760in}{1.909868in}}%
\pgfpathlineto{\pgfqpoint{7.405579in}{1.919685in}}%
\pgfpathlineto{\pgfqpoint{7.368590in}{1.930820in}}%
\pgfpathlineto{\pgfqpoint{7.332942in}{1.943227in}}%
\pgfpathlineto{\pgfqpoint{7.298777in}{1.956857in}}%
\pgfpathlineto{\pgfqpoint{7.266231in}{1.971655in}}%
\pgfpathlineto{\pgfqpoint{7.235434in}{1.987563in}}%
\pgfpathlineto{\pgfqpoint{7.206508in}{2.004517in}}%
\pgfpathlineto{\pgfqpoint{7.179569in}{2.022450in}}%
\pgfpathlineto{\pgfqpoint{7.154724in}{2.041290in}}%
\pgfpathlineto{\pgfqpoint{7.132072in}{2.060961in}}%
\pgfpathlineto{\pgfqpoint{7.111704in}{2.081387in}}%
\pgfpathlineto{\pgfqpoint{7.093701in}{2.102484in}}%
\pgfpathlineto{\pgfqpoint{7.078134in}{2.124170in}}%
\pgfpathlineto{\pgfqpoint{7.065065in}{2.146357in}}%
\pgfpathlineto{\pgfqpoint{7.054547in}{2.168957in}}%
\pgfpathlineto{\pgfqpoint{7.046622in}{2.191880in}}%
\pgfpathlineto{\pgfqpoint{7.041321in}{2.215035in}}%
\pgfpathlineto{\pgfqpoint{7.038665in}{2.238330in}}%
\pgfpathlineto{\pgfqpoint{7.038333in}{2.250000in}}%
\pgfpathlineto{\pgfqpoint{7.038333in}{2.250000in}}%
\pgfusepath{stroke}%
\end{pgfscope}%
\begin{pgfscope}%
\pgfpathrectangle{\pgfqpoint{5.486220in}{0.585980in}}{\pgfqpoint{4.437388in}{3.328041in}}%
\pgfusepath{clip}%
\pgfsetbuttcap%
\pgfsetroundjoin%
\pgfsetlinewidth{1.505625pt}%
\definecolor{currentstroke}{rgb}{0.665129,0.138566,0.585582}%
\pgfsetstrokecolor{currentstroke}%
\pgfsetstrokeopacity{0.800000}%
\pgfsetdash{{5.550000pt}{2.400000pt}}{0.000000pt}%
\pgfpathmoveto{\pgfqpoint{6.888103in}{2.250000in}}%
\pgfpathlineto{\pgfqpoint{6.889731in}{2.287838in}}%
\pgfpathlineto{\pgfqpoint{6.894608in}{2.325525in}}%
\pgfpathlineto{\pgfqpoint{6.902716in}{2.362911in}}%
\pgfpathlineto{\pgfqpoint{6.914022in}{2.399847in}}%
\pgfpathlineto{\pgfqpoint{6.928480in}{2.436186in}}%
\pgfpathlineto{\pgfqpoint{6.946033in}{2.471782in}}%
\pgfpathlineto{\pgfqpoint{6.966612in}{2.506495in}}%
\pgfpathlineto{\pgfqpoint{6.990133in}{2.540185in}}%
\pgfpathlineto{\pgfqpoint{7.016504in}{2.572718in}}%
\pgfpathlineto{\pgfqpoint{7.045619in}{2.603964in}}%
\pgfpathlineto{\pgfqpoint{7.077362in}{2.633800in}}%
\pgfpathlineto{\pgfqpoint{7.111607in}{2.662106in}}%
\pgfpathlineto{\pgfqpoint{7.148217in}{2.688769in}}%
\pgfpathlineto{\pgfqpoint{7.187046in}{2.713683in}}%
\pgfpathlineto{\pgfqpoint{7.227939in}{2.736748in}}%
\pgfpathlineto{\pgfqpoint{7.270734in}{2.757874in}}%
\pgfpathlineto{\pgfqpoint{7.315260in}{2.776974in}}%
\pgfpathlineto{\pgfqpoint{7.361338in}{2.793974in}}%
\pgfpathlineto{\pgfqpoint{7.408787in}{2.808806in}}%
\pgfpathlineto{\pgfqpoint{7.457416in}{2.821410in}}%
\pgfpathlineto{\pgfqpoint{7.507031in}{2.831736in}}%
\pgfpathlineto{\pgfqpoint{7.557435in}{2.839743in}}%
\pgfpathlineto{\pgfqpoint{7.608428in}{2.845400in}}%
\pgfpathlineto{\pgfqpoint{7.659804in}{2.848683in}}%
\pgfpathlineto{\pgfqpoint{7.711361in}{2.849579in}}%
\pgfpathlineto{\pgfqpoint{7.762892in}{2.848085in}}%
\pgfpathlineto{\pgfqpoint{7.814192in}{2.844207in}}%
\pgfpathlineto{\pgfqpoint{7.865056in}{2.837961in}}%
\pgfpathlineto{\pgfqpoint{7.915282in}{2.829370in}}%
\pgfpathlineto{\pgfqpoint{7.964669in}{2.818471in}}%
\pgfpathlineto{\pgfqpoint{8.013021in}{2.805305in}}%
\pgfpathlineto{\pgfqpoint{8.060144in}{2.789925in}}%
\pgfpathlineto{\pgfqpoint{8.105852in}{2.772393in}}%
\pgfpathlineto{\pgfqpoint{8.149961in}{2.752779in}}%
\pgfpathlineto{\pgfqpoint{8.192296in}{2.731160in}}%
\pgfpathlineto{\pgfqpoint{8.232688in}{2.707624in}}%
\pgfpathlineto{\pgfqpoint{8.270977in}{2.682263in}}%
\pgfpathlineto{\pgfqpoint{8.307009in}{2.655179in}}%
\pgfpathlineto{\pgfqpoint{8.340641in}{2.626480in}}%
\pgfpathlineto{\pgfqpoint{8.371738in}{2.596280in}}%
\pgfpathlineto{\pgfqpoint{8.400178in}{2.564700in}}%
\pgfpathlineto{\pgfqpoint{8.425846in}{2.531865in}}%
\pgfpathlineto{\pgfqpoint{8.448640in}{2.497907in}}%
\pgfpathlineto{\pgfqpoint{8.468469in}{2.462961in}}%
\pgfpathlineto{\pgfqpoint{8.485255in}{2.427165in}}%
\pgfpathlineto{\pgfqpoint{8.498930in}{2.390663in}}%
\pgfpathlineto{\pgfqpoint{8.509439in}{2.353601in}}%
\pgfpathlineto{\pgfqpoint{8.516742in}{2.316126in}}%
\pgfpathlineto{\pgfqpoint{8.520809in}{2.278387in}}%
\pgfpathlineto{\pgfqpoint{8.521623in}{2.240535in}}%
\pgfpathlineto{\pgfqpoint{8.519181in}{2.202720in}}%
\pgfpathlineto{\pgfqpoint{8.513494in}{2.165094in}}%
\pgfpathlineto{\pgfqpoint{8.504583in}{2.127807in}}%
\pgfpathlineto{\pgfqpoint{8.492485in}{2.091006in}}%
\pgfpathlineto{\pgfqpoint{8.477247in}{2.054840in}}%
\pgfpathlineto{\pgfqpoint{8.458930in}{2.019451in}}%
\pgfpathlineto{\pgfqpoint{8.437608in}{1.984982in}}%
\pgfpathlineto{\pgfqpoint{8.413365in}{1.951569in}}%
\pgfpathlineto{\pgfqpoint{8.386298in}{1.919345in}}%
\pgfpathlineto{\pgfqpoint{8.356514in}{1.888440in}}%
\pgfpathlineto{\pgfqpoint{8.324133in}{1.858975in}}%
\pgfpathlineto{\pgfqpoint{8.289284in}{1.831070in}}%
\pgfpathlineto{\pgfqpoint{8.252105in}{1.804835in}}%
\pgfpathlineto{\pgfqpoint{8.212745in}{1.780374in}}%
\pgfpathlineto{\pgfqpoint{8.171361in}{1.757785in}}%
\pgfpathlineto{\pgfqpoint{8.128117in}{1.737158in}}%
\pgfpathlineto{\pgfqpoint{8.083186in}{1.718576in}}%
\pgfpathlineto{\pgfqpoint{8.036748in}{1.702112in}}%
\pgfpathlineto{\pgfqpoint{7.988986in}{1.687832in}}%
\pgfpathlineto{\pgfqpoint{7.940093in}{1.675793in}}%
\pgfpathlineto{\pgfqpoint{7.890261in}{1.666043in}}%
\pgfpathlineto{\pgfqpoint{7.839691in}{1.658621in}}%
\pgfpathlineto{\pgfqpoint{7.788584in}{1.653556in}}%
\pgfpathlineto{\pgfqpoint{7.737143in}{1.650869in}}%
\pgfpathlineto{\pgfqpoint{7.685573in}{1.650571in}}%
\pgfpathlineto{\pgfqpoint{7.634081in}{1.652661in}}%
\pgfpathlineto{\pgfqpoint{7.582871in}{1.657133in}}%
\pgfpathlineto{\pgfqpoint{7.532147in}{1.663968in}}%
\pgfpathlineto{\pgfqpoint{7.482112in}{1.673140in}}%
\pgfpathlineto{\pgfqpoint{7.432966in}{1.684610in}}%
\pgfpathlineto{\pgfqpoint{7.384903in}{1.698335in}}%
\pgfpathlineto{\pgfqpoint{7.338116in}{1.714259in}}%
\pgfpathlineto{\pgfqpoint{7.292791in}{1.732318in}}%
\pgfpathlineto{\pgfqpoint{7.249110in}{1.752441in}}%
\pgfpathlineto{\pgfqpoint{7.207245in}{1.774547in}}%
\pgfpathlineto{\pgfqpoint{7.167363in}{1.798549in}}%
\pgfpathlineto{\pgfqpoint{7.129625in}{1.824350in}}%
\pgfpathlineto{\pgfqpoint{7.094180in}{1.851849in}}%
\pgfpathlineto{\pgfqpoint{7.061170in}{1.880934in}}%
\pgfpathlineto{\pgfqpoint{7.030725in}{1.911490in}}%
\pgfpathlineto{\pgfqpoint{7.002969in}{1.943396in}}%
\pgfpathlineto{\pgfqpoint{6.978010in}{1.976524in}}%
\pgfpathlineto{\pgfqpoint{6.955949in}{2.010742in}}%
\pgfpathlineto{\pgfqpoint{6.936874in}{2.045914in}}%
\pgfpathlineto{\pgfqpoint{6.920860in}{2.081900in}}%
\pgfpathlineto{\pgfqpoint{6.907972in}{2.118555in}}%
\pgfpathlineto{\pgfqpoint{6.898260in}{2.155735in}}%
\pgfpathlineto{\pgfqpoint{6.891764in}{2.193290in}}%
\pgfpathlineto{\pgfqpoint{6.888510in}{2.231072in}}%
\pgfpathlineto{\pgfqpoint{6.888103in}{2.250000in}}%
\pgfpathlineto{\pgfqpoint{6.888103in}{2.250000in}}%
\pgfusepath{stroke}%
\end{pgfscope}%
\begin{pgfscope}%
\pgfpathrectangle{\pgfqpoint{5.486220in}{0.585980in}}{\pgfqpoint{4.437388in}{3.328041in}}%
\pgfusepath{clip}%
\pgfsetbuttcap%
\pgfsetroundjoin%
\pgfsetlinewidth{1.505625pt}%
\definecolor{currentstroke}{rgb}{0.798216,0.280197,0.469538}%
\pgfsetstrokecolor{currentstroke}%
\pgfsetstrokeopacity{0.800000}%
\pgfsetdash{{5.550000pt}{2.400000pt}}{0.000000pt}%
\pgfpathmoveto{\pgfqpoint{6.657455in}{2.250000in}}%
\pgfpathlineto{\pgfqpoint{6.659543in}{2.306072in}}%
\pgfpathlineto{\pgfqpoint{6.665798in}{2.361920in}}%
\pgfpathlineto{\pgfqpoint{6.676195in}{2.417323in}}%
\pgfpathlineto{\pgfqpoint{6.690693in}{2.472058in}}%
\pgfpathlineto{\pgfqpoint{6.709233in}{2.525908in}}%
\pgfpathlineto{\pgfqpoint{6.731743in}{2.578659in}}%
\pgfpathlineto{\pgfqpoint{6.758133in}{2.630099in}}%
\pgfpathlineto{\pgfqpoint{6.788296in}{2.680024in}}%
\pgfpathlineto{\pgfqpoint{6.822113in}{2.728235in}}%
\pgfpathlineto{\pgfqpoint{6.859450in}{2.774539in}}%
\pgfpathlineto{\pgfqpoint{6.900157in}{2.818752in}}%
\pgfpathlineto{\pgfqpoint{6.944071in}{2.860699in}}%
\pgfpathlineto{\pgfqpoint{6.991019in}{2.900210in}}%
\pgfpathlineto{\pgfqpoint{7.040812in}{2.937130in}}%
\pgfpathlineto{\pgfqpoint{7.093253in}{2.971311in}}%
\pgfpathlineto{\pgfqpoint{7.148132in}{3.002616in}}%
\pgfpathlineto{\pgfqpoint{7.205231in}{3.030922in}}%
\pgfpathlineto{\pgfqpoint{7.264321in}{3.056114in}}%
\pgfpathlineto{\pgfqpoint{7.325168in}{3.078093in}}%
\pgfpathlineto{\pgfqpoint{7.387528in}{3.096771in}}%
\pgfpathlineto{\pgfqpoint{7.451154in}{3.112073in}}%
\pgfpathlineto{\pgfqpoint{7.515791in}{3.123939in}}%
\pgfpathlineto{\pgfqpoint{7.581182in}{3.132321in}}%
\pgfpathlineto{\pgfqpoint{7.647067in}{3.137186in}}%
\pgfpathlineto{\pgfqpoint{7.713182in}{3.138514in}}%
\pgfpathlineto{\pgfqpoint{7.779264in}{3.136301in}}%
\pgfpathlineto{\pgfqpoint{7.845049in}{3.130554in}}%
\pgfpathlineto{\pgfqpoint{7.910276in}{3.121298in}}%
\pgfpathlineto{\pgfqpoint{7.974685in}{3.108568in}}%
\pgfpathlineto{\pgfqpoint{8.038018in}{3.092415in}}%
\pgfpathlineto{\pgfqpoint{8.100023in}{3.072904in}}%
\pgfpathlineto{\pgfqpoint{8.160453in}{3.050113in}}%
\pgfpathlineto{\pgfqpoint{8.219067in}{3.024133in}}%
\pgfpathlineto{\pgfqpoint{8.275632in}{2.995067in}}%
\pgfpathlineto{\pgfqpoint{8.329921in}{2.963030in}}%
\pgfpathlineto{\pgfqpoint{8.381719in}{2.928152in}}%
\pgfpathlineto{\pgfqpoint{8.430820in}{2.890570in}}%
\pgfpathlineto{\pgfqpoint{8.477026in}{2.850434in}}%
\pgfpathlineto{\pgfqpoint{8.520155in}{2.807905in}}%
\pgfpathlineto{\pgfqpoint{8.560033in}{2.763152in}}%
\pgfpathlineto{\pgfqpoint{8.596504in}{2.716353in}}%
\pgfpathlineto{\pgfqpoint{8.629420in}{2.667696in}}%
\pgfpathlineto{\pgfqpoint{8.658650in}{2.617373in}}%
\pgfpathlineto{\pgfqpoint{8.684079in}{2.565586in}}%
\pgfpathlineto{\pgfqpoint{8.705604in}{2.512541in}}%
\pgfpathlineto{\pgfqpoint{8.723141in}{2.458449in}}%
\pgfpathlineto{\pgfqpoint{8.736618in}{2.403526in}}%
\pgfpathlineto{\pgfqpoint{8.745983in}{2.347992in}}%
\pgfpathlineto{\pgfqpoint{8.751198in}{2.292066in}}%
\pgfpathlineto{\pgfqpoint{8.752242in}{2.235973in}}%
\pgfpathlineto{\pgfqpoint{8.749111in}{2.179936in}}%
\pgfpathlineto{\pgfqpoint{8.741817in}{2.124178in}}%
\pgfpathlineto{\pgfqpoint{8.730391in}{2.068922in}}%
\pgfpathlineto{\pgfqpoint{8.714876in}{2.014388in}}%
\pgfpathlineto{\pgfqpoint{8.695335in}{1.960793in}}%
\pgfpathlineto{\pgfqpoint{8.671846in}{1.908350in}}%
\pgfpathlineto{\pgfqpoint{8.644503in}{1.857270in}}%
\pgfpathlineto{\pgfqpoint{8.613414in}{1.807755in}}%
\pgfpathlineto{\pgfqpoint{8.578704in}{1.760003in}}%
\pgfpathlineto{\pgfqpoint{8.540511in}{1.714205in}}%
\pgfpathlineto{\pgfqpoint{8.498986in}{1.670542in}}%
\pgfpathlineto{\pgfqpoint{8.454296in}{1.629189in}}%
\pgfpathlineto{\pgfqpoint{8.406619in}{1.590311in}}%
\pgfpathlineto{\pgfqpoint{8.356145in}{1.554062in}}%
\pgfpathlineto{\pgfqpoint{8.303075in}{1.520588in}}%
\pgfpathlineto{\pgfqpoint{8.247620in}{1.490021in}}%
\pgfpathlineto{\pgfqpoint{8.190002in}{1.462484in}}%
\pgfpathlineto{\pgfqpoint{8.130450in}{1.438086in}}%
\pgfpathlineto{\pgfqpoint{8.069202in}{1.416925in}}%
\pgfpathlineto{\pgfqpoint{8.006501in}{1.399085in}}%
\pgfpathlineto{\pgfqpoint{7.942599in}{1.384636in}}%
\pgfpathlineto{\pgfqpoint{7.877749in}{1.373637in}}%
\pgfpathlineto{\pgfqpoint{7.812210in}{1.366132in}}%
\pgfpathlineto{\pgfqpoint{7.746243in}{1.362150in}}%
\pgfpathlineto{\pgfqpoint{7.680112in}{1.361707in}}%
\pgfpathlineto{\pgfqpoint{7.614079in}{1.364805in}}%
\pgfpathlineto{\pgfqpoint{7.548409in}{1.371432in}}%
\pgfpathlineto{\pgfqpoint{7.483362in}{1.381561in}}%
\pgfpathlineto{\pgfqpoint{7.419199in}{1.395152in}}%
\pgfpathlineto{\pgfqpoint{7.356174in}{1.412151in}}%
\pgfpathlineto{\pgfqpoint{7.294540in}{1.432489in}}%
\pgfpathlineto{\pgfqpoint{7.234541in}{1.456086in}}%
\pgfpathlineto{\pgfqpoint{7.176418in}{1.482849in}}%
\pgfpathlineto{\pgfqpoint{7.120401in}{1.512669in}}%
\pgfpathlineto{\pgfqpoint{7.066715in}{1.545428in}}%
\pgfpathlineto{\pgfqpoint{7.015572in}{1.580996in}}%
\pgfpathlineto{\pgfqpoint{6.967177in}{1.619231in}}%
\pgfpathlineto{\pgfqpoint{6.921724in}{1.659980in}}%
\pgfpathlineto{\pgfqpoint{6.879392in}{1.703082in}}%
\pgfpathlineto{\pgfqpoint{6.840351in}{1.748363in}}%
\pgfpathlineto{\pgfqpoint{6.804756in}{1.795644in}}%
\pgfpathlineto{\pgfqpoint{6.772750in}{1.844737in}}%
\pgfpathlineto{\pgfqpoint{6.744459in}{1.895445in}}%
\pgfpathlineto{\pgfqpoint{6.719998in}{1.947566in}}%
\pgfpathlineto{\pgfqpoint{6.699462in}{2.000893in}}%
\pgfpathlineto{\pgfqpoint{6.682934in}{2.055212in}}%
\pgfpathlineto{\pgfqpoint{6.670481in}{2.110309in}}%
\pgfpathlineto{\pgfqpoint{6.662150in}{2.165962in}}%
\pgfpathlineto{\pgfqpoint{6.657977in}{2.221950in}}%
\pgfpathlineto{\pgfqpoint{6.657455in}{2.250000in}}%
\pgfpathlineto{\pgfqpoint{6.657455in}{2.250000in}}%
\pgfusepath{stroke}%
\end{pgfscope}%
\begin{pgfscope}%
\pgfpathrectangle{\pgfqpoint{5.486220in}{0.585980in}}{\pgfqpoint{4.437388in}{3.328041in}}%
\pgfusepath{clip}%
\pgfsetbuttcap%
\pgfsetroundjoin%
\pgfsetlinewidth{1.505625pt}%
\definecolor{currentstroke}{rgb}{0.901807,0.425087,0.359688}%
\pgfsetstrokecolor{currentstroke}%
\pgfsetstrokeopacity{0.800000}%
\pgfsetdash{{5.550000pt}{2.400000pt}}{0.000000pt}%
\pgfpathmoveto{\pgfqpoint{6.323681in}{2.250000in}}%
\pgfpathlineto{\pgfqpoint{6.324369in}{2.289933in}}%
\pgfpathlineto{\pgfqpoint{6.326434in}{2.329827in}}%
\pgfpathlineto{\pgfqpoint{6.329873in}{2.369640in}}%
\pgfpathlineto{\pgfqpoint{6.334682in}{2.409335in}}%
\pgfpathlineto{\pgfqpoint{6.340857in}{2.448871in}}%
\pgfpathlineto{\pgfqpoint{6.348392in}{2.488208in}}%
\pgfpathlineto{\pgfqpoint{6.357279in}{2.527308in}}%
\pgfpathlineto{\pgfqpoint{6.367510in}{2.566132in}}%
\pgfpathlineto{\pgfqpoint{6.379074in}{2.604640in}}%
\pgfpathlineto{\pgfqpoint{6.391959in}{2.642795in}}%
\pgfpathlineto{\pgfqpoint{6.406153in}{2.680558in}}%
\pgfpathlineto{\pgfqpoint{6.421642in}{2.717893in}}%
\pgfpathlineto{\pgfqpoint{6.438410in}{2.754760in}}%
\pgfpathlineto{\pgfqpoint{6.456440in}{2.791125in}}%
\pgfpathlineto{\pgfqpoint{6.475715in}{2.826950in}}%
\pgfpathlineto{\pgfqpoint{6.496215in}{2.862200in}}%
\pgfpathlineto{\pgfqpoint{6.517920in}{2.896840in}}%
\pgfpathlineto{\pgfqpoint{6.540808in}{2.930835in}}%
\pgfpathlineto{\pgfqpoint{6.564857in}{2.964152in}}%
\pgfpathlineto{\pgfqpoint{6.590042in}{2.996756in}}%
\pgfpathlineto{\pgfqpoint{6.616338in}{3.028616in}}%
\pgfpathlineto{\pgfqpoint{6.643720in}{3.059700in}}%
\pgfpathlineto{\pgfqpoint{6.672159in}{3.089977in}}%
\pgfpathlineto{\pgfqpoint{6.701628in}{3.119417in}}%
\pgfpathlineto{\pgfqpoint{6.732097in}{3.147990in}}%
\pgfpathlineto{\pgfqpoint{6.763536in}{3.175668in}}%
\pgfpathlineto{\pgfqpoint{6.795913in}{3.202423in}}%
\pgfpathlineto{\pgfqpoint{6.829196in}{3.228228in}}%
\pgfpathlineto{\pgfqpoint{6.863352in}{3.253059in}}%
\pgfpathlineto{\pgfqpoint{6.898347in}{3.276889in}}%
\pgfpathlineto{\pgfqpoint{6.934146in}{3.299697in}}%
\pgfpathlineto{\pgfqpoint{6.970713in}{3.321457in}}%
\pgfpathlineto{\pgfqpoint{7.008012in}{3.342150in}}%
\pgfpathlineto{\pgfqpoint{7.046006in}{3.361754in}}%
\pgfpathlineto{\pgfqpoint{7.084657in}{3.380250in}}%
\pgfpathlineto{\pgfqpoint{7.123926in}{3.397619in}}%
\pgfpathlineto{\pgfqpoint{7.163774in}{3.413844in}}%
\pgfpathlineto{\pgfqpoint{7.204161in}{3.428909in}}%
\pgfpathlineto{\pgfqpoint{7.245048in}{3.442798in}}%
\pgfpathlineto{\pgfqpoint{7.286393in}{3.455499in}}%
\pgfpathlineto{\pgfqpoint{7.328155in}{3.466998in}}%
\pgfpathlineto{\pgfqpoint{7.370293in}{3.477284in}}%
\pgfpathlineto{\pgfqpoint{7.412764in}{3.486347in}}%
\pgfpathlineto{\pgfqpoint{7.455527in}{3.494177in}}%
\pgfpathlineto{\pgfqpoint{7.498538in}{3.500767in}}%
\pgfpathlineto{\pgfqpoint{7.541755in}{3.506110in}}%
\pgfpathlineto{\pgfqpoint{7.585135in}{3.510201in}}%
\pgfpathlineto{\pgfqpoint{7.628634in}{3.513036in}}%
\pgfpathlineto{\pgfqpoint{7.672209in}{3.514612in}}%
\pgfpathlineto{\pgfqpoint{7.715816in}{3.514927in}}%
\pgfpathlineto{\pgfqpoint{7.759413in}{3.513982in}}%
\pgfpathlineto{\pgfqpoint{7.802955in}{3.511776in}}%
\pgfpathlineto{\pgfqpoint{7.846400in}{3.508313in}}%
\pgfpathlineto{\pgfqpoint{7.889704in}{3.503595in}}%
\pgfpathlineto{\pgfqpoint{7.932823in}{3.497628in}}%
\pgfpathlineto{\pgfqpoint{7.975715in}{3.490417in}}%
\pgfpathlineto{\pgfqpoint{8.018338in}{3.481969in}}%
\pgfpathlineto{\pgfqpoint{8.060647in}{3.472294in}}%
\pgfpathlineto{\pgfqpoint{8.102603in}{3.461400in}}%
\pgfpathlineto{\pgfqpoint{8.144162in}{3.449298in}}%
\pgfpathlineto{\pgfqpoint{8.185283in}{3.436001in}}%
\pgfpathlineto{\pgfqpoint{8.225925in}{3.421522in}}%
\pgfpathlineto{\pgfqpoint{8.266047in}{3.405875in}}%
\pgfpathlineto{\pgfqpoint{8.305611in}{3.389076in}}%
\pgfpathlineto{\pgfqpoint{8.344576in}{3.371141in}}%
\pgfpathlineto{\pgfqpoint{8.382903in}{3.352089in}}%
\pgfpathlineto{\pgfqpoint{8.420554in}{3.331938in}}%
\pgfpathlineto{\pgfqpoint{8.457492in}{3.310709in}}%
\pgfpathlineto{\pgfqpoint{8.493679in}{3.288422in}}%
\pgfpathlineto{\pgfqpoint{8.529081in}{3.265101in}}%
\pgfpathlineto{\pgfqpoint{8.563660in}{3.240767in}}%
\pgfpathlineto{\pgfqpoint{8.597384in}{3.215446in}}%
\pgfpathlineto{\pgfqpoint{8.630218in}{3.189162in}}%
\pgfpathlineto{\pgfqpoint{8.662130in}{3.161942in}}%
\pgfpathlineto{\pgfqpoint{8.693088in}{3.133813in}}%
\pgfpathlineto{\pgfqpoint{8.723060in}{3.104804in}}%
\pgfpathlineto{\pgfqpoint{8.752018in}{3.074942in}}%
\pgfpathlineto{\pgfqpoint{8.779932in}{3.044257in}}%
\pgfpathlineto{\pgfqpoint{8.806774in}{3.012781in}}%
\pgfpathlineto{\pgfqpoint{8.832518in}{2.980545in}}%
\pgfpathlineto{\pgfqpoint{8.857138in}{2.947580in}}%
\pgfpathlineto{\pgfqpoint{8.880610in}{2.913921in}}%
\pgfpathlineto{\pgfqpoint{8.902909in}{2.879599in}}%
\pgfpathlineto{\pgfqpoint{8.924014in}{2.844649in}}%
\pgfpathlineto{\pgfqpoint{8.943904in}{2.809107in}}%
\pgfpathlineto{\pgfqpoint{8.962559in}{2.773008in}}%
\pgfpathlineto{\pgfqpoint{8.979961in}{2.736387in}}%
\pgfpathlineto{\pgfqpoint{8.996091in}{2.699281in}}%
\pgfpathlineto{\pgfqpoint{9.010934in}{2.661728in}}%
\pgfpathlineto{\pgfqpoint{9.024475in}{2.623764in}}%
\pgfpathlineto{\pgfqpoint{9.036701in}{2.585428in}}%
\pgfpathlineto{\pgfqpoint{9.047600in}{2.546757in}}%
\pgfpathlineto{\pgfqpoint{9.057160in}{2.507790in}}%
\pgfpathlineto{\pgfqpoint{9.065372in}{2.468566in}}%
\pgfpathlineto{\pgfqpoint{9.072228in}{2.429125in}}%
\pgfpathlineto{\pgfqpoint{9.077721in}{2.389505in}}%
\pgfpathlineto{\pgfqpoint{9.081845in}{2.349746in}}%
\pgfpathlineto{\pgfqpoint{9.084597in}{2.309887in}}%
\pgfpathlineto{\pgfqpoint{9.085974in}{2.269969in}}%
\pgfpathlineto{\pgfqpoint{9.085974in}{2.230031in}}%
\pgfpathlineto{\pgfqpoint{9.084597in}{2.190113in}}%
\pgfpathlineto{\pgfqpoint{9.081845in}{2.150254in}}%
\pgfpathlineto{\pgfqpoint{9.077721in}{2.110495in}}%
\pgfpathlineto{\pgfqpoint{9.072228in}{2.070875in}}%
\pgfpathlineto{\pgfqpoint{9.065372in}{2.031434in}}%
\pgfpathlineto{\pgfqpoint{9.057160in}{1.992210in}}%
\pgfpathlineto{\pgfqpoint{9.047600in}{1.953243in}}%
\pgfpathlineto{\pgfqpoint{9.036701in}{1.914572in}}%
\pgfpathlineto{\pgfqpoint{9.024475in}{1.876236in}}%
\pgfpathlineto{\pgfqpoint{9.010934in}{1.838272in}}%
\pgfpathlineto{\pgfqpoint{8.996091in}{1.800719in}}%
\pgfpathlineto{\pgfqpoint{8.979961in}{1.763613in}}%
\pgfpathlineto{\pgfqpoint{8.962559in}{1.726992in}}%
\pgfpathlineto{\pgfqpoint{8.943904in}{1.690893in}}%
\pgfpathlineto{\pgfqpoint{8.924014in}{1.655351in}}%
\pgfpathlineto{\pgfqpoint{8.902909in}{1.620401in}}%
\pgfpathlineto{\pgfqpoint{8.880610in}{1.586079in}}%
\pgfpathlineto{\pgfqpoint{8.857138in}{1.552420in}}%
\pgfpathlineto{\pgfqpoint{8.832518in}{1.519455in}}%
\pgfpathlineto{\pgfqpoint{8.806774in}{1.487219in}}%
\pgfpathlineto{\pgfqpoint{8.779932in}{1.455743in}}%
\pgfpathlineto{\pgfqpoint{8.752018in}{1.425058in}}%
\pgfpathlineto{\pgfqpoint{8.723060in}{1.395196in}}%
\pgfpathlineto{\pgfqpoint{8.693088in}{1.366187in}}%
\pgfpathlineto{\pgfqpoint{8.662130in}{1.338058in}}%
\pgfpathlineto{\pgfqpoint{8.630218in}{1.310838in}}%
\pgfpathlineto{\pgfqpoint{8.597384in}{1.284554in}}%
\pgfpathlineto{\pgfqpoint{8.563660in}{1.259233in}}%
\pgfpathlineto{\pgfqpoint{8.529081in}{1.234899in}}%
\pgfpathlineto{\pgfqpoint{8.493679in}{1.211578in}}%
\pgfpathlineto{\pgfqpoint{8.457492in}{1.189291in}}%
\pgfpathlineto{\pgfqpoint{8.420554in}{1.168062in}}%
\pgfpathlineto{\pgfqpoint{8.382903in}{1.147911in}}%
\pgfpathlineto{\pgfqpoint{8.344576in}{1.128859in}}%
\pgfpathlineto{\pgfqpoint{8.305611in}{1.110924in}}%
\pgfpathlineto{\pgfqpoint{8.266047in}{1.094125in}}%
\pgfpathlineto{\pgfqpoint{8.225925in}{1.078478in}}%
\pgfpathlineto{\pgfqpoint{8.185283in}{1.063999in}}%
\pgfpathlineto{\pgfqpoint{8.144162in}{1.050702in}}%
\pgfpathlineto{\pgfqpoint{8.102603in}{1.038600in}}%
\pgfpathlineto{\pgfqpoint{8.060647in}{1.027706in}}%
\pgfpathlineto{\pgfqpoint{8.018338in}{1.018031in}}%
\pgfpathlineto{\pgfqpoint{7.975715in}{1.009583in}}%
\pgfpathlineto{\pgfqpoint{7.932823in}{1.002372in}}%
\pgfpathlineto{\pgfqpoint{7.889704in}{0.996405in}}%
\pgfpathlineto{\pgfqpoint{7.846400in}{0.991687in}}%
\pgfpathlineto{\pgfqpoint{7.802955in}{0.988224in}}%
\pgfpathlineto{\pgfqpoint{7.759413in}{0.986018in}}%
\pgfpathlineto{\pgfqpoint{7.715816in}{0.985073in}}%
\pgfpathlineto{\pgfqpoint{7.672209in}{0.985388in}}%
\pgfpathlineto{\pgfqpoint{7.628634in}{0.986964in}}%
\pgfpathlineto{\pgfqpoint{7.585135in}{0.989799in}}%
\pgfpathlineto{\pgfqpoint{7.541755in}{0.993890in}}%
\pgfpathlineto{\pgfqpoint{7.498538in}{0.999233in}}%
\pgfpathlineto{\pgfqpoint{7.455527in}{1.005823in}}%
\pgfpathlineto{\pgfqpoint{7.412764in}{1.013653in}}%
\pgfpathlineto{\pgfqpoint{7.370293in}{1.022716in}}%
\pgfpathlineto{\pgfqpoint{7.328155in}{1.033002in}}%
\pgfpathlineto{\pgfqpoint{7.286393in}{1.044501in}}%
\pgfpathlineto{\pgfqpoint{7.245048in}{1.057202in}}%
\pgfpathlineto{\pgfqpoint{7.204161in}{1.071091in}}%
\pgfpathlineto{\pgfqpoint{7.163774in}{1.086156in}}%
\pgfpathlineto{\pgfqpoint{7.123926in}{1.102381in}}%
\pgfpathlineto{\pgfqpoint{7.084657in}{1.119750in}}%
\pgfpathlineto{\pgfqpoint{7.046006in}{1.138246in}}%
\pgfpathlineto{\pgfqpoint{7.008012in}{1.157850in}}%
\pgfpathlineto{\pgfqpoint{6.970713in}{1.178543in}}%
\pgfpathlineto{\pgfqpoint{6.934146in}{1.200303in}}%
\pgfpathlineto{\pgfqpoint{6.898347in}{1.223111in}}%
\pgfpathlineto{\pgfqpoint{6.863352in}{1.246941in}}%
\pgfpathlineto{\pgfqpoint{6.829196in}{1.271772in}}%
\pgfpathlineto{\pgfqpoint{6.795913in}{1.297577in}}%
\pgfpathlineto{\pgfqpoint{6.763536in}{1.324332in}}%
\pgfpathlineto{\pgfqpoint{6.732097in}{1.352010in}}%
\pgfpathlineto{\pgfqpoint{6.701628in}{1.380583in}}%
\pgfpathlineto{\pgfqpoint{6.672159in}{1.410023in}}%
\pgfpathlineto{\pgfqpoint{6.643720in}{1.440300in}}%
\pgfpathlineto{\pgfqpoint{6.616338in}{1.471384in}}%
\pgfpathlineto{\pgfqpoint{6.590042in}{1.503244in}}%
\pgfpathlineto{\pgfqpoint{6.564857in}{1.535848in}}%
\pgfpathlineto{\pgfqpoint{6.540808in}{1.569165in}}%
\pgfpathlineto{\pgfqpoint{6.517920in}{1.603160in}}%
\pgfpathlineto{\pgfqpoint{6.496215in}{1.637800in}}%
\pgfpathlineto{\pgfqpoint{6.475715in}{1.673050in}}%
\pgfpathlineto{\pgfqpoint{6.456440in}{1.708875in}}%
\pgfpathlineto{\pgfqpoint{6.438410in}{1.745240in}}%
\pgfpathlineto{\pgfqpoint{6.421642in}{1.782107in}}%
\pgfpathlineto{\pgfqpoint{6.406153in}{1.819442in}}%
\pgfpathlineto{\pgfqpoint{6.391959in}{1.857205in}}%
\pgfpathlineto{\pgfqpoint{6.379074in}{1.895360in}}%
\pgfpathlineto{\pgfqpoint{6.367510in}{1.933868in}}%
\pgfpathlineto{\pgfqpoint{6.357279in}{1.972692in}}%
\pgfpathlineto{\pgfqpoint{6.348392in}{2.011792in}}%
\pgfpathlineto{\pgfqpoint{6.340857in}{2.051129in}}%
\pgfpathlineto{\pgfqpoint{6.334682in}{2.090665in}}%
\pgfpathlineto{\pgfqpoint{6.329873in}{2.130360in}}%
\pgfpathlineto{\pgfqpoint{6.326434in}{2.170173in}}%
\pgfpathlineto{\pgfqpoint{6.324369in}{2.210067in}}%
\pgfpathlineto{\pgfqpoint{6.323681in}{2.250000in}}%
\pgfpathlineto{\pgfqpoint{6.323681in}{2.250000in}}%
\pgfusepath{stroke}%
\end{pgfscope}%
\begin{pgfscope}%
\pgfpathrectangle{\pgfqpoint{5.486220in}{0.585980in}}{\pgfqpoint{4.437388in}{3.328041in}}%
\pgfusepath{clip}%
\pgfsetbuttcap%
\pgfsetroundjoin%
\pgfsetlinewidth{1.505625pt}%
\definecolor{currentstroke}{rgb}{0.973416,0.585761,0.251540}%
\pgfsetstrokecolor{currentstroke}%
\pgfsetstrokeopacity{0.800000}%
\pgfsetdash{{5.550000pt}{2.400000pt}}{0.000000pt}%
\pgfpathmoveto{\pgfqpoint{5.853920in}{2.250000in}}%
\pgfpathlineto{\pgfqpoint{5.854842in}{2.305748in}}%
\pgfpathlineto{\pgfqpoint{5.857609in}{2.361440in}}%
\pgfpathlineto{\pgfqpoint{5.862217in}{2.417021in}}%
\pgfpathlineto{\pgfqpoint{5.868662in}{2.472436in}}%
\pgfpathlineto{\pgfqpoint{5.876938in}{2.527629in}}%
\pgfpathlineto{\pgfqpoint{5.887035in}{2.582546in}}%
\pgfpathlineto{\pgfqpoint{5.898945in}{2.637130in}}%
\pgfpathlineto{\pgfqpoint{5.912655in}{2.691329in}}%
\pgfpathlineto{\pgfqpoint{5.928152in}{2.745088in}}%
\pgfpathlineto{\pgfqpoint{5.945419in}{2.798354in}}%
\pgfpathlineto{\pgfqpoint{5.964441in}{2.851073in}}%
\pgfpathlineto{\pgfqpoint{5.985197in}{2.903192in}}%
\pgfpathlineto{\pgfqpoint{6.007668in}{2.954661in}}%
\pgfpathlineto{\pgfqpoint{6.031830in}{3.005427in}}%
\pgfpathlineto{\pgfqpoint{6.057661in}{3.055440in}}%
\pgfpathlineto{\pgfqpoint{6.085133in}{3.104650in}}%
\pgfpathlineto{\pgfqpoint{6.114220in}{3.153009in}}%
\pgfpathlineto{\pgfqpoint{6.144893in}{3.200467in}}%
\pgfpathlineto{\pgfqpoint{6.177120in}{3.246977in}}%
\pgfpathlineto{\pgfqpoint{6.210871in}{3.292494in}}%
\pgfpathlineto{\pgfqpoint{6.246111in}{3.336972in}}%
\pgfpathlineto{\pgfqpoint{6.282805in}{3.380366in}}%
\pgfpathlineto{\pgfqpoint{6.320917in}{3.422634in}}%
\pgfpathlineto{\pgfqpoint{6.360408in}{3.463732in}}%
\pgfpathlineto{\pgfqpoint{6.401239in}{3.503621in}}%
\pgfpathlineto{\pgfqpoint{6.443370in}{3.542260in}}%
\pgfpathlineto{\pgfqpoint{6.486759in}{3.579611in}}%
\pgfpathlineto{\pgfqpoint{6.531362in}{3.615636in}}%
\pgfpathlineto{\pgfqpoint{6.577134in}{3.650300in}}%
\pgfpathlineto{\pgfqpoint{6.624031in}{3.683569in}}%
\pgfpathlineto{\pgfqpoint{6.672005in}{3.715408in}}%
\pgfpathlineto{\pgfqpoint{6.721009in}{3.745787in}}%
\pgfpathlineto{\pgfqpoint{6.770994in}{3.774674in}}%
\pgfpathlineto{\pgfqpoint{6.821910in}{3.802042in}}%
\pgfpathlineto{\pgfqpoint{6.873706in}{3.827863in}}%
\pgfpathlineto{\pgfqpoint{6.926330in}{3.852110in}}%
\pgfpathlineto{\pgfqpoint{6.979731in}{3.874761in}}%
\pgfpathlineto{\pgfqpoint{7.033854in}{3.895792in}}%
\pgfpathlineto{\pgfqpoint{7.116263in}{3.924020in}}%
\pgfpathmoveto{\pgfqpoint{8.294252in}{3.924020in}}%
\pgfpathlineto{\pgfqpoint{8.348657in}{3.905694in}}%
\pgfpathlineto{\pgfqpoint{8.403122in}{3.885480in}}%
\pgfpathlineto{\pgfqpoint{8.456891in}{3.863637in}}%
\pgfpathlineto{\pgfqpoint{8.509910in}{3.840185in}}%
\pgfpathlineto{\pgfqpoint{8.562126in}{3.815147in}}%
\pgfpathlineto{\pgfqpoint{8.613489in}{3.788550in}}%
\pgfpathlineto{\pgfqpoint{8.663945in}{3.760419in}}%
\pgfpathlineto{\pgfqpoint{8.713446in}{3.730782in}}%
\pgfpathlineto{\pgfqpoint{8.761941in}{3.699669in}}%
\pgfpathlineto{\pgfqpoint{8.809382in}{3.667111in}}%
\pgfpathlineto{\pgfqpoint{8.855723in}{3.633141in}}%
\pgfpathlineto{\pgfqpoint{8.900916in}{3.597791in}}%
\pgfpathlineto{\pgfqpoint{8.944917in}{3.561099in}}%
\pgfpathlineto{\pgfqpoint{8.987682in}{3.523099in}}%
\pgfpathlineto{\pgfqpoint{9.029169in}{3.483830in}}%
\pgfpathlineto{\pgfqpoint{9.069335in}{3.443332in}}%
\pgfpathlineto{\pgfqpoint{9.108141in}{3.401644in}}%
\pgfpathlineto{\pgfqpoint{9.145549in}{3.358807in}}%
\pgfpathlineto{\pgfqpoint{9.181520in}{3.314866in}}%
\pgfpathlineto{\pgfqpoint{9.216020in}{3.269863in}}%
\pgfpathlineto{\pgfqpoint{9.249013in}{3.223843in}}%
\pgfpathlineto{\pgfqpoint{9.280467in}{3.176853in}}%
\pgfpathlineto{\pgfqpoint{9.310351in}{3.128939in}}%
\pgfpathlineto{\pgfqpoint{9.338634in}{3.080149in}}%
\pgfpathlineto{\pgfqpoint{9.365289in}{3.030531in}}%
\pgfpathlineto{\pgfqpoint{9.390288in}{2.980135in}}%
\pgfpathlineto{\pgfqpoint{9.413608in}{2.929011in}}%
\pgfpathlineto{\pgfqpoint{9.435224in}{2.877211in}}%
\pgfpathlineto{\pgfqpoint{9.455115in}{2.824785in}}%
\pgfpathlineto{\pgfqpoint{9.473262in}{2.771786in}}%
\pgfpathlineto{\pgfqpoint{9.489646in}{2.718267in}}%
\pgfpathlineto{\pgfqpoint{9.504251in}{2.664281in}}%
\pgfpathlineto{\pgfqpoint{9.517063in}{2.609883in}}%
\pgfpathlineto{\pgfqpoint{9.528068in}{2.555125in}}%
\pgfpathlineto{\pgfqpoint{9.537256in}{2.500064in}}%
\pgfpathlineto{\pgfqpoint{9.544617in}{2.444753in}}%
\pgfpathlineto{\pgfqpoint{9.550144in}{2.389248in}}%
\pgfpathlineto{\pgfqpoint{9.553832in}{2.333604in}}%
\pgfpathlineto{\pgfqpoint{9.555677in}{2.277877in}}%
\pgfpathlineto{\pgfqpoint{9.555677in}{2.222123in}}%
\pgfpathlineto{\pgfqpoint{9.553832in}{2.166396in}}%
\pgfpathlineto{\pgfqpoint{9.550144in}{2.110752in}}%
\pgfpathlineto{\pgfqpoint{9.544617in}{2.055247in}}%
\pgfpathlineto{\pgfqpoint{9.537256in}{1.999936in}}%
\pgfpathlineto{\pgfqpoint{9.528068in}{1.944875in}}%
\pgfpathlineto{\pgfqpoint{9.517063in}{1.890117in}}%
\pgfpathlineto{\pgfqpoint{9.504251in}{1.835719in}}%
\pgfpathlineto{\pgfqpoint{9.489646in}{1.781733in}}%
\pgfpathlineto{\pgfqpoint{9.473262in}{1.728214in}}%
\pgfpathlineto{\pgfqpoint{9.455115in}{1.675215in}}%
\pgfpathlineto{\pgfqpoint{9.435224in}{1.622789in}}%
\pgfpathlineto{\pgfqpoint{9.413608in}{1.570989in}}%
\pgfpathlineto{\pgfqpoint{9.390288in}{1.519865in}}%
\pgfpathlineto{\pgfqpoint{9.365289in}{1.469469in}}%
\pgfpathlineto{\pgfqpoint{9.338634in}{1.419851in}}%
\pgfpathlineto{\pgfqpoint{9.310351in}{1.371061in}}%
\pgfpathlineto{\pgfqpoint{9.280467in}{1.323147in}}%
\pgfpathlineto{\pgfqpoint{9.249013in}{1.276157in}}%
\pgfpathlineto{\pgfqpoint{9.216020in}{1.230137in}}%
\pgfpathlineto{\pgfqpoint{9.181520in}{1.185134in}}%
\pgfpathlineto{\pgfqpoint{9.145549in}{1.141193in}}%
\pgfpathlineto{\pgfqpoint{9.108141in}{1.098356in}}%
\pgfpathlineto{\pgfqpoint{9.069335in}{1.056668in}}%
\pgfpathlineto{\pgfqpoint{9.029169in}{1.016170in}}%
\pgfpathlineto{\pgfqpoint{8.987682in}{0.976901in}}%
\pgfpathlineto{\pgfqpoint{8.944917in}{0.938901in}}%
\pgfpathlineto{\pgfqpoint{8.900916in}{0.902209in}}%
\pgfpathlineto{\pgfqpoint{8.855723in}{0.866859in}}%
\pgfpathlineto{\pgfqpoint{8.809382in}{0.832889in}}%
\pgfpathlineto{\pgfqpoint{8.761941in}{0.800331in}}%
\pgfpathlineto{\pgfqpoint{8.713446in}{0.769218in}}%
\pgfpathlineto{\pgfqpoint{8.663945in}{0.739581in}}%
\pgfpathlineto{\pgfqpoint{8.613489in}{0.711450in}}%
\pgfpathlineto{\pgfqpoint{8.562126in}{0.684853in}}%
\pgfpathlineto{\pgfqpoint{8.509910in}{0.659815in}}%
\pgfpathlineto{\pgfqpoint{8.456891in}{0.636363in}}%
\pgfpathlineto{\pgfqpoint{8.403122in}{0.614520in}}%
\pgfpathlineto{\pgfqpoint{8.348657in}{0.594306in}}%
\pgfpathlineto{\pgfqpoint{8.294252in}{0.575980in}}%
\pgfpathmoveto{\pgfqpoint{7.116263in}{0.575980in}}%
\pgfpathlineto{\pgfqpoint{7.088646in}{0.584817in}}%
\pgfpathlineto{\pgfqpoint{7.033854in}{0.604208in}}%
\pgfpathlineto{\pgfqpoint{6.979731in}{0.625239in}}%
\pgfpathlineto{\pgfqpoint{6.926330in}{0.647890in}}%
\pgfpathlineto{\pgfqpoint{6.873706in}{0.672137in}}%
\pgfpathlineto{\pgfqpoint{6.821910in}{0.697958in}}%
\pgfpathlineto{\pgfqpoint{6.770994in}{0.725326in}}%
\pgfpathlineto{\pgfqpoint{6.721009in}{0.754213in}}%
\pgfpathlineto{\pgfqpoint{6.672005in}{0.784592in}}%
\pgfpathlineto{\pgfqpoint{6.624031in}{0.816431in}}%
\pgfpathlineto{\pgfqpoint{6.577134in}{0.849700in}}%
\pgfpathlineto{\pgfqpoint{6.531362in}{0.884364in}}%
\pgfpathlineto{\pgfqpoint{6.486759in}{0.920389in}}%
\pgfpathlineto{\pgfqpoint{6.443370in}{0.957740in}}%
\pgfpathlineto{\pgfqpoint{6.401239in}{0.996379in}}%
\pgfpathlineto{\pgfqpoint{6.360408in}{1.036268in}}%
\pgfpathlineto{\pgfqpoint{6.320917in}{1.077366in}}%
\pgfpathlineto{\pgfqpoint{6.282805in}{1.119634in}}%
\pgfpathlineto{\pgfqpoint{6.246111in}{1.163028in}}%
\pgfpathlineto{\pgfqpoint{6.210871in}{1.207506in}}%
\pgfpathlineto{\pgfqpoint{6.177120in}{1.253023in}}%
\pgfpathlineto{\pgfqpoint{6.144893in}{1.299533in}}%
\pgfpathlineto{\pgfqpoint{6.114220in}{1.346991in}}%
\pgfpathlineto{\pgfqpoint{6.085133in}{1.395350in}}%
\pgfpathlineto{\pgfqpoint{6.057661in}{1.444560in}}%
\pgfpathlineto{\pgfqpoint{6.031830in}{1.494573in}}%
\pgfpathlineto{\pgfqpoint{6.007668in}{1.545339in}}%
\pgfpathlineto{\pgfqpoint{5.985197in}{1.596808in}}%
\pgfpathlineto{\pgfqpoint{5.964441in}{1.648927in}}%
\pgfpathlineto{\pgfqpoint{5.945419in}{1.701646in}}%
\pgfpathlineto{\pgfqpoint{5.928152in}{1.754912in}}%
\pgfpathlineto{\pgfqpoint{5.912655in}{1.808671in}}%
\pgfpathlineto{\pgfqpoint{5.898945in}{1.862870in}}%
\pgfpathlineto{\pgfqpoint{5.887035in}{1.917454in}}%
\pgfpathlineto{\pgfqpoint{5.876938in}{1.972371in}}%
\pgfpathlineto{\pgfqpoint{5.868662in}{2.027564in}}%
\pgfpathlineto{\pgfqpoint{5.862217in}{2.082979in}}%
\pgfpathlineto{\pgfqpoint{5.857609in}{2.138560in}}%
\pgfpathlineto{\pgfqpoint{5.854842in}{2.194252in}}%
\pgfpathlineto{\pgfqpoint{5.853920in}{2.250000in}}%
\pgfpathlineto{\pgfqpoint{5.853920in}{2.250000in}}%
\pgfusepath{stroke}%
\end{pgfscope}%
\begin{pgfscope}%
\pgfpathrectangle{\pgfqpoint{5.486220in}{0.585980in}}{\pgfqpoint{4.437388in}{3.328041in}}%
\pgfusepath{clip}%
\pgfsetbuttcap%
\pgfsetroundjoin%
\pgfsetlinewidth{1.505625pt}%
\definecolor{currentstroke}{rgb}{0.993033,0.771720,0.154808}%
\pgfsetstrokecolor{currentstroke}%
\pgfsetstrokeopacity{0.800000}%
\pgfsetdash{{5.550000pt}{2.400000pt}}{0.000000pt}%
\pgfpathmoveto{\pgfqpoint{5.476220in}{3.360835in}}%
\pgfpathlineto{\pgfqpoint{5.477429in}{3.363228in}}%
\pgfpathlineto{\pgfqpoint{5.514578in}{3.431244in}}%
\pgfpathlineto{\pgfqpoint{5.553911in}{3.498081in}}%
\pgfpathlineto{\pgfqpoint{5.595387in}{3.563675in}}%
\pgfpathlineto{\pgfqpoint{5.638967in}{3.627959in}}%
\pgfpathlineto{\pgfqpoint{5.684606in}{3.690870in}}%
\pgfpathlineto{\pgfqpoint{5.732259in}{3.752344in}}%
\pgfpathlineto{\pgfqpoint{5.781878in}{3.812321in}}%
\pgfpathlineto{\pgfqpoint{5.833414in}{3.870740in}}%
\pgfpathlineto{\pgfqpoint{5.883504in}{3.924020in}}%
\pgfpathmoveto{\pgfqpoint{9.526001in}{3.924020in}}%
\pgfpathlineto{\pgfqpoint{9.549942in}{3.899348in}}%
\pgfpathlineto{\pgfqpoint{9.602418in}{3.841729in}}%
\pgfpathlineto{\pgfqpoint{9.653002in}{3.782523in}}%
\pgfpathlineto{\pgfqpoint{9.701644in}{3.721790in}}%
\pgfpathlineto{\pgfqpoint{9.748295in}{3.659590in}}%
\pgfpathlineto{\pgfqpoint{9.792910in}{3.595985in}}%
\pgfpathlineto{\pgfqpoint{9.835444in}{3.531038in}}%
\pgfpathlineto{\pgfqpoint{9.875854in}{3.464814in}}%
\pgfpathlineto{\pgfqpoint{9.914099in}{3.397379in}}%
\pgfpathlineto{\pgfqpoint{9.933608in}{3.360261in}}%
\pgfpathmoveto{\pgfqpoint{9.933608in}{1.139739in}}%
\pgfpathlineto{\pgfqpoint{9.914099in}{1.102621in}}%
\pgfpathlineto{\pgfqpoint{9.875854in}{1.035186in}}%
\pgfpathlineto{\pgfqpoint{9.835444in}{0.968962in}}%
\pgfpathlineto{\pgfqpoint{9.792910in}{0.904015in}}%
\pgfpathlineto{\pgfqpoint{9.748295in}{0.840410in}}%
\pgfpathlineto{\pgfqpoint{9.701644in}{0.778210in}}%
\pgfpathlineto{\pgfqpoint{9.653002in}{0.717477in}}%
\pgfpathlineto{\pgfqpoint{9.602418in}{0.658271in}}%
\pgfpathlineto{\pgfqpoint{9.526001in}{0.575980in}}%
\pgfpathmoveto{\pgfqpoint{5.883504in}{0.575980in}}%
\pgfpathlineto{\pgfqpoint{5.833414in}{0.629260in}}%
\pgfpathlineto{\pgfqpoint{5.781878in}{0.687679in}}%
\pgfpathlineto{\pgfqpoint{5.732259in}{0.747656in}}%
\pgfpathlineto{\pgfqpoint{5.684606in}{0.809130in}}%
\pgfpathlineto{\pgfqpoint{5.638967in}{0.872041in}}%
\pgfpathlineto{\pgfqpoint{5.595387in}{0.936325in}}%
\pgfpathlineto{\pgfqpoint{5.553911in}{1.001919in}}%
\pgfpathlineto{\pgfqpoint{5.514578in}{1.068756in}}%
\pgfpathlineto{\pgfqpoint{5.476220in}{1.139165in}}%
\pgfpathlineto{\pgfqpoint{5.476220in}{1.139165in}}%
\pgfusepath{stroke}%
\end{pgfscope}%
\begin{pgfscope}%
\pgfpathrectangle{\pgfqpoint{5.486220in}{0.585980in}}{\pgfqpoint{4.437388in}{3.328041in}}%
\pgfusepath{clip}%
\pgfsetbuttcap%
\pgfsetroundjoin%
\pgfsetlinewidth{1.505625pt}%
\definecolor{currentstroke}{rgb}{0.940015,0.975158,0.131326}%
\pgfsetstrokecolor{currentstroke}%
\pgfsetstrokeopacity{0.800000}%
\pgfsetdash{{5.550000pt}{2.400000pt}}{0.000000pt}%
\pgfpathmoveto{\pgfqpoint{0.000000in}{0.000000in}}%
\pgfusepath{stroke}%
\end{pgfscope}%
\begin{pgfscope}%
\pgfpathrectangle{\pgfqpoint{5.486220in}{0.585980in}}{\pgfqpoint{4.437388in}{3.328041in}}%
\pgfusepath{clip}%
\pgfsetbuttcap%
\pgfsetroundjoin%
\definecolor{currentfill}{rgb}{0.000000,0.000000,0.000000}%
\pgfsetfillcolor{currentfill}%
\pgfsetfillopacity{0.900000}%
\pgfsetlinewidth{1.003750pt}%
\definecolor{currentstroke}{rgb}{0.000000,0.000000,0.000000}%
\pgfsetstrokecolor{currentstroke}%
\pgfsetstrokeopacity{0.900000}%
\pgfsetdash{}{0pt}%
\pgfsys@defobject{currentmarker}{\pgfqpoint{-0.001736in}{-0.001736in}}{\pgfqpoint{0.001736in}{0.001736in}}{%
\pgfpathmoveto{\pgfqpoint{0.000000in}{-0.001736in}}%
\pgfpathcurveto{\pgfqpoint{0.000460in}{-0.001736in}}{\pgfqpoint{0.000902in}{-0.001553in}}{\pgfqpoint{0.001228in}{-0.001228in}}%
\pgfpathcurveto{\pgfqpoint{0.001553in}{-0.000902in}}{\pgfqpoint{0.001736in}{-0.000460in}}{\pgfqpoint{0.001736in}{0.000000in}}%
\pgfpathcurveto{\pgfqpoint{0.001736in}{0.000460in}}{\pgfqpoint{0.001553in}{0.000902in}}{\pgfqpoint{0.001228in}{0.001228in}}%
\pgfpathcurveto{\pgfqpoint{0.000902in}{0.001553in}}{\pgfqpoint{0.000460in}{0.001736in}}{\pgfqpoint{0.000000in}{0.001736in}}%
\pgfpathcurveto{\pgfqpoint{-0.000460in}{0.001736in}}{\pgfqpoint{-0.000902in}{0.001553in}}{\pgfqpoint{-0.001228in}{0.001228in}}%
\pgfpathcurveto{\pgfqpoint{-0.001553in}{0.000902in}}{\pgfqpoint{-0.001736in}{0.000460in}}{\pgfqpoint{-0.001736in}{0.000000in}}%
\pgfpathcurveto{\pgfqpoint{-0.001736in}{-0.000460in}}{\pgfqpoint{-0.001553in}{-0.000902in}}{\pgfqpoint{-0.001228in}{-0.001228in}}%
\pgfpathcurveto{\pgfqpoint{-0.000902in}{-0.001553in}}{\pgfqpoint{-0.000460in}{-0.001736in}}{\pgfqpoint{0.000000in}{-0.001736in}}%
\pgfpathlineto{\pgfqpoint{0.000000in}{-0.001736in}}%
\pgfpathclose%
\pgfusepath{stroke,fill}%
}%
\begin{pgfscope}%
\pgfsys@transformshift{8.317965in}{2.190776in}%
\pgfsys@useobject{currentmarker}{}%
\end{pgfscope}%
\begin{pgfscope}%
\pgfsys@transformshift{8.313631in}{2.181824in}%
\pgfsys@useobject{currentmarker}{}%
\end{pgfscope}%
\begin{pgfscope}%
\pgfsys@transformshift{8.308690in}{2.172940in}%
\pgfsys@useobject{currentmarker}{}%
\end{pgfscope}%
\begin{pgfscope}%
\pgfsys@transformshift{8.303147in}{2.164132in}%
\pgfsys@useobject{currentmarker}{}%
\end{pgfscope}%
\begin{pgfscope}%
\pgfsys@transformshift{8.297008in}{2.155411in}%
\pgfsys@useobject{currentmarker}{}%
\end{pgfscope}%
\begin{pgfscope}%
\pgfsys@transformshift{8.290279in}{2.146783in}%
\pgfsys@useobject{currentmarker}{}%
\end{pgfscope}%
\begin{pgfscope}%
\pgfsys@transformshift{8.282966in}{2.138259in}%
\pgfsys@useobject{currentmarker}{}%
\end{pgfscope}%
\begin{pgfscope}%
\pgfsys@transformshift{8.275077in}{2.129846in}%
\pgfsys@useobject{currentmarker}{}%
\end{pgfscope}%
\begin{pgfscope}%
\pgfsys@transformshift{8.266620in}{2.121552in}%
\pgfsys@useobject{currentmarker}{}%
\end{pgfscope}%
\begin{pgfscope}%
\pgfsys@transformshift{8.257602in}{2.113387in}%
\pgfsys@useobject{currentmarker}{}%
\end{pgfscope}%
\begin{pgfscope}%
\pgfsys@transformshift{8.248034in}{2.105358in}%
\pgfsys@useobject{currentmarker}{}%
\end{pgfscope}%
\begin{pgfscope}%
\pgfsys@transformshift{8.237925in}{2.097473in}%
\pgfsys@useobject{currentmarker}{}%
\end{pgfscope}%
\begin{pgfscope}%
\pgfsys@transformshift{8.227284in}{2.089740in}%
\pgfsys@useobject{currentmarker}{}%
\end{pgfscope}%
\begin{pgfscope}%
\pgfsys@transformshift{8.216122in}{2.082166in}%
\pgfsys@useobject{currentmarker}{}%
\end{pgfscope}%
\begin{pgfscope}%
\pgfsys@transformshift{8.204451in}{2.074761in}%
\pgfsys@useobject{currentmarker}{}%
\end{pgfscope}%
\begin{pgfscope}%
\pgfsys@transformshift{8.192282in}{2.067529in}%
\pgfsys@useobject{currentmarker}{}%
\end{pgfscope}%
\begin{pgfscope}%
\pgfsys@transformshift{8.179627in}{2.060480in}%
\pgfsys@useobject{currentmarker}{}%
\end{pgfscope}%
\begin{pgfscope}%
\pgfsys@transformshift{8.166498in}{2.053620in}%
\pgfsys@useobject{currentmarker}{}%
\end{pgfscope}%
\begin{pgfscope}%
\pgfsys@transformshift{8.152910in}{2.046955in}%
\pgfsys@useobject{currentmarker}{}%
\end{pgfscope}%
\begin{pgfscope}%
\pgfsys@transformshift{8.321588in}{2.189168in}%
\pgfsys@useobject{currentmarker}{}%
\end{pgfscope}%
\begin{pgfscope}%
\pgfsys@transformshift{8.317228in}{2.179973in}%
\pgfsys@useobject{currentmarker}{}%
\end{pgfscope}%
\begin{pgfscope}%
\pgfsys@transformshift{8.312258in}{2.170848in}%
\pgfsys@useobject{currentmarker}{}%
\end{pgfscope}%
\begin{pgfscope}%
\pgfsys@transformshift{8.306682in}{2.161801in}%
\pgfsys@useobject{currentmarker}{}%
\end{pgfscope}%
\begin{pgfscope}%
\pgfsys@transformshift{8.300507in}{2.152843in}%
\pgfsys@useobject{currentmarker}{}%
\end{pgfscope}%
\begin{pgfscope}%
\pgfsys@transformshift{8.293738in}{2.143981in}%
\pgfsys@useobject{currentmarker}{}%
\end{pgfscope}%
\begin{pgfscope}%
\pgfsys@transformshift{8.286382in}{2.135225in}%
\pgfsys@useobject{currentmarker}{}%
\end{pgfscope}%
\begin{pgfscope}%
\pgfsys@transformshift{8.278446in}{2.126583in}%
\pgfsys@useobject{currentmarker}{}%
\end{pgfscope}%
\begin{pgfscope}%
\pgfsys@transformshift{8.269939in}{2.118065in}%
\pgfsys@useobject{currentmarker}{}%
\end{pgfscope}%
\begin{pgfscope}%
\pgfsys@transformshift{8.260868in}{2.109678in}%
\pgfsys@useobject{currentmarker}{}%
\end{pgfscope}%
\begin{pgfscope}%
\pgfsys@transformshift{8.251244in}{2.101431in}%
\pgfsys@useobject{currentmarker}{}%
\end{pgfscope}%
\begin{pgfscope}%
\pgfsys@transformshift{8.241074in}{2.093332in}%
\pgfsys@useobject{currentmarker}{}%
\end{pgfscope}%
\begin{pgfscope}%
\pgfsys@transformshift{8.230370in}{2.085389in}%
\pgfsys@useobject{currentmarker}{}%
\end{pgfscope}%
\begin{pgfscope}%
\pgfsys@transformshift{8.219143in}{2.077610in}%
\pgfsys@useobject{currentmarker}{}%
\end{pgfscope}%
\begin{pgfscope}%
\pgfsys@transformshift{8.207403in}{2.070003in}%
\pgfsys@useobject{currentmarker}{}%
\end{pgfscope}%
\begin{pgfscope}%
\pgfsys@transformshift{8.195162in}{2.062575in}%
\pgfsys@useobject{currentmarker}{}%
\end{pgfscope}%
\begin{pgfscope}%
\pgfsys@transformshift{8.182432in}{2.055335in}%
\pgfsys@useobject{currentmarker}{}%
\end{pgfscope}%
\begin{pgfscope}%
\pgfsys@transformshift{8.169226in}{2.048288in}%
\pgfsys@useobject{currentmarker}{}%
\end{pgfscope}%
\begin{pgfscope}%
\pgfsys@transformshift{8.155557in}{2.041442in}%
\pgfsys@useobject{currentmarker}{}%
\end{pgfscope}%
\begin{pgfscope}%
\pgfsys@transformshift{8.325308in}{2.187551in}%
\pgfsys@useobject{currentmarker}{}%
\end{pgfscope}%
\begin{pgfscope}%
\pgfsys@transformshift{8.320922in}{2.178111in}%
\pgfsys@useobject{currentmarker}{}%
\end{pgfscope}%
\begin{pgfscope}%
\pgfsys@transformshift{8.315922in}{2.168743in}%
\pgfsys@useobject{currentmarker}{}%
\end{pgfscope}%
\begin{pgfscope}%
\pgfsys@transformshift{8.310313in}{2.159456in}%
\pgfsys@useobject{currentmarker}{}%
\end{pgfscope}%
\begin{pgfscope}%
\pgfsys@transformshift{8.304100in}{2.150259in}%
\pgfsys@useobject{currentmarker}{}%
\end{pgfscope}%
\begin{pgfscope}%
\pgfsys@transformshift{8.297290in}{2.141162in}%
\pgfsys@useobject{currentmarker}{}%
\end{pgfscope}%
\begin{pgfscope}%
\pgfsys@transformshift{8.289890in}{2.132173in}%
\pgfsys@useobject{currentmarker}{}%
\end{pgfscope}%
\begin{pgfscope}%
\pgfsys@transformshift{8.281906in}{2.123302in}%
\pgfsys@useobject{currentmarker}{}%
\end{pgfscope}%
\begin{pgfscope}%
\pgfsys@transformshift{8.273347in}{2.114557in}%
\pgfsys@useobject{currentmarker}{}%
\end{pgfscope}%
\begin{pgfscope}%
\pgfsys@transformshift{8.264222in}{2.105947in}%
\pgfsys@useobject{currentmarker}{}%
\end{pgfscope}%
\begin{pgfscope}%
\pgfsys@transformshift{8.254539in}{2.097480in}%
\pgfsys@useobject{currentmarker}{}%
\end{pgfscope}%
\begin{pgfscope}%
\pgfsys@transformshift{8.244309in}{2.089166in}%
\pgfsys@useobject{currentmarker}{}%
\end{pgfscope}%
\begin{pgfscope}%
\pgfsys@transformshift{8.233540in}{2.081012in}%
\pgfsys@useobject{currentmarker}{}%
\end{pgfscope}%
\begin{pgfscope}%
\pgfsys@transformshift{8.222245in}{2.073026in}%
\pgfsys@useobject{currentmarker}{}%
\end{pgfscope}%
\begin{pgfscope}%
\pgfsys@transformshift{8.210434in}{2.065217in}%
\pgfsys@useobject{currentmarker}{}%
\end{pgfscope}%
\begin{pgfscope}%
\pgfsys@transformshift{8.198119in}{2.057592in}%
\pgfsys@useobject{currentmarker}{}%
\end{pgfscope}%
\begin{pgfscope}%
\pgfsys@transformshift{8.185312in}{2.050159in}%
\pgfsys@useobject{currentmarker}{}%
\end{pgfscope}%
\begin{pgfscope}%
\pgfsys@transformshift{8.172027in}{2.042925in}%
\pgfsys@useobject{currentmarker}{}%
\end{pgfscope}%
\begin{pgfscope}%
\pgfsys@transformshift{8.158276in}{2.035897in}%
\pgfsys@useobject{currentmarker}{}%
\end{pgfscope}%
\begin{pgfscope}%
\pgfsys@transformshift{8.329126in}{2.185923in}%
\pgfsys@useobject{currentmarker}{}%
\end{pgfscope}%
\begin{pgfscope}%
\pgfsys@transformshift{8.324713in}{2.176238in}%
\pgfsys@useobject{currentmarker}{}%
\end{pgfscope}%
\begin{pgfscope}%
\pgfsys@transformshift{8.319682in}{2.166626in}%
\pgfsys@useobject{currentmarker}{}%
\end{pgfscope}%
\begin{pgfscope}%
\pgfsys@transformshift{8.314038in}{2.157097in}%
\pgfsys@useobject{currentmarker}{}%
\end{pgfscope}%
\begin{pgfscope}%
\pgfsys@transformshift{8.307787in}{2.147660in}%
\pgfsys@useobject{currentmarker}{}%
\end{pgfscope}%
\begin{pgfscope}%
\pgfsys@transformshift{8.300936in}{2.138326in}%
\pgfsys@useobject{currentmarker}{}%
\end{pgfscope}%
\begin{pgfscope}%
\pgfsys@transformshift{8.293490in}{2.129103in}%
\pgfsys@useobject{currentmarker}{}%
\end{pgfscope}%
\begin{pgfscope}%
\pgfsys@transformshift{8.285457in}{2.120000in}%
\pgfsys@useobject{currentmarker}{}%
\end{pgfscope}%
\begin{pgfscope}%
\pgfsys@transformshift{8.276846in}{2.111027in}%
\pgfsys@useobject{currentmarker}{}%
\end{pgfscope}%
\begin{pgfscope}%
\pgfsys@transformshift{8.267664in}{2.102193in}%
\pgfsys@useobject{currentmarker}{}%
\end{pgfscope}%
\begin{pgfscope}%
\pgfsys@transformshift{8.257922in}{2.093506in}%
\pgfsys@useobject{currentmarker}{}%
\end{pgfscope}%
\begin{pgfscope}%
\pgfsys@transformshift{8.247628in}{2.084975in}%
\pgfsys@useobject{currentmarker}{}%
\end{pgfscope}%
\begin{pgfscope}%
\pgfsys@transformshift{8.236793in}{2.076608in}%
\pgfsys@useobject{currentmarker}{}%
\end{pgfscope}%
\begin{pgfscope}%
\pgfsys@transformshift{8.225429in}{2.068414in}%
\pgfsys@useobject{currentmarker}{}%
\end{pgfscope}%
\begin{pgfscope}%
\pgfsys@transformshift{8.213545in}{2.060402in}%
\pgfsys@useobject{currentmarker}{}%
\end{pgfscope}%
\begin{pgfscope}%
\pgfsys@transformshift{8.201154in}{2.052578in}%
\pgfsys@useobject{currentmarker}{}%
\end{pgfscope}%
\begin{pgfscope}%
\pgfsys@transformshift{8.188269in}{2.044951in}%
\pgfsys@useobject{currentmarker}{}%
\end{pgfscope}%
\begin{pgfscope}%
\pgfsys@transformshift{8.174902in}{2.037528in}%
\pgfsys@useobject{currentmarker}{}%
\end{pgfscope}%
\begin{pgfscope}%
\pgfsys@transformshift{8.161066in}{2.030318in}%
\pgfsys@useobject{currentmarker}{}%
\end{pgfscope}%
\begin{pgfscope}%
\pgfsys@transformshift{8.333042in}{2.184286in}%
\pgfsys@useobject{currentmarker}{}%
\end{pgfscope}%
\begin{pgfscope}%
\pgfsys@transformshift{8.328602in}{2.174353in}%
\pgfsys@useobject{currentmarker}{}%
\end{pgfscope}%
\begin{pgfscope}%
\pgfsys@transformshift{8.323539in}{2.164495in}%
\pgfsys@useobject{currentmarker}{}%
\end{pgfscope}%
\begin{pgfscope}%
\pgfsys@transformshift{8.317860in}{2.154722in}%
\pgfsys@useobject{currentmarker}{}%
\end{pgfscope}%
\begin{pgfscope}%
\pgfsys@transformshift{8.311570in}{2.145045in}%
\pgfsys@useobject{currentmarker}{}%
\end{pgfscope}%
\begin{pgfscope}%
\pgfsys@transformshift{8.304675in}{2.135472in}%
\pgfsys@useobject{currentmarker}{}%
\end{pgfscope}%
\begin{pgfscope}%
\pgfsys@transformshift{8.297183in}{2.126013in}%
\pgfsys@useobject{currentmarker}{}%
\end{pgfscope}%
\begin{pgfscope}%
\pgfsys@transformshift{8.289099in}{2.116678in}%
\pgfsys@useobject{currentmarker}{}%
\end{pgfscope}%
\begin{pgfscope}%
\pgfsys@transformshift{8.280434in}{2.107476in}%
\pgfsys@useobject{currentmarker}{}%
\end{pgfscope}%
\begin{pgfscope}%
\pgfsys@transformshift{8.271195in}{2.098416in}%
\pgfsys@useobject{currentmarker}{}%
\end{pgfscope}%
\begin{pgfscope}%
\pgfsys@transformshift{8.261392in}{2.089507in}%
\pgfsys@useobject{currentmarker}{}%
\end{pgfscope}%
\begin{pgfscope}%
\pgfsys@transformshift{8.251033in}{2.080758in}%
\pgfsys@useobject{currentmarker}{}%
\end{pgfscope}%
\begin{pgfscope}%
\pgfsys@transformshift{8.240131in}{2.072177in}%
\pgfsys@useobject{currentmarker}{}%
\end{pgfscope}%
\begin{pgfscope}%
\pgfsys@transformshift{8.228695in}{2.063774in}%
\pgfsys@useobject{currentmarker}{}%
\end{pgfscope}%
\begin{pgfscope}%
\pgfsys@transformshift{8.216736in}{2.055557in}%
\pgfsys@useobject{currentmarker}{}%
\end{pgfscope}%
\begin{pgfscope}%
\pgfsys@transformshift{8.204268in}{2.047533in}%
\pgfsys@useobject{currentmarker}{}%
\end{pgfscope}%
\begin{pgfscope}%
\pgfsys@transformshift{8.191302in}{2.039711in}%
\pgfsys@useobject{currentmarker}{}%
\end{pgfscope}%
\begin{pgfscope}%
\pgfsys@transformshift{8.177851in}{2.032099in}%
\pgfsys@useobject{currentmarker}{}%
\end{pgfscope}%
\begin{pgfscope}%
\pgfsys@transformshift{8.163928in}{2.024704in}%
\pgfsys@useobject{currentmarker}{}%
\end{pgfscope}%
\begin{pgfscope}%
\pgfsys@transformshift{8.337058in}{2.182638in}%
\pgfsys@useobject{currentmarker}{}%
\end{pgfscope}%
\begin{pgfscope}%
\pgfsys@transformshift{8.332589in}{2.172456in}%
\pgfsys@useobject{currentmarker}{}%
\end{pgfscope}%
\begin{pgfscope}%
\pgfsys@transformshift{8.327494in}{2.162351in}%
\pgfsys@useobject{currentmarker}{}%
\end{pgfscope}%
\begin{pgfscope}%
\pgfsys@transformshift{8.321779in}{2.152333in}%
\pgfsys@useobject{currentmarker}{}%
\end{pgfscope}%
\begin{pgfscope}%
\pgfsys@transformshift{8.315448in}{2.142413in}%
\pgfsys@useobject{currentmarker}{}%
\end{pgfscope}%
\begin{pgfscope}%
\pgfsys@transformshift{8.308509in}{2.132600in}%
\pgfsys@useobject{currentmarker}{}%
\end{pgfscope}%
\begin{pgfscope}%
\pgfsys@transformshift{8.300969in}{2.122904in}%
\pgfsys@useobject{currentmarker}{}%
\end{pgfscope}%
\begin{pgfscope}%
\pgfsys@transformshift{8.292834in}{2.113335in}%
\pgfsys@useobject{currentmarker}{}%
\end{pgfscope}%
\begin{pgfscope}%
\pgfsys@transformshift{8.284113in}{2.103902in}%
\pgfsys@useobject{currentmarker}{}%
\end{pgfscope}%
\begin{pgfscope}%
\pgfsys@transformshift{8.274815in}{2.094615in}%
\pgfsys@useobject{currentmarker}{}%
\end{pgfscope}%
\begin{pgfscope}%
\pgfsys@transformshift{8.264949in}{2.085482in}%
\pgfsys@useobject{currentmarker}{}%
\end{pgfscope}%
\begin{pgfscope}%
\pgfsys@transformshift{8.254525in}{2.076514in}%
\pgfsys@useobject{currentmarker}{}%
\end{pgfscope}%
\begin{pgfscope}%
\pgfsys@transformshift{8.243552in}{2.067718in}%
\pgfsys@useobject{currentmarker}{}%
\end{pgfscope}%
\begin{pgfscope}%
\pgfsys@transformshift{8.232043in}{2.059104in}%
\pgfsys@useobject{currentmarker}{}%
\end{pgfscope}%
\begin{pgfscope}%
\pgfsys@transformshift{8.220008in}{2.050681in}%
\pgfsys@useobject{currentmarker}{}%
\end{pgfscope}%
\begin{pgfscope}%
\pgfsys@transformshift{8.207460in}{2.042456in}%
\pgfsys@useobject{currentmarker}{}%
\end{pgfscope}%
\begin{pgfscope}%
\pgfsys@transformshift{8.194411in}{2.034438in}%
\pgfsys@useobject{currentmarker}{}%
\end{pgfscope}%
\begin{pgfscope}%
\pgfsys@transformshift{8.180874in}{2.026635in}%
\pgfsys@useobject{currentmarker}{}%
\end{pgfscope}%
\begin{pgfscope}%
\pgfsys@transformshift{8.166862in}{2.019054in}%
\pgfsys@useobject{currentmarker}{}%
\end{pgfscope}%
\begin{pgfscope}%
\pgfsys@transformshift{8.341173in}{2.180979in}%
\pgfsys@useobject{currentmarker}{}%
\end{pgfscope}%
\begin{pgfscope}%
\pgfsys@transformshift{8.336675in}{2.170547in}%
\pgfsys@useobject{currentmarker}{}%
\end{pgfscope}%
\begin{pgfscope}%
\pgfsys@transformshift{8.331547in}{2.160193in}%
\pgfsys@useobject{currentmarker}{}%
\end{pgfscope}%
\begin{pgfscope}%
\pgfsys@transformshift{8.325795in}{2.149929in}%
\pgfsys@useobject{currentmarker}{}%
\end{pgfscope}%
\begin{pgfscope}%
\pgfsys@transformshift{8.319423in}{2.139764in}%
\pgfsys@useobject{currentmarker}{}%
\end{pgfscope}%
\begin{pgfscope}%
\pgfsys@transformshift{8.312439in}{2.129710in}%
\pgfsys@useobject{currentmarker}{}%
\end{pgfscope}%
\begin{pgfscope}%
\pgfsys@transformshift{8.304849in}{2.119775in}%
\pgfsys@useobject{currentmarker}{}%
\end{pgfscope}%
\begin{pgfscope}%
\pgfsys@transformshift{8.296662in}{2.109970in}%
\pgfsys@useobject{currentmarker}{}%
\end{pgfscope}%
\begin{pgfscope}%
\pgfsys@transformshift{8.287884in}{2.100305in}%
\pgfsys@useobject{currentmarker}{}%
\end{pgfscope}%
\begin{pgfscope}%
\pgfsys@transformshift{8.278526in}{2.090789in}%
\pgfsys@useobject{currentmarker}{}%
\end{pgfscope}%
\begin{pgfscope}%
\pgfsys@transformshift{8.268595in}{2.081432in}%
\pgfsys@useobject{currentmarker}{}%
\end{pgfscope}%
\begin{pgfscope}%
\pgfsys@transformshift{8.258103in}{2.072242in}%
\pgfsys@useobject{currentmarker}{}%
\end{pgfscope}%
\begin{pgfscope}%
\pgfsys@transformshift{8.247059in}{2.063230in}%
\pgfsys@useobject{currentmarker}{}%
\end{pgfscope}%
\begin{pgfscope}%
\pgfsys@transformshift{8.235475in}{2.054404in}%
\pgfsys@useobject{currentmarker}{}%
\end{pgfscope}%
\begin{pgfscope}%
\pgfsys@transformshift{8.223362in}{2.045773in}%
\pgfsys@useobject{currentmarker}{}%
\end{pgfscope}%
\begin{pgfscope}%
\pgfsys@transformshift{8.210732in}{2.037346in}%
\pgfsys@useobject{currentmarker}{}%
\end{pgfscope}%
\begin{pgfscope}%
\pgfsys@transformshift{8.197598in}{2.029131in}%
\pgfsys@useobject{currentmarker}{}%
\end{pgfscope}%
\begin{pgfscope}%
\pgfsys@transformshift{8.183973in}{2.021135in}%
\pgfsys@useobject{currentmarker}{}%
\end{pgfscope}%
\begin{pgfscope}%
\pgfsys@transformshift{8.169870in}{2.013368in}%
\pgfsys@useobject{currentmarker}{}%
\end{pgfscope}%
\begin{pgfscope}%
\pgfsys@transformshift{8.345389in}{2.179310in}%
\pgfsys@useobject{currentmarker}{}%
\end{pgfscope}%
\begin{pgfscope}%
\pgfsys@transformshift{8.340861in}{2.168625in}%
\pgfsys@useobject{currentmarker}{}%
\end{pgfscope}%
\begin{pgfscope}%
\pgfsys@transformshift{8.335699in}{2.158021in}%
\pgfsys@useobject{currentmarker}{}%
\end{pgfscope}%
\begin{pgfscope}%
\pgfsys@transformshift{8.329908in}{2.147508in}%
\pgfsys@useobject{currentmarker}{}%
\end{pgfscope}%
\begin{pgfscope}%
\pgfsys@transformshift{8.323495in}{2.137098in}%
\pgfsys@useobject{currentmarker}{}%
\end{pgfscope}%
\begin{pgfscope}%
\pgfsys@transformshift{8.316464in}{2.126800in}%
\pgfsys@useobject{currentmarker}{}%
\end{pgfscope}%
\begin{pgfscope}%
\pgfsys@transformshift{8.308825in}{2.116625in}%
\pgfsys@useobject{currentmarker}{}%
\end{pgfscope}%
\begin{pgfscope}%
\pgfsys@transformshift{8.300583in}{2.106583in}%
\pgfsys@useobject{currentmarker}{}%
\end{pgfscope}%
\begin{pgfscope}%
\pgfsys@transformshift{8.291747in}{2.096684in}%
\pgfsys@useobject{currentmarker}{}%
\end{pgfscope}%
\begin{pgfscope}%
\pgfsys@transformshift{8.282326in}{2.086938in}%
\pgfsys@useobject{currentmarker}{}%
\end{pgfscope}%
\begin{pgfscope}%
\pgfsys@transformshift{8.272330in}{2.077355in}%
\pgfsys@useobject{currentmarker}{}%
\end{pgfscope}%
\begin{pgfscope}%
\pgfsys@transformshift{8.261768in}{2.067943in}%
\pgfsys@useobject{currentmarker}{}%
\end{pgfscope}%
\begin{pgfscope}%
\pgfsys@transformshift{8.250651in}{2.058713in}%
\pgfsys@useobject{currentmarker}{}%
\end{pgfscope}%
\begin{pgfscope}%
\pgfsys@transformshift{8.238990in}{2.049674in}%
\pgfsys@useobject{currentmarker}{}%
\end{pgfscope}%
\begin{pgfscope}%
\pgfsys@transformshift{8.226797in}{2.040834in}%
\pgfsys@useobject{currentmarker}{}%
\end{pgfscope}%
\begin{pgfscope}%
\pgfsys@transformshift{8.214083in}{2.032203in}%
\pgfsys@useobject{currentmarker}{}%
\end{pgfscope}%
\begin{pgfscope}%
\pgfsys@transformshift{8.200862in}{2.023789in}%
\pgfsys@useobject{currentmarker}{}%
\end{pgfscope}%
\begin{pgfscope}%
\pgfsys@transformshift{8.187147in}{2.015600in}%
\pgfsys@useobject{currentmarker}{}%
\end{pgfscope}%
\begin{pgfscope}%
\pgfsys@transformshift{8.172951in}{2.007645in}%
\pgfsys@useobject{currentmarker}{}%
\end{pgfscope}%
\begin{pgfscope}%
\pgfsys@transformshift{8.349706in}{2.177630in}%
\pgfsys@useobject{currentmarker}{}%
\end{pgfscope}%
\begin{pgfscope}%
\pgfsys@transformshift{8.345147in}{2.166690in}%
\pgfsys@useobject{currentmarker}{}%
\end{pgfscope}%
\begin{pgfscope}%
\pgfsys@transformshift{8.339951in}{2.155834in}%
\pgfsys@useobject{currentmarker}{}%
\end{pgfscope}%
\begin{pgfscope}%
\pgfsys@transformshift{8.334121in}{2.145072in}%
\pgfsys@useobject{currentmarker}{}%
\end{pgfscope}%
\begin{pgfscope}%
\pgfsys@transformshift{8.327664in}{2.134414in}%
\pgfsys@useobject{currentmarker}{}%
\end{pgfscope}%
\begin{pgfscope}%
\pgfsys@transformshift{8.320586in}{2.123871in}%
\pgfsys@useobject{currentmarker}{}%
\end{pgfscope}%
\begin{pgfscope}%
\pgfsys@transformshift{8.312895in}{2.113455in}%
\pgfsys@useobject{currentmarker}{}%
\end{pgfscope}%
\begin{pgfscope}%
\pgfsys@transformshift{8.304597in}{2.103174in}%
\pgfsys@useobject{currentmarker}{}%
\end{pgfscope}%
\begin{pgfscope}%
\pgfsys@transformshift{8.295702in}{2.093040in}%
\pgfsys@useobject{currentmarker}{}%
\end{pgfscope}%
\begin{pgfscope}%
\pgfsys@transformshift{8.286218in}{2.083062in}%
\pgfsys@useobject{currentmarker}{}%
\end{pgfscope}%
\begin{pgfscope}%
\pgfsys@transformshift{8.276154in}{2.073250in}%
\pgfsys@useobject{currentmarker}{}%
\end{pgfscope}%
\begin{pgfscope}%
\pgfsys@transformshift{8.265521in}{2.063615in}%
\pgfsys@useobject{currentmarker}{}%
\end{pgfscope}%
\begin{pgfscope}%
\pgfsys@transformshift{8.254330in}{2.054166in}%
\pgfsys@useobject{currentmarker}{}%
\end{pgfscope}%
\begin{pgfscope}%
\pgfsys@transformshift{8.242590in}{2.044911in}%
\pgfsys@useobject{currentmarker}{}%
\end{pgfscope}%
\begin{pgfscope}%
\pgfsys@transformshift{8.230315in}{2.035861in}%
\pgfsys@useobject{currentmarker}{}%
\end{pgfscope}%
\begin{pgfscope}%
\pgfsys@transformshift{8.217515in}{2.027025in}%
\pgfsys@useobject{currentmarker}{}%
\end{pgfscope}%
\begin{pgfscope}%
\pgfsys@transformshift{8.204205in}{2.018411in}%
\pgfsys@useobject{currentmarker}{}%
\end{pgfscope}%
\begin{pgfscope}%
\pgfsys@transformshift{8.190397in}{2.010028in}%
\pgfsys@useobject{currentmarker}{}%
\end{pgfscope}%
\begin{pgfscope}%
\pgfsys@transformshift{8.176105in}{2.001884in}%
\pgfsys@useobject{currentmarker}{}%
\end{pgfscope}%
\begin{pgfscope}%
\pgfsys@transformshift{8.354125in}{2.175938in}%
\pgfsys@useobject{currentmarker}{}%
\end{pgfscope}%
\begin{pgfscope}%
\pgfsys@transformshift{8.349535in}{2.164743in}%
\pgfsys@useobject{currentmarker}{}%
\end{pgfscope}%
\begin{pgfscope}%
\pgfsys@transformshift{8.344303in}{2.153633in}%
\pgfsys@useobject{currentmarker}{}%
\end{pgfscope}%
\begin{pgfscope}%
\pgfsys@transformshift{8.338433in}{2.142619in}%
\pgfsys@useobject{currentmarker}{}%
\end{pgfscope}%
\begin{pgfscope}%
\pgfsys@transformshift{8.331932in}{2.131712in}%
\pgfsys@useobject{currentmarker}{}%
\end{pgfscope}%
\begin{pgfscope}%
\pgfsys@transformshift{8.324805in}{2.120923in}%
\pgfsys@useobject{currentmarker}{}%
\end{pgfscope}%
\begin{pgfscope}%
\pgfsys@transformshift{8.317061in}{2.110262in}%
\pgfsys@useobject{currentmarker}{}%
\end{pgfscope}%
\begin{pgfscope}%
\pgfsys@transformshift{8.308707in}{2.099741in}%
\pgfsys@useobject{currentmarker}{}%
\end{pgfscope}%
\begin{pgfscope}%
\pgfsys@transformshift{8.299751in}{2.089370in}%
\pgfsys@useobject{currentmarker}{}%
\end{pgfscope}%
\begin{pgfscope}%
\pgfsys@transformshift{8.290202in}{2.079159in}%
\pgfsys@useobject{currentmarker}{}%
\end{pgfscope}%
\begin{pgfscope}%
\pgfsys@transformshift{8.280069in}{2.069118in}%
\pgfsys@useobject{currentmarker}{}%
\end{pgfscope}%
\begin{pgfscope}%
\pgfsys@transformshift{8.269363in}{2.059258in}%
\pgfsys@useobject{currentmarker}{}%
\end{pgfscope}%
\begin{pgfscope}%
\pgfsys@transformshift{8.258095in}{2.049587in}%
\pgfsys@useobject{currentmarker}{}%
\end{pgfscope}%
\begin{pgfscope}%
\pgfsys@transformshift{8.246275in}{2.040117in}%
\pgfsys@useobject{currentmarker}{}%
\end{pgfscope}%
\begin{pgfscope}%
\pgfsys@transformshift{8.233915in}{2.030855in}%
\pgfsys@useobject{currentmarker}{}%
\end{pgfscope}%
\begin{pgfscope}%
\pgfsys@transformshift{8.221028in}{2.021812in}%
\pgfsys@useobject{currentmarker}{}%
\end{pgfscope}%
\begin{pgfscope}%
\pgfsys@transformshift{8.207627in}{2.012997in}%
\pgfsys@useobject{currentmarker}{}%
\end{pgfscope}%
\begin{pgfscope}%
\pgfsys@transformshift{8.193724in}{2.004417in}%
\pgfsys@useobject{currentmarker}{}%
\end{pgfscope}%
\begin{pgfscope}%
\pgfsys@transformshift{8.179334in}{1.996083in}%
\pgfsys@useobject{currentmarker}{}%
\end{pgfscope}%
\begin{pgfscope}%
\pgfsys@transformshift{8.358646in}{2.174234in}%
\pgfsys@useobject{currentmarker}{}%
\end{pgfscope}%
\begin{pgfscope}%
\pgfsys@transformshift{8.354024in}{2.162781in}%
\pgfsys@useobject{currentmarker}{}%
\end{pgfscope}%
\begin{pgfscope}%
\pgfsys@transformshift{8.348755in}{2.151416in}%
\pgfsys@useobject{currentmarker}{}%
\end{pgfscope}%
\begin{pgfscope}%
\pgfsys@transformshift{8.342845in}{2.140149in}%
\pgfsys@useobject{currentmarker}{}%
\end{pgfscope}%
\begin{pgfscope}%
\pgfsys@transformshift{8.336298in}{2.128991in}%
\pgfsys@useobject{currentmarker}{}%
\end{pgfscope}%
\begin{pgfscope}%
\pgfsys@transformshift{8.329122in}{2.117954in}%
\pgfsys@useobject{currentmarker}{}%
\end{pgfscope}%
\begin{pgfscope}%
\pgfsys@transformshift{8.321324in}{2.107048in}%
\pgfsys@useobject{currentmarker}{}%
\end{pgfscope}%
\begin{pgfscope}%
\pgfsys@transformshift{8.312912in}{2.096285in}%
\pgfsys@useobject{currentmarker}{}%
\end{pgfscope}%
\begin{pgfscope}%
\pgfsys@transformshift{8.303893in}{2.085675in}%
\pgfsys@useobject{currentmarker}{}%
\end{pgfscope}%
\begin{pgfscope}%
\pgfsys@transformshift{8.294278in}{2.075229in}%
\pgfsys@useobject{currentmarker}{}%
\end{pgfscope}%
\begin{pgfscope}%
\pgfsys@transformshift{8.284074in}{2.064957in}%
\pgfsys@useobject{currentmarker}{}%
\end{pgfscope}%
\begin{pgfscope}%
\pgfsys@transformshift{8.273294in}{2.054870in}%
\pgfsys@useobject{currentmarker}{}%
\end{pgfscope}%
\begin{pgfscope}%
\pgfsys@transformshift{8.261947in}{2.044977in}%
\pgfsys@useobject{currentmarker}{}%
\end{pgfscope}%
\begin{pgfscope}%
\pgfsys@transformshift{8.250045in}{2.035289in}%
\pgfsys@useobject{currentmarker}{}%
\end{pgfscope}%
\begin{pgfscope}%
\pgfsys@transformshift{8.237599in}{2.025814in}%
\pgfsys@useobject{currentmarker}{}%
\end{pgfscope}%
\begin{pgfscope}%
\pgfsys@transformshift{8.224622in}{2.016563in}%
\pgfsys@useobject{currentmarker}{}%
\end{pgfscope}%
\begin{pgfscope}%
\pgfsys@transformshift{8.211128in}{2.007545in}%
\pgfsys@useobject{currentmarker}{}%
\end{pgfscope}%
\begin{pgfscope}%
\pgfsys@transformshift{8.197128in}{1.998768in}%
\pgfsys@useobject{currentmarker}{}%
\end{pgfscope}%
\begin{pgfscope}%
\pgfsys@transformshift{8.182638in}{1.990242in}%
\pgfsys@useobject{currentmarker}{}%
\end{pgfscope}%
\begin{pgfscope}%
\pgfsys@transformshift{8.363270in}{2.172518in}%
\pgfsys@useobject{currentmarker}{}%
\end{pgfscope}%
\begin{pgfscope}%
\pgfsys@transformshift{8.358616in}{2.160807in}%
\pgfsys@useobject{currentmarker}{}%
\end{pgfscope}%
\begin{pgfscope}%
\pgfsys@transformshift{8.353310in}{2.149184in}%
\pgfsys@useobject{currentmarker}{}%
\end{pgfscope}%
\begin{pgfscope}%
\pgfsys@transformshift{8.347357in}{2.137661in}%
\pgfsys@useobject{currentmarker}{}%
\end{pgfscope}%
\begin{pgfscope}%
\pgfsys@transformshift{8.340765in}{2.126251in}%
\pgfsys@useobject{currentmarker}{}%
\end{pgfscope}%
\begin{pgfscope}%
\pgfsys@transformshift{8.333538in}{2.114964in}%
\pgfsys@useobject{currentmarker}{}%
\end{pgfscope}%
\begin{pgfscope}%
\pgfsys@transformshift{8.325685in}{2.103811in}%
\pgfsys@useobject{currentmarker}{}%
\end{pgfscope}%
\begin{pgfscope}%
\pgfsys@transformshift{8.317213in}{2.092804in}%
\pgfsys@useobject{currentmarker}{}%
\end{pgfscope}%
\begin{pgfscope}%
\pgfsys@transformshift{8.308130in}{2.081954in}%
\pgfsys@useobject{currentmarker}{}%
\end{pgfscope}%
\begin{pgfscope}%
\pgfsys@transformshift{8.298447in}{2.071272in}%
\pgfsys@useobject{currentmarker}{}%
\end{pgfscope}%
\begin{pgfscope}%
\pgfsys@transformshift{8.288171in}{2.060767in}%
\pgfsys@useobject{currentmarker}{}%
\end{pgfscope}%
\begin{pgfscope}%
\pgfsys@transformshift{8.277315in}{2.050452in}%
\pgfsys@useobject{currentmarker}{}%
\end{pgfscope}%
\begin{pgfscope}%
\pgfsys@transformshift{8.265887in}{2.040335in}%
\pgfsys@useobject{currentmarker}{}%
\end{pgfscope}%
\begin{pgfscope}%
\pgfsys@transformshift{8.253901in}{2.030427in}%
\pgfsys@useobject{currentmarker}{}%
\end{pgfscope}%
\begin{pgfscope}%
\pgfsys@transformshift{8.241367in}{2.020738in}%
\pgfsys@useobject{currentmarker}{}%
\end{pgfscope}%
\begin{pgfscope}%
\pgfsys@transformshift{8.228299in}{2.011277in}%
\pgfsys@useobject{currentmarker}{}%
\end{pgfscope}%
\begin{pgfscope}%
\pgfsys@transformshift{8.214708in}{2.002055in}%
\pgfsys@useobject{currentmarker}{}%
\end{pgfscope}%
\begin{pgfscope}%
\pgfsys@transformshift{8.200610in}{1.993080in}%
\pgfsys@useobject{currentmarker}{}%
\end{pgfscope}%
\begin{pgfscope}%
\pgfsys@transformshift{8.186017in}{1.984360in}%
\pgfsys@useobject{currentmarker}{}%
\end{pgfscope}%
\begin{pgfscope}%
\pgfsys@transformshift{8.367998in}{2.170791in}%
\pgfsys@useobject{currentmarker}{}%
\end{pgfscope}%
\begin{pgfscope}%
\pgfsys@transformshift{8.363310in}{2.158818in}%
\pgfsys@useobject{currentmarker}{}%
\end{pgfscope}%
\begin{pgfscope}%
\pgfsys@transformshift{8.357966in}{2.146935in}%
\pgfsys@useobject{currentmarker}{}%
\end{pgfscope}%
\begin{pgfscope}%
\pgfsys@transformshift{8.351971in}{2.135156in}%
\pgfsys@useobject{currentmarker}{}%
\end{pgfscope}%
\begin{pgfscope}%
\pgfsys@transformshift{8.345331in}{2.123491in}%
\pgfsys@useobject{currentmarker}{}%
\end{pgfscope}%
\begin{pgfscope}%
\pgfsys@transformshift{8.338053in}{2.111952in}%
\pgfsys@useobject{currentmarker}{}%
\end{pgfscope}%
\begin{pgfscope}%
\pgfsys@transformshift{8.330143in}{2.100551in}%
\pgfsys@useobject{currentmarker}{}%
\end{pgfscope}%
\begin{pgfscope}%
\pgfsys@transformshift{8.321610in}{2.089299in}%
\pgfsys@useobject{currentmarker}{}%
\end{pgfscope}%
\begin{pgfscope}%
\pgfsys@transformshift{8.312463in}{2.078207in}%
\pgfsys@useobject{currentmarker}{}%
\end{pgfscope}%
\begin{pgfscope}%
\pgfsys@transformshift{8.302709in}{2.067286in}%
\pgfsys@useobject{currentmarker}{}%
\end{pgfscope}%
\begin{pgfscope}%
\pgfsys@transformshift{8.292360in}{2.056548in}%
\pgfsys@useobject{currentmarker}{}%
\end{pgfscope}%
\begin{pgfscope}%
\pgfsys@transformshift{8.281425in}{2.046002in}%
\pgfsys@useobject{currentmarker}{}%
\end{pgfscope}%
\begin{pgfscope}%
\pgfsys@transformshift{8.269916in}{2.035659in}%
\pgfsys@useobject{currentmarker}{}%
\end{pgfscope}%
\begin{pgfscope}%
\pgfsys@transformshift{8.257844in}{2.025530in}%
\pgfsys@useobject{currentmarker}{}%
\end{pgfscope}%
\begin{pgfscope}%
\pgfsys@transformshift{8.245220in}{2.015625in}%
\pgfsys@useobject{currentmarker}{}%
\end{pgfscope}%
\begin{pgfscope}%
\pgfsys@transformshift{8.232057in}{2.005954in}%
\pgfsys@useobject{currentmarker}{}%
\end{pgfscope}%
\begin{pgfscope}%
\pgfsys@transformshift{8.218370in}{1.996526in}%
\pgfsys@useobject{currentmarker}{}%
\end{pgfscope}%
\begin{pgfscope}%
\pgfsys@transformshift{8.204170in}{1.987350in}%
\pgfsys@useobject{currentmarker}{}%
\end{pgfscope}%
\begin{pgfscope}%
\pgfsys@transformshift{8.189473in}{1.978437in}%
\pgfsys@useobject{currentmarker}{}%
\end{pgfscope}%
\begin{pgfscope}%
\pgfsys@transformshift{8.372831in}{2.169050in}%
\pgfsys@useobject{currentmarker}{}%
\end{pgfscope}%
\begin{pgfscope}%
\pgfsys@transformshift{8.368109in}{2.156814in}%
\pgfsys@useobject{currentmarker}{}%
\end{pgfscope}%
\begin{pgfscope}%
\pgfsys@transformshift{8.362726in}{2.144671in}%
\pgfsys@useobject{currentmarker}{}%
\end{pgfscope}%
\begin{pgfscope}%
\pgfsys@transformshift{8.356687in}{2.132633in}%
\pgfsys@useobject{currentmarker}{}%
\end{pgfscope}%
\begin{pgfscope}%
\pgfsys@transformshift{8.349999in}{2.120712in}%
\pgfsys@useobject{currentmarker}{}%
\end{pgfscope}%
\begin{pgfscope}%
\pgfsys@transformshift{8.342667in}{2.108919in}%
\pgfsys@useobject{currentmarker}{}%
\end{pgfscope}%
\begin{pgfscope}%
\pgfsys@transformshift{8.334700in}{2.097268in}%
\pgfsys@useobject{currentmarker}{}%
\end{pgfscope}%
\begin{pgfscope}%
\pgfsys@transformshift{8.326105in}{2.085768in}%
\pgfsys@useobject{currentmarker}{}%
\end{pgfscope}%
\begin{pgfscope}%
\pgfsys@transformshift{8.316891in}{2.074432in}%
\pgfsys@useobject{currentmarker}{}%
\end{pgfscope}%
\begin{pgfscope}%
\pgfsys@transformshift{8.307066in}{2.063272in}%
\pgfsys@useobject{currentmarker}{}%
\end{pgfscope}%
\begin{pgfscope}%
\pgfsys@transformshift{8.296642in}{2.052297in}%
\pgfsys@useobject{currentmarker}{}%
\end{pgfscope}%
\begin{pgfscope}%
\pgfsys@transformshift{8.285627in}{2.041520in}%
\pgfsys@useobject{currentmarker}{}%
\end{pgfscope}%
\begin{pgfscope}%
\pgfsys@transformshift{8.274034in}{2.030950in}%
\pgfsys@useobject{currentmarker}{}%
\end{pgfscope}%
\begin{pgfscope}%
\pgfsys@transformshift{8.261874in}{2.020599in}%
\pgfsys@useobject{currentmarker}{}%
\end{pgfscope}%
\begin{pgfscope}%
\pgfsys@transformshift{8.249158in}{2.010476in}%
\pgfsys@useobject{currentmarker}{}%
\end{pgfscope}%
\begin{pgfscope}%
\pgfsys@transformshift{8.235899in}{2.000592in}%
\pgfsys@useobject{currentmarker}{}%
\end{pgfscope}%
\begin{pgfscope}%
\pgfsys@transformshift{8.222112in}{1.990957in}%
\pgfsys@useobject{currentmarker}{}%
\end{pgfscope}%
\begin{pgfscope}%
\pgfsys@transformshift{8.207809in}{1.981580in}%
\pgfsys@useobject{currentmarker}{}%
\end{pgfscope}%
\begin{pgfscope}%
\pgfsys@transformshift{8.193004in}{1.972470in}%
\pgfsys@useobject{currentmarker}{}%
\end{pgfscope}%
\begin{pgfscope}%
\pgfsys@transformshift{8.377769in}{2.167297in}%
\pgfsys@useobject{currentmarker}{}%
\end{pgfscope}%
\begin{pgfscope}%
\pgfsys@transformshift{8.373013in}{2.154796in}%
\pgfsys@useobject{currentmarker}{}%
\end{pgfscope}%
\begin{pgfscope}%
\pgfsys@transformshift{8.367590in}{2.142390in}%
\pgfsys@useobject{currentmarker}{}%
\end{pgfscope}%
\begin{pgfscope}%
\pgfsys@transformshift{8.361506in}{2.130091in}%
\pgfsys@useobject{currentmarker}{}%
\end{pgfscope}%
\begin{pgfscope}%
\pgfsys@transformshift{8.354768in}{2.117912in}%
\pgfsys@useobject{currentmarker}{}%
\end{pgfscope}%
\begin{pgfscope}%
\pgfsys@transformshift{8.347382in}{2.105864in}%
\pgfsys@useobject{currentmarker}{}%
\end{pgfscope}%
\begin{pgfscope}%
\pgfsys@transformshift{8.339356in}{2.093960in}%
\pgfsys@useobject{currentmarker}{}%
\end{pgfscope}%
\begin{pgfscope}%
\pgfsys@transformshift{8.330698in}{2.082212in}%
\pgfsys@useobject{currentmarker}{}%
\end{pgfscope}%
\begin{pgfscope}%
\pgfsys@transformshift{8.321415in}{2.070630in}%
\pgfsys@useobject{currentmarker}{}%
\end{pgfscope}%
\begin{pgfscope}%
\pgfsys@transformshift{8.311518in}{2.059228in}%
\pgfsys@useobject{currentmarker}{}%
\end{pgfscope}%
\begin{pgfscope}%
\pgfsys@transformshift{8.301017in}{2.048016in}%
\pgfsys@useobject{currentmarker}{}%
\end{pgfscope}%
\begin{pgfscope}%
\pgfsys@transformshift{8.289921in}{2.037005in}%
\pgfsys@useobject{currentmarker}{}%
\end{pgfscope}%
\begin{pgfscope}%
\pgfsys@transformshift{8.278242in}{2.026206in}%
\pgfsys@useobject{currentmarker}{}%
\end{pgfscope}%
\begin{pgfscope}%
\pgfsys@transformshift{8.265991in}{2.015631in}%
\pgfsys@useobject{currentmarker}{}%
\end{pgfscope}%
\begin{pgfscope}%
\pgfsys@transformshift{8.253182in}{2.005289in}%
\pgfsys@useobject{currentmarker}{}%
\end{pgfscope}%
\begin{pgfscope}%
\pgfsys@transformshift{8.239825in}{1.995191in}%
\pgfsys@useobject{currentmarker}{}%
\end{pgfscope}%
\begin{pgfscope}%
\pgfsys@transformshift{8.225936in}{1.985347in}%
\pgfsys@useobject{currentmarker}{}%
\end{pgfscope}%
\begin{pgfscope}%
\pgfsys@transformshift{8.211527in}{1.975767in}%
\pgfsys@useobject{currentmarker}{}%
\end{pgfscope}%
\begin{pgfscope}%
\pgfsys@transformshift{8.196613in}{1.966460in}%
\pgfsys@useobject{currentmarker}{}%
\end{pgfscope}%
\begin{pgfscope}%
\pgfsys@transformshift{8.382814in}{2.165531in}%
\pgfsys@useobject{currentmarker}{}%
\end{pgfscope}%
\begin{pgfscope}%
\pgfsys@transformshift{8.378021in}{2.152763in}%
\pgfsys@useobject{currentmarker}{}%
\end{pgfscope}%
\begin{pgfscope}%
\pgfsys@transformshift{8.372558in}{2.140092in}%
\pgfsys@useobject{currentmarker}{}%
\end{pgfscope}%
\begin{pgfscope}%
\pgfsys@transformshift{8.366429in}{2.127530in}%
\pgfsys@useobject{currentmarker}{}%
\end{pgfscope}%
\begin{pgfscope}%
\pgfsys@transformshift{8.359640in}{2.115091in}%
\pgfsys@useobject{currentmarker}{}%
\end{pgfscope}%
\begin{pgfscope}%
\pgfsys@transformshift{8.352199in}{2.102786in}%
\pgfsys@useobject{currentmarker}{}%
\end{pgfscope}%
\begin{pgfscope}%
\pgfsys@transformshift{8.344113in}{2.090628in}%
\pgfsys@useobject{currentmarker}{}%
\end{pgfscope}%
\begin{pgfscope}%
\pgfsys@transformshift{8.335389in}{2.078628in}%
\pgfsys@useobject{currentmarker}{}%
\end{pgfscope}%
\begin{pgfscope}%
\pgfsys@transformshift{8.326037in}{2.066800in}%
\pgfsys@useobject{currentmarker}{}%
\end{pgfscope}%
\begin{pgfscope}%
\pgfsys@transformshift{8.316066in}{2.055154in}%
\pgfsys@useobject{currentmarker}{}%
\end{pgfscope}%
\begin{pgfscope}%
\pgfsys@transformshift{8.305486in}{2.043702in}%
\pgfsys@useobject{currentmarker}{}%
\end{pgfscope}%
\begin{pgfscope}%
\pgfsys@transformshift{8.294307in}{2.032456in}%
\pgfsys@useobject{currentmarker}{}%
\end{pgfscope}%
\begin{pgfscope}%
\pgfsys@transformshift{8.282540in}{2.021427in}%
\pgfsys@useobject{currentmarker}{}%
\end{pgfscope}%
\begin{pgfscope}%
\pgfsys@transformshift{8.270198in}{2.010626in}%
\pgfsys@useobject{currentmarker}{}%
\end{pgfscope}%
\begin{pgfscope}%
\pgfsys@transformshift{8.257292in}{2.000063in}%
\pgfsys@useobject{currentmarker}{}%
\end{pgfscope}%
\begin{pgfscope}%
\pgfsys@transformshift{8.243836in}{1.989749in}%
\pgfsys@useobject{currentmarker}{}%
\end{pgfscope}%
\begin{pgfscope}%
\pgfsys@transformshift{8.229842in}{1.979695in}%
\pgfsys@useobject{currentmarker}{}%
\end{pgfscope}%
\begin{pgfscope}%
\pgfsys@transformshift{8.215325in}{1.969910in}%
\pgfsys@useobject{currentmarker}{}%
\end{pgfscope}%
\begin{pgfscope}%
\pgfsys@transformshift{8.200299in}{1.960405in}%
\pgfsys@useobject{currentmarker}{}%
\end{pgfscope}%
\begin{pgfscope}%
\pgfsys@transformshift{8.387965in}{2.163752in}%
\pgfsys@useobject{currentmarker}{}%
\end{pgfscope}%
\begin{pgfscope}%
\pgfsys@transformshift{8.383136in}{2.150715in}%
\pgfsys@useobject{currentmarker}{}%
\end{pgfscope}%
\begin{pgfscope}%
\pgfsys@transformshift{8.377631in}{2.137777in}%
\pgfsys@useobject{currentmarker}{}%
\end{pgfscope}%
\begin{pgfscope}%
\pgfsys@transformshift{8.371456in}{2.124950in}%
\pgfsys@useobject{currentmarker}{}%
\end{pgfscope}%
\begin{pgfscope}%
\pgfsys@transformshift{8.364616in}{2.112249in}%
\pgfsys@useobject{currentmarker}{}%
\end{pgfscope}%
\begin{pgfscope}%
\pgfsys@transformshift{8.357118in}{2.099685in}%
\pgfsys@useobject{currentmarker}{}%
\end{pgfscope}%
\begin{pgfscope}%
\pgfsys@transformshift{8.348970in}{2.087270in}%
\pgfsys@useobject{currentmarker}{}%
\end{pgfscope}%
\begin{pgfscope}%
\pgfsys@transformshift{8.340180in}{2.075018in}%
\pgfsys@useobject{currentmarker}{}%
\end{pgfscope}%
\begin{pgfscope}%
\pgfsys@transformshift{8.330757in}{2.062940in}%
\pgfsys@useobject{currentmarker}{}%
\end{pgfscope}%
\begin{pgfscope}%
\pgfsys@transformshift{8.320710in}{2.051049in}%
\pgfsys@useobject{currentmarker}{}%
\end{pgfscope}%
\begin{pgfscope}%
\pgfsys@transformshift{8.310050in}{2.039356in}%
\pgfsys@useobject{currentmarker}{}%
\end{pgfscope}%
\begin{pgfscope}%
\pgfsys@transformshift{8.298786in}{2.027873in}%
\pgfsys@useobject{currentmarker}{}%
\end{pgfscope}%
\begin{pgfscope}%
\pgfsys@transformshift{8.286930in}{2.016612in}%
\pgfsys@useobject{currentmarker}{}%
\end{pgfscope}%
\begin{pgfscope}%
\pgfsys@transformshift{8.274494in}{2.005583in}%
\pgfsys@useobject{currentmarker}{}%
\end{pgfscope}%
\begin{pgfscope}%
\pgfsys@transformshift{8.261490in}{1.994797in}%
\pgfsys@useobject{currentmarker}{}%
\end{pgfscope}%
\begin{pgfscope}%
\pgfsys@transformshift{8.247931in}{1.984267in}%
\pgfsys@useobject{currentmarker}{}%
\end{pgfscope}%
\begin{pgfscope}%
\pgfsys@transformshift{8.233831in}{1.974001in}%
\pgfsys@useobject{currentmarker}{}%
\end{pgfscope}%
\begin{pgfscope}%
\pgfsys@transformshift{8.219204in}{1.964010in}%
\pgfsys@useobject{currentmarker}{}%
\end{pgfscope}%
\begin{pgfscope}%
\pgfsys@transformshift{8.204064in}{1.954304in}%
\pgfsys@useobject{currentmarker}{}%
\end{pgfscope}%
\begin{pgfscope}%
\pgfsys@transformshift{8.393225in}{2.161958in}%
\pgfsys@useobject{currentmarker}{}%
\end{pgfscope}%
\begin{pgfscope}%
\pgfsys@transformshift{8.388358in}{2.148650in}%
\pgfsys@useobject{currentmarker}{}%
\end{pgfscope}%
\begin{pgfscope}%
\pgfsys@transformshift{8.382811in}{2.135443in}%
\pgfsys@useobject{currentmarker}{}%
\end{pgfscope}%
\begin{pgfscope}%
\pgfsys@transformshift{8.376588in}{2.122351in}%
\pgfsys@useobject{currentmarker}{}%
\end{pgfscope}%
\begin{pgfscope}%
\pgfsys@transformshift{8.369695in}{2.109385in}%
\pgfsys@useobject{currentmarker}{}%
\end{pgfscope}%
\begin{pgfscope}%
\pgfsys@transformshift{8.362140in}{2.096560in}%
\pgfsys@useobject{currentmarker}{}%
\end{pgfscope}%
\begin{pgfscope}%
\pgfsys@transformshift{8.353929in}{2.083887in}%
\pgfsys@useobject{currentmarker}{}%
\end{pgfscope}%
\begin{pgfscope}%
\pgfsys@transformshift{8.345072in}{2.071380in}%
\pgfsys@useobject{currentmarker}{}%
\end{pgfscope}%
\begin{pgfscope}%
\pgfsys@transformshift{8.335576in}{2.059051in}%
\pgfsys@useobject{currentmarker}{}%
\end{pgfscope}%
\begin{pgfscope}%
\pgfsys@transformshift{8.325452in}{2.046913in}%
\pgfsys@useobject{currentmarker}{}%
\end{pgfscope}%
\begin{pgfscope}%
\pgfsys@transformshift{8.314709in}{2.034977in}%
\pgfsys@useobject{currentmarker}{}%
\end{pgfscope}%
\begin{pgfscope}%
\pgfsys@transformshift{8.303358in}{2.023255in}%
\pgfsys@useobject{currentmarker}{}%
\end{pgfscope}%
\begin{pgfscope}%
\pgfsys@transformshift{8.291411in}{2.011760in}%
\pgfsys@useobject{currentmarker}{}%
\end{pgfscope}%
\begin{pgfscope}%
\pgfsys@transformshift{8.278879in}{2.000501in}%
\pgfsys@useobject{currentmarker}{}%
\end{pgfscope}%
\begin{pgfscope}%
\pgfsys@transformshift{8.265775in}{1.989492in}%
\pgfsys@useobject{currentmarker}{}%
\end{pgfscope}%
\begin{pgfscope}%
\pgfsys@transformshift{8.252112in}{1.978742in}%
\pgfsys@useobject{currentmarker}{}%
\end{pgfscope}%
\begin{pgfscope}%
\pgfsys@transformshift{8.237904in}{1.968263in}%
\pgfsys@useobject{currentmarker}{}%
\end{pgfscope}%
\begin{pgfscope}%
\pgfsys@transformshift{8.223164in}{1.958064in}%
\pgfsys@useobject{currentmarker}{}%
\end{pgfscope}%
\begin{pgfscope}%
\pgfsys@transformshift{8.207907in}{1.948157in}%
\pgfsys@useobject{currentmarker}{}%
\end{pgfscope}%
\begin{pgfscope}%
\pgfsys@transformshift{8.398593in}{2.160152in}%
\pgfsys@useobject{currentmarker}{}%
\end{pgfscope}%
\begin{pgfscope}%
\pgfsys@transformshift{8.393688in}{2.146570in}%
\pgfsys@useobject{currentmarker}{}%
\end{pgfscope}%
\begin{pgfscope}%
\pgfsys@transformshift{8.388098in}{2.133092in}%
\pgfsys@useobject{currentmarker}{}%
\end{pgfscope}%
\begin{pgfscope}%
\pgfsys@transformshift{8.381826in}{2.119731in}%
\pgfsys@useobject{currentmarker}{}%
\end{pgfscope}%
\begin{pgfscope}%
\pgfsys@transformshift{8.374879in}{2.106499in}%
\pgfsys@useobject{currentmarker}{}%
\end{pgfscope}%
\begin{pgfscope}%
\pgfsys@transformshift{8.367265in}{2.093410in}%
\pgfsys@useobject{currentmarker}{}%
\end{pgfscope}%
\begin{pgfscope}%
\pgfsys@transformshift{8.358991in}{2.080478in}%
\pgfsys@useobject{currentmarker}{}%
\end{pgfscope}%
\begin{pgfscope}%
\pgfsys@transformshift{8.350064in}{2.067714in}%
\pgfsys@useobject{currentmarker}{}%
\end{pgfscope}%
\begin{pgfscope}%
\pgfsys@transformshift{8.340494in}{2.055132in}%
\pgfsys@useobject{currentmarker}{}%
\end{pgfscope}%
\begin{pgfscope}%
\pgfsys@transformshift{8.330291in}{2.042745in}%
\pgfsys@useobject{currentmarker}{}%
\end{pgfscope}%
\begin{pgfscope}%
\pgfsys@transformshift{8.319465in}{2.030564in}%
\pgfsys@useobject{currentmarker}{}%
\end{pgfscope}%
\begin{pgfscope}%
\pgfsys@transformshift{8.308025in}{2.018601in}%
\pgfsys@useobject{currentmarker}{}%
\end{pgfscope}%
\begin{pgfscope}%
\pgfsys@transformshift{8.295985in}{2.006870in}%
\pgfsys@useobject{currentmarker}{}%
\end{pgfscope}%
\begin{pgfscope}%
\pgfsys@transformshift{8.283355in}{1.995381in}%
\pgfsys@useobject{currentmarker}{}%
\end{pgfscope}%
\begin{pgfscope}%
\pgfsys@transformshift{8.270149in}{1.984145in}%
\pgfsys@useobject{currentmarker}{}%
\end{pgfscope}%
\begin{pgfscope}%
\pgfsys@transformshift{8.256379in}{1.973175in}%
\pgfsys@useobject{currentmarker}{}%
\end{pgfscope}%
\begin{pgfscope}%
\pgfsys@transformshift{8.242060in}{1.962480in}%
\pgfsys@useobject{currentmarker}{}%
\end{pgfscope}%
\begin{pgfscope}%
\pgfsys@transformshift{8.227205in}{1.952072in}%
\pgfsys@useobject{currentmarker}{}%
\end{pgfscope}%
\begin{pgfscope}%
\pgfsys@transformshift{8.211830in}{1.941961in}%
\pgfsys@useobject{currentmarker}{}%
\end{pgfscope}%
\begin{pgfscope}%
\pgfsys@transformshift{8.404070in}{2.158330in}%
\pgfsys@useobject{currentmarker}{}%
\end{pgfscope}%
\begin{pgfscope}%
\pgfsys@transformshift{8.399127in}{2.144474in}%
\pgfsys@useobject{currentmarker}{}%
\end{pgfscope}%
\begin{pgfscope}%
\pgfsys@transformshift{8.393492in}{2.130723in}%
\pgfsys@useobject{currentmarker}{}%
\end{pgfscope}%
\begin{pgfscope}%
\pgfsys@transformshift{8.387171in}{2.117090in}%
\pgfsys@useobject{currentmarker}{}%
\end{pgfscope}%
\begin{pgfscope}%
\pgfsys@transformshift{8.380170in}{2.103590in}%
\pgfsys@useobject{currentmarker}{}%
\end{pgfscope}%
\begin{pgfscope}%
\pgfsys@transformshift{8.372495in}{2.090236in}%
\pgfsys@useobject{currentmarker}{}%
\end{pgfscope}%
\begin{pgfscope}%
\pgfsys@transformshift{8.364155in}{2.077042in}%
\pgfsys@useobject{currentmarker}{}%
\end{pgfscope}%
\begin{pgfscope}%
\pgfsys@transformshift{8.355158in}{2.064019in}%
\pgfsys@useobject{currentmarker}{}%
\end{pgfscope}%
\begin{pgfscope}%
\pgfsys@transformshift{8.345513in}{2.051183in}%
\pgfsys@useobject{currentmarker}{}%
\end{pgfscope}%
\begin{pgfscope}%
\pgfsys@transformshift{8.335229in}{2.038544in}%
\pgfsys@useobject{currentmarker}{}%
\end{pgfscope}%
\begin{pgfscope}%
\pgfsys@transformshift{8.324317in}{2.026116in}%
\pgfsys@useobject{currentmarker}{}%
\end{pgfscope}%
\begin{pgfscope}%
\pgfsys@transformshift{8.312788in}{2.013911in}%
\pgfsys@useobject{currentmarker}{}%
\end{pgfscope}%
\begin{pgfscope}%
\pgfsys@transformshift{8.300652in}{2.001942in}%
\pgfsys@useobject{currentmarker}{}%
\end{pgfscope}%
\begin{pgfscope}%
\pgfsys@transformshift{8.287923in}{1.990220in}%
\pgfsys@useobject{currentmarker}{}%
\end{pgfscope}%
\begin{pgfscope}%
\pgfsys@transformshift{8.274612in}{1.978756in}%
\pgfsys@useobject{currentmarker}{}%
\end{pgfscope}%
\begin{pgfscope}%
\pgfsys@transformshift{8.260734in}{1.967564in}%
\pgfsys@useobject{currentmarker}{}%
\end{pgfscope}%
\begin{pgfscope}%
\pgfsys@transformshift{8.246301in}{1.956652in}%
\pgfsys@useobject{currentmarker}{}%
\end{pgfscope}%
\begin{pgfscope}%
\pgfsys@transformshift{8.231329in}{1.946034in}%
\pgfsys@useobject{currentmarker}{}%
\end{pgfscope}%
\begin{pgfscope}%
\pgfsys@transformshift{8.215833in}{1.935718in}%
\pgfsys@useobject{currentmarker}{}%
\end{pgfscope}%
\begin{pgfscope}%
\pgfsys@transformshift{8.409658in}{2.156495in}%
\pgfsys@useobject{currentmarker}{}%
\end{pgfscope}%
\begin{pgfscope}%
\pgfsys@transformshift{8.404675in}{2.142361in}%
\pgfsys@useobject{currentmarker}{}%
\end{pgfscope}%
\begin{pgfscope}%
\pgfsys@transformshift{8.398995in}{2.128334in}%
\pgfsys@useobject{currentmarker}{}%
\end{pgfscope}%
\begin{pgfscope}%
\pgfsys@transformshift{8.392624in}{2.114429in}%
\pgfsys@useobject{currentmarker}{}%
\end{pgfscope}%
\begin{pgfscope}%
\pgfsys@transformshift{8.385566in}{2.100659in}%
\pgfsys@useobject{currentmarker}{}%
\end{pgfscope}%
\begin{pgfscope}%
\pgfsys@transformshift{8.377831in}{2.087037in}%
\pgfsys@useobject{currentmarker}{}%
\end{pgfscope}%
\begin{pgfscope}%
\pgfsys@transformshift{8.369424in}{2.073578in}%
\pgfsys@useobject{currentmarker}{}%
\end{pgfscope}%
\begin{pgfscope}%
\pgfsys@transformshift{8.360355in}{2.060295in}%
\pgfsys@useobject{currentmarker}{}%
\end{pgfscope}%
\begin{pgfscope}%
\pgfsys@transformshift{8.350633in}{2.047201in}%
\pgfsys@useobject{currentmarker}{}%
\end{pgfscope}%
\begin{pgfscope}%
\pgfsys@transformshift{8.340267in}{2.034310in}%
\pgfsys@useobject{currentmarker}{}%
\end{pgfscope}%
\begin{pgfscope}%
\pgfsys@transformshift{8.329267in}{2.021633in}%
\pgfsys@useobject{currentmarker}{}%
\end{pgfscope}%
\begin{pgfscope}%
\pgfsys@transformshift{8.317646in}{2.009184in}%
\pgfsys@useobject{currentmarker}{}%
\end{pgfscope}%
\begin{pgfscope}%
\pgfsys@transformshift{8.305413in}{1.996975in}%
\pgfsys@useobject{currentmarker}{}%
\end{pgfscope}%
\begin{pgfscope}%
\pgfsys@transformshift{8.292582in}{1.985018in}%
\pgfsys@useobject{currentmarker}{}%
\end{pgfscope}%
\begin{pgfscope}%
\pgfsys@transformshift{8.279165in}{1.973325in}%
\pgfsys@useobject{currentmarker}{}%
\end{pgfscope}%
\begin{pgfscope}%
\pgfsys@transformshift{8.265176in}{1.961908in}%
\pgfsys@useobject{currentmarker}{}%
\end{pgfscope}%
\begin{pgfscope}%
\pgfsys@transformshift{8.250628in}{1.950778in}%
\pgfsys@useobject{currentmarker}{}%
\end{pgfscope}%
\begin{pgfscope}%
\pgfsys@transformshift{8.235537in}{1.939947in}%
\pgfsys@useobject{currentmarker}{}%
\end{pgfscope}%
\begin{pgfscope}%
\pgfsys@transformshift{8.219916in}{1.929424in}%
\pgfsys@useobject{currentmarker}{}%
\end{pgfscope}%
\begin{pgfscope}%
\pgfsys@transformshift{8.415357in}{2.154644in}%
\pgfsys@useobject{currentmarker}{}%
\end{pgfscope}%
\begin{pgfscope}%
\pgfsys@transformshift{8.410334in}{2.140231in}%
\pgfsys@useobject{currentmarker}{}%
\end{pgfscope}%
\begin{pgfscope}%
\pgfsys@transformshift{8.404608in}{2.125926in}%
\pgfsys@useobject{currentmarker}{}%
\end{pgfscope}%
\begin{pgfscope}%
\pgfsys@transformshift{8.398185in}{2.111746in}%
\pgfsys@useobject{currentmarker}{}%
\end{pgfscope}%
\begin{pgfscope}%
\pgfsys@transformshift{8.391070in}{2.097703in}%
\pgfsys@useobject{currentmarker}{}%
\end{pgfscope}%
\begin{pgfscope}%
\pgfsys@transformshift{8.383272in}{2.083812in}%
\pgfsys@useobject{currentmarker}{}%
\end{pgfscope}%
\begin{pgfscope}%
\pgfsys@transformshift{8.374798in}{2.070087in}%
\pgfsys@useobject{currentmarker}{}%
\end{pgfscope}%
\begin{pgfscope}%
\pgfsys@transformshift{8.365655in}{2.056541in}%
\pgfsys@useobject{currentmarker}{}%
\end{pgfscope}%
\begin{pgfscope}%
\pgfsys@transformshift{8.355854in}{2.043188in}%
\pgfsys@useobject{currentmarker}{}%
\end{pgfscope}%
\begin{pgfscope}%
\pgfsys@transformshift{8.345405in}{2.030041in}%
\pgfsys@useobject{currentmarker}{}%
\end{pgfscope}%
\begin{pgfscope}%
\pgfsys@transformshift{8.334316in}{2.017114in}%
\pgfsys@useobject{currentmarker}{}%
\end{pgfscope}%
\begin{pgfscope}%
\pgfsys@transformshift{8.322601in}{2.004418in}%
\pgfsys@useobject{currentmarker}{}%
\end{pgfscope}%
\begin{pgfscope}%
\pgfsys@transformshift{8.310269in}{1.991967in}%
\pgfsys@useobject{currentmarker}{}%
\end{pgfscope}%
\begin{pgfscope}%
\pgfsys@transformshift{8.297334in}{1.979774in}%
\pgfsys@useobject{currentmarker}{}%
\end{pgfscope}%
\begin{pgfscope}%
\pgfsys@transformshift{8.283809in}{1.967850in}%
\pgfsys@useobject{currentmarker}{}%
\end{pgfscope}%
\begin{pgfscope}%
\pgfsys@transformshift{8.269707in}{1.956207in}%
\pgfsys@useobject{currentmarker}{}%
\end{pgfscope}%
\begin{pgfscope}%
\pgfsys@transformshift{8.255041in}{1.944857in}%
\pgfsys@useobject{currentmarker}{}%
\end{pgfscope}%
\begin{pgfscope}%
\pgfsys@transformshift{8.239827in}{1.933811in}%
\pgfsys@useobject{currentmarker}{}%
\end{pgfscope}%
\begin{pgfscope}%
\pgfsys@transformshift{8.224080in}{1.923080in}%
\pgfsys@useobject{currentmarker}{}%
\end{pgfscope}%
\begin{pgfscope}%
\pgfsys@transformshift{8.421168in}{2.152779in}%
\pgfsys@useobject{currentmarker}{}%
\end{pgfscope}%
\begin{pgfscope}%
\pgfsys@transformshift{8.416104in}{2.138083in}%
\pgfsys@useobject{currentmarker}{}%
\end{pgfscope}%
\begin{pgfscope}%
\pgfsys@transformshift{8.410331in}{2.123499in}%
\pgfsys@useobject{currentmarker}{}%
\end{pgfscope}%
\begin{pgfscope}%
\pgfsys@transformshift{8.403855in}{2.109041in}%
\pgfsys@useobject{currentmarker}{}%
\end{pgfscope}%
\begin{pgfscope}%
\pgfsys@transformshift{8.396683in}{2.094724in}%
\pgfsys@useobject{currentmarker}{}%
\end{pgfscope}%
\begin{pgfscope}%
\pgfsys@transformshift{8.388821in}{2.080561in}%
\pgfsys@useobject{currentmarker}{}%
\end{pgfscope}%
\begin{pgfscope}%
\pgfsys@transformshift{8.380277in}{2.066567in}%
\pgfsys@useobject{currentmarker}{}%
\end{pgfscope}%
\begin{pgfscope}%
\pgfsys@transformshift{8.371060in}{2.052757in}%
\pgfsys@useobject{currentmarker}{}%
\end{pgfscope}%
\begin{pgfscope}%
\pgfsys@transformshift{8.361179in}{2.039142in}%
\pgfsys@useobject{currentmarker}{}%
\end{pgfscope}%
\begin{pgfscope}%
\pgfsys@transformshift{8.350643in}{2.025738in}%
\pgfsys@useobject{currentmarker}{}%
\end{pgfscope}%
\begin{pgfscope}%
\pgfsys@transformshift{8.339465in}{2.012558in}%
\pgfsys@useobject{currentmarker}{}%
\end{pgfscope}%
\begin{pgfscope}%
\pgfsys@transformshift{8.327653in}{1.999614in}%
\pgfsys@useobject{currentmarker}{}%
\end{pgfscope}%
\begin{pgfscope}%
\pgfsys@transformshift{8.315221in}{1.986920in}%
\pgfsys@useobject{currentmarker}{}%
\end{pgfscope}%
\begin{pgfscope}%
\pgfsys@transformshift{8.302180in}{1.974487in}%
\pgfsys@useobject{currentmarker}{}%
\end{pgfscope}%
\begin{pgfscope}%
\pgfsys@transformshift{8.288544in}{1.962330in}%
\pgfsys@useobject{currentmarker}{}%
\end{pgfscope}%
\begin{pgfscope}%
\pgfsys@transformshift{8.274326in}{1.950459in}%
\pgfsys@useobject{currentmarker}{}%
\end{pgfscope}%
\begin{pgfscope}%
\pgfsys@transformshift{8.259541in}{1.938887in}%
\pgfsys@useobject{currentmarker}{}%
\end{pgfscope}%
\begin{pgfscope}%
\pgfsys@transformshift{8.244203in}{1.927625in}%
\pgfsys@useobject{currentmarker}{}%
\end{pgfscope}%
\begin{pgfscope}%
\pgfsys@transformshift{8.228327in}{1.916685in}%
\pgfsys@useobject{currentmarker}{}%
\end{pgfscope}%
\begin{pgfscope}%
\pgfsys@transformshift{8.427092in}{2.150898in}%
\pgfsys@useobject{currentmarker}{}%
\end{pgfscope}%
\begin{pgfscope}%
\pgfsys@transformshift{8.421986in}{2.135918in}%
\pgfsys@useobject{currentmarker}{}%
\end{pgfscope}%
\begin{pgfscope}%
\pgfsys@transformshift{8.416166in}{2.121052in}%
\pgfsys@useobject{currentmarker}{}%
\end{pgfscope}%
\begin{pgfscope}%
\pgfsys@transformshift{8.409636in}{2.106314in}%
\pgfsys@useobject{currentmarker}{}%
\end{pgfscope}%
\begin{pgfscope}%
\pgfsys@transformshift{8.402405in}{2.091720in}%
\pgfsys@useobject{currentmarker}{}%
\end{pgfscope}%
\begin{pgfscope}%
\pgfsys@transformshift{8.394477in}{2.077283in}%
\pgfsys@useobject{currentmarker}{}%
\end{pgfscope}%
\begin{pgfscope}%
\pgfsys@transformshift{8.385863in}{2.063019in}%
\pgfsys@useobject{currentmarker}{}%
\end{pgfscope}%
\begin{pgfscope}%
\pgfsys@transformshift{8.376570in}{2.048941in}%
\pgfsys@useobject{currentmarker}{}%
\end{pgfscope}%
\begin{pgfscope}%
\pgfsys@transformshift{8.366607in}{2.035063in}%
\pgfsys@useobject{currentmarker}{}%
\end{pgfscope}%
\begin{pgfscope}%
\pgfsys@transformshift{8.355984in}{2.021400in}%
\pgfsys@useobject{currentmarker}{}%
\end{pgfscope}%
\begin{pgfscope}%
\pgfsys@transformshift{8.344713in}{2.007964in}%
\pgfsys@useobject{currentmarker}{}%
\end{pgfscope}%
\begin{pgfscope}%
\pgfsys@transformshift{8.332804in}{1.994770in}%
\pgfsys@useobject{currentmarker}{}%
\end{pgfscope}%
\begin{pgfscope}%
\pgfsys@transformshift{8.320269in}{1.981830in}%
\pgfsys@useobject{currentmarker}{}%
\end{pgfscope}%
\begin{pgfscope}%
\pgfsys@transformshift{8.307120in}{1.969158in}%
\pgfsys@useobject{currentmarker}{}%
\end{pgfscope}%
\begin{pgfscope}%
\pgfsys@transformshift{8.293371in}{1.956765in}%
\pgfsys@useobject{currentmarker}{}%
\end{pgfscope}%
\begin{pgfscope}%
\pgfsys@transformshift{8.279036in}{1.944665in}%
\pgfsys@useobject{currentmarker}{}%
\end{pgfscope}%
\begin{pgfscope}%
\pgfsys@transformshift{8.264128in}{1.932869in}%
\pgfsys@useobject{currentmarker}{}%
\end{pgfscope}%
\begin{pgfscope}%
\pgfsys@transformshift{8.248663in}{1.921389in}%
\pgfsys@useobject{currentmarker}{}%
\end{pgfscope}%
\begin{pgfscope}%
\pgfsys@transformshift{8.232656in}{1.910237in}%
\pgfsys@useobject{currentmarker}{}%
\end{pgfscope}%
\begin{pgfscope}%
\pgfsys@transformshift{8.433130in}{2.149002in}%
\pgfsys@useobject{currentmarker}{}%
\end{pgfscope}%
\begin{pgfscope}%
\pgfsys@transformshift{8.427982in}{2.133735in}%
\pgfsys@useobject{currentmarker}{}%
\end{pgfscope}%
\begin{pgfscope}%
\pgfsys@transformshift{8.422113in}{2.118584in}%
\pgfsys@useobject{currentmarker}{}%
\end{pgfscope}%
\begin{pgfscope}%
\pgfsys@transformshift{8.415529in}{2.103565in}%
\pgfsys@useobject{currentmarker}{}%
\end{pgfscope}%
\begin{pgfscope}%
\pgfsys@transformshift{8.408236in}{2.088691in}%
\pgfsys@useobject{currentmarker}{}%
\end{pgfscope}%
\begin{pgfscope}%
\pgfsys@transformshift{8.400243in}{2.073978in}%
\pgfsys@useobject{currentmarker}{}%
\end{pgfscope}%
\begin{pgfscope}%
\pgfsys@transformshift{8.391556in}{2.059441in}%
\pgfsys@useobject{currentmarker}{}%
\end{pgfscope}%
\begin{pgfscope}%
\pgfsys@transformshift{8.439283in}{2.147089in}%
\pgfsys@useobject{currentmarker}{}%
\end{pgfscope}%
\begin{pgfscope}%
\pgfsys@transformshift{8.434091in}{2.131534in}%
\pgfsys@useobject{currentmarker}{}%
\end{pgfscope}%
\begin{pgfscope}%
\pgfsys@transformshift{8.428173in}{2.116096in}%
\pgfsys@useobject{currentmarker}{}%
\end{pgfscope}%
\begin{pgfscope}%
\pgfsys@transformshift{8.421533in}{2.100792in}%
\pgfsys@useobject{currentmarker}{}%
\end{pgfscope}%
\begin{pgfscope}%
\pgfsys@transformshift{8.414179in}{2.085637in}%
\pgfsys@useobject{currentmarker}{}%
\end{pgfscope}%
\begin{pgfscope}%
\pgfsys@transformshift{8.406118in}{2.070645in}%
\pgfsys@useobject{currentmarker}{}%
\end{pgfscope}%
\begin{pgfscope}%
\pgfsys@transformshift{8.397358in}{2.055833in}%
\pgfsys@useobject{currentmarker}{}%
\end{pgfscope}%
\begin{pgfscope}%
\pgfsys@transformshift{8.445552in}{2.145161in}%
\pgfsys@useobject{currentmarker}{}%
\end{pgfscope}%
\begin{pgfscope}%
\pgfsys@transformshift{8.440316in}{2.129314in}%
\pgfsys@useobject{currentmarker}{}%
\end{pgfscope}%
\begin{pgfscope}%
\pgfsys@transformshift{8.434347in}{2.113587in}%
\pgfsys@useobject{currentmarker}{}%
\end{pgfscope}%
\begin{pgfscope}%
\pgfsys@transformshift{8.427650in}{2.097996in}%
\pgfsys@useobject{currentmarker}{}%
\end{pgfscope}%
\begin{pgfscope}%
\pgfsys@transformshift{8.420234in}{2.082557in}%
\pgfsys@useobject{currentmarker}{}%
\end{pgfscope}%
\begin{pgfscope}%
\pgfsys@transformshift{8.412104in}{2.067284in}%
\pgfsys@useobject{currentmarker}{}%
\end{pgfscope}%
\begin{pgfscope}%
\pgfsys@transformshift{8.403269in}{2.052194in}%
\pgfsys@useobject{currentmarker}{}%
\end{pgfscope}%
\begin{pgfscope}%
\pgfsys@transformshift{8.451938in}{2.143216in}%
\pgfsys@useobject{currentmarker}{}%
\end{pgfscope}%
\begin{pgfscope}%
\pgfsys@transformshift{8.446656in}{2.127074in}%
\pgfsys@useobject{currentmarker}{}%
\end{pgfscope}%
\begin{pgfscope}%
\pgfsys@transformshift{8.440636in}{2.111056in}%
\pgfsys@useobject{currentmarker}{}%
\end{pgfscope}%
\begin{pgfscope}%
\pgfsys@transformshift{8.433882in}{2.095176in}%
\pgfsys@useobject{currentmarker}{}%
\end{pgfscope}%
\begin{pgfscope}%
\pgfsys@transformshift{8.426401in}{2.079450in}%
\pgfsys@useobject{currentmarker}{}%
\end{pgfscope}%
\begin{pgfscope}%
\pgfsys@transformshift{8.418201in}{2.063894in}%
\pgfsys@useobject{currentmarker}{}%
\end{pgfscope}%
\begin{pgfscope}%
\pgfsys@transformshift{8.409290in}{2.048524in}%
\pgfsys@useobject{currentmarker}{}%
\end{pgfscope}%
\begin{pgfscope}%
\pgfsys@transformshift{8.458441in}{2.141254in}%
\pgfsys@useobject{currentmarker}{}%
\end{pgfscope}%
\begin{pgfscope}%
\pgfsys@transformshift{8.453114in}{2.124816in}%
\pgfsys@useobject{currentmarker}{}%
\end{pgfscope}%
\begin{pgfscope}%
\pgfsys@transformshift{8.447041in}{2.108503in}%
\pgfsys@useobject{currentmarker}{}%
\end{pgfscope}%
\begin{pgfscope}%
\pgfsys@transformshift{8.440228in}{2.092331in}%
\pgfsys@useobject{currentmarker}{}%
\end{pgfscope}%
\begin{pgfscope}%
\pgfsys@transformshift{8.432683in}{2.076316in}%
\pgfsys@useobject{currentmarker}{}%
\end{pgfscope}%
\begin{pgfscope}%
\pgfsys@transformshift{8.424411in}{2.060475in}%
\pgfsys@useobject{currentmarker}{}%
\end{pgfscope}%
\begin{pgfscope}%
\pgfsys@transformshift{8.415423in}{2.044822in}%
\pgfsys@useobject{currentmarker}{}%
\end{pgfscope}%
\begin{pgfscope}%
\pgfsys@transformshift{8.465064in}{2.139274in}%
\pgfsys@useobject{currentmarker}{}%
\end{pgfscope}%
\begin{pgfscope}%
\pgfsys@transformshift{8.459690in}{2.122537in}%
\pgfsys@useobject{currentmarker}{}%
\end{pgfscope}%
\begin{pgfscope}%
\pgfsys@transformshift{8.453564in}{2.105928in}%
\pgfsys@useobject{currentmarker}{}%
\end{pgfscope}%
\begin{pgfscope}%
\pgfsys@transformshift{8.446691in}{2.089461in}%
\pgfsys@useobject{currentmarker}{}%
\end{pgfscope}%
\begin{pgfscope}%
\pgfsys@transformshift{8.439079in}{2.073155in}%
\pgfsys@useobject{currentmarker}{}%
\end{pgfscope}%
\begin{pgfscope}%
\pgfsys@transformshift{8.430735in}{2.057025in}%
\pgfsys@useobject{currentmarker}{}%
\end{pgfscope}%
\begin{pgfscope}%
\pgfsys@transformshift{8.421667in}{2.041088in}%
\pgfsys@useobject{currentmarker}{}%
\end{pgfscope}%
\begin{pgfscope}%
\pgfsys@transformshift{8.471807in}{2.137278in}%
\pgfsys@useobject{currentmarker}{}%
\end{pgfscope}%
\begin{pgfscope}%
\pgfsys@transformshift{8.466385in}{2.120239in}%
\pgfsys@useobject{currentmarker}{}%
\end{pgfscope}%
\begin{pgfscope}%
\pgfsys@transformshift{8.460204in}{2.103330in}%
\pgfsys@useobject{currentmarker}{}%
\end{pgfscope}%
\begin{pgfscope}%
\pgfsys@transformshift{8.453271in}{2.086567in}%
\pgfsys@useobject{currentmarker}{}%
\end{pgfscope}%
\begin{pgfscope}%
\pgfsys@transformshift{8.445591in}{2.069966in}%
\pgfsys@useobject{currentmarker}{}%
\end{pgfscope}%
\begin{pgfscope}%
\pgfsys@transformshift{8.437173in}{2.053546in}%
\pgfsys@useobject{currentmarker}{}%
\end{pgfscope}%
\begin{pgfscope}%
\pgfsys@transformshift{8.428025in}{2.037321in}%
\pgfsys@useobject{currentmarker}{}%
\end{pgfscope}%
\begin{pgfscope}%
\pgfsys@transformshift{8.478670in}{2.135263in}%
\pgfsys@useobject{currentmarker}{}%
\end{pgfscope}%
\begin{pgfscope}%
\pgfsys@transformshift{8.473200in}{2.117920in}%
\pgfsys@useobject{currentmarker}{}%
\end{pgfscope}%
\begin{pgfscope}%
\pgfsys@transformshift{8.466964in}{2.100709in}%
\pgfsys@useobject{currentmarker}{}%
\end{pgfscope}%
\begin{pgfscope}%
\pgfsys@transformshift{8.459968in}{2.083646in}%
\pgfsys@useobject{currentmarker}{}%
\end{pgfscope}%
\begin{pgfscope}%
\pgfsys@transformshift{8.452220in}{2.066749in}%
\pgfsys@useobject{currentmarker}{}%
\end{pgfscope}%
\begin{pgfscope}%
\pgfsys@transformshift{8.443727in}{2.050035in}%
\pgfsys@useobject{currentmarker}{}%
\end{pgfscope}%
\begin{pgfscope}%
\pgfsys@transformshift{8.434497in}{2.033520in}%
\pgfsys@useobject{currentmarker}{}%
\end{pgfscope}%
\begin{pgfscope}%
\pgfsys@transformshift{8.485656in}{2.133231in}%
\pgfsys@useobject{currentmarker}{}%
\end{pgfscope}%
\begin{pgfscope}%
\pgfsys@transformshift{8.480137in}{2.115580in}%
\pgfsys@useobject{currentmarker}{}%
\end{pgfscope}%
\begin{pgfscope}%
\pgfsys@transformshift{8.473844in}{2.098064in}%
\pgfsys@useobject{currentmarker}{}%
\end{pgfscope}%
\begin{pgfscope}%
\pgfsys@transformshift{8.466785in}{2.080699in}%
\pgfsys@useobject{currentmarker}{}%
\end{pgfscope}%
\begin{pgfscope}%
\pgfsys@transformshift{8.458967in}{2.063503in}%
\pgfsys@useobject{currentmarker}{}%
\end{pgfscope}%
\begin{pgfscope}%
\pgfsys@transformshift{8.450397in}{2.046492in}%
\pgfsys@useobject{currentmarker}{}%
\end{pgfscope}%
\begin{pgfscope}%
\pgfsys@transformshift{8.441084in}{2.029685in}%
\pgfsys@useobject{currentmarker}{}%
\end{pgfscope}%
\begin{pgfscope}%
\pgfsys@transformshift{8.492765in}{2.131180in}%
\pgfsys@useobject{currentmarker}{}%
\end{pgfscope}%
\begin{pgfscope}%
\pgfsys@transformshift{8.487195in}{2.113219in}%
\pgfsys@useobject{currentmarker}{}%
\end{pgfscope}%
\begin{pgfscope}%
\pgfsys@transformshift{8.480846in}{2.095395in}%
\pgfsys@useobject{currentmarker}{}%
\end{pgfscope}%
\begin{pgfscope}%
\pgfsys@transformshift{8.473722in}{2.077725in}%
\pgfsys@useobject{currentmarker}{}%
\end{pgfscope}%
\begin{pgfscope}%
\pgfsys@transformshift{8.465833in}{2.060227in}%
\pgfsys@useobject{currentmarker}{}%
\end{pgfscope}%
\begin{pgfscope}%
\pgfsys@transformshift{8.457185in}{2.042918in}%
\pgfsys@useobject{currentmarker}{}%
\end{pgfscope}%
\begin{pgfscope}%
\pgfsys@transformshift{8.447787in}{2.025815in}%
\pgfsys@useobject{currentmarker}{}%
\end{pgfscope}%
\begin{pgfscope}%
\pgfsys@transformshift{8.499999in}{2.129110in}%
\pgfsys@useobject{currentmarker}{}%
\end{pgfscope}%
\begin{pgfscope}%
\pgfsys@transformshift{8.494378in}{2.110837in}%
\pgfsys@useobject{currentmarker}{}%
\end{pgfscope}%
\begin{pgfscope}%
\pgfsys@transformshift{8.487970in}{2.092702in}%
\pgfsys@useobject{currentmarker}{}%
\end{pgfscope}%
\begin{pgfscope}%
\pgfsys@transformshift{8.480781in}{2.074724in}%
\pgfsys@useobject{currentmarker}{}%
\end{pgfscope}%
\begin{pgfscope}%
\pgfsys@transformshift{8.472819in}{2.056921in}%
\pgfsys@useobject{currentmarker}{}%
\end{pgfscope}%
\begin{pgfscope}%
\pgfsys@transformshift{8.464092in}{2.039311in}%
\pgfsys@useobject{currentmarker}{}%
\end{pgfscope}%
\begin{pgfscope}%
\pgfsys@transformshift{8.454607in}{2.021910in}%
\pgfsys@useobject{currentmarker}{}%
\end{pgfscope}%
\begin{pgfscope}%
\pgfsys@transformshift{8.507357in}{2.127021in}%
\pgfsys@useobject{currentmarker}{}%
\end{pgfscope}%
\begin{pgfscope}%
\pgfsys@transformshift{8.501684in}{2.108432in}%
\pgfsys@useobject{currentmarker}{}%
\end{pgfscope}%
\begin{pgfscope}%
\pgfsys@transformshift{8.495217in}{2.089984in}%
\pgfsys@useobject{currentmarker}{}%
\end{pgfscope}%
\begin{pgfscope}%
\pgfsys@transformshift{8.487962in}{2.071696in}%
\pgfsys@useobject{currentmarker}{}%
\end{pgfscope}%
\begin{pgfscope}%
\pgfsys@transformshift{8.479926in}{2.053585in}%
\pgfsys@useobject{currentmarker}{}%
\end{pgfscope}%
\begin{pgfscope}%
\pgfsys@transformshift{8.471118in}{2.035670in}%
\pgfsys@useobject{currentmarker}{}%
\end{pgfscope}%
\begin{pgfscope}%
\pgfsys@transformshift{8.461546in}{2.017969in}%
\pgfsys@useobject{currentmarker}{}%
\end{pgfscope}%
\begin{pgfscope}%
\pgfsys@transformshift{8.514843in}{2.124913in}%
\pgfsys@useobject{currentmarker}{}%
\end{pgfscope}%
\begin{pgfscope}%
\pgfsys@transformshift{8.509117in}{2.106005in}%
\pgfsys@useobject{currentmarker}{}%
\end{pgfscope}%
\begin{pgfscope}%
\pgfsys@transformshift{8.502589in}{2.087241in}%
\pgfsys@useobject{currentmarker}{}%
\end{pgfscope}%
\begin{pgfscope}%
\pgfsys@transformshift{8.495266in}{2.068639in}%
\pgfsys@useobject{currentmarker}{}%
\end{pgfscope}%
\begin{pgfscope}%
\pgfsys@transformshift{8.487156in}{2.050218in}%
\pgfsys@useobject{currentmarker}{}%
\end{pgfscope}%
\begin{pgfscope}%
\pgfsys@transformshift{8.478266in}{2.031996in}%
\pgfsys@useobject{currentmarker}{}%
\end{pgfscope}%
\begin{pgfscope}%
\pgfsys@transformshift{8.468604in}{2.013991in}%
\pgfsys@useobject{currentmarker}{}%
\end{pgfscope}%
\begin{pgfscope}%
\pgfsys@transformshift{8.522456in}{2.122785in}%
\pgfsys@useobject{currentmarker}{}%
\end{pgfscope}%
\begin{pgfscope}%
\pgfsys@transformshift{8.516676in}{2.103556in}%
\pgfsys@useobject{currentmarker}{}%
\end{pgfscope}%
\begin{pgfscope}%
\pgfsys@transformshift{8.510087in}{2.084472in}%
\pgfsys@useobject{currentmarker}{}%
\end{pgfscope}%
\begin{pgfscope}%
\pgfsys@transformshift{8.502696in}{2.065554in}%
\pgfsys@useobject{currentmarker}{}%
\end{pgfscope}%
\begin{pgfscope}%
\pgfsys@transformshift{8.494509in}{2.046819in}%
\pgfsys@useobject{currentmarker}{}%
\end{pgfscope}%
\begin{pgfscope}%
\pgfsys@transformshift{8.485535in}{2.028287in}%
\pgfsys@useobject{currentmarker}{}%
\end{pgfscope}%
\begin{pgfscope}%
\pgfsys@transformshift{8.475783in}{2.009976in}%
\pgfsys@useobject{currentmarker}{}%
\end{pgfscope}%
\begin{pgfscope}%
\pgfsys@transformshift{8.530199in}{2.120637in}%
\pgfsys@useobject{currentmarker}{}%
\end{pgfscope}%
\begin{pgfscope}%
\pgfsys@transformshift{8.524364in}{2.101083in}%
\pgfsys@useobject{currentmarker}{}%
\end{pgfscope}%
\begin{pgfscope}%
\pgfsys@transformshift{8.517713in}{2.081677in}%
\pgfsys@useobject{currentmarker}{}%
\end{pgfscope}%
\begin{pgfscope}%
\pgfsys@transformshift{8.510251in}{2.062440in}%
\pgfsys@useobject{currentmarker}{}%
\end{pgfscope}%
\begin{pgfscope}%
\pgfsys@transformshift{8.501987in}{2.043389in}%
\pgfsys@useobject{currentmarker}{}%
\end{pgfscope}%
\begin{pgfscope}%
\pgfsys@transformshift{8.492928in}{2.024544in}%
\pgfsys@useobject{currentmarker}{}%
\end{pgfscope}%
\begin{pgfscope}%
\pgfsys@transformshift{8.483083in}{2.005924in}%
\pgfsys@useobject{currentmarker}{}%
\end{pgfscope}%
\begin{pgfscope}%
\pgfsys@transformshift{8.538071in}{2.118469in}%
\pgfsys@useobject{currentmarker}{}%
\end{pgfscope}%
\begin{pgfscope}%
\pgfsys@transformshift{8.532181in}{2.098587in}%
\pgfsys@useobject{currentmarker}{}%
\end{pgfscope}%
\begin{pgfscope}%
\pgfsys@transformshift{8.525466in}{2.078856in}%
\pgfsys@useobject{currentmarker}{}%
\end{pgfscope}%
\begin{pgfscope}%
\pgfsys@transformshift{8.517933in}{2.059296in}%
\pgfsys@useobject{currentmarker}{}%
\end{pgfscope}%
\begin{pgfscope}%
\pgfsys@transformshift{8.509590in}{2.039925in}%
\pgfsys@useobject{currentmarker}{}%
\end{pgfscope}%
\begin{pgfscope}%
\pgfsys@transformshift{8.500445in}{2.020765in}%
\pgfsys@useobject{currentmarker}{}%
\end{pgfscope}%
\begin{pgfscope}%
\pgfsys@transformshift{8.490506in}{2.001832in}%
\pgfsys@useobject{currentmarker}{}%
\end{pgfscope}%
\begin{pgfscope}%
\pgfsys@transformshift{8.546075in}{2.116279in}%
\pgfsys@useobject{currentmarker}{}%
\end{pgfscope}%
\begin{pgfscope}%
\pgfsys@transformshift{8.540128in}{2.096067in}%
\pgfsys@useobject{currentmarker}{}%
\end{pgfscope}%
\begin{pgfscope}%
\pgfsys@transformshift{8.533349in}{2.076007in}%
\pgfsys@useobject{currentmarker}{}%
\end{pgfscope}%
\begin{pgfscope}%
\pgfsys@transformshift{8.525744in}{2.056121in}%
\pgfsys@useobject{currentmarker}{}%
\end{pgfscope}%
\begin{pgfscope}%
\pgfsys@transformshift{8.517321in}{2.036429in}%
\pgfsys@useobject{currentmarker}{}%
\end{pgfscope}%
\begin{pgfscope}%
\pgfsys@transformshift{8.508087in}{2.016949in}%
\pgfsys@useobject{currentmarker}{}%
\end{pgfscope}%
\begin{pgfscope}%
\pgfsys@transformshift{8.498054in}{1.997702in}%
\pgfsys@useobject{currentmarker}{}%
\end{pgfscope}%
\begin{pgfscope}%
\pgfsys@transformshift{8.554212in}{2.114069in}%
\pgfsys@useobject{currentmarker}{}%
\end{pgfscope}%
\begin{pgfscope}%
\pgfsys@transformshift{8.548208in}{2.093522in}%
\pgfsys@useobject{currentmarker}{}%
\end{pgfscope}%
\begin{pgfscope}%
\pgfsys@transformshift{8.541363in}{2.073131in}%
\pgfsys@useobject{currentmarker}{}%
\end{pgfscope}%
\begin{pgfscope}%
\pgfsys@transformshift{8.533684in}{2.052917in}%
\pgfsys@useobject{currentmarker}{}%
\end{pgfscope}%
\begin{pgfscope}%
\pgfsys@transformshift{8.525179in}{2.032899in}%
\pgfsys@useobject{currentmarker}{}%
\end{pgfscope}%
\begin{pgfscope}%
\pgfsys@transformshift{8.515857in}{2.013097in}%
\pgfsys@useobject{currentmarker}{}%
\end{pgfscope}%
\begin{pgfscope}%
\pgfsys@transformshift{8.505726in}{1.993531in}%
\pgfsys@useobject{currentmarker}{}%
\end{pgfscope}%
\begin{pgfscope}%
\pgfsys@transformshift{8.562483in}{2.111837in}%
\pgfsys@useobject{currentmarker}{}%
\end{pgfscope}%
\begin{pgfscope}%
\pgfsys@transformshift{8.556420in}{2.090953in}%
\pgfsys@useobject{currentmarker}{}%
\end{pgfscope}%
\begin{pgfscope}%
\pgfsys@transformshift{8.549509in}{2.070227in}%
\pgfsys@useobject{currentmarker}{}%
\end{pgfscope}%
\begin{pgfscope}%
\pgfsys@transformshift{8.541755in}{2.049681in}%
\pgfsys@useobject{currentmarker}{}%
\end{pgfscope}%
\begin{pgfscope}%
\pgfsys@transformshift{8.533167in}{2.029334in}%
\pgfsys@useobject{currentmarker}{}%
\end{pgfscope}%
\begin{pgfscope}%
\pgfsys@transformshift{8.523754in}{2.009207in}%
\pgfsys@useobject{currentmarker}{}%
\end{pgfscope}%
\begin{pgfscope}%
\pgfsys@transformshift{8.513525in}{1.989321in}%
\pgfsys@useobject{currentmarker}{}%
\end{pgfscope}%
\begin{pgfscope}%
\pgfsys@transformshift{8.570889in}{2.109584in}%
\pgfsys@useobject{currentmarker}{}%
\end{pgfscope}%
\begin{pgfscope}%
\pgfsys@transformshift{8.564767in}{2.088359in}%
\pgfsys@useobject{currentmarker}{}%
\end{pgfscope}%
\begin{pgfscope}%
\pgfsys@transformshift{8.557788in}{2.067295in}%
\pgfsys@useobject{currentmarker}{}%
\end{pgfscope}%
\begin{pgfscope}%
\pgfsys@transformshift{8.549958in}{2.046414in}%
\pgfsys@useobject{currentmarker}{}%
\end{pgfscope}%
\begin{pgfscope}%
\pgfsys@transformshift{8.541286in}{2.025735in}%
\pgfsys@useobject{currentmarker}{}%
\end{pgfscope}%
\begin{pgfscope}%
\pgfsys@transformshift{8.531781in}{2.005280in}%
\pgfsys@useobject{currentmarker}{}%
\end{pgfscope}%
\begin{pgfscope}%
\pgfsys@transformshift{8.521451in}{1.985069in}%
\pgfsys@useobject{currentmarker}{}%
\end{pgfscope}%
\begin{pgfscope}%
\pgfsys@transformshift{8.579432in}{2.107308in}%
\pgfsys@useobject{currentmarker}{}%
\end{pgfscope}%
\begin{pgfscope}%
\pgfsys@transformshift{8.573249in}{2.085739in}%
\pgfsys@useobject{currentmarker}{}%
\end{pgfscope}%
\begin{pgfscope}%
\pgfsys@transformshift{8.566201in}{2.064334in}%
\pgfsys@useobject{currentmarker}{}%
\end{pgfscope}%
\begin{pgfscope}%
\pgfsys@transformshift{8.558295in}{2.043114in}%
\pgfsys@useobject{currentmarker}{}%
\end{pgfscope}%
\begin{pgfscope}%
\pgfsys@transformshift{8.549537in}{2.022100in}%
\pgfsys@useobject{currentmarker}{}%
\end{pgfscope}%
\begin{pgfscope}%
\pgfsys@transformshift{8.539938in}{2.001314in}%
\pgfsys@useobject{currentmarker}{}%
\end{pgfscope}%
\begin{pgfscope}%
\pgfsys@transformshift{8.529506in}{1.980775in}%
\pgfsys@useobject{currentmarker}{}%
\end{pgfscope}%
\begin{pgfscope}%
\pgfsys@transformshift{8.588113in}{2.105010in}%
\pgfsys@useobject{currentmarker}{}%
\end{pgfscope}%
\begin{pgfscope}%
\pgfsys@transformshift{8.581869in}{2.083093in}%
\pgfsys@useobject{currentmarker}{}%
\end{pgfscope}%
\begin{pgfscope}%
\pgfsys@transformshift{8.574751in}{2.061344in}%
\pgfsys@useobject{currentmarker}{}%
\end{pgfscope}%
\begin{pgfscope}%
\pgfsys@transformshift{8.566766in}{2.039782in}%
\pgfsys@useobject{currentmarker}{}%
\end{pgfscope}%
\begin{pgfscope}%
\pgfsys@transformshift{8.557921in}{2.018430in}%
\pgfsys@useobject{currentmarker}{}%
\end{pgfscope}%
\begin{pgfscope}%
\pgfsys@transformshift{8.548227in}{1.997308in}%
\pgfsys@useobject{currentmarker}{}%
\end{pgfscope}%
\begin{pgfscope}%
\pgfsys@transformshift{8.537691in}{1.976439in}%
\pgfsys@useobject{currentmarker}{}%
\end{pgfscope}%
\begin{pgfscope}%
\pgfsys@transformshift{8.596933in}{2.102689in}%
\pgfsys@useobject{currentmarker}{}%
\end{pgfscope}%
\begin{pgfscope}%
\pgfsys@transformshift{8.590627in}{2.080422in}%
\pgfsys@useobject{currentmarker}{}%
\end{pgfscope}%
\begin{pgfscope}%
\pgfsys@transformshift{8.583438in}{2.058323in}%
\pgfsys@useobject{currentmarker}{}%
\end{pgfscope}%
\begin{pgfscope}%
\pgfsys@transformshift{8.575373in}{2.036416in}%
\pgfsys@useobject{currentmarker}{}%
\end{pgfscope}%
\begin{pgfscope}%
\pgfsys@transformshift{8.566440in}{2.014722in}%
\pgfsys@useobject{currentmarker}{}%
\end{pgfscope}%
\begin{pgfscope}%
\pgfsys@transformshift{8.556649in}{1.993263in}%
\pgfsys@useobject{currentmarker}{}%
\end{pgfscope}%
\begin{pgfscope}%
\pgfsys@transformshift{8.546008in}{1.972059in}%
\pgfsys@useobject{currentmarker}{}%
\end{pgfscope}%
\begin{pgfscope}%
\pgfsys@transformshift{8.605894in}{2.100344in}%
\pgfsys@useobject{currentmarker}{}%
\end{pgfscope}%
\begin{pgfscope}%
\pgfsys@transformshift{8.599525in}{2.077723in}%
\pgfsys@useobject{currentmarker}{}%
\end{pgfscope}%
\begin{pgfscope}%
\pgfsys@transformshift{8.592263in}{2.055273in}%
\pgfsys@useobject{currentmarker}{}%
\end{pgfscope}%
\begin{pgfscope}%
\pgfsys@transformshift{8.584117in}{2.033017in}%
\pgfsys@useobject{currentmarker}{}%
\end{pgfscope}%
\begin{pgfscope}%
\pgfsys@transformshift{8.575095in}{2.010978in}%
\pgfsys@useobject{currentmarker}{}%
\end{pgfscope}%
\begin{pgfscope}%
\pgfsys@transformshift{8.565205in}{1.989177in}%
\pgfsys@useobject{currentmarker}{}%
\end{pgfscope}%
\begin{pgfscope}%
\pgfsys@transformshift{8.554458in}{1.967636in}%
\pgfsys@useobject{currentmarker}{}%
\end{pgfscope}%
\begin{pgfscope}%
\pgfsys@transformshift{8.614998in}{2.097976in}%
\pgfsys@useobject{currentmarker}{}%
\end{pgfscope}%
\begin{pgfscope}%
\pgfsys@transformshift{8.608564in}{2.074997in}%
\pgfsys@useobject{currentmarker}{}%
\end{pgfscope}%
\begin{pgfscope}%
\pgfsys@transformshift{8.601229in}{2.052192in}%
\pgfsys@useobject{currentmarker}{}%
\end{pgfscope}%
\begin{pgfscope}%
\pgfsys@transformshift{8.593001in}{2.029584in}%
\pgfsys@useobject{currentmarker}{}%
\end{pgfscope}%
\begin{pgfscope}%
\pgfsys@transformshift{8.583887in}{2.007196in}%
\pgfsys@useobject{currentmarker}{}%
\end{pgfscope}%
\begin{pgfscope}%
\pgfsys@transformshift{8.573897in}{1.985050in}%
\pgfsys@useobject{currentmarker}{}%
\end{pgfscope}%
\begin{pgfscope}%
\pgfsys@transformshift{8.563041in}{1.963168in}%
\pgfsys@useobject{currentmarker}{}%
\end{pgfscope}%
\begin{pgfscope}%
\pgfsys@transformshift{8.624245in}{2.095584in}%
\pgfsys@useobject{currentmarker}{}%
\end{pgfscope}%
\begin{pgfscope}%
\pgfsys@transformshift{8.617745in}{2.072243in}%
\pgfsys@useobject{currentmarker}{}%
\end{pgfscope}%
\begin{pgfscope}%
\pgfsys@transformshift{8.610336in}{2.049079in}%
\pgfsys@useobject{currentmarker}{}%
\end{pgfscope}%
\begin{pgfscope}%
\pgfsys@transformshift{8.602024in}{2.026116in}%
\pgfsys@useobject{currentmarker}{}%
\end{pgfscope}%
\begin{pgfscope}%
\pgfsys@transformshift{8.592818in}{2.003376in}%
\pgfsys@useobject{currentmarker}{}%
\end{pgfscope}%
\begin{pgfscope}%
\pgfsys@transformshift{8.582727in}{1.980881in}%
\pgfsys@useobject{currentmarker}{}%
\end{pgfscope}%
\begin{pgfscope}%
\pgfsys@transformshift{8.571760in}{1.958655in}%
\pgfsys@useobject{currentmarker}{}%
\end{pgfscope}%
\begin{pgfscope}%
\pgfsys@transformshift{8.633637in}{2.093168in}%
\pgfsys@useobject{currentmarker}{}%
\end{pgfscope}%
\begin{pgfscope}%
\pgfsys@transformshift{8.627071in}{2.069462in}%
\pgfsys@useobject{currentmarker}{}%
\end{pgfscope}%
\begin{pgfscope}%
\pgfsys@transformshift{8.619586in}{2.045935in}%
\pgfsys@useobject{currentmarker}{}%
\end{pgfscope}%
\begin{pgfscope}%
\pgfsys@transformshift{8.611189in}{2.022613in}%
\pgfsys@useobject{currentmarker}{}%
\end{pgfscope}%
\begin{pgfscope}%
\pgfsys@transformshift{8.601889in}{1.999516in}%
\pgfsys@useobject{currentmarker}{}%
\end{pgfscope}%
\begin{pgfscope}%
\pgfsys@transformshift{8.591694in}{1.976670in}%
\pgfsys@useobject{currentmarker}{}%
\end{pgfscope}%
\begin{pgfscope}%
\pgfsys@transformshift{8.580616in}{1.954096in}%
\pgfsys@useobject{currentmarker}{}%
\end{pgfscope}%
\begin{pgfscope}%
\pgfsys@transformshift{8.643175in}{2.090727in}%
\pgfsys@useobject{currentmarker}{}%
\end{pgfscope}%
\begin{pgfscope}%
\pgfsys@transformshift{8.636542in}{2.066651in}%
\pgfsys@useobject{currentmarker}{}%
\end{pgfscope}%
\begin{pgfscope}%
\pgfsys@transformshift{8.628980in}{2.042759in}%
\pgfsys@useobject{currentmarker}{}%
\end{pgfscope}%
\begin{pgfscope}%
\pgfsys@transformshift{8.620497in}{2.019073in}%
\pgfsys@useobject{currentmarker}{}%
\end{pgfscope}%
\begin{pgfscope}%
\pgfsys@transformshift{8.611101in}{1.995618in}%
\pgfsys@useobject{currentmarker}{}%
\end{pgfscope}%
\begin{pgfscope}%
\pgfsys@transformshift{8.600802in}{1.972415in}%
\pgfsys@useobject{currentmarker}{}%
\end{pgfscope}%
\begin{pgfscope}%
\pgfsys@transformshift{8.589610in}{1.949490in}%
\pgfsys@useobject{currentmarker}{}%
\end{pgfscope}%
\begin{pgfscope}%
\pgfsys@transformshift{8.652862in}{2.088260in}%
\pgfsys@useobject{currentmarker}{}%
\end{pgfscope}%
\begin{pgfscope}%
\pgfsys@transformshift{8.646160in}{2.063812in}%
\pgfsys@useobject{currentmarker}{}%
\end{pgfscope}%
\begin{pgfscope}%
\pgfsys@transformshift{8.638520in}{2.039550in}%
\pgfsys@useobject{currentmarker}{}%
\end{pgfscope}%
\begin{pgfscope}%
\pgfsys@transformshift{8.629950in}{2.015497in}%
\pgfsys@useobject{currentmarker}{}%
\end{pgfscope}%
\begin{pgfscope}%
\pgfsys@transformshift{8.620457in}{1.991679in}%
\pgfsys@useobject{currentmarker}{}%
\end{pgfscope}%
\begin{pgfscope}%
\pgfsys@transformshift{8.610052in}{1.968117in}%
\pgfsys@useobject{currentmarker}{}%
\end{pgfscope}%
\begin{pgfscope}%
\pgfsys@transformshift{8.598744in}{1.944837in}%
\pgfsys@useobject{currentmarker}{}%
\end{pgfscope}%
\begin{pgfscope}%
\pgfsys@transformshift{8.662698in}{2.085769in}%
\pgfsys@useobject{currentmarker}{}%
\end{pgfscope}%
\begin{pgfscope}%
\pgfsys@transformshift{8.655927in}{2.060944in}%
\pgfsys@useobject{currentmarker}{}%
\end{pgfscope}%
\begin{pgfscope}%
\pgfsys@transformshift{8.648208in}{2.036308in}%
\pgfsys@useobject{currentmarker}{}%
\end{pgfscope}%
\begin{pgfscope}%
\pgfsys@transformshift{8.639548in}{2.011885in}%
\pgfsys@useobject{currentmarker}{}%
\end{pgfscope}%
\begin{pgfscope}%
\pgfsys@transformshift{8.629957in}{1.987699in}%
\pgfsys@useobject{currentmarker}{}%
\end{pgfscope}%
\begin{pgfscope}%
\pgfsys@transformshift{8.619444in}{1.963774in}%
\pgfsys@useobject{currentmarker}{}%
\end{pgfscope}%
\begin{pgfscope}%
\pgfsys@transformshift{8.608019in}{1.940135in}%
\pgfsys@useobject{currentmarker}{}%
\end{pgfscope}%
\begin{pgfscope}%
\pgfsys@transformshift{8.672686in}{2.083251in}%
\pgfsys@useobject{currentmarker}{}%
\end{pgfscope}%
\begin{pgfscope}%
\pgfsys@transformshift{8.665844in}{2.058046in}%
\pgfsys@useobject{currentmarker}{}%
\end{pgfscope}%
\begin{pgfscope}%
\pgfsys@transformshift{8.658044in}{2.033032in}%
\pgfsys@useobject{currentmarker}{}%
\end{pgfscope}%
\begin{pgfscope}%
\pgfsys@transformshift{8.649294in}{2.008234in}%
\pgfsys@useobject{currentmarker}{}%
\end{pgfscope}%
\begin{pgfscope}%
\pgfsys@transformshift{8.639603in}{1.983678in}%
\pgfsys@useobject{currentmarker}{}%
\end{pgfscope}%
\begin{pgfscope}%
\pgfsys@transformshift{8.628980in}{1.959387in}%
\pgfsys@useobject{currentmarker}{}%
\end{pgfscope}%
\begin{pgfscope}%
\pgfsys@transformshift{8.617436in}{1.935385in}%
\pgfsys@useobject{currentmarker}{}%
\end{pgfscope}%
\begin{pgfscope}%
\pgfsys@transformshift{8.682826in}{2.080707in}%
\pgfsys@useobject{currentmarker}{}%
\end{pgfscope}%
\begin{pgfscope}%
\pgfsys@transformshift{8.675912in}{2.055117in}%
\pgfsys@useobject{currentmarker}{}%
\end{pgfscope}%
\begin{pgfscope}%
\pgfsys@transformshift{8.668031in}{2.029722in}%
\pgfsys@useobject{currentmarker}{}%
\end{pgfscope}%
\begin{pgfscope}%
\pgfsys@transformshift{8.659189in}{2.004546in}%
\pgfsys@useobject{currentmarker}{}%
\end{pgfscope}%
\begin{pgfscope}%
\pgfsys@transformshift{8.649397in}{1.979614in}%
\pgfsys@useobject{currentmarker}{}%
\end{pgfscope}%
\begin{pgfscope}%
\pgfsys@transformshift{8.638662in}{1.954953in}%
\pgfsys@useobject{currentmarker}{}%
\end{pgfscope}%
\begin{pgfscope}%
\pgfsys@transformshift{8.626997in}{1.930585in}%
\pgfsys@useobject{currentmarker}{}%
\end{pgfscope}%
\begin{pgfscope}%
\pgfsys@transformshift{8.693121in}{2.078136in}%
\pgfsys@useobject{currentmarker}{}%
\end{pgfscope}%
\begin{pgfscope}%
\pgfsys@transformshift{8.686134in}{2.052158in}%
\pgfsys@useobject{currentmarker}{}%
\end{pgfscope}%
\begin{pgfscope}%
\pgfsys@transformshift{8.678170in}{2.026377in}%
\pgfsys@useobject{currentmarker}{}%
\end{pgfscope}%
\begin{pgfscope}%
\pgfsys@transformshift{8.669235in}{2.000818in}%
\pgfsys@useobject{currentmarker}{}%
\end{pgfscope}%
\begin{pgfscope}%
\pgfsys@transformshift{8.659339in}{1.975509in}%
\pgfsys@useobject{currentmarker}{}%
\end{pgfscope}%
\begin{pgfscope}%
\pgfsys@transformshift{8.648492in}{1.950472in}%
\pgfsys@useobject{currentmarker}{}%
\end{pgfscope}%
\begin{pgfscope}%
\pgfsys@transformshift{8.636704in}{1.925735in}%
\pgfsys@useobject{currentmarker}{}%
\end{pgfscope}%
\begin{pgfscope}%
\pgfsys@transformshift{8.703571in}{2.075538in}%
\pgfsys@useobject{currentmarker}{}%
\end{pgfscope}%
\begin{pgfscope}%
\pgfsys@transformshift{8.696511in}{2.049167in}%
\pgfsys@useobject{currentmarker}{}%
\end{pgfscope}%
\begin{pgfscope}%
\pgfsys@transformshift{8.688462in}{2.022996in}%
\pgfsys@useobject{currentmarker}{}%
\end{pgfscope}%
\begin{pgfscope}%
\pgfsys@transformshift{8.679433in}{1.997052in}%
\pgfsys@useobject{currentmarker}{}%
\end{pgfscope}%
\begin{pgfscope}%
\pgfsys@transformshift{8.669432in}{1.971359in}%
\pgfsys@useobject{currentmarker}{}%
\end{pgfscope}%
\begin{pgfscope}%
\pgfsys@transformshift{8.658470in}{1.945945in}%
\pgfsys@useobject{currentmarker}{}%
\end{pgfscope}%
\begin{pgfscope}%
\pgfsys@transformshift{8.646558in}{1.920833in}%
\pgfsys@useobject{currentmarker}{}%
\end{pgfscope}%
\begin{pgfscope}%
\pgfsys@transformshift{8.714179in}{2.072913in}%
\pgfsys@useobject{currentmarker}{}%
\end{pgfscope}%
\begin{pgfscope}%
\pgfsys@transformshift{8.707044in}{2.046145in}%
\pgfsys@useobject{currentmarker}{}%
\end{pgfscope}%
\begin{pgfscope}%
\pgfsys@transformshift{8.698910in}{2.019580in}%
\pgfsys@useobject{currentmarker}{}%
\end{pgfscope}%
\begin{pgfscope}%
\pgfsys@transformshift{8.689785in}{1.993245in}%
\pgfsys@useobject{currentmarker}{}%
\end{pgfscope}%
\begin{pgfscope}%
\pgfsys@transformshift{8.679678in}{1.967166in}%
\pgfsys@useobject{currentmarker}{}%
\end{pgfscope}%
\begin{pgfscope}%
\pgfsys@transformshift{8.668600in}{1.941369in}%
\pgfsys@useobject{currentmarker}{}%
\end{pgfscope}%
\begin{pgfscope}%
\pgfsys@transformshift{8.656560in}{1.915880in}%
\pgfsys@useobject{currentmarker}{}%
\end{pgfscope}%
\begin{pgfscope}%
\pgfsys@transformshift{8.724947in}{2.070260in}%
\pgfsys@useobject{currentmarker}{}%
\end{pgfscope}%
\begin{pgfscope}%
\pgfsys@transformshift{8.717735in}{2.043091in}%
\pgfsys@useobject{currentmarker}{}%
\end{pgfscope}%
\begin{pgfscope}%
\pgfsys@transformshift{8.709514in}{2.016128in}%
\pgfsys@useobject{currentmarker}{}%
\end{pgfscope}%
\begin{pgfscope}%
\pgfsys@transformshift{8.700292in}{1.989398in}%
\pgfsys@useobject{currentmarker}{}%
\end{pgfscope}%
\begin{pgfscope}%
\pgfsys@transformshift{8.690077in}{1.962929in}%
\pgfsys@useobject{currentmarker}{}%
\end{pgfscope}%
\begin{pgfscope}%
\pgfsys@transformshift{8.678881in}{1.936745in}%
\pgfsys@useobject{currentmarker}{}%
\end{pgfscope}%
\begin{pgfscope}%
\pgfsys@transformshift{8.666713in}{1.910874in}%
\pgfsys@useobject{currentmarker}{}%
\end{pgfscope}%
\begin{pgfscope}%
\pgfsys@transformshift{8.735875in}{2.067578in}%
\pgfsys@useobject{currentmarker}{}%
\end{pgfscope}%
\begin{pgfscope}%
\pgfsys@transformshift{8.728586in}{2.040003in}%
\pgfsys@useobject{currentmarker}{}%
\end{pgfscope}%
\begin{pgfscope}%
\pgfsys@transformshift{8.720277in}{2.012639in}%
\pgfsys@useobject{currentmarker}{}%
\end{pgfscope}%
\begin{pgfscope}%
\pgfsys@transformshift{8.710956in}{1.985510in}%
\pgfsys@useobject{currentmarker}{}%
\end{pgfscope}%
\begin{pgfscope}%
\pgfsys@transformshift{8.700632in}{1.958645in}%
\pgfsys@useobject{currentmarker}{}%
\end{pgfscope}%
\begin{pgfscope}%
\pgfsys@transformshift{8.689316in}{1.932071in}%
\pgfsys@useobject{currentmarker}{}%
\end{pgfscope}%
\begin{pgfscope}%
\pgfsys@transformshift{8.677018in}{1.905814in}%
\pgfsys@useobject{currentmarker}{}%
\end{pgfscope}%
\begin{pgfscope}%
\pgfsys@transformshift{8.746966in}{2.064867in}%
\pgfsys@useobject{currentmarker}{}%
\end{pgfscope}%
\begin{pgfscope}%
\pgfsys@transformshift{8.739599in}{2.036883in}%
\pgfsys@useobject{currentmarker}{}%
\end{pgfscope}%
\begin{pgfscope}%
\pgfsys@transformshift{8.731201in}{2.009112in}%
\pgfsys@useobject{currentmarker}{}%
\end{pgfscope}%
\begin{pgfscope}%
\pgfsys@transformshift{8.721779in}{1.981580in}%
\pgfsys@useobject{currentmarker}{}%
\end{pgfscope}%
\begin{pgfscope}%
\pgfsys@transformshift{8.711344in}{1.954316in}%
\pgfsys@useobject{currentmarker}{}%
\end{pgfscope}%
\begin{pgfscope}%
\pgfsys@transformshift{8.699906in}{1.927347in}%
\pgfsys@useobject{currentmarker}{}%
\end{pgfscope}%
\begin{pgfscope}%
\pgfsys@transformshift{8.687476in}{1.900700in}%
\pgfsys@useobject{currentmarker}{}%
\end{pgfscope}%
\begin{pgfscope}%
\pgfsys@transformshift{8.758222in}{2.062128in}%
\pgfsys@useobject{currentmarker}{}%
\end{pgfscope}%
\begin{pgfscope}%
\pgfsys@transformshift{8.750775in}{2.033729in}%
\pgfsys@useobject{currentmarker}{}%
\end{pgfscope}%
\begin{pgfscope}%
\pgfsys@transformshift{8.742286in}{2.005547in}%
\pgfsys@useobject{currentmarker}{}%
\end{pgfscope}%
\begin{pgfscope}%
\pgfsys@transformshift{8.732763in}{1.977608in}%
\pgfsys@useobject{currentmarker}{}%
\end{pgfscope}%
\begin{pgfscope}%
\pgfsys@transformshift{8.722215in}{1.949941in}%
\pgfsys@useobject{currentmarker}{}%
\end{pgfscope}%
\begin{pgfscope}%
\pgfsys@transformshift{8.710653in}{1.922572in}%
\pgfsys@useobject{currentmarker}{}%
\end{pgfscope}%
\begin{pgfscope}%
\pgfsys@transformshift{8.698089in}{1.895530in}%
\pgfsys@useobject{currentmarker}{}%
\end{pgfscope}%
\begin{pgfscope}%
\pgfsys@transformshift{8.769644in}{2.059358in}%
\pgfsys@useobject{currentmarker}{}%
\end{pgfscope}%
\begin{pgfscope}%
\pgfsys@transformshift{8.762116in}{2.030541in}%
\pgfsys@useobject{currentmarker}{}%
\end{pgfscope}%
\begin{pgfscope}%
\pgfsys@transformshift{8.753535in}{2.001943in}%
\pgfsys@useobject{currentmarker}{}%
\end{pgfscope}%
\begin{pgfscope}%
\pgfsys@transformshift{8.743908in}{1.973593in}%
\pgfsys@useobject{currentmarker}{}%
\end{pgfscope}%
\begin{pgfscope}%
\pgfsys@transformshift{8.733246in}{1.945517in}%
\pgfsys@useobject{currentmarker}{}%
\end{pgfscope}%
\begin{pgfscope}%
\pgfsys@transformshift{8.721559in}{1.917746in}%
\pgfsys@useobject{currentmarker}{}%
\end{pgfscope}%
\begin{pgfscope}%
\pgfsys@transformshift{8.708858in}{1.890305in}%
\pgfsys@useobject{currentmarker}{}%
\end{pgfscope}%
\begin{pgfscope}%
\pgfsys@transformshift{8.781233in}{2.056559in}%
\pgfsys@useobject{currentmarker}{}%
\end{pgfscope}%
\begin{pgfscope}%
\pgfsys@transformshift{8.773624in}{2.027319in}%
\pgfsys@useobject{currentmarker}{}%
\end{pgfscope}%
\begin{pgfscope}%
\pgfsys@transformshift{8.764949in}{1.998301in}%
\pgfsys@useobject{currentmarker}{}%
\end{pgfscope}%
\begin{pgfscope}%
\pgfsys@transformshift{8.755218in}{1.969534in}%
\pgfsys@useobject{currentmarker}{}%
\end{pgfscope}%
\begin{pgfscope}%
\pgfsys@transformshift{8.744440in}{1.941046in}%
\pgfsys@useobject{currentmarker}{}%
\end{pgfscope}%
\begin{pgfscope}%
\pgfsys@transformshift{8.732626in}{1.912867in}%
\pgfsys@useobject{currentmarker}{}%
\end{pgfscope}%
\begin{pgfscope}%
\pgfsys@transformshift{8.719787in}{1.885023in}%
\pgfsys@useobject{currentmarker}{}%
\end{pgfscope}%
\begin{pgfscope}%
\pgfsys@transformshift{8.792993in}{2.053729in}%
\pgfsys@useobject{currentmarker}{}%
\end{pgfscope}%
\begin{pgfscope}%
\pgfsys@transformshift{8.785301in}{2.024061in}%
\pgfsys@useobject{currentmarker}{}%
\end{pgfscope}%
\begin{pgfscope}%
\pgfsys@transformshift{8.776531in}{1.994619in}%
\pgfsys@useobject{currentmarker}{}%
\end{pgfscope}%
\begin{pgfscope}%
\pgfsys@transformshift{8.766694in}{1.965431in}%
\pgfsys@useobject{currentmarker}{}%
\end{pgfscope}%
\begin{pgfscope}%
\pgfsys@transformshift{8.755798in}{1.936526in}%
\pgfsys@useobject{currentmarker}{}%
\end{pgfscope}%
\begin{pgfscope}%
\pgfsys@transformshift{8.743854in}{1.907935in}%
\pgfsys@useobject{currentmarker}{}%
\end{pgfscope}%
\begin{pgfscope}%
\pgfsys@transformshift{8.730875in}{1.879684in}%
\pgfsys@useobject{currentmarker}{}%
\end{pgfscope}%
\begin{pgfscope}%
\pgfsys@transformshift{8.804925in}{2.050868in}%
\pgfsys@useobject{currentmarker}{}%
\end{pgfscope}%
\begin{pgfscope}%
\pgfsys@transformshift{8.797148in}{2.020768in}%
\pgfsys@useobject{currentmarker}{}%
\end{pgfscope}%
\begin{pgfscope}%
\pgfsys@transformshift{8.788282in}{1.990896in}%
\pgfsys@useobject{currentmarker}{}%
\end{pgfscope}%
\begin{pgfscope}%
\pgfsys@transformshift{8.778337in}{1.961283in}%
\pgfsys@useobject{currentmarker}{}%
\end{pgfscope}%
\begin{pgfscope}%
\pgfsys@transformshift{8.767321in}{1.931957in}%
\pgfsys@useobject{currentmarker}{}%
\end{pgfscope}%
\begin{pgfscope}%
\pgfsys@transformshift{8.755247in}{1.902948in}%
\pgfsys@useobject{currentmarker}{}%
\end{pgfscope}%
\begin{pgfscope}%
\pgfsys@transformshift{8.742125in}{1.874286in}%
\pgfsys@useobject{currentmarker}{}%
\end{pgfscope}%
\begin{pgfscope}%
\pgfsys@transformshift{8.817030in}{2.047975in}%
\pgfsys@useobject{currentmarker}{}%
\end{pgfscope}%
\begin{pgfscope}%
\pgfsys@transformshift{8.809167in}{2.017438in}%
\pgfsys@useobject{currentmarker}{}%
\end{pgfscope}%
\begin{pgfscope}%
\pgfsys@transformshift{8.800204in}{1.987132in}%
\pgfsys@useobject{currentmarker}{}%
\end{pgfscope}%
\begin{pgfscope}%
\pgfsys@transformshift{8.790149in}{1.957089in}%
\pgfsys@useobject{currentmarker}{}%
\end{pgfscope}%
\begin{pgfscope}%
\pgfsys@transformshift{8.779013in}{1.927337in}%
\pgfsys@useobject{currentmarker}{}%
\end{pgfscope}%
\begin{pgfscope}%
\pgfsys@transformshift{8.766805in}{1.897908in}%
\pgfsys@useobject{currentmarker}{}%
\end{pgfscope}%
\begin{pgfscope}%
\pgfsys@transformshift{8.753539in}{1.868829in}%
\pgfsys@useobject{currentmarker}{}%
\end{pgfscope}%
\begin{pgfscope}%
\pgfsys@transformshift{8.829310in}{2.045051in}%
\pgfsys@useobject{currentmarker}{}%
\end{pgfscope}%
\begin{pgfscope}%
\pgfsys@transformshift{8.821361in}{2.014072in}%
\pgfsys@useobject{currentmarker}{}%
\end{pgfscope}%
\begin{pgfscope}%
\pgfsys@transformshift{8.812299in}{1.983327in}%
\pgfsys@useobject{currentmarker}{}%
\end{pgfscope}%
\begin{pgfscope}%
\pgfsys@transformshift{8.802133in}{1.952849in}%
\pgfsys@useobject{currentmarker}{}%
\end{pgfscope}%
\begin{pgfscope}%
\pgfsys@transformshift{8.790873in}{1.922667in}%
\pgfsys@useobject{currentmarker}{}%
\end{pgfscope}%
\begin{pgfscope}%
\pgfsys@transformshift{8.778531in}{1.892811in}%
\pgfsys@useobject{currentmarker}{}%
\end{pgfscope}%
\begin{pgfscope}%
\pgfsys@transformshift{8.765119in}{1.863311in}%
\pgfsys@useobject{currentmarker}{}%
\end{pgfscope}%
\begin{pgfscope}%
\pgfsys@transformshift{8.841768in}{2.042094in}%
\pgfsys@useobject{currentmarker}{}%
\end{pgfscope}%
\begin{pgfscope}%
\pgfsys@transformshift{8.833731in}{2.010668in}%
\pgfsys@useobject{currentmarker}{}%
\end{pgfscope}%
\begin{pgfscope}%
\pgfsys@transformshift{8.824569in}{1.979480in}%
\pgfsys@useobject{currentmarker}{}%
\end{pgfscope}%
\begin{pgfscope}%
\pgfsys@transformshift{8.814290in}{1.948562in}%
\pgfsys@useobject{currentmarker}{}%
\end{pgfscope}%
\begin{pgfscope}%
\pgfsys@transformshift{8.802906in}{1.917945in}%
\pgfsys@useobject{currentmarker}{}%
\end{pgfscope}%
\begin{pgfscope}%
\pgfsys@transformshift{8.790427in}{1.887658in}%
\pgfsys@useobject{currentmarker}{}%
\end{pgfscope}%
\begin{pgfscope}%
\pgfsys@transformshift{8.776866in}{1.857733in}%
\pgfsys@useobject{currentmarker}{}%
\end{pgfscope}%
\begin{pgfscope}%
\pgfsys@transformshift{8.854406in}{2.039105in}%
\pgfsys@useobject{currentmarker}{}%
\end{pgfscope}%
\begin{pgfscope}%
\pgfsys@transformshift{8.846279in}{2.007227in}%
\pgfsys@useobject{currentmarker}{}%
\end{pgfscope}%
\begin{pgfscope}%
\pgfsys@transformshift{8.837015in}{1.975591in}%
\pgfsys@useobject{currentmarker}{}%
\end{pgfscope}%
\begin{pgfscope}%
\pgfsys@transformshift{8.826622in}{1.944228in}%
\pgfsys@useobject{currentmarker}{}%
\end{pgfscope}%
\begin{pgfscope}%
\pgfsys@transformshift{8.815111in}{1.913170in}%
\pgfsys@useobject{currentmarker}{}%
\end{pgfscope}%
\begin{pgfscope}%
\pgfsys@transformshift{8.802493in}{1.882448in}%
\pgfsys@useobject{currentmarker}{}%
\end{pgfscope}%
\begin{pgfscope}%
\pgfsys@transformshift{8.788782in}{1.852092in}%
\pgfsys@useobject{currentmarker}{}%
\end{pgfscope}%
\begin{pgfscope}%
\pgfsys@transformshift{8.867225in}{2.036082in}%
\pgfsys@useobject{currentmarker}{}%
\end{pgfscope}%
\begin{pgfscope}%
\pgfsys@transformshift{8.859007in}{2.003747in}%
\pgfsys@useobject{currentmarker}{}%
\end{pgfscope}%
\begin{pgfscope}%
\pgfsys@transformshift{8.849640in}{1.971658in}%
\pgfsys@useobject{currentmarker}{}%
\end{pgfscope}%
\begin{pgfscope}%
\pgfsys@transformshift{8.839131in}{1.939845in}%
\pgfsys@useobject{currentmarker}{}%
\end{pgfscope}%
\begin{pgfscope}%
\pgfsys@transformshift{8.827492in}{1.908342in}%
\pgfsys@useobject{currentmarker}{}%
\end{pgfscope}%
\begin{pgfscope}%
\pgfsys@transformshift{8.814733in}{1.877180in}%
\pgfsys@useobject{currentmarker}{}%
\end{pgfscope}%
\begin{pgfscope}%
\pgfsys@transformshift{8.800869in}{1.846389in}%
\pgfsys@useobject{currentmarker}{}%
\end{pgfscope}%
\begin{pgfscope}%
\pgfsys@transformshift{8.880227in}{2.033026in}%
\pgfsys@useobject{currentmarker}{}%
\end{pgfscope}%
\begin{pgfscope}%
\pgfsys@transformshift{8.871918in}{2.000229in}%
\pgfsys@useobject{currentmarker}{}%
\end{pgfscope}%
\begin{pgfscope}%
\pgfsys@transformshift{8.862445in}{1.967681in}%
\pgfsys@useobject{currentmarker}{}%
\end{pgfscope}%
\begin{pgfscope}%
\pgfsys@transformshift{8.851819in}{1.935414in}%
\pgfsys@useobject{currentmarker}{}%
\end{pgfscope}%
\begin{pgfscope}%
\pgfsys@transformshift{8.840050in}{1.903461in}%
\pgfsys@useobject{currentmarker}{}%
\end{pgfscope}%
\begin{pgfscope}%
\pgfsys@transformshift{8.827148in}{1.871853in}%
\pgfsys@useobject{currentmarker}{}%
\end{pgfscope}%
\begin{pgfscope}%
\pgfsys@transformshift{8.813129in}{1.840622in}%
\pgfsys@useobject{currentmarker}{}%
\end{pgfscope}%
\begin{pgfscope}%
\pgfsys@transformshift{8.893415in}{2.029935in}%
\pgfsys@useobject{currentmarker}{}%
\end{pgfscope}%
\begin{pgfscope}%
\pgfsys@transformshift{8.885012in}{1.996671in}%
\pgfsys@useobject{currentmarker}{}%
\end{pgfscope}%
\begin{pgfscope}%
\pgfsys@transformshift{8.875434in}{1.963659in}%
\pgfsys@useobject{currentmarker}{}%
\end{pgfscope}%
\begin{pgfscope}%
\pgfsys@transformshift{8.864688in}{1.930933in}%
\pgfsys@useobject{currentmarker}{}%
\end{pgfscope}%
\begin{pgfscope}%
\pgfsys@transformshift{8.852787in}{1.898524in}%
\pgfsys@useobject{currentmarker}{}%
\end{pgfscope}%
\begin{pgfscope}%
\pgfsys@transformshift{8.839741in}{1.866466in}%
\pgfsys@useobject{currentmarker}{}%
\end{pgfscope}%
\begin{pgfscope}%
\pgfsys@transformshift{8.825564in}{1.834791in}%
\pgfsys@useobject{currentmarker}{}%
\end{pgfscope}%
\begin{pgfscope}%
\pgfsys@transformshift{8.906790in}{2.026809in}%
\pgfsys@useobject{currentmarker}{}%
\end{pgfscope}%
\begin{pgfscope}%
\pgfsys@transformshift{8.898293in}{1.993073in}%
\pgfsys@useobject{currentmarker}{}%
\end{pgfscope}%
\begin{pgfscope}%
\pgfsys@transformshift{8.888607in}{1.959592in}%
\pgfsys@useobject{currentmarker}{}%
\end{pgfscope}%
\begin{pgfscope}%
\pgfsys@transformshift{8.877740in}{1.926401in}%
\pgfsys@useobject{currentmarker}{}%
\end{pgfscope}%
\begin{pgfscope}%
\pgfsys@transformshift{8.865705in}{1.893533in}%
\pgfsys@useobject{currentmarker}{}%
\end{pgfscope}%
\begin{pgfscope}%
\pgfsys@transformshift{8.852512in}{1.861019in}%
\pgfsys@useobject{currentmarker}{}%
\end{pgfscope}%
\begin{pgfscope}%
\pgfsys@transformshift{8.838175in}{1.828894in}%
\pgfsys@useobject{currentmarker}{}%
\end{pgfscope}%
\begin{pgfscope}%
\pgfsys@transformshift{8.920355in}{2.023649in}%
\pgfsys@useobject{currentmarker}{}%
\end{pgfscope}%
\begin{pgfscope}%
\pgfsys@transformshift{8.911763in}{1.989434in}%
\pgfsys@useobject{currentmarker}{}%
\end{pgfscope}%
\begin{pgfscope}%
\pgfsys@transformshift{8.901967in}{1.955480in}%
\pgfsys@useobject{currentmarker}{}%
\end{pgfscope}%
\begin{pgfscope}%
\pgfsys@transformshift{8.890978in}{1.921818in}%
\pgfsys@useobject{currentmarker}{}%
\end{pgfscope}%
\begin{pgfscope}%
\pgfsys@transformshift{8.878806in}{1.888484in}%
\pgfsys@useobject{currentmarker}{}%
\end{pgfscope}%
\begin{pgfscope}%
\pgfsys@transformshift{8.865465in}{1.855511in}%
\pgfsys@useobject{currentmarker}{}%
\end{pgfscope}%
\begin{pgfscope}%
\pgfsys@transformshift{8.850966in}{1.822930in}%
\pgfsys@useobject{currentmarker}{}%
\end{pgfscope}%
\begin{pgfscope}%
\pgfsys@transformshift{8.934112in}{2.020452in}%
\pgfsys@useobject{currentmarker}{}%
\end{pgfscope}%
\begin{pgfscope}%
\pgfsys@transformshift{8.925422in}{1.985755in}%
\pgfsys@useobject{currentmarker}{}%
\end{pgfscope}%
\begin{pgfscope}%
\pgfsys@transformshift{8.915515in}{1.951321in}%
\pgfsys@useobject{currentmarker}{}%
\end{pgfscope}%
\begin{pgfscope}%
\pgfsys@transformshift{8.904402in}{1.917184in}%
\pgfsys@useobject{currentmarker}{}%
\end{pgfscope}%
\begin{pgfscope}%
\pgfsys@transformshift{8.892093in}{1.883379in}%
\pgfsys@useobject{currentmarker}{}%
\end{pgfscope}%
\begin{pgfscope}%
\pgfsys@transformshift{8.878600in}{1.849940in}%
\pgfsys@useobject{currentmarker}{}%
\end{pgfscope}%
\begin{pgfscope}%
\pgfsys@transformshift{8.863938in}{1.816900in}%
\pgfsys@useobject{currentmarker}{}%
\end{pgfscope}%
\begin{pgfscope}%
\pgfsys@transformshift{8.948063in}{2.017220in}%
\pgfsys@useobject{currentmarker}{}%
\end{pgfscope}%
\begin{pgfscope}%
\pgfsys@transformshift{8.939275in}{1.982033in}%
\pgfsys@useobject{currentmarker}{}%
\end{pgfscope}%
\begin{pgfscope}%
\pgfsys@transformshift{8.929255in}{1.947114in}%
\pgfsys@useobject{currentmarker}{}%
\end{pgfscope}%
\begin{pgfscope}%
\pgfsys@transformshift{8.918016in}{1.912497in}%
\pgfsys@useobject{currentmarker}{}%
\end{pgfscope}%
\begin{pgfscope}%
\pgfsys@transformshift{8.905567in}{1.878216in}%
\pgfsys@useobject{currentmarker}{}%
\end{pgfscope}%
\begin{pgfscope}%
\pgfsys@transformshift{8.891921in}{1.844306in}%
\pgfsys@useobject{currentmarker}{}%
\end{pgfscope}%
\begin{pgfscope}%
\pgfsys@transformshift{8.877092in}{1.810800in}%
\pgfsys@useobject{currentmarker}{}%
\end{pgfscope}%
\begin{pgfscope}%
\pgfsys@transformshift{8.962210in}{2.013950in}%
\pgfsys@useobject{currentmarker}{}%
\end{pgfscope}%
\begin{pgfscope}%
\pgfsys@transformshift{8.953322in}{1.978270in}%
\pgfsys@useobject{currentmarker}{}%
\end{pgfscope}%
\begin{pgfscope}%
\pgfsys@transformshift{8.943188in}{1.942860in}%
\pgfsys@useobject{currentmarker}{}%
\end{pgfscope}%
\begin{pgfscope}%
\pgfsys@transformshift{8.931821in}{1.907757in}%
\pgfsys@useobject{currentmarker}{}%
\end{pgfscope}%
\begin{pgfscope}%
\pgfsys@transformshift{8.919230in}{1.872995in}%
\pgfsys@useobject{currentmarker}{}%
\end{pgfscope}%
\begin{pgfscope}%
\pgfsys@transformshift{8.905429in}{1.838608in}%
\pgfsys@useobject{currentmarker}{}%
\end{pgfscope}%
\begin{pgfscope}%
\pgfsys@transformshift{8.890432in}{1.804632in}%
\pgfsys@useobject{currentmarker}{}%
\end{pgfscope}%
\begin{pgfscope}%
\pgfsys@transformshift{8.976556in}{2.010644in}%
\pgfsys@useobject{currentmarker}{}%
\end{pgfscope}%
\begin{pgfscope}%
\pgfsys@transformshift{8.967566in}{1.974463in}%
\pgfsys@useobject{currentmarker}{}%
\end{pgfscope}%
\begin{pgfscope}%
\pgfsys@transformshift{8.957317in}{1.938558in}%
\pgfsys@useobject{currentmarker}{}%
\end{pgfscope}%
\begin{pgfscope}%
\pgfsys@transformshift{8.945820in}{1.902963in}%
\pgfsys@useobject{currentmarker}{}%
\end{pgfscope}%
\begin{pgfscope}%
\pgfsys@transformshift{8.933086in}{1.867713in}%
\pgfsys@useobject{currentmarker}{}%
\end{pgfscope}%
\begin{pgfscope}%
\pgfsys@transformshift{8.919127in}{1.832845in}%
\pgfsys@useobject{currentmarker}{}%
\end{pgfscope}%
\begin{pgfscope}%
\pgfsys@transformshift{8.903958in}{1.798393in}%
\pgfsys@useobject{currentmarker}{}%
\end{pgfscope}%
\end{pgfscope}%
\begin{pgfscope}%
\pgfpathrectangle{\pgfqpoint{5.486220in}{0.585980in}}{\pgfqpoint{4.437388in}{3.328041in}}%
\pgfusepath{clip}%
\pgfsetbuttcap%
\pgfsetroundjoin%
\definecolor{currentfill}{rgb}{1.000000,0.000000,0.000000}%
\pgfsetfillcolor{currentfill}%
\pgfsetlinewidth{1.003750pt}%
\definecolor{currentstroke}{rgb}{1.000000,0.000000,0.000000}%
\pgfsetstrokecolor{currentstroke}%
\pgfsetdash{}{0pt}%
\pgfsys@defobject{currentmarker}{\pgfqpoint{-0.034722in}{-0.034722in}}{\pgfqpoint{0.034722in}{0.034722in}}{%
\pgfpathmoveto{\pgfqpoint{0.000000in}{-0.034722in}}%
\pgfpathcurveto{\pgfqpoint{0.009208in}{-0.034722in}}{\pgfqpoint{0.018041in}{-0.031064in}}{\pgfqpoint{0.024552in}{-0.024552in}}%
\pgfpathcurveto{\pgfqpoint{0.031064in}{-0.018041in}}{\pgfqpoint{0.034722in}{-0.009208in}}{\pgfqpoint{0.034722in}{0.000000in}}%
\pgfpathcurveto{\pgfqpoint{0.034722in}{0.009208in}}{\pgfqpoint{0.031064in}{0.018041in}}{\pgfqpoint{0.024552in}{0.024552in}}%
\pgfpathcurveto{\pgfqpoint{0.018041in}{0.031064in}}{\pgfqpoint{0.009208in}{0.034722in}}{\pgfqpoint{0.000000in}{0.034722in}}%
\pgfpathcurveto{\pgfqpoint{-0.009208in}{0.034722in}}{\pgfqpoint{-0.018041in}{0.031064in}}{\pgfqpoint{-0.024552in}{0.024552in}}%
\pgfpathcurveto{\pgfqpoint{-0.031064in}{0.018041in}}{\pgfqpoint{-0.034722in}{0.009208in}}{\pgfqpoint{-0.034722in}{0.000000in}}%
\pgfpathcurveto{\pgfqpoint{-0.034722in}{-0.009208in}}{\pgfqpoint{-0.031064in}{-0.018041in}}{\pgfqpoint{-0.024552in}{-0.024552in}}%
\pgfpathcurveto{\pgfqpoint{-0.018041in}{-0.031064in}}{\pgfqpoint{-0.009208in}{-0.034722in}}{\pgfqpoint{0.000000in}{-0.034722in}}%
\pgfpathlineto{\pgfqpoint{0.000000in}{-0.034722in}}%
\pgfpathclose%
\pgfusepath{stroke,fill}%
}%
\begin{pgfscope}%
\pgfsys@transformshift{7.150240in}{2.250000in}%
\pgfsys@useobject{currentmarker}{}%
\end{pgfscope}%
\begin{pgfscope}%
\pgfsys@transformshift{8.259587in}{2.250000in}%
\pgfsys@useobject{currentmarker}{}%
\end{pgfscope}%
\end{pgfscope}%
\begin{pgfscope}%
\pgfsetrectcap%
\pgfsetmiterjoin%
\pgfsetlinewidth{0.803000pt}%
\definecolor{currentstroke}{rgb}{0.000000,0.000000,0.000000}%
\pgfsetstrokecolor{currentstroke}%
\pgfsetdash{}{0pt}%
\pgfpathmoveto{\pgfqpoint{5.486220in}{0.585980in}}%
\pgfpathlineto{\pgfqpoint{5.486220in}{3.914020in}}%
\pgfusepath{stroke}%
\end{pgfscope}%
\begin{pgfscope}%
\pgfsetrectcap%
\pgfsetmiterjoin%
\pgfsetlinewidth{0.803000pt}%
\definecolor{currentstroke}{rgb}{0.000000,0.000000,0.000000}%
\pgfsetstrokecolor{currentstroke}%
\pgfsetdash{}{0pt}%
\pgfpathmoveto{\pgfqpoint{9.923608in}{0.585980in}}%
\pgfpathlineto{\pgfqpoint{9.923608in}{3.914020in}}%
\pgfusepath{stroke}%
\end{pgfscope}%
\begin{pgfscope}%
\pgfsetrectcap%
\pgfsetmiterjoin%
\pgfsetlinewidth{0.803000pt}%
\definecolor{currentstroke}{rgb}{0.000000,0.000000,0.000000}%
\pgfsetstrokecolor{currentstroke}%
\pgfsetdash{}{0pt}%
\pgfpathmoveto{\pgfqpoint{5.486220in}{0.585980in}}%
\pgfpathlineto{\pgfqpoint{9.923608in}{0.585980in}}%
\pgfusepath{stroke}%
\end{pgfscope}%
\begin{pgfscope}%
\pgfsetrectcap%
\pgfsetmiterjoin%
\pgfsetlinewidth{0.803000pt}%
\definecolor{currentstroke}{rgb}{0.000000,0.000000,0.000000}%
\pgfsetstrokecolor{currentstroke}%
\pgfsetdash{}{0pt}%
\pgfpathmoveto{\pgfqpoint{5.486220in}{3.914020in}}%
\pgfpathlineto{\pgfqpoint{9.923608in}{3.914020in}}%
\pgfusepath{stroke}%
\end{pgfscope}%
\begin{pgfscope}%
\definecolor{textcolor}{rgb}{0.000000,0.000000,0.000000}%
\pgfsetstrokecolor{textcolor}%
\pgfsetfillcolor{textcolor}%
\pgftext[x=7.704914in,y=3.997354in,,base]{\color{textcolor}{\rmfamily\fontsize{12.000000}{14.400000}\selectfont\catcode`\^=\active\def^{\ifmmode\sp\else\^{}\fi}\catcode`\%=\active\def%{\%}w-plane: $w = \cos(z)$}}%
\end{pgfscope}%
\begin{pgfscope}%
\pgfsetbuttcap%
\pgfsetmiterjoin%
\definecolor{currentfill}{rgb}{1.000000,1.000000,1.000000}%
\pgfsetfillcolor{currentfill}%
\pgfsetfillopacity{0.800000}%
\pgfsetlinewidth{1.003750pt}%
\definecolor{currentstroke}{rgb}{0.800000,0.800000,0.800000}%
\pgfsetstrokecolor{currentstroke}%
\pgfsetstrokeopacity{0.800000}%
\pgfsetdash{}{0pt}%
\pgfpathmoveto{\pgfqpoint{5.583442in}{3.594576in}}%
\pgfpathlineto{\pgfqpoint{6.612439in}{3.594576in}}%
\pgfpathquadraticcurveto{\pgfqpoint{6.640217in}{3.594576in}}{\pgfqpoint{6.640217in}{3.622354in}}%
\pgfpathlineto{\pgfqpoint{6.640217in}{3.816798in}}%
\pgfpathquadraticcurveto{\pgfqpoint{6.640217in}{3.844576in}}{\pgfqpoint{6.612439in}{3.844576in}}%
\pgfpathlineto{\pgfqpoint{5.583442in}{3.844576in}}%
\pgfpathquadraticcurveto{\pgfqpoint{5.555664in}{3.844576in}}{\pgfqpoint{5.555664in}{3.816798in}}%
\pgfpathlineto{\pgfqpoint{5.555664in}{3.622354in}}%
\pgfpathquadraticcurveto{\pgfqpoint{5.555664in}{3.594576in}}{\pgfqpoint{5.583442in}{3.594576in}}%
\pgfpathlineto{\pgfqpoint{5.583442in}{3.594576in}}%
\pgfpathclose%
\pgfusepath{stroke,fill}%
\end{pgfscope}%
\begin{pgfscope}%
\pgfsetbuttcap%
\pgfsetroundjoin%
\definecolor{currentfill}{rgb}{1.000000,0.000000,0.000000}%
\pgfsetfillcolor{currentfill}%
\pgfsetlinewidth{1.003750pt}%
\definecolor{currentstroke}{rgb}{1.000000,0.000000,0.000000}%
\pgfsetstrokecolor{currentstroke}%
\pgfsetdash{}{0pt}%
\pgfsys@defobject{currentmarker}{\pgfqpoint{-0.034722in}{-0.034722in}}{\pgfqpoint{0.034722in}{0.034722in}}{%
\pgfpathmoveto{\pgfqpoint{0.000000in}{-0.034722in}}%
\pgfpathcurveto{\pgfqpoint{0.009208in}{-0.034722in}}{\pgfqpoint{0.018041in}{-0.031064in}}{\pgfqpoint{0.024552in}{-0.024552in}}%
\pgfpathcurveto{\pgfqpoint{0.031064in}{-0.018041in}}{\pgfqpoint{0.034722in}{-0.009208in}}{\pgfqpoint{0.034722in}{0.000000in}}%
\pgfpathcurveto{\pgfqpoint{0.034722in}{0.009208in}}{\pgfqpoint{0.031064in}{0.018041in}}{\pgfqpoint{0.024552in}{0.024552in}}%
\pgfpathcurveto{\pgfqpoint{0.018041in}{0.031064in}}{\pgfqpoint{0.009208in}{0.034722in}}{\pgfqpoint{0.000000in}{0.034722in}}%
\pgfpathcurveto{\pgfqpoint{-0.009208in}{0.034722in}}{\pgfqpoint{-0.018041in}{0.031064in}}{\pgfqpoint{-0.024552in}{0.024552in}}%
\pgfpathcurveto{\pgfqpoint{-0.031064in}{0.018041in}}{\pgfqpoint{-0.034722in}{0.009208in}}{\pgfqpoint{-0.034722in}{0.000000in}}%
\pgfpathcurveto{\pgfqpoint{-0.034722in}{-0.009208in}}{\pgfqpoint{-0.031064in}{-0.018041in}}{\pgfqpoint{-0.024552in}{-0.024552in}}%
\pgfpathcurveto{\pgfqpoint{-0.018041in}{-0.031064in}}{\pgfqpoint{-0.009208in}{-0.034722in}}{\pgfqpoint{0.000000in}{-0.034722in}}%
\pgfpathlineto{\pgfqpoint{0.000000in}{-0.034722in}}%
\pgfpathclose%
\pgfusepath{stroke,fill}%
}%
\begin{pgfscope}%
\pgfsys@transformshift{5.750109in}{3.733465in}%
\pgfsys@useobject{currentmarker}{}%
\end{pgfscope}%
\end{pgfscope}%
\begin{pgfscope}%
\definecolor{textcolor}{rgb}{0.000000,0.000000,0.000000}%
\pgfsetstrokecolor{textcolor}%
\pgfsetfillcolor{textcolor}%
\pgftext[x=6.000109in,y=3.684854in,left,base]{\color{textcolor}{\rmfamily\fontsize{10.000000}{12.000000}\selectfont\catcode`\^=\active\def^{\ifmmode\sp\else\^{}\fi}\catcode`\%=\active\def%{\%}Foci ($\pm 1$)}}%
\end{pgfscope}%
\end{pgfpicture}%
\makeatother%
\endgroup%
}
    \caption{Mapping $z \mapsto \cos(z)$. Vertical lines $x=c$ map to hyperbolas (foci $\pm 1$), and horizontal lines $y=d$ map to ellipses (foci $\pm 1$).}
    \label{fig:mapping_cos}
  \end{figure}
\end{solution}

\section{Exercise 1.116}
Calculate $i^i$.

\begin{solution}
  \[
    i^i = (e^{\pi i / 2})^{i} = e^{\frac{1}{2} \pi (i \cdot i)} =
    e^{-\frac{\pi}{2}}
  \]
\end{solution}

\section{Exercise 1.117*}
We have been augmenting sets of numbers: $\mathbb{N} \subset
\mathbb{Z} \subset \mathbb{Q} \subset \mathbb{R} \subset \mathbb{C}$.
For $n \in \mathbb{Z}$ and $x, z \in X$ where $X$ is a set in the
first column, is $X$ closed under the operations in the first row?
\begin{center}
  \begin{tabular}{|c|c|c|c|c|}
    \hline
    & $+/-$ & $\times/\div$ & $(\cdot)^n$ & $e^z, z^x$ \\ \hline
    $\mathbb{N}$ & & & & \\ \hline
    $\mathbb{Z}$ & & & & \\ \hline
    $\mathbb{Q}$ & & & & \\ \hline
    $\mathbb{R}$ & & & & \\ \hline
    $\mathbb{C}$ & & & & \\ \hline
  \end{tabular}
\end{center}
Use the above table to tell a story of marching from $\mathbb{N}$ to
$\mathbb{C}$.

\begin{solution}
  We fill in the table as follows:
  \begin{center}
    \begin{tabular}{|c|c|c|c|c|}
      \hline
      & $+/-$ & $\times/\div$ & $(\cdot)^n$ & $e^z, z^x$ \\ \hline
      $\mathbb{N}$ & No & No & No & No \\ \hline
      $\mathbb{Z}$ & Yes & No & No & No \\ \hline
      $\mathbb{Q}$ & Yes & Yes & Yes & No \\ \hline
      $\mathbb{R}$ & Yes & Yes & Yes & No \\ \hline
      $\mathbb{C}$ & Yes & Yes & Yes & Yes \\ \hline
    \end{tabular}
  \end{center}

  Firstly we construct $\N$ by von Neumann's construction based on
  ZFC, where each
  natural number is defined as the set of all smaller natural numbers:
  \[
    0 = \varnothing, n + 1 = n \cup \{n\}.
  \]
  Then we construct $\Z$ by taking formal differences of natural numbers:
  \[
    \Z = \{(a, b) : a, b \in \N\} / \sim, \quad (a, b) \sim (c, d)
    \iff a + d = b + c.
  \]
  This makes $\Z$ an abelian group under addition, but it is not closed under
  multiplication or exponentiation. Next, we construct $\Q$ by taking formal
  quotients of integers:
  \[
    \Q = \{(a, b) : a \in \Z, b \in \Z \setminus \{0\}\} / \sim,
    \quad (a, b) \sim (c, d) \iff ad = bc.
  \]
  This makes $\Q$ a field, but it is not closed under exponentiation.
  Next, we construct $\R$ by taking equivalence classes of rational
  Cauchy sequences, which makes $\R$ a complete ordered field, but it
  is not closed under exponentiation. Finally, we construct $\C$ by
  taking ordered pairs of real numbers with the usual addition and
  multiplication, which makes $\C$ an algebraically closed field that is
  closed under exponentiation.

\end{solution}

\section{Exercise 1.120}
For an $\mathbb{R}$-linear map $A : \mathbb{C} \to \mathbb{C}$, the
following statements are equivalent:
\begin{enumerate}[(a)]
  \item $\exists \ell \in \mathbb{C}$ s.t. $\forall z \in \mathbb{C},
    A(z) = \ell z$.
  \item $A$ is $\mathbb{C}$-linear.
  \item $A(i) = iA(1)$.
  \item The matrix of $A$ with respect to the canonical basis $1 =
    (1, 0)$ and $i = (0, 1)$ has the special form
    \[
      \begin{bmatrix}
        \alpha & -\beta \\
        \beta & \alpha
      \end{bmatrix}, \quad \alpha, \beta \in \mathbb{R}.
    \]
\end{enumerate}

\begin{solution}
  (a) $\Rightarrow$ (b): Suppose $A(z) = \ell z$ for some $\ell \in \C$.
  To show $A$ is $\C$-linear, we need to verify that for any
  $c \in \C$ and $z \in \C$, we have $A(cz) = c \cdot A(z)$.

  Write $\ell = \alpha + i\beta$ with $\alpha, \beta \in \R$, and
  write $c = a + ib$, $z = x + iy$ with $a, b, x, y \in \R$. Then
  \[
    cz = (a+ib)(x+iy) = (ax - by) + i(ay + bx).
  \]
  Applying $A$ to $cz$, we obtain
  \[
    \begin{aligned}
      A(cz) &= \ell \cdot cz \\
      &= (\alpha + i\beta)\big((ax - by) + i(ay + bx)\big) \\
      &= \alpha(ax - by) - \beta(ay + bx)
      + i\big(\beta(ax - by) + \alpha(ay + bx)\big) \\
      &= (\alpha a - \beta b)x + (-\alpha b - \beta a)y
      + i\big((\beta a + \alpha b)x + (-\beta b + \alpha a)y\big).
    \end{aligned}
  \]

  On the other hand,
  \[
    \begin{aligned}
      c \cdot A(z) &= (a + ib) \cdot (\ell z) \\
      &= (a + ib) \cdot \big((\alpha + i\beta)(x + iy)\big) \\
      &= (a + ib) \cdot \big((\alpha x - \beta y) + i(\beta x +
      \alpha y)\big) \\
      &= a(\alpha x - \beta y) - b(\beta x + \alpha y)
      + i\big(b(\alpha x - \beta y) + a(\beta x + \alpha y)\big) \\
      &= (\alpha a - \beta b)x + (-\alpha b - \beta a)y
      + i\big((\beta a + \alpha b)x + (-\beta b + \alpha a)y\big).
    \end{aligned}
  \]

  The two expressions are identical, so $A(cz) = c \cdot A(z)$.
  Hence $A$ is $\C$-linear.

  (b) $\Rightarrow$ (c): If $A$ is $\C$-linear, then $A(i) = A(i \cdot 1) =
  i A(1)$.

  (c) $\Rightarrow$ (d): Let $A(1) = \alpha + i\beta$. Then $A(i) = iA(1)
  = i(\alpha + i\beta) = -\beta + i\alpha$. Thus the matrix of $A$ with
  respect to the basis $(1, i)$ is
  \[
    \begin{bmatrix}
      \alpha & -\beta \\
      \beta & \alpha
    \end{bmatrix}.
  \]

  (d) $\Rightarrow$ (a): If the matrix of $A$ has the form in (d), then
  for any $z = x + iy$, we have
  \[
    A(z) = A(x + iy) = xA(1) + yA(i) = (\alpha x - \beta y) + i(\beta x +
    \alpha y).
  \]
  Let $\ell = \alpha + i\beta$. Then $\ell z = (\alpha + i\beta)(x + iy)
  = (\alpha x - \beta y) + i(\beta x + \alpha y)$. Thus $A(z) = \ell z$.
\end{solution}

\end{document}
